\documentclass[11pt,envcountchap,pdf]{svmono}
 
\usepackage[pdftex]{graphicx}


%\usepackage{algorithm2e} 

\usepackage{tikz}
\usetikzlibrary{arrows,%
                petri,%
                topaths,%
                positioning}%
\usepackage{tkz-berge}



\usepackage{listings}
\lstset{ 
  language = Python
}

\lstloadlanguages{Python}







% Define box and box title style
\tikzstyle{mybox} = [draw=black, fill=lightgray!20,
    rectangle,  inner sep=10pt, inner ysep=5pt]
\tikzstyle{fancytitle} =[fill=lightgray!40]

\newcommand{\fancybox}[2][Title of the box]{%
  \begin{center}
    \begin{tikzpicture}
      \node [mybox] (box){%       
        \begin{minipage}{0.8\textwidth}
          #1\\[1ex]          
          #2
        \end{minipage}
      };
%\node[fancytitle, right=10pt] at (box.north west) {#1};
%\node[fancytitle, rounded corners] at (box.east) {$\clubsuit$};
\end{tikzpicture}%
\end{center}
}





\newtheorem{assumption}{Assumption}


\usepackage[T1]{fontenc}
\usepackage[utf8]{inputenc} 
\usepackage[american]{babel}

% \usepackage{pstricks}
% \usepackage{pst-node}
% \usepackage{pst-coil}

% \newgray{vlg}{.85}
% \newgray{vvlg}{.95}


\usepackage{colortbl}
\usepackage{caption}
\usepackage{subcaption}
\usepackage{amssymb}
\usepackage{amsmath}



\usepackage{mathrsfs}
\usepackage{psfrag}
\usepackage{hyperref}

     
\usepackage{enumerate}

\usepackage{utf8math}


\usepackage{xspace}
\newcommand{\plw}{PLW\xspace}
\newcommand{\psv}{PSV\xspace}
\newcommand{\ip}{IP\xspace}

%\usepackage{showlabels}
\usepackage{tabularx}



%Math Operators

\DeclareMathOperator{\size}{size}
\DeclareMathOperator{\conv}{conv}
\newcommand{\SV}{\mathrm{SV}}
\newcommand{\eps}{{\varepsilon}}
\newcommand{\beps}{\mathbf{\epsilon}}
\newcommand{\bigO}{O}
\newcommand{\cut}{\mathrm{cut}}
\newcommand{\LLL}{\mathrm{LLL}}
\newcommand{\setR}{\mathbb{R}}
\newcommand{\setZ}{\mathbb{Z}}
\newcommand{\setQ}{\mathbb{Q}}
\newcommand{\setC}{\mathbb{C}}
\newcommand{\setN}{\mathbb{N}}
\newcommand{\wt}[1]{\widetilde{#1}}
\newcommand{\opt}{{\sc 0/1-opt}\xspace}
\newcommand{\aug}{{\sc 0/1-aug}\xspace}
\newcommand{\psep}{{\sc 0/1-psep}\xspace}
\newcommand{\sep}{{\sc 0/1-sep}\xspace}
\newcommand{\fopt}{{\sc 0/1-testopt\xspace} }

\newcommand{\hpp}{\mathrm{HPP}}
\newcommand{\nodes}{\mathcal{V}}
\newcommand{\vol}{\mathrm{vol}}
\newcommand{\diag}{\mathrm{diag}}
\newcommand{\arcs}{\mathcal{A}}
\newcommand{\edges}{\mathcal{E}}
\newcommand{\paths}{\mathscr{P}}
\newcommand{\cycles}{\mathcal{C}}




\newcommand{\K}{{\mathcal K}}
\newcommand{\A}{{A}}
\newcommand{\B}{{B}}
\newcommand{\T}{\mathscr{T}}
\newcommand{\eE}{\mathscr{E}}
\newcommand{\eS}{\mathscr{S}}
\newcommand{\eP}{\mathscr{P}}
\newcommand{\eM}{\mathscr{M}}



\newcommand{\transp}{^{\mathrm{T}}}

\newcommand{\smallmat}[1]{\left( \begin{smallmatrix} #1 \end{smallmatrix}\right)}

\newcommand{\mat}[1]{ \begin{pmatrix} #1 \end{pmatrix}}
\newcommand{\smat}[1]{ \big(\begin{smallmatrix} #1 \end{smallmatrix}\big)}

\newcommand{\pc}{\mathscr{P}}
\newcommand{\ob}{\mathscr{O}}
\newcommand{\odds}{\mathscr{W}}
\newcommand{\up}{\mathscr{U}}
\newcommand{\ef}{\mathscr{F}}
\newcommand{\eh}{\mathscr{H}}
\newcommand{\ev}{\mathscr{V}}
\newcommand{\ec}{\mathscr{C}}
\newcommand{\eu}{\mathscr{U}}

\newcommand{\lex}{\mathrm{lex}}

\renewcommand{\leq}{\leqslant}
\renewcommand{\geq}{\geqslant}





\newcommand{\R}{\setR}
\newcommand{\Z}{\setZ}
\newcommand{\N}{\setN}



\newcommand{\linhull}{\mathrm{lin.hull}}
\newcommand{\affhull}{\mathrm{affine.hull}}
\newcommand{\charcone}{\mathrm{char.cone}}
\newcommand{\cone}{\mathrm{cone}}
\newcommand{\rank}{\mathrm{rank}}
\newcommand{\wb}[1]{\overline{#1}}





\newenvironment{dproblem}[2]{
\begin{center}
\psframebox{
  \begin{minipage}{4in}  
      
    {\sc #1 \\}
    #2
  \end{minipage}
}
\end{center}}



\date{\today}
\author{Friedrich Eisenbrand} 
\title{Linear Optimization}
\subtitle{These are notes of my course Discrete Optimization. They are constantly updated. }


\includeonly{algorithms,convex}

\begin{document}
\maketitle

%\Extrachap{Linear Optimization} 

%\maketitle
\tableofcontents 

\spdefaulttheorem{algorithm}{Algorithm}{}{}
\smartqed


\definecolor{Gray}{gray}{0.6} 
\newcolumntype{A}{%
>{\centering}X} 
\newcolumntype{B}{%
>{\columncolor{Gray}\centering}X} 
\newcolumntype{R}{%
>{\columncolor{red}\centering}X} 
\newcolumntype{G}{%
>{\columncolor{Gray}\centering}c} 

\newcolumntype{E}{%
>{\columncolor{green}\centering}X} 

\newcolumntype{C}{>{\bfseries\centering\arraybackslash}X}

\chapter*{Preface} 

A central topic in linear algebra is the theory around  linear equations 
\begin{equation} 
\label{p:eq:1}
  A \, x = b
\end{equation}
where $A ∈ ℝ^{m×n}$ is a given matrix and $b ∈ ℝ^m$ is a given vector. Every student of mathematics learns learns to appreciate \emph{Gaussian elimination} which transforms the system~\eqref{p:eq:1}  into an equivalent system  
\begin{displaymath}
  {A}' x = {b}'
\end{displaymath}
that is in row-echelon form. This means that ${A}' ∈ ℝ^{m×n}$ is such that 
  for each $i ∈ \{ 1,\dots , m-1\}$
  \begin{displaymath}
    \min\{ j : {a}'_{ij} \neq 0\} <  \min\{ j : {a}'_{i+1j} \neq 0\},
  \end{displaymath}
see Figure~\ref{fig:8}. 


\begin{figure}
  \centering
 \includegraphics[height=5cm]{figures/RowEchelon.jpg}
  \caption{A schematic image of a matrix in row-echelon form} 
\label{fig:8}
\end{figure}


The system~\eqref{p:eq:1} has a solution if and only if ${b}'$ has only zero's in those components that correspond to the zero-rows of ${A}'$ and a solution is readily computed. In fact, Gaussian elimination also provides an invertible matrix $Q ∈ ℝ^{m ×m}$ such that $Q\cdot A = {A}'$ and $Q \cdot b = {b}'$. If the $i$-th component of ${b}'$ is nonzero and the $i$-th row of ${A}'$ consists of zeros only, then with $q$ being the column vector corresponding to the $i$-th row of $Q$ one has 
\begin{displaymath}
  q^T A = 0 \text{ and } q^Tb \neq 0. 
\end{displaymath}
Thus there is a convenient \emph{certificate} of the fact that \eqref{p:eq:1} is not solvable. 


\begin{theorem}
  \label{thr:8}
  A linear system~\eqref{p:eq:1} is not solvable if and only if there exists a vector $q ∈ ℝ^m$ such that 
  \begin{displaymath}
    q^T A = 0 \text{ and } q^T b \neq 0. 
  \end{displaymath}
\end{theorem}

In this course, we will now also deal with systems of \emph{linear inequalities} 
\begin{equation}
  \label{p:ineq}
  A \,x ≤ b,
\end{equation}
where $A ∈ ℝ^{m × n}$ and $b ∈ ℝ^m$. A \emph{solution} to such a system is a vector $x^* ∈ ℝ^m$ such that
\begin{equation}
  \label{p:sol}
 \text{for each } i=1,\dots,m \quad : \quad   a_{i1} x^*_1 + \cdots + a_{in} x^*_n \leq b_i.  
\end{equation}
If a solution exists, the system~\eqref{p:ineq} is called \emph{feasible} otherwise it is called \emph{infeasible}. 
We will ask analogous questions for  systems of linear inequalities as we did for linear equations. Can we efficiently find a solution of~\eqref{p:ineq}?  Is there a simple certificate to convince somebody of the infeasibility of \eqref{p:ineq}? 

To shed a bit of light on this second question, let us consider a vector $λ ∈ ℝ^m_{≥0}$. If $x^*$ is a solution, then $x^*$ also satisfies the inequality 
\begin{equation}
  \label{p:impli}
  λ^T A \, x ≤ λ^T b.
\end{equation}
If there exists a $λ \in ℝ^m_{≥0}$ such that 
\begin{equation}
  \label{p:Farkas}
  λ^T A = 0^T \text{ and } λ^T b = -1,
\end{equation}
then this $λ$ \emph{certifies} that~\eqref{p:ineq} is infeasible. In the case of infeasibility, can such a $λ≥0$ always be found? Among the many results presented in  this course,  we will show that the answer is ``yes''. 
\begin{theorem}[Farkas' lemma]
  \label{thr:7}
  A system of linear inequalities~\eqref{p:ineq} is infeasible if and only if there exists a $λ ∈ ℝ^m_{≥0}$ such that 
  \begin{displaymath}
    λ^T A = 0 \text{ and } λ^T b = -1.
  \end{displaymath}
\end{theorem}


Central to this course is the \emph{linear programming} or \emph{linear optimization} problem.  There, one is given a system of linear inequalities~\eqref{p:ineq} together with a linear objective function $f(x) = c^Tx$ for some fixed $c ∈ℝ^n$. The task is to solve 
\begin{equation}
  \label{eq:1}
  \max\{ c^Tx : x ∈ ℝ^n, \, Ax \leq b\}, 
\end{equation}
or in other words, to find an $x^* ∈ ℝ^n$ that is a solution of \eqref{p:ineq} such that $c^T x^* ≥ c^T y^*$ for each solution $y^*$ of \eqref{p:ineq}, if such an $x^*$ exists. 
The linear optimization problem is one of the most important types of mathematical optimization problems. It has numerous applications in science and engineering and, in particular, modern fields like \emph{machine learning}.

In this course, we will learn  the theory of linear optimization and develop \emph{efficient algorithms} to solve linear programming problems. This means that we do not content ourselves with an algorithm that is correct, but we want this algorithm to be capable to solve large scale problems within a short time. 

\emph{Efficient algorithms} is a subfield of computational mathematics on its own. Before we begin with our considerations on linear programming, we will have to learn the basics of it. 






\section*{Exercises}

\begin{enumerate}[1)]
\item Provide a certificate as in Theorem~\ref{thr:8} of the unsolvability of the linear equation 
  \begin{displaymath}
    \begin{pmatrix}
      2 & 1 & 0 \\
      5 & 4 & 1 \\
      7 & 5 & 1
    \end{pmatrix} \,
    \begin{pmatrix}
      x_1 \\ x_2 \\ x_3 
    \end{pmatrix} =
    \begin{pmatrix}
      1\\2\\4
    \end{pmatrix}
  \end{displaymath}
\item Show the ``if'' direction of the Farkas' lemma (Theorem~\ref{thr:7}). 
\end{enumerate}




\chapter{Algorithms and running time analysis} 
\label{cha:runn-time-analys}



An \emph{algorithm} executes a set of \emph{instructions}
used in common programming languages like  arithmetic
operations, comparisons or read/write instructions. The sequence of these instructions is controlled by \emph{loops and conditionals} like {\tt if, while, for} etc. 

Each of these instructions requires \emph{time}. The \emph{running time} of an algorithm is the \emph{number of instructions} that the algorithm performs. This number depends on the \emph{input} of the algorithm. 


\begin{example}
\label{ex-a-1}
  Consider the following algorithm to compute the product of two $n × n$ matrices $A,B ∈ ℚ^{n ×n}$: 
    \begin{tabbing}            
    {\bf for} \= $i=1,\dots,n$ \\ 
              \> {\bf for} \= $j=1,\dots,n$ \\
              \> \> $c_{ij} := 0$ \\
              \> \> {\bf for} \= $k=1,\dots,n$ \\
              \> \> \> $c_{ij} := c_{ij} + a_{ik} ⋅ a_{kj}$
            \end{tabbing} 
%
 The number of additions that are carried out is $n^2 ⋅ (n-1)$ and the number of multiplications is $n^3$.
 The number of store-instructions is $n^2 ⋅(n+1)$. The number of read-instructions is of similar magnitude. 
\end{example}


The above example shows that an \emph{exact counting} is sometimes
tedious. Looking at the algorithm however, you quickly agree that
there exists \emph{some constant} $c \in \R_{>0}$ such that the
algorithm performs at most $c ⋅ n^3$ instructions. 

In the analysis of algorithms, one does usually not care so much about the constant $c$ above in the beginning. There are sub-fields of algorithms where this constant however matters. Especially for algorithms on large data sets, where access to external data is costly. However, this is another story that does not concern us here. When we analyze algorithms, we are interested in the \emph{asymptotic running time}. 

\begin{definition}[$O$-notation]  

 Let $T,f: \setN \to \setR_{\geq0}$ be functions. We say 
  \begin{itemize}
  \item \emph{$T(n) = O(f(n))$}, if there exist positive  constants
    $n_o\in \setN$ and  $c\in\setR_{>0}$ 
    with $$T(n) \leq c \cdot f(n) \text{ for all }n\geq n_0.$$  
  \item \emph{$T(n) = \Omega(f(n))$}, if there exist constants   $n_o\in \setN$
    and  $c\in\setR_{>0}$ 
    with $$T(n) \geq c \cdot f(n)\text{ for all } n\geq n_0.$$
  \item \emph{$T(n) = \Theta(f(n))$}  if $$T(n)=O(f(n))\text{ and }
     T(n) = \Omega(f(n)).$$
  \end{itemize}
\end{definition}


\begin{example}
  The function $T(n)=2n^2 + 3n +1$ is in $O(n^2)$, since for all
  $n\geq1$ one has $2n^2 + 3n + 1 \leq 6n^2$. Here $n_0 = 1$ and $c =
  6$.  Similarly $T(n) = \Omega(n^2)$, since for each $n\geq1$ one has $2n^2 + 3n
  +1\geq n^2$. Thus $T(n)$ is in $\Theta(n^2)$. 
\end{example}

We measure the running time of algorithms in terms of the \emph{length
  of the input}.  The matrices $A$ and $B $ that are the input of the matrix-multiplication
algorithm of Example~\ref{ex-a-1} consist of $n^2$ 
numbers each. The algorithm runs in time $O(n^3 = (n^2)^{3/2})$. 


What does it mean for an algorithm to be efficient? For us, this will
mean that it runs in \emph{polynomial time}. As a first definition of
polynomial time algorithm, we say that an algorithm runs in
\emph{polynomial time}, if there exists a constant $k$ such that the
algorithm runs in time $O(n^k)$, where $n$ is the length of the input
of the algorithm.

However, we recall the \emph{binary} representation of natural numbers $n ∈ ℕ$. A sequence of \emph{bits} $a_0,\dots,a_{k-1}$ with $a_j ∈ \{0,1\}$ for $0 ≤j≤k-1$ represents the number 
\begin{displaymath}
  ∑_{j=0}^{k-1} a_j ⋅ 2^j. 
\end{displaymath}  
Conversely, each positive natural number $n≥1$ has the binary representation that is found recursively by the following process. 
If $n=1$, then its representation is $a_0 = 1$. 
If $n>1$ and  is even, then the sequence representing $n$ is 
\begin{displaymath}
  0,b_0,\dots,b_{k-1},
\end{displaymath}
where $b_0,\dots,b_{k-1}$ is the representation of $n/2$. If $n>1$ and n is odd, then the sequence representing $n$ is 
\begin{displaymath}
  1,b_0,\dots,b_{k-1},
\end{displaymath}
where $b_0,\dots,b_{k-1}$ is the representation of $⌊n/2⌋$. This creates a representation with leading bit one, i.e. $a_{k-1}=1$. 
By deleting \emph{leading zeros}, i.e., ensuring $a_{k-1}=1$, the representation of a natural positive number is unique, (see exercise~\ref{cha:runn-time-analys}.\ref{alg:ex2}). 



\begin{example}
  This example will make us revise the our first definition of a polynomial-time algorithm. 
The input of this algorithm is a list of characters. Let us say that the list has $n$ elements. 
 
 \label{ex-a-3}
  \begin{tabbing}
    Input: A list $L$ or characters \\
    $s :=2$ \\
    {\bf for} \= $c ∈ L$:\\
              \> $s:= s ⋅ s$ \\
    {\bf return} $s$         
  \end{tabbing}
  Then clearly, the algorithm carries out a polynomial, even linear,
  number of operations in the input, if we consider an arithmetic
  operation also as a basic operation that can be carried out in
  constant time. However, the algorithm squares $2$
  repeatedly, $n$
  times to be precise. Thus, at the end, the variable $s$
  holds the number $2^{2^n}$.
  The number of \emph{bits} in the binary representation  that the return value
  is $2^n$
  which is \emph{exponential} in the input length. 
\end{example}

\begin{definition}
  \label{def:2}
  The \emph{size} of an integer $x$
  is $\size(x) = \lceil\log( |x| +1)\rceil$
  and for $x \in \setQ$,
  $\size(x) = \size(p)+\size(q)$,
  where $x = p/q$ with $p,q \in \setZ$, $q\geq1$ and $\gcd(p,q)=1$.   
\end{definition}
Here $\gcd(p,q)$ denotes the \emph{greatest common divisor} of $p$ and $q$ which we define precisely below. 
Thus the size of a number is asymptotically equal to the number of
bits that are needed to store the number.  We now provide the
definition of what a polynomial-time algorithm is.
\begin{definition}
  \label{def:1}
  An algorithm is \emph{polynomial time}, if there exists a
  constant $k$
  such that the algorithm performs $O(n^k)$
  operations on rational numbers whose size is bounded by
  $O(n^k)$.
  Here $n$
  is the number of bits that encode the input of the algorithm. We say that the algorithm runs in time $O(n^k)$. 
\end{definition}

We now use this definition to analyze the famous \emph{Euclidean} algorithm that computes the \emph{greatest common divisor} of two integers. 

 For $a,b \in \setZ$, $b \neq0$ we say $b$ \emph{divides} $a$ if there
    exists an $x \in \setZ$ such that $a = b \cdot  x$. We 
    write    $b\mid a$.  
   For $a,b,c \in \setZ$, if $c\mid a$ and $c\mid b$, then $c$ is a
    \emph{common divisor} of $a$ and $b$. 
    If at least one of the two integers $a$ and $b$ is non-zero,
    then there exists a \emph{greatest common divisor} of $a$ and
    $b$. It is denoted by \emph{$\gcd(a,b)$}. 
    

    How do we compute the greatest common divisor efficiently? The
    following is called \emph{division with remainder}.  For
    $a,b \in \setZ$
    with $b >0$ there exist unique integers $q,r \in \setZ$ with
    \begin{displaymath}
      a = q\cdot b + r, \text{ and }  0\leq r <b. 
    \end{displaymath}
    Now clearly, 
    for $a,b \in \setZ$ with $b >0$ and $q,r \in \setZ$ as above one has
    $\gcd(a,b) = \gcd(b,r)$.  
  This gives rise to the famous \emph{Euclidean algorithm}. 
  \begin{algorithm}[Euclidean algorithm]
    \label{alg:1}
    \begin{tabbing}
      Input: Integers $a ≥ b ≥0$ not both equal to zero \\
      Output: The greatest common divisor $\gcd(a,b)$\\
      {\bf if}  $(b=0)$ {\bf return} $a$ \\
      {\bf else} \= \\
      \> Compute $q,r ∈ ℕ$ with \=  $b > r ≥ 0$ and $a = q⋅b +r$ \\
      \> \> (division with remainder)\\
      \> {\bf return} $\gcd(b,r)$ 
    \end{tabbing}
  \end{algorithm}  
%
  \begin{theorem}
    The Euclidean algorithm runs in time $O(n)$. 
  \end{theorem}
  \begin{proof}
    Suppose that $a$ and $b$ have at most $n$ bits each.  Clearly, the
    numbers in the course of the algorithm have at most $n$
    bits. Furthermore, if $a \geq b$, then $r \leq a/2$, where $r$ is
    the remainder of the division of $a$ by $b$. Thus each second
    iteration, the first parameter of the input has one bit less. Thus
    the number of operations is bounded by $O(n)$.
  \end{proof}
  
\begin{example}
\label{exe:det}
The \emph{determinant} of a matrix $A ∈ ℝ^{n ×n}$ can be computed by the recursive formula 
\begin{displaymath}
  \det(A) = ∑_{i=1}^n (-1)^{1+j}a_{1j} \det(A_{1j}),
\end{displaymath}
where $A_{1j}$ is the $(n-1)×(n-1)$ matrix that is obtained from $A$ by deleting its  first row and $j$-th column.  This yields the following algorithm. 

\begin{tabbing}
  Input: $A ∈ ℚ^{n  ×n}$ \\
  Output: $\det(A)$ \\
  
  {\bf if} \= $(n=1)$ \\
           \> {\bf return} $a_{11}$ \\
  {\bf else} \\
           \> $d:=0$  \\
           \> {\bf for } \= $j=1,\dots,n$ \\
           \>            \> $d:= (-1)^{1+j}⋅ \det(A_{1j}) +d$\\
           \> {\bf return} $d$   
\end{tabbing}
This algorithm is \emph{recursive}. This basically means that a tree is constructed, where nodes of the tree correspond to recursive calls to the algorithm $\det(⋅)$ including the \emph{root} call. Any node that corresponds to a matrix with $k>1$ rows and columns is expanded by adding $k$ \emph{child nodes} corresponding to the recursive calls that are executed. Once the lower-level nodes of the tree cannot expanded anymore, one can evaluate the tree, in this case the determinant, in a bottom-up fashion. See 
Let $T(n)$ denote the number of basic operations that the algorithm performs. Then $T(1) ≥1$ and 
\begin{displaymath}
  T(n) ≥ n ⋅ T(n-1), 
\end{displaymath}
which shows that $T(n) >= n! = 2^{\Omega({n \log n})}$ which is \emph{exponential} in $n$. 
\end{example}


\begin{figure}
  \centering
  \includegraphics[height=5cm]{figures/DeterminantAlgoSlow.jpg}
  \caption{An example of the recursion tree of the algorithm from Example~\ref{exe:det}. The tree corresponds to the run of the algorithm on input $\protect\smat{         3 & 7 & 5 \\ 2&9&8 \\ 1 & 3 & 3  }.$
} 
\end{figure}


\section*{Exercises} 


\begin{enumerate}
\item Find the binary representation of $134$. 
\item Show that the binary representation with leading bit one of a positive natural number is
 unique. \label{alg:ex2}
\item Show that there are $n$-bit numbers $a,b ∈ ℕ$ such that the  Euclidean algorithm on input $a$ and $b$ performs  $\Omega(n)$ arithmetic operations. \emph{Hint: Fibonacci numbers} 
\item Show $n! =  2^{\Omega({n \log n})}$. 
\item 
Let $A ∈ ℝ^{n ×n}$  and suppose that the $n^2$ components of $A$ are pairwise different.
Suppose that $B$ is a matrix that can be obtained from $A$ by deleting the first $k$ rows and some of the $k$ columns of $A$. How many (recursive) calls of the form $\det(B)$ does the algorithm of Example~\ref{exe:det} create? 
\item Let $A ∈ ℝ^{n ×n}$  and suppose that the $n^2$ components of $A$ are pairwise different. How many different submatrices can be obtained from $A$ by deleting the first $k$ rows and some set of $k$ columns? Conclude that the algorithm  of Example~\ref{exe:det} remains exponential, even if it does not expand repeated subcalls. 
\item Complete the  algorithm below such  that it adds two natural numbers  in binary representation  $a_0,\dots,a_{l-1}$,  $b_0,\dots,b_{l-1}$. What is the asymptotic running time (number of basic operations) of your algorithm? Can there be an asymptotically  faster algorithm? 

  \begin{tabbing}
    Input:~~ \= Two natural numbers $a$ and $b$  in their binary representation\\ 
    \>   $a_0,\dots,a_{l-1}$, $b_0,\dots,b_{l-1}$. \\
    Output: \> The binary representation $c_0,\dots,c_l$ of $a+b$ \\
\pushtabs 	
\\
   $\mathrm{carry} := 0$ \\
   {\bf for} \= $i=0, \dots, l-1$ \\
             \> $c_i = \mathrm{carry} + a_i + b_i \pmod{2} $\\
             \> $\mbox{carry} := $ \\
   $c_l := $  \\
   {\bf return} $c_0,\dots ,c_l$ \\   
\poptabs
  \end{tabbing} \label{item:ex:7}

\end{enumerate}

\section{Analysis of Gaussian elimination} 
\label{sec:analys-gauss-elim}

We recall Gaussian elimination. 

\begin{algorithm}[Gaussian elimination]
  \label{alg:3}
  \begin{tabbing}
    Input: $A ∈ ℚ^{m ×n}$  \\
    Output: \= $A'$ in row echelon form such that there exists an invertible \\ 
            \> $Q ∈ ℚ^{m × m}$ such that $Q⋅A = A'$ . \\
            \pushtabs 
\\
$A' := A$ \\
$i := 1$\\
{\bf while} \=  ($i≤m$)  \\
\> find \emph{minimal} $1 ≤ j ≤n$ such that there exists $k≥i$ such that $a'_{kj} ≠ 0$ \\
\> If no such element exists, then {\bf stop} \\
            \> swap rows $i$ and $k$ in $A'$ \\
            \>  {\bf for} \= $k = i+1,\dots, m$ \\
            \>            \> subtract $(a'_{kj}/a'_{ij})$ times row $i$ from row $k$ in $A'$  \\
\> $i:=i+1$ 
\poptabs        
  \end{tabbing}
\end{algorithm}

\noindent 
We can easily prove correctness of the algorithm. First of all, the algorithm does only perform elementary row-operations of the form 
\begin{enumerate}[i)]
\item swap two rows 
\item subtract a multiple of one row from \emph{another} row. 
\end{enumerate}
This means that the resulting matrix $A'$ can be obtained via 
\begin{displaymath}
  A' = Q ⋅ A
\end{displaymath}
with a non-singular $Q ∈ ℚ^{m × m}$. On the other hand we have the following invariant. 
\begin{quote}
  After each iteration of the while-loop, the matrix $H$ obtained from the first  first $j$ columns of $A'$ is  in row-echelon form and rows $i,i+1,\dots,m$ of $H$  are entirely zero, see Figure\ref{fig:9}.  
\end{quote}

\begin{figure}  
   \centering
  \includegraphics[height=5cm]{figures/Echelon.jpg}
  \caption{Gaussian elimination: The matrix $A'$ before  the $i$-th iteration of the while loop. \label{fig:9}}
\end{figure}


How many arithmetic operations does Gaussian elimination perform? Subtracting a multiple of one row from another row in $A'$ can be done in time $O(n)$ Thus the number of operations that are performed within the {\bf for}-loop are $O(m ⋅n)$ in total. There are $O(m)$ iterations through the while-loop. All-together this shows that Gaussian elimination performs $O(m^2 ⋅n)$ iterations. But how large can the numbers grow in the course of the Gaussian algorithm? Could it be that numbers have to be manipulated, whose binary encoding length is not polynomial in the total encoding length of the matrix $A$? Luckily the answer to this question is ``No''. To provide this answer, we have to show the following. 


\begin{theorem}[Hadamard bound] 
  Let $A \in \R^{n \times n}$ be non-singular. Then 
  \begin{displaymath}
    |\det(A)| \leq \prod_{i=1}^n \|a_i\|_2 \leq n^{n/2} \cdot B^n, 
  \end{displaymath}
  where $B$ is upper bound on absolute values of entries of $A$.
\end{theorem}

\begin{proof}
  The Gram-Schmidt orthogonalization of $A$ yields a factorization 
  \begin{displaymath}
    A = Q \cdot R,
  \end{displaymath}
where $R$ is an upper triangular matrix with ones on the diagonal. The matrix $Q$ has orthogonal columns, where the length of the $i$-th column $q^{(i)}$ is upper bounded by the length of the $i$-th column of $A$. 
The assertion follows from 
\begin{displaymath}
  \det(A)^2 = \det(Q)^2 = \det(Q^T) \det(Q) = \prod_i \|q^{(i)}\|^2. 
\end{displaymath}
\end{proof}



\begin{corollary}
\label{co:11}
  If $A\in \setZ^{n\times n}$ is integral and each entry in absolute
  value is bounded by $B$, then $\size(\det(A)) = O(n \log n+ n \cdot
  \size(B))$.
\end{corollary}
% 


\begin{corollary}
  \label{co:10}
   If $A\in \setQ^{n\times n}$ is a rational matrix and $φ$ is an upper bound on the size of each component of $A$, then $\size(\det(A)) = O( n^3 \cdot
  φ)$.
\end{corollary}

\begin{proof}
Suppose that $a_{ij} = p_{ij}/q_{ij}$, where $p_{ij}$ and $q_{ij}$ are integers with $gcd(p_{ij},q_{ij})=1$ for each $i,j$. Then $(∏_{ij}q_{ij}) A$ is an integer matrix  and 
\begin{displaymath}
  \det(A) = (∏_{ij}q_{ij})^{-n} \det((∏_{ij}q_{ij}) A). 
\end{displaymath}
The size of $(∏_{ij}q_{ij})^{-n} $ is $O(n^3 φ)$ and the size of $\det((∏_{ij}q_{ij}) A$ is also $O(n^3 φ)$ thanks to Corollary~\ref{co:11} . 
\end{proof}


Now that we have shown that the determinant of a rational matrix is a number of polynomial encoding length, we can prove that Gaussian elimination is indeed a polynomial time algorithm.

\begin{theorem}
  \label{thr:9}
The   Gaussian algorithm runs in polynomial time on input $A ∈ ℤ^{m ×n}$.  More precisely, the rational numbers produced in the algorithm can be maintained to be  ratios of sub-determinants of $A'$ and are thus of polynomial binary encoding length. 
\end{theorem}

\begin{proof}
For the proof of this theorem, we assume that we have performed row and column swaps on $A∈ ℤ$ beforehand such that the pivot element in iteration $i≥1$    is in row $i$ and column $i$. This means that we never have to swap rows in $A'$, after the $i$-th iteration of the wile-loop the matrix $A'$ is of the form 
\begin{displaymath}
  A' =
  \begin{pmatrix}
    U &  E \\
    0 & D 
  \end{pmatrix}
\end{displaymath}
where $U ∈ ℚ^{i ×i}$ is upper triangular and non-singular and $D ∈ ℚ^{(m-i) × (n-i)}$ is the part that still has to be eliminated. 



  Consider an element $d_{k,j}$ of $D$ and  define the index sets
  $K = \{1,\dots,i\} ∪\{i+ k\}$ and $J = \{1,\dots,i\} + ∪ \{i+ j\}$ and the matrix
  $A_{KJ}$ as the one induced from $A$ by the rows and columns indexed
  by $K$ and $J$ respectively. A crucial observation is that
  \begin{displaymath}
   \det(A_{KJ}) = \det(A'_{KJ}) 
 \end{displaymath}
 holds. Likewise one has $\det(A_i) = \det(U)$, where $A_i$ is the
 matrix stemming from the first $i$ rows and columns of $A$.


 From this it follows that
 \begin{displaymath}
  d_{k,j} ⋅ \det(A_i)  =   \det(A_{KJ}) 
 \end{displaymath}
 and thus
 \begin{displaymath}
   d_{k,j}  =   \det(A_{KJ})   / \det(A_i). 
 \end{displaymath}
 We also write
 \begin{displaymath}
   d_{k,j}  =   r_{ij}   / \det(A_i)
 \end{displaymath}
 with the suitable integer $r_{ij} = \det(A_{KJ})$. 
 Now the  pivot element is $d_{11}$ and the new element in position $k,j$ of $D$ becomes with $r_{11} = \det(A_{i+1})$ 
 % \begin{eqnarray*}\displaystyle
 %   d_{kj} - d_{k1}/d_{11} ⋅  d_{1j} & = & \frac{\det(A_{KJ})} {\det(A_{i})} - \frac{ \det(A_{K1})  }{\det(A_{i+1})} ⋅ \frac{\det(A_{1j})}{\det(A_{i})}  \\
 %   & = &  \frac{\frac{\det(A_{KJ}) \det(A_{i+1}) - \det(A_{K1}) \det(A_{1J})}{  {\det(A_i)}}} {\det(A_{i+1})}
 % \end{eqnarray*}

  \begin{eqnarray*}\displaystyle
   d_{kj} - d_{k1}/d_{11} ⋅  d_{1j} & = & \frac{r_{kj}} {\det(A_{i})} - \frac{r_{k1}  }{r_{11}} ⋅ \frac{r_{1j}}{\det(A_{i})}  \\
                                    & = &  \frac{(r_{kj} r_{11} - r_{k1} r_{1j}) / \det(A_{i})}
                                          {\det(A_{i+1})}
 \end{eqnarray*}
 The numerator is an integer that can be computed with basic integer arithmetic operations. All integers are bounded by $n^{n/2} B^n$, where $B$ is the largest binary encoding length of an entry of $A$. This shows the claim. 
 
\end{proof}






\section*{Exercises} 
\label{sec:exercises}

\begin{enumerate}
\item \label{item:3} Show that the matrix $Q ∈ ℚ^{m × m}$ that transforms $A \in ℚ^{m × n}$ into $A'$ in the Gaussian algorithm via $Q ⋅A = A'$ has entries that are od polynomial size in the binary encoding length of $A$. 
\end{enumerate}


\section{Fast matrix multiplication} 
\label{sec:fast-matr-mult}


We conclude this chapter on the analysis of algorithms with a result
of Volker Strassen~\cite{strassen1969gaussian} who showed that two
$n ×n$
matrices can be multiplied in time (number of arithmetic operations)
$O(n^{2.805})$.
This algorithm was published in 1969 and it showed as well that a
matrix can be inverted within the same timebound, see~\cite{AHU74}. 

Our task is to compute the product 

\begin{equation}
  \label{eq:str:1}
C =   A \cdot B  
\end{equation}
for two $n ×n$ matrices. We have seen that straightforward matrix-multiplication (Example~\ref{ex-a-1}) requires $O(n^3)$ arithmetic operations. 
We can assume, by padding $A$ and $B$ with zeroes, that $n = 2^\ell$ for some $\ell ∈ ℕ_0$. 

If we split the matrices $A$ and $B$ into 4 $n/2 × n/2$ matrices 
\begin{equation}
  \label{eq:str:3}
  A =
  \begin{pmatrix}
    A_{11} & A_{12} \\ 
    A_{21} & A_{22} 
  \end{pmatrix} \, \text{ and }  B =
  \begin{pmatrix}
    B_{11} & B_{12} \\ 
    B_{21} & B_{22} 
  \end{pmatrix}
\end{equation}
Then 
\begin{displaymath}
  \begin{pmatrix}
    C_{11} & C_{12} \\
    C_{21} & C_{22}
  \end{pmatrix}
   =
  \begin{pmatrix}
    A_{11}\cdot B_{11} + A_{12}\cdot B_{21} & \, &  A_{11}\cdot B_{12} + A_{12}\cdot B_{22}  \\
        A_{21}\cdot B_{11} + A_{22}\cdot B_{21} &\, &     A_{21}\cdot B_{12} + A_{22}\cdot B_{22} 
  \end{pmatrix}.
\end{displaymath}
Thus wan can reduce \emph{one}   multiplication of two $n ×n$ matrices to $8$ multiplications of two $n/2 × n/2$ matrices. For the running time, one obtains the recurrence 
\begin{equation}
  \label{eq:stra:4}
  T(n ) = 8 \cdot T(n/2) + Θ(n^2) 
\end{equation}
which leads again to $T(n) = Θ(n^3)$, see exercise \ref{item:str:4}). 

Strassen discovered that one can preprocess the input matrices in such a way that the correct result can be retrieved from the result  $7$ multiplications of $n/2 × n/2$ matrices. The preprocessing and retrival time is $O(n)$ which yields the recursion 
\begin{equation}
  \label{eq:str:5}
  T(n) = 7 \cdot T(n/2) + O(n^2). 
\end{equation}

The idea is to compute the $7$ matrices
\begin{eqnarray*}
  M_1 & = & (A_{11} + A_{22}) \cdot (B_{11}+ B_{22}) \\
  M_2 & = & (A_{21} + A_{22}) \cdot B_{11} \\
  M_3 & = & A_{11} \cdot (B_{12} - B_{22}) \\
  M_4 & = & A_{22} \cdot (B_{21} - B_{11}) \\
  M_5 & = & (A_{11} + A_{12})\cdot B_{22} \\
  M_6 & = & (A_{21}-A_{11}) \cdot (B_{11}+B_{12}) \\
  M_7 & = & (A_{12}-A_{22}) \cdot  (B_{21} + B_{22})\\. 
\end{eqnarray*}
This amounts to $O(n^2)$ arithmetic operations and $7$ multiplications of $n/2 × n/2$ matrices. From $M_1,\dots,M_7$ one can retrieve 
\begin{eqnarray*}
C_{11} & = & M_1 + M_4 - M_5 + M_7\\
C_{12} & = & M_3 + M_5 \\
C_{21} & = & M_2 + M_4\\
C_{22} & = & M_1 - M_2 + M_3 + M_6.  
\end{eqnarray*}
again with $O(n^2)$ arithmetic operations. 


\begin{algorithm}[Fast matrix multiplication ($\mathrm{FMM}$)]

  \begin{tabbing}
    Input: Two $n ×n$ matrices $A$ and $B$ \\
    Output: $C = \mathrm{FMM}(A,B)$, the product $ A\cdot B$ \\

    {\bf if}  $n=1$ return $a_{11} ⋅b_{11}$ \\
    {\bf else} \=  \\
               \>  $M_1 = \mathrm{FMM} (A_{11} + A_{22} , B_{11}+ B_{22}) $ \\
               \> $M_2 = \mathrm{FMM}(A_{21} + A_{22},  B_{11})$ \\
               \> $M_3  = \mathrm{FMM} A_{11} , B_{12} - B_{22}) $\\
               \> $M_4  = \mathrm{FMM}( A_{22} , (B_{21} - B_{22})$ \\
               \> $M_5  = \mathrm{FMM}(A_{11} + A_{12}, B_{22})$ \\
               \> $M_6  = \mathrm{FMM} (A_{21}-A_{11}, B_{11}+B_{12}) $\\
               \> $M_7  = \mathrm{FMM} (A_{12}-A_{22}, B_{21} + B_{22})$\\ 
               \> Compute the matrices $C_{11}, C_{12}, C_{21}, C_{22}$ from $M_1,\dots,M_7$ \\
               \> {\bf return} $C$
  \end{tabbing}
\end{algorithm}

  \begin{figure}
    \centering
     \includegraphics[height=8cm]{figures/Strassen2.pdf} 
    \caption{The analysis of the Strassen algorithm. }
    \label{fig:str:1}
  \end{figure}


  \begin{theorem}[Strassen]
    \label{thr:10}
    Two $n × n$ matrices can be multiplied in time (number of arithmetic operations) $O(n^{2+ \log_2(7/4)})$. 
  \end{theorem}



  \begin{proof}
    See Figure~\ref{fig:str:1}. 

  \end{proof}

\section*{Exercises}

\begin{enumerate}
\item Suppose we are given three $n×n$ matrices $A,B,C ∈ ℤ^{n×n}$ and we want to test whether $A ⋅ B = C$ holds. We could multiply $A$ and $B$ and then compare the result with $C$. This would amount to running time (number of arithmetic operations) of $O(n^3)$ with the standard matrix-multiplication algorithm. 

We now show how to perform an efficient \emph{randomized test}. Suppose that you can draw a vector $v ∈ \{0,1\}^n$  i.i.d. at random in time $O(n)$. The idea is then to compute the product $B ⋅v$ and then the product $A ⋅ (B ⋅v)$ and afterwards $C ⋅ v$, all in time $O(n^2)$. Show the following. 

\begin{enumerate}
\item If $A⋅B \neq C$, then $P(A ⋅ (B ⋅v) = C⋅v) ≤1/2$. 
\item Let $v_1,\dots,v_k∈ \{0,1\}^n$ be i.i.d. at random and suppose that  $A⋅B \neq C$. The probability of the event: $A ⋅ (B ⋅v_i) = C⋅v_i$ for each $i=1,\dots,k$ is bounded by $1/2^k$. 
\item Conclude that there is an algorithm that runs in time $O(k⋅n^2)$ which tests whether $A⋅B = C$ holds. The probability that the algorithm gives the wrong result is bounded by $1/2^k$. 
\end{enumerate}
\item Show that the recursion \eqref{eq:stra:4} has the solution $T(n) = Θ (n^3)$.  \label{item:str:4}
\item Describe an algorithm that multiplies two $n$-bit integers in time $O(n^2)$. You may use the algorithm to add two $n$-bit integers from exercise~\ref{cha:runn-time-analys}.\ref{item:ex:7}.  
\item Suppose $n = 2^\ell $  and $a,b ∈ ℕ$ are two $n$-bit integers. Consider the numbers $a_h$ and $a_l$ which are represented by the first $n/2$ bits and the last $n/2$ bits of $a$ respectively. Likewise the numbers $b_h$ and $b_l$ are the numbers represented by the first half and the second half of the bit-representation of $b$. 
  \begin{enumerate}[i)]
   \item Show $a = a_h ⋅ 2^{n/2} + a_l$ and  $b = b_h ⋅ 2^{n/2} + b_l$
  \item Show $a⋅b = a_h ⋅ b_h ⋅ 2^n + (a_h ⋅ b_l + a_l ⋅ b_h) ⋅2^{n/2} + a_l ⋅b_l$ 
  \item Conclude very carefully that two $n$-bit
    numbers can be multiplied by resorting to three multiplications of
    $n/2$-bit numbers and $O(n)$ basic  operations.
  \item Conclude that two $n$-bit
    numbers can be computed in time $O(n^{\log_2(3)})$
    elementary bit operations.
  \end{enumerate}

% \item Gaussian elimination in $Q(x)$. Suppose that $A ∈ ℤ[x]^{m ×n}$ is given, i.e. component $ij$ is a polynomial $f_{ij} (x) ∈ℤ[x]$. We will analyse the complexity of Gaussian elimination on this matrix. For simplicity, we assume that no swapping of rows and columns is necessary and stage $k$ of Gaussian elimination has resulted in a the transformation of $A$ as
%   \begin{displaymath}
%     \begin{pmatrix}
%       R & C \\
%       0 & D
%     \end{pmatrix},
%   \end{displaymath}
%   where $R ∈ℚ(x)^{ k ×k}$ is upper triangular and non-singular and $D ∈ ℚ(x)^{ {(n-k) ×(m-k)}$ is the part that still needs to be further eliminated. 
%   \begin{enumerate}
%   \item Show that the element $d_{ij}$ can be written as
%     \begin{displaymath}
%       d_{ij} = \det(A_{I,J}) / \det{A_k}.
%     \end{displaymath}
%     Here $I = \{1,\dots,k\} +i$ and $J = \{1,\dots,k\} +k$ and $A_{I,J}$ is the submatrix of $A$ induced by the rows/columns  indexed by $I$ and $J$ respectively and $A_k$ is the $k$-th principal minor of $A$.

%   \item Let $D'$ result  from $D$ by pivoting with $d_{11}$, i.e. $d'_{ij} = d_{ij} - (d_{i1}/d_{11} ) d_{1j}$ 
%   \end{enumerate}



\end{enumerate}

%%% Local Variables:
%%% mode: latex
%%% TeX-master: "lecture"
%%% End:


\chapter{Linear programming}
\label{cha:introduction-1}

We start by giving some examples of linear programs and how they are
used in practice. 
 





\section{Softdrink production}
\label{sec:healthy-low-priced}

Imagine that you own a company that produces the two softdrinks,
\emph{Spring} and \emph{Nebsi}. These softdrinks are a mixture of
\emph{water}, an ingredient \emph{A} and ingredient \emph{B}.
The recipes for Spring and Nebsi are different. Also, the profit for the two drinks is not the same. Those are as follows. 

The use  of ingredients A and B and the profit  per $100l$ are as follows. 
\begin{center}
  \begin{tabular}{c|c|c|c}
    &  A &  B & Profit  \\\hline
    Spring      & $3l$          & $8l$ & $100$ CHF\\\hline 
    Nebsi      & $6l$           & $4l$ & $125$ CHF
  \end{tabular}
\end{center}
%
While the supply of water is unlimited, your company has only $30l$ of
ingredient~$A$ and $44l$ of ingredient~$B$.

At the end of the production day, the local wholesaler picks up the
drinks in two barrels. The capacity of the barrel for Spring is $500l$
while the barrel for Nebsi has a capacity of $400l$.

As the manager of your small company, your goal is to come up with a
\emph{production plan} that maximizes your profit. A production plan
is a two-dimensional vector $(x_1,x_2) \in \R^2$ which means that you
will produce $x_1 \cdot 100 l$ of Spring and $x_2 \cdot 100 l$ of
Nebsi. A production plan is \emph{feasible} if the produced drinks fit
into the respective barrels and not more of A and B is used than what
is on stock. Clearly, $(5,4)$ is not a feasible production plan, as
this would require $39l$ of ingredient A which exceeds the capacity.





A feasible production plan that maximizes your profit can be found
with the help of a \emph{linear program}, a central object of study in
this course.

  \begin{equation}\label{eq:1-4}
  \begin{array}[]{l rcl}
    \text{max.} & 100\cdot x_1 + 125 \cdot x_2 \\
    \text{ s.t.:} &  3\cdot x_1 + 6 \cdot x_2 & \leq & 30 \\
    &    8\cdot x_1 + 4 \cdot x_2 & \leq & 44 \\
    & x_1 & \leq & 5\\
    & x_2 & \leq & 4\\
     & x_1 & \geq & 0\\
    & x_2 & \geq & 0\\
  \end{array}
\end{equation}

One has to maximize a linear \emph{objective function}, in this case
$f(x_1,x_2) = 100\cdot x_1 + 125 \cdot x_2$ where $(x_1,x_2)$
satisfies linear inequalities. The linear inequalities $x_1\geq 0$ and
$x_2 \geq 0$ reflect the fact that only positive amounts can be
produced, while $x_1 \leq 5$ reflects the barrel capacity for
Spring. The linear inequality $3\cdot x_1 + 6 \cdot x_2 \leq 30$
reflects the amount of $30l$ of ingredient A that is on stock.

We can now make a drawing of all feasible production plans. 


\begin{figure}
  \centering
  
  \begin{tikzpicture}[scale=.5]       
     
          \filldraw[fill=green!20,draw=green!20!](0,0) -- (0,4) -- (2,4) --
          (4,3) -- (5,1) -- (5,0) -- (0,0); 
     
          
          \draw [-,draw=gray] (11,-.5) -- (-1.5,5.75) ;
          \draw [-,draw=gray] (-1.5,4) -- (11,4) ;
          \draw [-,draw=gray] (5,-.5) -- (5,6) ;
          \draw [-,draw=gray] (2.5,6) -- (5.75,-.5) ;
                              
          
          
          \draw[->] (-1.5,0) -- (11,0) node[below right] {$x_1$}; \draw[->]
          (0,-.5) -- (0,6) node[left] {$x_2$};
          
          
          \foreach \x in {1,...,10}
          \draw (\x cm,1pt) -- (\x cm,-1pt) node[anchor=north] {$\x$};
          \foreach \y in {1,...,5}
          \draw (1pt,\y cm) -- (-1pt,\y cm) node[anchor=east] {$\y$};     
        \end{tikzpicture} 

  \caption{The feasible production plans are the green area. }
  \label{fig:2}
\end{figure}

The set of points $(x_1,x_2)$ that have objective function $\beta$ is the line 
\begin{displaymath}
  \{(x_1,x_2) \in \R^2 \colon 100\cdot x_1 + 125 \cdot x_2 = \beta \} 
\end{displaymath}
and our task is now to find the largest value for $\beta$ such that the corresponding line still intersects the set of feasible production plans. 


\begin{figure}
  \centering
  
  \begin{tikzpicture}[scale=.5]       
     
          \filldraw[fill=green!20,draw=green!20!](0,0) -- (0,4) -- (2,4) --
          (4,3) -- (5,1) -- (5,0) -- (0,0); 
     
          
          \draw [-,draw=gray] (11,-.5) -- (-1.5,5.75) ;
          \draw [-,draw=gray] (-1.5,4) -- (11,4) ;
          \draw [-,draw=gray] (5,-.5) -- (5,6) ;
          \draw [-,draw=gray] (2.5,6) -- (5.75,-.5) ;
                              
          
          
          \draw[->] (-1.5,0) -- (11,0) node[below right] {$x_1$}; \draw[->]
          (0,-.5) -- (0,6) node[left] {$x_2$};
          
          {
          \draw[draw=blue] (-1.50000000000000,
          7.40000000000000)node[left]{\small  
            \color{blue}{$  \beta = 775$}} --
          (8.37500000000000, -0.500000000000000) ; }
 
          {
          \draw[draw = blue] 
          (-1.50000000000000, 5.20000000000000)node[left]{\small 
            \color{blue}{$  \beta = 500$}}
          --
          (5.62500000000000, -0.500000000000000) ; %beta = 500 
}

    
          {      
          \draw[draw = blue] 
          (-1.50000000000000, 3.20000000000000)node[left]{\small 
            \color{blue}{$  \beta = 250$}}
          -- (3.12500000000000, -0.500000000000000); % beta = 250 

        }
         
        {
        \filldraw [red] (4,3) circle (3pt)node[above right] {$(4,3)$};
}
          
          \foreach \x in {1,...,10}
          \draw (\x cm,1pt) -- (\x cm,-1pt) node[anchor=north] {$\x$};
          \foreach \y in {1,...,5}
          \draw (1pt,\y cm) -- (-1pt,\y cm) node[anchor=east] {$\y$};     
        \end{tikzpicture} 

  \caption{The optimal production plan is $(4,3)$.}
  \label{fig:intro3}
\end{figure}
Figure~\ref{fig:intro3} reveals that $(4,3)$ is an optimal production plan and that the maximum profit that the manager can achieve is $775$. 



\section{Proving optimality}
\label{sec:proving-optimality}


How can the manager be convinced that $(4,3)$ is an optimal production
plan? Maybe he has made a mistake in his drawing or with his
calculations and $(4,3)$ is not optimal. What we will see now is a
very important principle of linear programming. There is a simple way
to prove optimality of solutions that we will explore later on. 

Inspecting the drawing, one can see that there are two inequalities that $(4,3)$ satisfies with equality, namely the inequalities 
\begin{eqnarray}
\label{eq:1-2}
    3\cdot x_1 + 6 \cdot x_2 & \leq & 30 \\
     8\cdot x_1 + 4 \cdot x_2 & \leq & 44.  \label{eq:1-3}
\end{eqnarray}
Clearly all feasible production plans satisfy these inequalities and
inspecting Figure~\ref{fig:intro3} it seems clear that $(4,3)$ is an
optimal solution of the optimization problem~\eqref{eq:1-4} where each
linear inequality but the inequalities~\eqref{eq:1-2} and
\eqref{eq:1-3} have been removed. 

What now follows is a very important technique that we will apply later on again in greater generality. Since each feasible production plan satisfies the inequalities ~\eqref{eq:1-2} and
\eqref{eq:1-3} it satisfies also these inequalities, after they have been multiplied by $50/3$ and $25/4$ respectively. In fact the inequalities~\eqref{eq:1-2} and
\eqref{eq:1-3} are equivalent to the following two inequalities 
\begin{eqnarray}
\label{eq:1-5}
    50\cdot x_1 + 100 \cdot x_2 & \leq & 500 \\
     50\cdot x_1 + 25 \cdot x_2 & \leq & 275.  \label{eq:1-6}
\end{eqnarray}
By adding up the inequalities~\eqref{eq:1-5} and \eqref{eq:1-6} we obtain the inequality 
\begin{eqnarray}
  \label{eq:1-7}
  100\cdot x_1 + 125 \cdot x_2 & \leq & 775
\end{eqnarray}
which in turn is also satisfied by each feasible production plan. The
left-hand-side of inequality~\eqref{eq:1-7} is the objective function
and $775$ is the value of the objective function evaluated at
$(4,3)$. Thus each feasible production plan yields a profit of at most
$775$ which is the profit yielded by $(4,3)$. This shows that $(4,3)$
is optimal. 





\section{Linear Programs}
\label{sec:linear-programming}
We use the following notation. For a matrix $A \in \setR^{m\times n}$, $i\in
\{1,\ldots,m\}$ and $ j \in \{ 1,\ldots,n\}$ we denote 
the  $i$-th row of  $A$ by  $a_i$ and the $j$-th column of $A$ by
$a^j$. With  $A_{ij}$ we denote the element of $A$ which is in the $i$-th
row and $j$-th column of $A$. For a vector $v \in \setR^m$ and $i\in
\{1,\ldots,m\}$ we denote the $i$-th element of $v$ by $v_i$. 



\begin{definition}
\label{def:9}
Let  $A \in \setR^{m\times n}$ be a matrix,  $b \in \setR^{m}$ and $c \in \setR^n$ be
vectors and  $I_\geq,I_\leq,I_= \subseteq \{1,\ldots,m\}$
and $J_\geq,J_\leq \subseteq\{1,\ldots,n\}$ be index sets. A \emph{linear program (LP)}
consists of 
\begin{enumerate}[i)]
\item a linear \emph{objective function}
  \begin{displaymath}
    \begin{array}{c}
      \max c^T x \\
      \text{or } \min c^Tx
    \end{array}
  \end{displaymath}
\item linear \emph{constraints} 
  \begin{displaymath}
    \begin{array}{c}
      a_i^T x \geq b_i, i \in {I_\geq} \\ 
      a_j^T x \leq b_j, j \in {I_\leq} \\ 
      a_k^T x = b_k, k \in {I_=} 
    \end{array}
  \end{displaymath}
\item and  \emph{bounds on the variables} 
  \begin{displaymath}
    \begin{array}{c}
      x_j \geq0, \, j \in J_\geq \\
      x_j \leq0, \, j \in J_\leq. 
    \end{array}
  \end{displaymath}
\end{enumerate}

\end{definition}




  Notice that we can re-write the objective function $\min c^Tx$ as
  $\max -c^Tx $. Similarly, the  constraints
  $a_i^T x \geq b_i, i \in {I_\geq} $ are equivalent to the constrains 
  $-a_i^T x \leq -b_i, i \in {I_\geq}$. Also the constraints  $ a_k^T x =
  b_k, k \in {I_=}$ can be replaced by the constraints  $a_k^T x \leq
  b_k,\, -a_k^T x \leq  -b_k,\,   k \in {I_=}$. 
  A  lower bound $x_j \geq0$ can be written as $-e_j^Tx\leq0$, where $e_j$ is
  the $j$-th unit vector which has zeroes in every component, except
  for the $j$-th component, which is $1$. Similarly an upper bound
  $x_j\leq0$ can be written as $e_j^Tx\leq0$. 

  All-together, a linear program as in Definition~\ref{def:9}  can
  always be written as 
  $$\max\{c^Tx\colon \wt{A}x\leq \wt{b}, x\in \setR^n\}$$ with a suitable matrix
  $\wt{A} \in \setR^{m\times n}$ and a   suitable vector $\wt{b}\in \setR^m$. This
  representation has a name. 



\begin{definition}
  \label{def:11}
  A linear program is in \emph{inequality standard form}, if it is of
  the form
  \begin{displaymath}
    \max \{c^Tx \colon Ax\leq b, \, x \in \setR^n\}
  \end{displaymath}
  for some matrix $A\in \setR^{m\times n}$ and some vector $b \in \setR^m$. 
\end{definition}

\begin{definition}
\label{def:10}
A point $x^* \in \setR^n$ is called \emph{feasible}, if $x^*$
satisfies all constraints and bounds on the variables. If there are
feasible solutions of a linear program, then the linear program is
called \emph{feasible} itself. A linear program is \emph{bounded} if
there exists a constant $M \in \setR$ such for all feasible $x^* \in
\setR^n$ $c^Tx^* \leq M$, if the linear program is a maximization
problem and $c^Tx^*\geq M$, if the linear program is a minimization
problem.  A feasible solution $x^*$ is an optimal solution if
$c^Tx^*\geq c^Ty^*$ for all feasible $y^*$ if the linear program is a
maximization problem and $c^Tx^* \leq c^Ty^*$ if the linear program is
a minimization problem.
\end{definition}


\begin{figure}[htbp]
  \begin{center}
    \definecolor{ffqqqq}{rgb}{1.,0.,0.}
    \begin{tikzpicture}[scale=0.4,line cap=round,line join=round,>=triangle 45,x=1.0cm,y=1.0cm]
      \clip(-5.543788,-3) rectangle (38.525622,7.562352);
      \fill[fill=black,fill opacity=0.1] (9.48,1.62) -- (12.7,3.44) --
      (15.64,0.26) -- (13.9,-2.58) -- (10.,-2.) -- cycle;
      \fill[fill=black,fill opacity=0.1] (0.66,3.22) -- (3.64,-2.64) --
      (6.76,3.18) -- cycle; \fill[fill=black,fill opacity=0.1]
      (-4.284662,3.17592620061) -- (-3.581894,3.170052) --
      (-3.581894,-3.067014) -- (-4.284662,-3.067014) -- cycle;
      \fill[fill=black,fill opacity=0.1] (-2.,3.1898147234) --
      (-1.239334,3.19443883283) -- (-1.239334,-3.067014) --
      (-2.,-3.067014) -- cycle; \draw (9.48,1.62)-- (12.7,3.44); \draw
      (12.7,3.44)-- (15.64,0.26); \draw (15.64,0.26)-- (13.9,-2.58); \draw
      (13.9,-2.58)-- (10.,-2.); \draw (10.,-2.)-- (9.48,1.62); \draw
      (0.66,3.22)-- (3.64,-2.64); \draw (3.64,-2.64)-- (6.76,3.18); \draw
      (-3.581894,3.170052)-- (-3.581894,-3.067014); \draw
      (-2.,-3.067014)-- (-2.,3.1898147234); \draw [->,color=ffqqqq]
      (17.99894,1.208158) -- (18.,6.);      
\end{tikzpicture}

\end{center}
\caption{With the objective function being to find the highest point,
  we have from left-to-right an infeasible linear program, an
  unbounded linear program and a bounded linear program.}
  \label{fig:123}
\end{figure}




We will see later that a feasible and bounded linear program has an
optimal solution.



\section{Fitting a line} 
\label{sec:fitting-line}

The following is an example which is well known in statistics. Suppose
that you  measure points $(x_i,y_i)\in \setR^2 \; i=1,\ldots,n$ and you are interested in a
linear function $y = a\cdot x +b$ that reflects the sample. One way to do
that is by minimizing the expression
\begin{equation}
  \label{eq:56}
  \sum_{i=1}^n (ax_i + b-y_i)^2 , 
\end{equation}
where $a,b\in \setR$ are the parameters of the line that we are looking
for. The number $(ax_i + b-y_i)^2$ is the square of the vertical
distance of the point $(x_i,y_i)$ from the line $ y = a\,x +b$. 

Instead of using the method of least-squares, we could also minimize
the following function, see also~\cite[Chapter 2.4]{1214763},
\begin{equation}
  \label{eq:57}
  \sum_{i=1}^n |ax_i + b-y_i|.
\end{equation}
This objective has the advantage to be slightly more robust towards
outliers. How can we model this as a linear program. The
trick is to use an extra variable $h_i$ which models the absolute
value of $ax_i+b - y_i$. 
\begin{equation}
  \label{eq:58}
  \begin{array}{lcr}
    \min & \sum_{i=1}^n h_i \\
    h_i & \geq & a x_i + b-y_i,  \, i=1,\ldots,n \\
    h_i & \geq & -(a x_i + b-y_i ),  \, i=1,\ldots,n \\    
  \end{array}
\end{equation}
The variables of this linear program are  $h_i$, $i=1,\ldots,n$, $a$ and
$b$. For a fixed $a\in \setR$ and $b \in \setR$ the optimal $h_i$'s will be
$h_i  = |ax_i + b-y_i|$ since the objective minimizes the sum of the
$h_i$'s. If one of the $h_i$'s was strictly larger than  $|ax_i + b-y_i|$,
then the objective could be improved by making it smaller. 

\begin{figure}
  \centering
  \begin{tikzpicture}[inner sep=0pt,thick,
     dot/.style={fill=black,circle,minimum size=3pt}]
     %\draw[help lines] (0,0) grid (7,4);
     % \draw [<->,thick] (0,4) node (yaxis) [above] {$y$}
     % |- (9,0) node (xaxis) [right] {$x$};
     \node[dot] (a) at (1,3) (1,1) {};
     \node[dot] (b) at (2.8,3.3) (2,2) {};
     \node[dot] (c) at (4,2) (1,2) {};
     \node[dot] (d) at (5.5,2.1) (1.25,0.25) {};
     \node[dot] (e) at (6,.2) (1.75,1.5) {};
     \node[dot] (f) at (7,1.3) {};
     \node[dot] (g) at (7.5,1)  {};
     \node[dot] (h) at (8.5,0.8)  {};
      (8.5,0.8)
      \draw[thick,blue] (1,3.5) -- (9,0.5) ;
 %     \node[draw=red, fit=(a) (b) (c) (d) (e)] {box};
  %    \node[draw,circle,fit=(a) (b) (c) (d) (e)] {};
   \end{tikzpicture}


  \caption{A line that minimizes the sum of the vertical distances.}
  \label{fig:4}
\end{figure}


\section{Linear Programming solvers and modeling languages}
\label{sec:line-progr-solv}

We will demonstrate now how to use a modeling language for linear
programming and a linear programming solver to find a fitting line, as
described in Section~\ref{sec:fitting-line} for the points 
\begin{displaymath}
   (1,3), (2.8,3.3),(4,2),(5.5,2.1),(6,0.2), (7,1.3), (7.5,1), (8.5,0.8) 
\end{displaymath}

There are two popular formats for linear programming problems which
are widely used by linear programming solvers, the \emph{lp-format}
and the \emph{mps-format}. Both are not easy to read. To facilitate
the modeling of a linear program, so-called modeling languages are
used. We demonstrate the use of the popular open source modeling
software called
\href{http://www.zib.de/koch/zimpl}{zimpl}~\cite{Koch2004}. Below you
see a way to model our fitting line linear program with zimpl:

{\small 
\begin{verbatim}
set I := {1 to 8};
param X[I] :=  <1> 1, <2> 2.8, <3> 4  , <4> 5.5, 
               <5> 6, <6> 7  , <7> 7.5, <8> 8.5 ;
param Y[I] :=  <1> 3  , <2> 3.3, <3> 2, <4> 2.1, 
               <5> 0.2, <6> 1.3, <7> 1, <8> 0.8 ;
var h[I] >= -infinity <= infinity;
var a    >= -infinity <= infinity ;
var b    >= -infinity <= infinity ;

minimize cost: sum <i> in I: h[i];

subto c1: forall <i> in I:      h[i] >=   ( a * X[i] + b -Y[i]);
subto c2: forall <i> in I:      h[i] >= - ( a * X[i] + b -Y[i]);
\end{verbatim}
}
%
\noindent
Zimpl  creates a linear program which is readable by linear
programming solvers like
\href{http://www2.isye.gatech.edu/~wcook/qsopt}{QSopt} or
\href{http://soplex.zib.de/}{SoPlex}. %  An optimal fitting-line
% w.r.t. the distance measure~\eqref{eq:57} is the line $y = -0.293333
% \cdot x + 3.293333$. It is depicted in figure~\ref{fig:1}. 

% \begin{figure}[htbp]
%   \begin{center}
    
%     \begin{pspicture}(0,0)(11,4)\showgrid
%       \psdot(1,3)
%       \psdot(2.8,3.3)
%       \psdot(4,2)
%       \psdot(5.5,2.1)
%       \psdot(6,.2)
%       \psdot(7,1.3)
%       \psdot(7.5,1)
%       \psdot(8.5,0.8)
      
%       \psline(0,3.293)(11,0)
%     \end{pspicture}
%   \end{center}
%   \caption{A set of points $\{ (1,3), (2.8,3.3),
%     (4,2),(5.5,2.1),(6,.2), (7,1.3), (7.5,1), (8.5,0.8)\}$ and the
%     line determined by linear program~\eqref{eq:58}.} 
% \end{figure}



\section{Linear programming for longer OLED-lifetime}



\emph{Organic Light Emitting Diodes} (OLEDs) are considered as the
display technology of the future and  more and more commercial
products are equipped 
with such displays as shown in Fig.~\ref{fig:displayandchip}. However,
the cheapest OLED technology suffers from short lifetimes.
 We will
show in this section how linear programming can be used to increase
the lifetime of such displays. 

\begin{figure}[ht]
  \centering \includegraphics[height=3cm]{figures/displayandchip.eps}
 \caption{Sample of a commercial OLED device with integrated driver chip}
 \label{fig:displayandchip}
\end{figure}



A (passive matrix) OLED display has a matrix structure consisting of $n$ rows
and $m$ columns.  At any crossover between a row and a column there is
a vertical diode which works as a pixel. 
The image itself is given as an integral non-negative $n \times m$ matrix
$(r_{ij}) \in [0,\ldots,\varrho]^{n \times m}$ representing its RGB values.
Consider the contacts for the rows and columns as switches.  For the
time the switch of row $i$ and column $j$ is closed, an electrical
current flows through the diode of pixel $(i,j)$ and it shines. Hence,
we can control the intensity of a pixel by the two quantities
\emph{electrical current} and \emph{time}. 
The value $r_{ij}$ determines the amount of time within the time frame
in which the switches $i$~and~$j$ have to be simultaneously closed.
At a sufficient high frame rate e.g. 50 Hz, the perception by the eye
is the average value of the light emitted by the pixel and one sees the image.

The traditional addressing scheme is row-by-row. This means that the
switch for the first row is closed for a certain time while the
switches for the columns are closed for the necessary amount of time
dictated by the entries $r_{1j},\, j=1,\ldots,m$. Consequently the first
row can be displayed in time $\max \{ r_{1j}: j = 1, \ldots, m \}$.
Then the second row is displayed and so on. With this addressing
scheme, the pixels are idle most of the time and then have to shine
with very high intensity. This puts the diodes under stress and is a
major cause of the short lifetime of the displays. 


How can this lifetime problem be dealt with? The main idea is to save
time, or equivalently to lower the maximum intensity, by displaying
several rows at once. 

Consider the schematic image on the left of Fig.~\ref{fig:example}.
Let us compute the amount of time which is
necessary to display the image with this addressing scheme.  The
maximum value of the entries in the first row is $238$. This is the
amount of time which is necessary to display the first row. After that
the second row is displayed in time $237$. In total the time which is
required to display the image is $238+237+234+232+229=1170$ time
units.

\begin{figure}[ht]
  \centering
  $$
\columnseprule5mm
\xdefinecolor{r1}{RGB}{109,238,28}
\xdefinecolor{r2}{RGB}{112,237,28}
\xdefinecolor{r3}{RGB}{150,234,25}
\xdefinecolor{r4}{RGB}{189,232,22}
\xdefinecolor{r5}{RGB}{227,229,19}
\begin{matrix}
\rowcolor{r1} 109 & 238 & 28 \\
\rowcolor{r2} 112 & 237 & 28 \\
\rowcolor{r3} 150 & 234 & 25 \\
\rowcolor{r4} 189 & 232 & 22 \\
\rowcolor{r5} 227 & 229 & 19
\end{matrix}
\>\> = \>\>
\xdefinecolor{m1}{RGB}{0,82,25}
\xdefinecolor{m2}{RGB}{112,155,3}
\xdefinecolor{m3}{RGB}{0,41,22}
\xdefinecolor{m4}{RGB}{189,191,0}
\begin{matrix}
\rowcolor{m1}    \color{white}0 & \color{white} 82 & \color{white}25 \\
\rowcolor{m1}    \color{white}0 & \color{white} 82 & \color{white}25 \\
\rowcolor{m3}    \color{white}0 & \color{white} 41 & \color{white}22 \\
\rowcolor{m3}    \color{white}0 & \color{white} 41 & \color{white}22 \\
\rowcolor{black} \color{white}0 & \color{white}  0 & \color{white} 0 \\
\end{matrix}
\>\> + \>\> 
\begin{matrix}
\rowcolor{black} \color{white} 0 & \color{white}  0 & \color{white} 0 \\
\rowcolor{m2} \color{white} 112 & \color{white} 155 & \color{white} 3 \\
\rowcolor{m2} \color{white} 112 & \color{white} 155 & \color{white} 3 \\
\rowcolor{m4} \color{white} 189 & \color{white} 191 & \color{white} 0 \\
\rowcolor{m4} \color{white} 189 & \color{white} 191 & \color{white} 0
\end{matrix}
\>\> + \>\> 
\xdefinecolor{s1}{RGB}{109,156,3}
\xdefinecolor{s2}{RGB}{0,0,0}
\xdefinecolor{s3}{RGB}{38,38,0}
\xdefinecolor{s4}{RGB}{0,0,0}
\xdefinecolor{s5}{RGB}{38,38,19}
\begin{matrix}
\rowcolor{s1} \color{white} 109 & \color{white} 156 & \color{white}   3 \\
\rowcolor{s2} \color{white}   0 & \color{white}   0 &  \color{white}  0 \\
\rowcolor{s3} \color{white}  38 & \color{white}  38 &  \color{white}  0 \\
\rowcolor{s4} \color{white}   0 &  \color{white}  0 &  \color{white}  0 \\
\rowcolor{s5} \color{white}  38 & \color{white}  38 & \color{white}  19
\end{matrix}
$$
  \caption{An example decomposition}
  \label{fig:example}
\end{figure}

\noindent Now consider the decomposition of the image as the sum of the three
images on the right of Fig.~\ref{fig:example}. In the first image, each
odd row is equal to its even successor. This means that we can close
the switches for rows $1$ and $2$ simultaneously, and these two equal
rows are displayed in $82$ time units.  Rows $3$ and $4$ can also be
displayed simultaneously which shows that the first image on the right
can be displayed in $82+41$ time units. The second image on the right
can be displayed in $155+191$ time units while the third image has to
be displayed traditionally. In total all three images, and thus the
original image on the left via this decomposition, can be displayed in
$82+41+155+191+156+38+38=701$ time units. This means that we could
reduce the necessary time via this decomposition by roughly
$40\%$. We could equally display the image in the original $1170$ time
units but reduce the peak intensity, or equally the maximum electrical current
through a diode by roughly $40\%$.


We now show how to model the time-optimal decomposition of an image as
a linear program. 
To decompose $R$ we need to find matrices
$F^{(1)}=(f_{ij}^{(1)})$ and $F^{(2)}=(f_{ij}^{(2)})$ where $F^{(1)}$
represents the singleline part and $F^{(2)}$ the two doubleline parts.
More precisely, the $i$-th row of matrix $F^{(2)}$ represents the
doubleline covering rows $i$ and $i+1$.  Since the overlay (addition)
of the subframes must be equal to the original image to get a valid
decomposition of $R$, the matrices $F^{(1)}$ and $F^{(2)}$ must
fulfill the constraint $f^{(1)}_{ij} + f^{(2)}_{i-1,j} +
f^{(2)}_{ij}=r_{ij}$ for $i=1,\dots,n$ and $j=1,\dots,m$, where we now
and in the following use the convention to simply omit terms with
indices running out of bounds.  Since we cannot produce ``negative''
light we require also non-negativity of the variables
$f^{(\alpha)}_{ij}\geq0$.  The goal is to find an integral decomposition
that minimizes
\begin{equation*}
\sum_{i=1}^n \max \{ f^{(1)}_{ij}: 1 \leq j \leq m \} + \sum_{i=1}^{n-1} \max \{
f^{(2)}_{ij}: 1 \leq j \leq m \}\enspace.
\end{equation*}
This problem can be formulated as a  linear program by
replacing the objective by $\sum_{i=1}^n u^{(1)}_i +
\sum_{i=1}^{n-1}u^{(2)}_i$ and by adding the constraints $ f^{(\alpha)}_{ij} \leq
u^{(\alpha)}_i$.  This yields
\begin{align}
  \min~ & \sum_{i=1}^n u^{(1)}_i+\sum_{i=1}^{n-1}u^{(2)}_i\nonumber\\
  \text{s.t.}\quad & f^{(1)}_{ij} + f^{(2)}_{i-1,j} + f^{(2)}_{ij} = r_{ij}
  &&\text{for all $i,j$}\label{eqn:original}\\
  & f^{(\alpha)}_{ij} \leq u^{(\alpha)}_i
  &&\text{for all $i,j,\alpha$}\\
   & f^{(\alpha)}_{ij} \in\setR_{\geq 0}
  &&\text{for all $i,j,\alpha$} \nonumber\nonumber
\end{align}
%
Note that the objective does not contain the $f$-variables. By
decomposing images like this, the average lifetime of an OLED display
can be increased by roughly 100\%, see~\cite{EKX07}. 


\section*{Exercises}

\begin{enumerate}[1)]
\item A company produces and sells two different products.
Our goal is to determine the number of units of each product they should produce during one month,
assuming that there is an unlimited demand for the products,
but there are some constraints on production capacity and budget.

There are 20000 hours of machine time in the month.
Producing one unit takes 3 hours of machine time for the first product and 4 hours for the second product.
Material and other costs for producing one unit of the first product amount to 3CHF,
while producing one unit of the second product costs 2CHF.
The products are sold for 6CHF and 5CHF per unit, respectively.
The available budget for production is 4000CHF initially.
25\% of the income from selling the first product can be used immediately as additional budget for production,
and so can 28\% of the income from selling the second product.
\begin{enumerate}
\item Formulate a linear program to maximize the profit subject to the described constraints.
\item Solve the linear program graphically by drawing its set of feasible solutions and determining an optimal solution from the drawing.
\item Suppose the company could modernize their production line to get an additional 2000 machine hours for the cost of 400CHF.
  Would this investment pay off?
\end{enumerate}

\item A factory produces two different products.
To create one unit of product 1, it needs one unit of raw material $A$
and one unit of raw material $B$.
To create one unit of product 2, it needs one units of raw material $B$ and
two units of raw material $C$.
Raw material $B$ needs preprocessing before it can be used, which takes
one minute per unit. 
At most 20 hours of time is available per day for the preprocessing.
Raw materials of capacity at most 1200 can be delivered to the factory per day. 
One unit of raw material A, B and C has size 4, 3 and 2 respectively. 

At most 130 units of the first and 100 units of the second product can be sold
per day. The first product sells for 6 CHF per unit and the second one
for 9 CHF per unit.

Formulate the problem of maximizing turnover as a linear program in two variables and solve it.

\item Prove the following statement or give a counterexample:
The set of optimal solutions of a linear program is always finite.
\item Let (\ref{eq:lp-inequality}) be a linear program in inequality standard form, i.e.
\begin{equation} \label{eq:lp-inequality}
\max\{ c^T x \mid Ax \leq b, x \in \setR^n \}
\end{equation}
where $A \in \setR^{m \times n}$, $b\in \setR^m$, and $c\in\setR^n$.

Prove that there is an equivalent linear program (\ref{eq:lp-equality}) of the form
\begin{equation} \label{eq:lp-equality}
\max\{ \tilde c^T x \mid \tilde Ax = \tilde b, x \geq 0, x \in \setR^{\tilde n} \}
\end{equation}
where $\tilde A \in \setR^{\tilde m \times \tilde n}$, $\tilde b \in \setR^{\tilde m}$, and $\tilde c \in \setR^{\tilde n}$
are such that every feasible point of (\ref{eq:lp-inequality}) corresponds to a feasible point of (\ref{eq:lp-equality})
with the same objective function value and vice versa.

Linear programs of the form in (\ref{eq:lp-equality}) are said to be in \emph{equality standard form}.

\item  Model the linear program~\eqref{eqn:original} to decompose the EPFL logo with Zimpl. 
An incomplete model containing the encoding of the grayscale values of the logo can be found \href{http://disopt.epfl.ch/webdav/site/disopt/users/190205/public/logo\_dec.zmpl}{here}\footnote{http://disopt.epfl.ch/webdav/site/disopt/users/190205/public/logo\_dec.zmpl}.

Use an LP solver library of your choice to compute an optmal solution.

\end{enumerate}


% \bibliographystyle{abbrv}

% \bibliography{mybib,papers,books,my_publications}


%%% Local Variables: 
%%% mode: latex
%%% TeX-master: "lecture"
%%% End: 

\chapter{Polyhedra and convex sets}
\label{conv:cha:convex-sets}



\begin{definition}
  A polyhedron $P\subseteq\setR^n$ is a set of the form $P = \{ x \in
   \setR^n \colon Ax\leq b\}$ for some $A\in \setR^{m\times n}$ and
  some $b \in \setR^m$.
\end{definition}

We are interested in polyhedra, since the
 set of feasible  solutions  of a linear program $\max\{c^Tx \colon
Ax\leq b\}$ is a polyhedron. 


 
\begin{example}
  \label{cha:polyh-conv-sets}
Consider again the soft-drink production problem from chapter~\ref{sec:healthy-low-priced}. 
The corresponding set of feasible solutions is the polyhedron $P = \{x \in \R^2 \colon Ax \leq b\}$  with 
    \begin{displaymath}
      A =
      \begin{pmatrix}
        3 & 6 \\
        8 & 4 \\
        1 & 0 \\
        0 & 1 \\
        -1 & 0 \\
        0 & -1
      \end{pmatrix},
      \quad \text{ and } \quad b = 
      \begin{pmatrix}
        30 \\ 44 \\ 5 \\ 4 \\ 0 \\ 0
      \end{pmatrix}. 
    \end{displaymath}

    \begin{figure}
      \centering
      \begin{tikzpicture}[scale=.45]       
     
      \filldraw[fill=green!20,draw=green!20!](0,0) -- (0,4) -- (2,4) --
      (4,3) -- (5,1) -- (5,0) -- (0,0); 
      
      
      \draw [-,draw=gray] (7,1.5) -- (-1.5,5.75) ;
      \draw [-,draw=gray] (-1.5,4) -- (7,4) ;
      \draw [-,draw=gray] (5,-.5) -- (5,6) ;
      \draw [-,draw=gray] (2.5,6) -- (5.75,-.5) ;
      
      
          
      \draw[->] (-1.5,0) -- (8,0) node[below right] {$x_1$}; \draw[->]
      (0,-.5) -- (0,6) node[left] {$x_2$};
      
%      \draw[draw=blue] (-1.50000000000000,
%      7.40000000000000)node[left]{\small  
%        \color{blue}{$  \beta = 775$}} --
%      (8.37500000000000, -0.500000000000000) ; 
      
%      \filldraw [red] (4,3) circle (3pt)node[above right] {$(4,3)$};
      
          
%      \foreach \x in {1,...,7}
 %     \draw (\x cm,1pt) -- (\x cm,-1pt) node[anchor=north] {$\x$};
 %     \foreach \y in {1,...,5}
 %     \draw (1pt,\y cm) -- (-1pt,\y cm) node[anchor=east] {$\y$};     
    \end{tikzpicture} 
      \caption{The polyhedron of feasible solutions of the linear program \eqref{eq:1-4}.}
      \label{fig:1}
    \end{figure}

\end{example}








\begin{definition}
  \label{conv:def:2}
  A set $K\subseteq\setR^n$ is \emph{convex} if for each $u,v \in K$
  and $\lambda \in [0,1]$ the point $\lambda u+(1-\lambda)v$ is also
  contained in $K$. \end{definition}
\begin{figure}[htbp]
  \centering
  \includegraphics[height=3cm]{figures/exconv.pdf} 
  \caption{The set on the left is convex, the set on the right is  non-convex.}\label{conv:fig:3}
  
\end{figure}


A \emph{halfspace} is a set of solutions of one inequality $a^Tx \leq
\beta$ where $a \in \R^n$ and $\beta \in \R$, i.e., a set of the form
\begin{displaymath}
  \{ x \in \R^n \colon a^Tx \leq \beta\}. 
\end{displaymath} 
A \emph{hyperplane} is a set of the form 
\begin{displaymath}
  \{ x \in \R^n \colon a^Tx = \beta\}. 
\end{displaymath} 
It is easy to see that a halfspace is convex. Convexity is also maintained if convex sets are intersected. 

\begin{proposition}
  \label{lem:1}
  Let $I$ be an index set and $C_i \subseteq \R^n$ be convex sets for
  each $i \in I$, then $\cap_{i \in I}C_i$ is a convex set.
\end{proposition}

Consequently, the set of feasible solutions of a linear program $\{x
\in \R^n \colon Ax \leq b\}$ is a convex set. This is our motivation
to study properties of convex sets.


\section{Extreme points and vertices}
\label{sec:extr-points-vert}


\begin{definition}
  An inequality $a^Tx \leq \beta$ is \emph{valid} for a  set $K \subseteq \R^n$ if each $x^* \in K$ satisfies $a^Tx^* \leq \beta$.  If in addition $(a^Tx = β) \cap K \neq\emptyset$,
then $a^Tx\leq β$ is a \emph{supporting inequality} and $a^Tx = β$ is a
\emph{supporting hyperplane}. 
\end{definition}



% \begin{figure}
%   \centering  
%  \begin{tikzpicture}[scale=.45]       
     
%       \filldraw[fill=green!20,draw=green!20!](0,0) -- (0,4) -- (2,4) --
%       (4,3) -- (5,1) -- (5,0) -- (0,0); 
      
      
%       \draw [-,draw=gray] (7,1.5) -- (-1.5,5.75) ;
%       \draw [-,draw=gray] (-1.5,4) -- (7,4) ;
%       \draw [-,draw=gray] (5,-.5) -- (5,6) ;
%       \draw [-,draw=gray] (2.5,6) -- (5.75,-.5) ;
      
      
          
%       \draw[->] (-1.5,0) -- (8,0) node[below right] {$x_1$}; \draw[->]
%       (0,-.5) -- (0,6) node[left] {$x_2$};
% %\uncover<2->
% {      
%       \draw[draw=blue] (-1.50000000000000,
%       7.40000000000000)node[above]{\scriptsize  
%         \color{blue}{$  100 \cdot x_1 + 125 \cdot x_2 \leq 775$}} -- 
%       (8.37500000000000, -0.500000000000000) ; 
%       }
% % \uncover<5->
% {
      
% %       \filldraw [red] (4,3) circle (3pt)node[above right] {\scriptsize
% %         $(4,3)$};
%        }
          
% %      \foreach \x in {1,...,7}
%  %     \draw (\x cm,1pt) -- (\x cm,-1pt) node[anchor=north] {$\x$};
%  %     \foreach \y in {1,...,5}
%  %     \draw (1pt,\y cm) -- (-1pt,\y cm) node[anchor=east] {$\y$};     
%     \end{tikzpicture} 
    
%   \caption{The inequality $100\cdot x_1 + 125 \cdot x_2$ is valid for the set of feasible solutions of the linear program~\eqref{eq:1-4}. }
%   \label{fig:6}
% \end{figure}



\begin{definition}
  Let $K \subseteq \R^n$ be a convex set. A point $x^* \in K$ is an 
 \emph{extreme point} or \emph{vertex} of $K$ if there exists a valid inequality $a^Tx \leq \beta$ of $K$ such that 
 \begin{displaymath}
   \{x^*\} = K \cap \{x \in \R^n \colon a^Tx = \beta\}.  
 \end{displaymath}
In other words, if $x^*$ is the only point of $K$ that satisfies the valid inequality with equality. 
\end{definition}

\begin{figure}
  \centering
  \includegraphics[height=3cm]{figures/ExtremePoint.pdf} 
  \caption{An extreme point of a convex set.}
  \label{fig:5}
\end{figure}

We can now characterize the extreme points of polyhedra. In fact, there are only finitely many of them. 

\begin{theorem}
  \label{thr:1}
  Let $P = \{x \in \R^n \colon Ax \leq b\}$ be a polyhedron. A feasible point $x^*$ is an extreme point of $P$ if and only if there is a sub-system $A'x \leq b'$ of  $Ax \leq b$  such that
  \begin{enumerate}[i)]
  \item $x^*$ satisfies all inequalities of $A'x \leq b'$ with
    equality. \label{item:1}
  \item $A'$ has $n$ rows and $A'$ is non-singular. \label{item:2}
  \end{enumerate}  
\end{theorem}


\begin{proof}
  Let $A'x \leq b'$ be such a sub-system and consider the valid
  inequality $\mathbf{1}^T A' x \leq \mathbf{1}^T b'$.  Clearly $x^*$
  satisfies this inequality with equality. Any $y^* \in P$ that
  satisfies this inequality with equality must satisfy $A'x = b'$.
  Since $A'$ is non-singular, $x^*$ is the unique solution of $A'x =
  b'$ which means that $x^*$ is the unique point of $P$ that satisfies
  $\mathbf{1}^T A' x \leq \mathbf{1}^T b'$ with equality. 

  Assume now that there does not exist a sub-system $A'x \leq b'$ of
  $Ax \leq b$ with properties~\ref{item:1}) and \ref{item:2}). Denote
  the sub-system of inequalities that are satisfied by $x^*$ with
  equality by $\wt{A}x \leq \wt{b}$. Then $\rank(\wt{A}) < n$ and
  there exists a $d \neq 0 \in \R^n$ with $\wt{A}d = 0$. Consequently
  there exists an $\varepsilon >0$ such that $x^* \pm \varepsilon
  \cdot d \in P$.

  Clearly, any inequality that is satisfied by $x^*$ with equality and
  that is satisfied by $x^* \pm\varepsilon d$ is satisfied by $x^*
  \pm\varepsilon d$ with equality as well. This implies that $x^*$ is
  not an extreme point.
 
\end{proof}



The relevance of vertices for linear programming is reflected in the following theorem. 

\begin{theorem}
  \label{thr:2}
  If a linear program $\max\{c^Tx \colon x \in \R^n, \, Ax \leq b\}$
  is feasible and bounded and if $\rank(A) = n$, then the linear program has an optimal solution that is  an extreme point. 
\end{theorem}


\begin{proof}
  We use the following notation. If $x^*$ is a feasible solution then
  $A_{x^*} x \leq b_{x^*}$ is the subsystem of $Ax \leq b$ that is
  satisfied by $x^*$ with equality. The rank of $x^*$, $\rank(x^*)$ is
  the rank of $A_{x^*}$.  The following claim implies the assertion. 
  \begin{quote}
    If $x^*$ is feasible and $\rank(x^*) < n$, then there exists a
    $y^*$ with $c^T y^* \geq c^Tx^*$ and $\rank(y^*)>\rank(x^*)$.
  \end{quote}
  To prove this, let $d \neq 0 \in \R^n$ be a vector with $A_{x^*}d =
  0$.  We can assume $c^Td \geq 0$ by switching to $-d$ otherwise.

  If $c^Td >0$, then consider the points $x^* + \lambda d$ with
  $\lambda \geq 0$ and let $\lambda_{max}$ be maximal with the
  corresponding point feasible. Clearly $y^* = x^* + \lambda_{max} d$
  satisfies the condition of the claim. 

  Suppose now that  $c^T d = 0$. Then $A d \neq 0$ since $\rank(A) = n$. Let $\lambda_{max}$ be the maximum of the set $\{\lambda \geq 0 \colon A (x^* \pm \lambda d) \leq b\}$. Then $y^* = x^* + \lambda_{max}d$ or $y^* = x^* - \lambda_{max}d$  satisfies the condition of the claim. 
  
\end{proof}



\begin{corollary}
  \label{co:12}
  A linear program $\max \{ c^Tx : x ∈ ℝ^n, \, Ax ≤ b\}$ which is feasible and bounded has an optimal solution. 
\end{corollary}


\begin{proof}
  The linear program 
  \begin{displaymath}
    \max \{ c^Tx^+ -c^T x^- : x^+,x^- ∈ ℝ^n, x^+≥0, \, x^‐≥0, \, A(x^+-x^-) ≤ b\}
  \end{displaymath}
 is feasible and bounded and the rank of the constraint matrix is $2\,n$.  Thus, by Theorem~\ref{thr:2} possesses an optimal solution. $x^+,x^-$. This corresponds to an optimal solution $x^+-x^-$ of the linear program $\max \{ c^Tx : x ∈ ℝ^n, \, Ax ≤ b\}$. 
\end{proof}

\section{Linear, affine, conic  and convex hulls}
\label{conv:cha:linear-affine-convex}



We now describe convex sets that are generated by a set 
 $X\subseteq\setR^n$  of $n$-dimensional
vectors. The \emph{linear hull}, \emph{affine hull}, \emph{conic hull} and \emph{convex
  hull} of $X$ are defined as follows.
\begin{eqnarray}
  \linhull(X) & = & \{ \lambda_1 x_1+ \cdots + \lambda_t x_t \mid t\geq0,  \label{conv:eq:2}\\
   & & \quad \quad  x_1,\ldots,x_t
  \in  X, \, \lambda_1,\ldots,\lambda_t\in \setR\} \nonumber \\
  \affhull(X) & = & \{ \lambda_1 x_1+ \cdots + \lambda_t x_t \mid t\geq1, \label{conv:eq:4}\\
  & &   \quad \quad  x_1,\ldots,x_t \in  X, \, \sum_{i=1}^t \lambda_i = 1, \,
  \lambda_1,\ldots,\lambda_t\in \setR\} \nonumber   \\ 
  \cone(X)  & = & \{ \lambda_1 x_1+ \cdots + \lambda_t x_t \mid t\geq0, \label{conv:eq:5} \\
  &&     \quad \quad  x_1,\ldots,x_t \in  X,  \, \lambda_1,\ldots,\lambda_t\in
  \setR_{\geq0}\} \nonumber \\
  \conv(X) & = & \{ \lambda_1 x_1+ \cdots + \lambda_t x_t \mid t\geq1, \label{conv:eq:6} \\
  &&     \quad \quad  x_1,\ldots,x_t \in  X, \,  \sum_{i=1}^t \lambda_i = 1, \, \lambda_1,\ldots,\lambda_t\in
  \setR_{\geq0}\} \nonumber 
\end{eqnarray}



\begin{figure}[htbp]
  \begin{center}
    \includegraphics{figures/picture1.pdf}    
  \end{center}
  \caption{The convex hull of $7$ points in $\setR^2$. }\label{conv:fig:2}
\end{figure}


\begin{figure}[htbp]
  \begin{center}{
      \includegraphics{figures/picture2-1.pdf} 
      \hfill
      \includegraphics{figures/picture2-2.pdf} 
    }
    
  \end{center}
  \caption{Two points with their convex hull on the left and  their
    affine hull on the right. }\label{conv:fig:1}
\end{figure}





\begin{figure}[htbp]
  \begin{center}{
   \includegraphics{figures/picture3.pdf}

    }
    
  \end{center}
  \caption{Two points with their conic hull}\label{conv:fig:4}
\end{figure}



\begin{proposition}
  \label{conv:prop:1}
  Let $X\subseteq\setR^n$ and $x_0\in X$. One has 
  \begin{displaymath}
    \affhull(X) = x_0 + \linhull(X - x_0),
  \end{displaymath}
  where for $u \in \setR^n$ and $V\subseteq\setR^n$,   $u +V$ denotes the set
  $u+V = \{ u+v \mid v \in V\}$.
\end{proposition}


\begin{proof}
  We first show that each $x \in \affhull(X)$ is also an element of the
  set $ x_0 + \linhull(X-x_0)$ and then we show that each point $x \in
  x_0 + \linhull(X-x_0)$ is also an element of $\affhull(X)$.


  Let $x \in \affhull(X)$,i.e., there exists a natural number $t\geq1$
  and $\lambda_1,\ldots,\lambda_t \in \setR$, with $x = \lambda_1x_1+\cdots+\lambda_t x_t$ and
  $\sum_{i=1}^t\lambda_i=1$. Now
  \begin{eqnarray*}
    x & = &  x_0  - x_0 + \lambda_1 x_1+   \lambda_2x_2+\cdots+\lambda_t x_t \\
    & = & x_0  - \lambda_1x_0 -  \cdots-\lambda_t x_0 + \lambda_1 x_1+  \lambda_2x_2+\cdots+\lambda_t x_t \\
    & = & x_0 +\lambda_1(x_1-x_0)+\cdots+\lambda_t (x_t-x_0),
  \end{eqnarray*}
  which shows that $x \in x_0 + \linhull(X - x_0)$. 


  Suppose now that
  $x \in x_0 + \linhull(X - x_0)$. Then there exist $\lambda_1,\ldots,\lambda_t\in \setR$
  with $x = x_0 +\lambda_1(x_1-x_0)+\cdots+\lambda_t (x_t-x_0)$. With $\lambda_0 = 1
  -\sum_{i=1}^t \lambda_i$ one has $\sum_{i=0}^t \lambda_i =1$ and
  \begin{eqnarray*}
    x & = &  x_0 +\lambda_1(x_1-x_0)+\cdots+\lambda_t (x_t-x_0)\\
    & = & \lambda_0 x_0 + \cdots + \lambda_t x_t
  \end{eqnarray*}
  and thus that $x \in \affhull(X)$.  \qed 

\end{proof}






\begin{theorem}
  \label{conv:thr:1}
  Let $X\subseteq\setR^n$ be a set of points. The convex hull, $\conv(X)$,  of $X$
  is convex. 
\end{theorem}

\begin{proof}
  Let $u$ and $v$ be points in $\conv(X)$. This means that there exists
  a natural number $t\geq1$,  real numbers  $\alpha_i,\beta_i\geq0$, and points $x_i
  \in X$,  $i=1,\ldots,t$ with
  $\sum_{i=1}^t \alpha_i = \sum_{i=1}^t \beta_i =1$ with  $u =  \sum_{i=1}^t
  \alpha_ix_i$  and $v = \sum_{i=1}^t \beta_ix_i$. For $\lambda \in [0,1]$ one has 
  $\lambda\alpha_i+(1-\lambda)\beta_i\geq0$ for $i=1,\ldots,t$  and $\sum_{i=1}^t
  \left(\lambda\alpha_i+(1-\lambda)\beta_i \right) =1$. This shows that   
  \begin{displaymath}
    \lambda u + (1-\lambda) v  =  \sum \left(\lambda_i \alpha_i + (1- \lambda_i) \beta_i\right)
    x_i  \in \conv(X),
  \end{displaymath}
  and therefore that $\conv(X)$ is convex. 
\qed \end{proof}

\begin{theorem}
  \label{conv:thr:2}
  Let $X\subseteq\setR^n$ be a set of points. Each convex set $K$ containing $X$
  also contains $\conv(X)$. 
\end{theorem}

\begin{proof}
  Let $K$ be a convex set containing $X$, 
  and let $x_1,\ldots,x_t\in X$ and $\lambda_i\in \setR$ with $\lambda_i\geq0$, $i=1,\ldots,t$ and
  $\sum_{i=1}^t  \lambda_i =1$. We need to show that $u = \sum_{i=1}^t \lambda_i
  x_i$ is contained in  $K$.   This is true for $t\leq2$ by the definition
  of convex sets.  

  We argue by induction.  Suppose that $t\geq3$. If one of the $\lambda_i$ is
  equal   to $0$, then one can represent $u$ as a convex combination
  of $t-1$ points in $X$ and, by induction, $u \in K$. If    $t\geq3$,
  each $\lambda_i>0$ and $\sum_{i=1}^t \lambda_i = 1$, then one has $0<\lambda_i<1$ for
  $i=1,\ldots,t$ and thus we can write 
  \begin{displaymath}
    u = \lambda_1 x_1 + (1-\lambda_1) \sum_{i=2}^t \frac{\lambda_i}{1-\lambda_1} x_i. 
  \end{displaymath}
  One has $\lambda_i / (1-\lambda_1)>0$ and 
  \begin{displaymath}
    \sum_{i=2}^t \frac{\lambda_i}{1-\lambda_1} = 1,
  \end{displaymath}
  which means that the point $\sum_{i=2}^t \frac{\lambda_i}{1-\lambda_1} x_i$ is in
  $K$ by induction. Again, by the definition of convex sets, we
  conclude that $u$ lies in $K$. 
\qed \end{proof}

Theorem~\ref{conv:thr:2} implies that $\conv(X)$ is the intersection of all
convex sets containing $X$, i.e., 
\begin{displaymath}
  \conv(X) = \bigcap_{\substack{K \supseteq X\\ K \text{  convex}}} K. 
\end{displaymath}



\begin{definition}
  \label{conv:def:4}
  A set $C\subseteq\setR^n$ is a \emph{cone}, if it is convex and for each $c \in
  C$ and each $\lambda \in \setR_{\geq0}$ one has  $\lambda\cdot c \in C$. 
\end{definition}

Similarly to Theorem~\ref{conv:thr:1} and Theorem~\ref{conv:thr:2} one proves
the following. 
\begin{theorem}
  \label{conv:thr:8}
  For any $X\subseteq\setR^n$, the set $\cone(X)$ is a cone. 
\end{theorem}

\begin{theorem}
  \label{conv:thr:9}
  Let $X\subseteq\setR^n$ be a set of points. Each cone containing $X$ also
  contains $\cone(X)$. 
\end{theorem}

These theorems imply that $\cone(X)$ is the intersection of all
cones  containing $X$, i.e., 
\begin{displaymath}
  \cone(X) = \bigcap_{\substack{C \supseteq X\\ C \text{  is a cone}}} C. 
\end{displaymath}


\section{Radon's lemma and Carath\'eodory's theorem}
\label{conv:sec:radons-lemma-carath}

\begin{theorem}[Radon's lemma]
  \label{conv:thr:6}
  Let $A\subseteq\setR^n$ be  a set of $n+2$ points. There exist disjoint
  subsets $A_1,A_2\subseteq A$ with
  \begin{displaymath}
    \conv(A_1) \cap \conv(A_2) \neq \emptyset.
  \end{displaymath}
\end{theorem}

\begin{proof}
  Let $A = \{a_1,\ldots,a_{n+2}\}$. We embed these points into $\setR^{n+1}$
  by appending a $1$ in the $n+1$-st component, i.e., we construct 
  \begin{displaymath}
    A' = \left \{ \smat{a_1\\1}, \ldots,\smat{a_{n+2}\\1} \right\}\subseteq\setR^{n+1}.
  \end{displaymath}
  The set $A'$ consists of $n+2$ vectors in $\setR^{n+1}$. Those vectors
  are linearly dependent.  Let 
  \begin{equation}
    \label{conv:eq:1}
    0= \sum_{i=1}^{n+2} \lambda_i \smat{a_i\\1}
  \end{equation}
  be  a nontrivial linear representation of  $0$, i.e., not all $\lambda_i$
  are~$0$.  Furthermore, let $P = \{i \colon \lambda_i \geq0, \, i=1,\ldots,n+2\}$
  and $N = \{i   \colon \lambda_i <0, \, i=1,\ldots,n+2\}$. We claim that
  \begin{displaymath}
      \conv(\{a_i \colon i \in P\}) \cap \conv(\{a_i  \colon i \in N\} ) \neq \emptyset.
  \end{displaymath}
  It follows from~\eqref{conv:eq:1} and the fact that the $n+1$-st
  component of the vectors is $1$  that $\sum_{i \in P} \lambda_i = - \sum_{i \in N}
  \lambda_i  = s >0$. It follows also from~\eqref{conv:eq:1} that 
  \begin{displaymath}
     \sum_{i \in P} \lambda_i a_i = \sum_{i \in N} -\lambda_i a_i.
  \end{displaymath}
  The point $ u = \sum_{i \in P} (\lambda_i / s) \cdot  a_i  = \sum_{i \in N} (-\lambda_i
  /s) a_i $ is contained in $ \conv(\{a_i \colon i \in P\}) \cap \conv(\{a_i
  \colon i \in N\} )$, implying the claim.  \qed
  
\end{proof}


\begin{theorem}[Carath\'eodory's theorem]
  \label{conv:thr:7}
  Let $X\subseteq\setR^n$, then for each $x \in \cone(X)$ there exists a set
  $\wt{X}\subseteq X$ of cardinality at most $n$  such that $x \in
  \cone(\wt{X})$. The vectors in $\wt{X}$ are linearly independent. 
\end{theorem}


\begin{proof}
  Let $x \in \cone(X)$, then there exist $t \in \setN_+$, $x_i \in X$ and
  $\lambda_i\geq0$, $i=1,\ldots,t$, with  $ x = \sum_{i=1}^t \lambda_i x_i$.
  Suppose that $t\in \setN_+$ is minimal such that $x$ can be represented
  as   above. We claim that $t \leq n$. If $t\geq n+1$, then the $x_i$ are
  linearly dependent.  This means that there are $\mu_i \in \setR$, not all
  equal to $0$ with   
  \begin{equation}
    \label{conv:eq:13}
    \sum_{i=1}^t \mu_i x_i = 0.
  \end{equation}
 By multiplying each
  $\mu_i$ in \eqref{conv:eq:13} with   $-1$ if necessary, we can  assume
  that at least one of the  $\mu_i$ is strictly larger than $0$.  
  One
  has for each $\varepsilon \in \setR$  
  \begin{equation}
    \label{conv:eq:14}
    x = \sum_{i=1}^t (\lambda_i - \varepsilon \cdot \mu_i) x_i. 
  \end{equation}
  What is the largest $\varepsilon^*>0$ that we can pick for $\varepsilon$ such
  that~\eqref{conv:eq:14} is still a conic combination? We need to have
  \begin{equation}
    \label{conv:eq:15}
      \lambda_i - \varepsilon \cdot \mu_i \geq0, \,\text{  for each }i \in \{1,\ldots,t\}. 
  \end{equation} 
   Let $J$ be the set of indices $J = \{j  \colon j \in \{1,\ldots,t\}, \, \mu_j
   >0\}$. We observed that we can assume $J \neq \emptyset$. 
We have~\eqref{conv:eq:15} as long as 
  \begin{equation}
    \label{conv:eq:16}
    \varepsilon \leq  \lambda_j / \mu_j \text{ for each } j \in J.
  \end{equation}
This means that $\varepsilon^* = \min\{ \lambda_j / \mu_j \colon j \in J\}$.  Let $j^*\in J$
be an index where this minimum is attained.  
Since $\lambda_i - \varepsilon^* \cdot \mu_i \geq0$ for all $i=1,\ldots,t$ and since
$\lambda_{j^*}-\varepsilon^*\cdot\mu_{j^*}=0$, we have $x \in  \cone(\{x_1,\ldots,x_t\}\setminus
\{x_{j^*}\}$, which is a contradiction to the minimality of~$t$.
\qed
\end{proof}


\begin{corollary}[Carath\'eodory's theorem for convex hulls]
  \label{conv:co:1}
  Let $X\subseteq\setR^n$, then for each $x \in \conv(X)$ there exists a set
  $\wt{X}\subseteq X$ of cardinality at most $n+1$ such that $x \in
  \conv(\wt{X})$. 
\end{corollary}




\section{Separation theorem and Farkas' lemma}
\label{conv:sec:separ-theor-fark}


We recall a basic fact from analysis, see,
e.g.~\cite[Theorem~4.4.1]{MarsdenHoffman93}. 

\begin{theorem}
  \label{conv:thr:11}
  Let $X\subseteq\setR^n$ be compact and $f: X \to \setR$ be continuous. Then $f$ is
  bounded and there exist
  points $x_1,x_2 \in X$ with $f(x_1) = \sup\{ f(x) \colon  x \in X\}$ and
  $f(x_2) = \inf \{ f(x) \colon x \in X\}$. 
\end{theorem}


\begin{theorem}
  \label{conv:thr:10}
  Let $K\subseteq\setR^n$ be a closed  convex set and $x^* \in \setR^n \setminus K$, then there
  exists an inequality $a^Tx ≤ \beta$ such that $a^T y < \beta$ holds for all
  $y \in K$ and $a^Tx^*>\beta$. 
\end{theorem}

\begin{proof}
  Since the mapping $f(x) = \|x^*-x\|$ is continuous and since for any
  $k \in K$, $K \cap \{ x \in K \colon \|x^* - x\|\leq \|x^* - k\|\}$ is
  compact, there exists a point $k^* \in K$ with minimal distance to
  $x^*$. Consider the midpoint $m = 1/2(k^*+x^*)$ on the line-segment
  $\overline{k^*x^*}$   and the hyperplane $a^Tx = \beta$ with $\beta = a^Tm$ and $a =
  (x^* - k^*)$.%  Clearly, $a^T x^* = \beta - 1/2 \|k^*-x^*\|^2$ and $a^T k^*
  % = \beta +  1/2 \|k^*-x^*\|^2$.

  Clearly $a^T x^* > β$ and $a^T k^* < β$. Suppose now that there exists a point $k' ∈ K$ with $a^T k' ≥ β$. Since the line segment spanned by $k^*$ and $k'$ is contained in $K$, we can assume that $a^T k' = β$ holds. This means that $k' = k^* + a/2 + x'$, where $x'$ is orthogonal to $a$. Again, by convexity $ k^* + λ(a/2 + x') ∈ K$ for $0≤λ≤1$. But the square of the distance of such a point to $x^*$ is 
  \begin{eqnarray*}
    \|x^* - k^* - λ(a/2 + x')\|^2 & = & \|(1-λ/2) a - λ x'\|^2 \\
    & = & (1-λ/2)^2 \| a\|^2  + λ^2  \|x'\|^2 \\ 
  \end{eqnarray*}
  
  As a function of $\lambda$, this is decreasing 
  at $\lambda=0$. Thus there exists a point on the line segment between $k^*$ and $k'$ which is closer to $x^*$ than $k^*$ and this is a contradiction. 
  \qed 
\end{proof}


\begin{theorem}[Farkas' lemma]
  \label{conv:thr:12}
  Let $A \in \setR^{m\times n}$ be a matrix and $b \in \setR^m$ be a vector. The
  system $Ax = b, \,x\geq0$ has a solution if and only if for all $\lambda \in
  \setR^m$ with $\lambda^TA\geq0$ one has $\lambda^Tb \geq0$.  
\end{theorem}

\begin{proof}
  Suppose that $x^* \in \setR^n_{\geq0}$ satisfies $Ax^* = b$ and let $\lambda \in
  \setR^m$ with $\lambda^T A \geq0$. Then $\lambda^Tb = \lambda^TA x^* \geq0$, since
  $\lambda^TA\geq0$ and $x^*\geq0$. 

  Now suppose that $Ax = b, \,x\geq0$ does not have a solution. Then,
  with   $X\subseteq\setR^m$ being  the set of column vectors of $A$, 
  $b$ is not in $\cone(X)$. The set $\cone(X)$ is convex and
  closed, see exercise~\ref{conv:item:2}. Theorem~\ref{conv:thr:10} implies 
  that there is an inequality $\lambda^Tx \geq \beta$ such that $\lambda^Ty > \beta$ for
  each $y \in \cone(X)$ and $\lambda^Tb < \beta$. Since $0 \in \cone(X)$ it follows that $0 > \beta$ and thus that
  $\lambda^Tb<0$. 
  
\end{proof}

\section{Decomposition theorem for polyhedra}
\label{sec:decomp-theor-polyh}

In the following we use the notation $P(A,b) = \{x \in\R^n \colon Ax \leq b\}$ for the polyhedron that is defined by $Ax \leq b$. We prove the Minkowski-Weyl theorem in this section that shows that polyhedra can be decomposed into the Minkowski sum of a polytope and a cone. 

\begin{definition}
\label{po:def:5}
An inequality $a^Tx\leq\beta$ is called an \emph{implicit equality} of
$Ax\leq b$ if each $x^* \in P(A,b)$ satisfies $a^Tx^* = \beta$. We denote the
subsystem  consisting of implicit equalities of $Ax\leq b$ by $A^=x\leq b^=$
and the subsystem consisting of the other inequalities by
$A^\leq x\leq b^\leq$. A constraint is \emph{redundant} if its removal from
$Ax\leq b$ does not change the set of feasible solution of $Ax\leq b$.  
\end{definition}

In the following, a vector $x$ satisfies $Ax < b$ if and only if
$a_i^T x < b_i$ for all $1\leq i\leq m$, where $a_1$,\ldots,$a_m$ are the rows of $A$.

\begin{lemma}
  \label{po:lem:3}
  Let $P(A,b)$ be a non-empty polyhedron.
  Then there exists an $x \in P(A,b)$ with $A^\leq x<b^\leq$. 
\end{lemma}
\begin{proof}
  Suppose that the inequalities in  $A^\leq x\leq b^\leq$ are $a_1^Tx\leq\beta_1
  ,\ldots,a_k^Tx\leq\beta_k$. For each $1\leq i\leq k$ there exists an $x_i \in P$ with
  $a_i^Tx_i<\beta_i$. Thus the point $x = 1/k (x_1+\cdots+x_k)$ is a point of
  $P(A,b)$ satisfying  $A^\leq x<b^\leq$. 
\end{proof}


\begin{lemma}
  \label{po:lem:2}
  Let $Ax\leq b$ be a system of inequalities. One has 
  \begin{displaymath}
    \affhull(P(A,b)) = \{ x \in \setR^n \mid A^=x = b^=\} = \{ x \in \setR^n \mid A^=x\leq b^=\}.
  \end{displaymath}
\end{lemma}
\begin{proof}
  Let $x_1,\ldots,x_t \in P(A,b)$ and suppose that $a^Tx\leq\beta$ is an
  implicit equality. Then since $a^Tx_i = \beta$ one has
  $a^T(\sum_{j=1}^t\lambda_ix_i) = \beta$. Therefore the inclusions $\subseteq$
  follow. 

  Suppose now that $x_0$ satisfies $A^=x\leq b^=$. Let $x_1 \in P(A,b)$
  with $A^\leq x_1<b^\leq$. If $x_0=x_1$ then $x_0 \in P(A,b) \subseteq
  \affhull(P(A,b))$.  Otherwise the line segment between $x_0$ and
  $x_1$ contains more than one point in $P$ and thus $x_0 \in
  \affhull(P)$. 
\end{proof}








A nonempty set  $C\subseteq\setR^n$ is a \emph{cone} if $\lambda\, x + \mu\,y \in C$
for each $x,y\in C$ and $\lambda,\mu\in \setR_{\geq0}$. A cone  $C$ is \emph{polyhedral}
if $C = \{ x \in \setR^n \mid Ax\leq0\}$. A cone \emph{generated by} vectors
$x_1,\ldots,x_m \in \setR^n$ is a set of the form $C = \{ \sum_{i=1}^m \lambda_i x_i
\mid \lambda_i\in \setR_{\geq0}, \, i=1,\ldots,m\}$.   A point $x = \sum_{i=1}^m \lambda_i x_i$
with  $\lambda_i\in \setR_{\geq0}, \, i=1,\ldots,m$ is called a \emph{conic
  combination}   of the $x_1,\ldots,x_m$. The set of conic combinations of
$X$ is denoted by $\cone(X)$. 


\begin{theorem}[Farkas-Minkowsi-Weyl theorem]
  \label{po:thr:3}
  A convex cone is polyhedral if and only if it is finitely
  generated. 
\end{theorem}


\begin{proof}
  Suppose that $a_1,\ldots,a_m$ span $\setR^n$ and consider the cone $C = \{
  \sum_{i=1}^m \lambda_i a_i \mid \lambda_i\geq0, \, i=1,\ldots,m\}$.
  Let $b \notin C$.
  Then the system $A\lambda = b$, $\lambda \geq 0$ has no solution.
  By Theorem~\ref{conv:thr:12} (Farkas' lemma), this implies that there exists a $y\in\setR^n$
  such that $A^Ty \leq 0$ and $b^Ty > 0$.

  Suppose that the columns of $A$ which correspond to
  inequalities in $A^Ty\leq0$   that are satisfied  by $y$ with equality
  have rank $<n-1$.
  Denote these columns by $a_{i_1},\ldots,a_{i_k}$.  
  Then there exists a $v\neq0$ which is orthogonal to
  each of these columns and to $b$, i.e., $a_{i_j}^Tv = 0$ for each
  $j=1,\ldots,k$ and $b^Tv =0 $. 
  There also exists a column $a^*$ of $A$ which is not in the set
  $\{a_{i_1},\ldots,a_{i_k}\}$ such that $(a^*)^Tv>0$ since the columns of
  $A$ span $\setR^n$. Therefore there exists an $\epsilon>0$ such that 
  \begin{enumerate}[i)]
  \item $A^T(y + \epsilon \cdot v)\leq0$ 
  \item The subspace generated by the columns of $A$ which correspond
    to inequalities of $A^Tx\leq0$ which are satisfied by $y + \epsilon \cdot v$
    with equality strictly contains $\langle a_{i_1},\ldots,a_{i_k}\rangle$. 
  \end{enumerate}
  
  Notice that we have $b^Ty = b^T(y + \epsilon \cdot v)>0$. 

  Continuing this way, we obtain a solution of the form $y + u$ of
  $A^Tx\leq0$ such that one has $n-1$ linearly independent columns of $A$
  whose corresponding inequality in $A^Tx\leq0$ are satisfied with
  equality.   Thus we see that each $b$ which does
  not belong to $C$ can be separated from $C$ with an inequality of
  the form $c^Tx\leq0$  which
  is uniquely defined by $n-1$ linearly independent vectors from the set
  $a_1,\ldots,a_m$.  This shows that $C$ is polyhedral. 

  Suppose now that $a_1,\ldots,a_m$ do not span $\setR^n$. Then there  exist
  linearly independent vectors $d_1,\ldots,d_k$ such that each $d_i$ is
  orthogonal to each of the $a_1,\ldots,a_m$ and $a_1,\ldots,a_m,d_1,\ldots,d_k$
  spans $\setR^n$.   The cone generated by
  $a_1,\ldots,a_m,d_1,\ldots,d_k$ is polyhedral and thus of the form $Ax\leq0$
  with some matrix $A\in \setR^{m\times n}$. Suppose that  $\langle a_1,\ldots,a_m\rangle = \{x
  \in \setR^n \mid Ux = 0\}$.  Now $C = \{ x \in \setR^n \mid Ax\leq0, \, Ux = 0\}$
  and $C$ is polyhedral. 


  Now suppose that $C = \{ x \in \setR^n \mid a_1^Tx\leq0,\ldots,a_m^Tx\leq0\}$.
  The cone
  \[ C' := \cone(a_1,\ldots,a_m) = \{ \sum_{i=1}^m \lambda_i a_i \mid
  \lambda_i\geq0,\,i=1,\ldots,m\} \]
  is polyhedral and thus of the form $C' = \{ x
  \in \setR^n \mid b_1^Tx\leq0, \ldots,b_k^Tx\leq0\}$. Clearly,
  $\cone(b_1,\ldots,b_k)\subseteq C$  since $b_i^Ta_j\leq0$. Suppose now that $y \in
  C \textbackslash{} \cone(b_1,\ldots,b_k)$. Then, since $\cone(b_1,\ldots,b_k)$ is
  polyhedral, there exists a $w\in \setR^n$ with $w^Ty>0$ and $w^Tb_i\leq0$
  for each $i=1,\ldots,k$. From the latter we conclude that $w \in
  C'$. From $y \in C$ and $w \in C'$ we conclude $w^Ty\leq0$, which is a contradiction.
\end{proof}

A set of vectors $Q = \conv(X)$, where $X\subseteq\setR^n$ is finite is called a
\emph{polytope}. 

\begin{theorem}[Decomposition theorem for polyhedra]
  \label{po:thr:4}
  A set $P\subseteq\setR^n$ is a polyhedron if and only if $P = Q + C$ for some
  polytope $Q$ and a polyhedral cone $C$. 
\end{theorem}

\begin{proof}
  Suppose $P = \{ x \in \setR^n \mid Ax\leq b\}$ is a polyhedron. Consider the
  polyhedral cone 
  \begin{equation}
    \label{po:eq:6}
    \left\{\mat{x\\\lambda} \mid x \in \setR^n, \, \lambda \in \setR_{\geq0}; Ax -
      \lambda b\leq0\right\} 
  \end{equation}
  is generated by finitely many vectors $\mat{x_i\\\lambda_i}$,
  $i=1,\ldots,m$. By scaling with a positive number we may assume that
  each $\lambda_i\in \{0,1\}$.  Let $Q$ be the convex hull of the $x_i$ with
  $\lambda_i=1$ and let $C$ be the cone generated by the $x_i$ with
  $\lambda_i=0$. A point $x \in \setR^n$ is in $P$ if and only if $\mat{x\\1}$
  belongs to~\eqref{po:eq:6} and thus if and only if 
  \begin{displaymath}
    \mat{x\\1} \in
    \cone\left\{\mat{x_1\\\lambda_1},\ldots,\mat{x_m\\\lambda_m}\right\}. 
  \end{displaymath}
  Therefore $P = Q + C$. 


  Suppose now that $P = Q+C$ for some polytope $Q$ and a polyhedral
  cone $C$ with $Q = \conv(x_1,\ldots,x_m)$ and $C = \cone(y_1,\ldots,y_t)$. A
  vector $x_0$ is in $P$ if and only if 
  \begin{equation}
    \label{po:eq:7}
    \mat{x_0\\1} \in \cone\left\{\mat{x_1\\1},\ldots, \mat{x_m\\1},\mat{y_1\\0},\ldots,\mat{y_t\\0}     \right\}
  \end{equation}
By Theorem~\ref{po:thr:3} \eqref{po:eq:7} is equal to 
\begin{equation}
\label{po:eq:12}
  \left\{ \mat{x\\\lambda} \mid Ax - \lambda b \leq0\right\}
  \end{equation}
for some matrix $A$ and vector $b$.  Thus $x_0\in P$  if and only if
$Ax_0\leq b$ and thus $P$ is a polyhedron.

\end{proof}




\begin{figure}[htbp]
  \begin{center}
    %\resizebox{4cm}{4cm}{%
    \includegraphics{figures/PicPolyhedra1.pdf}
  %}
  \end{center}
  \caption{A polyhedron and its decomposition into $Q$ and $C$\label{po:fig:decomp}}
\end{figure}




Let $P = \{ x \in \setR^n \mid Ax\leq b\}$. The \emph{characteristic cone} is 
$\charcone(P) =\{ y \mid y+x \in P \text{ for all } x \in P\} = \{y \mid Ay
\leq0\}$. One has
\begin{enumerate}[i)]
\item $y \in \charcone(P)$ if and only if there exists an $x \in P$ such
  that $x + \lambda\,y \in P$ for all $\lambda\geq0$ 
\item $P + \charcone(P)=P$
\item $P$ is bounded if and only if $\charcone(P)=\{0\}$. 
\item If the decomposition of $P$ is $P = Q +C$, then $C = \charcone(P)$. 
\end{enumerate}



The \emph{lineality space} of $P$ is defined as $\charcone(P) \cap -
\charcone(P)$. A polyhedron is \emph{pointed}, if its lineality space is
$\{0\}$.






%\subsection{Faces}
%\label{po:sec:faces}
\begin{definition}
  \label{def:f4}
  A set $F\subseteq\setR^n$ is called a \emph{face} of $P$ if there exists a
  valid inequality $c^Tx\leq\delta$ for $P$ with $F = P \cap (c^Tx = \delta)$. 
\end{definition}
 
\begin{lemma}
  \label{po:lem:4}
  A set $\emptyset \neq F \subseteq \setR^n$ is a face of $P$
  if and only if $F = \{ x \in P \mid A'x = b'\}$ for a subset $A'x\leq b'$ of $Ax\leq b$.
\end{lemma}

\begin{proof}
  Suppose that $F = \{ x \in P \mid A'x = b'\}$. Consider the vector $c =
  1^TA'$ and $\delta = 1^Tb'$. The inequality $c^Tx\leq \delta$ is valid for
  $P$. It is satisfied with equality by each $x \in F$. If $x' \in P\textbackslash{}
  F$, then there exists an inequality $a^Tx\leq\beta$ of $A'x\leq b'$ such
  that $ a^Tx' < \beta$ and consequently $c^Tx'<\delta$. 

  On the other hand, if $c^Tx\leq\delta$ defines the face $F$,
  then by the linear programming duality (see chapter~\ref{cha:duality})
  \[ \max\{ c^Tx \mid Ax\leq b \} = \min\{ b^T\lambda \mid A^T\lambda = c, \lambda \geq 0 \} \]
  there exists a $\lambda \in \setR^m_{\geq 0}$ such that $c=\lambda^TA$ and $\delta = \lambda^Tb$.
  Let $A'x\leq b'$ be the
  subsystem of $Ax\leq b$ which corresponds to strictly positive entries
  in $Ax\leq b$. One has $F = \{ x \in P \mid A'x = b'\}$. 
\end{proof}





A \emph{facet} of $P$ is an inclusion-wise maximal face $F$ of $P$
with $F\neq P$.  An inequality $a^Tx\leq\beta$ of $Ax\leq b$ is called
\emph{redundant} if $P(A,b) = P(A',b')$, where $A'x\leq b'$ is the system
stemming from $Ax\leq b$ by deleting $a^Tx\leq\beta$.  A system $Ax\leq b$ is
irredundant if $Ax\leq b$ does not contain a redundant inequality.

\begin{lemma}
  \label{po:lem:1}
  Let $Ax\leq b$ be an irredundant system. 
  Then a set $F\subseteq P$ is a facet if and only if it is
  of the form $F = \{ x \in P \mid a^Tx = \beta\}$ for an
  inequality $a^Tx\leq\beta$ of $A^\leq x\leq b^\leq$. 
\end{lemma}


\begin{proof}
  Let $F$ be a facet of $P$. Then $F = \{x \in P \mid c^Tx\leq\delta\}$ for a valid
  inequality $c^Tx\leq\delta$ of $P$. There exists a $\lambda \in \setR_{\geq0}^m$ with
  $c=\lambda^TA$ and $\delta=\lambda^Tb$.  There exists an inequality  $a^Tx\leq\beta$ of
  $A^\leq x\leq b^\leq$ whose corresponding entry in $\lambda$ is strictly
  positive. Clearly $F\subseteq\{x \in P \mid a^Tx=\beta\}\subset P$. Since $F$ is an
  inclusion-wise maximal face one has $F = \{x \in P \mid a^Tx=\beta\}$. 
  
  Let $F$ be of the form $F = \{ x \in P \mid a^Tx = \beta\}$ for an inequality
  $a^Tx\leq\beta$ of $A^\leq x\leq b^\leq$.  Clearly $F \neq \emptyset$ since the system $Ax\leq b$ is
  irredundant. If $F$ is not a facet, then $F\subseteq F'=\{ x \in P \mid a'^Tx = \beta'\}$
  with another inequality $a'^Tx\leq\beta'$ of $A^\leq x\leq b^\leq$. Let $x^*\in \setR^n$ be a point with
  $a^Tx^*>\beta$ and which satisfies all other inequalities of $Ax\leq b$. Such an $x^*$
  exists, since $Ax\leq b$ is irredundant.  Let $\wt{x}\in P$ with
  $A^\leq\wt{x}<b^\leq$.  There exists   a point $\wb{x}$ on the
  line-segment $\wb{\wt{x}x^*}$ with   $a^T\wb{x}=\beta$.  This point is then
  also in $F'$ and thus $a'^Tx = \beta'$ follows. This shows that
  $a'^Tx^*>\beta'$
  and thus $a^Tx\leq\beta$ can be
  removed from the system.  This is a contradiction to 
   $Ax\leq b$ being irredundant.  
\end{proof}



\begin{lemma}
  \label{po:lem:5}
  A face $F$ of $P(A,b)$ is inclusion-wise minimal if and only if it
  is of the form $F = \{ x \in \setR^n \mid A'x=b'\}$ for some subsystem
    $A'x\leq b'$ of $Ax\leq b$. 
\end{lemma}


\begin{proof}
  Let $F$ be a minimal face of $P$ and let $A'x\leq b'$ a the subsystem
  of inequalities of $Ax\leq b$ with $F = \{ x \in P \mid A'x=b'\}$. Suppose that
  $F \subset \{x\in\setR^n  \mid A'x=b'\}$ and let $x_1 \in \setR^n \textbackslash{} P$ satisfy
  $A'x_1=b'$ and  $x_2 \in F$.  There exists ``a first''  inequality $a^Tx\leq\beta$ of
  $Ax\leq b$ which is ``hit'' by the line-segment  $\wb{x_2x_1}$. Let $x^* =
  \wb{x_2x_1}\cap(a^Tx=\beta)$. Then $x^*\in F$ and thus $F \cap (a^Tx=\beta) \neq
  \emptyset$. But $F \supset F \cap (a^Tx=\beta)$ since $a^Tx \leq\beta$ is not an inequality of
  $A'x\leq b'$. This is a contradiction to the minimality of $F$. 

  
  Suppose that $F$ is a face with $F = \{x \in \setR^n \mid A'x= b'\}
  = \{x \in P \mid A'x=b'\}$ for a subsystem $A'x\leq b'$ of $Ax\leq
  b$. Suppose that there exists a face $\wt{F}$ of $P$ with $\emptyset
  \subset \wt{F}\subset F$. By Lemma~\ref{po:lem:4} $\wt{F} = \{ x \in
  P \mid A'x=b', A^*x=b^*\}$, where $A^*x\leq b^*$ is a sub-system of
  $Ax\leq b$ which contains an inequality $a^Tx\leq\beta$ such that
  there exists an $x_1,x_2 \in F$ with $a^Tx_1<\beta$ and
  $a^Tx_2\leq\beta$.  The line $\ell(x_1,x_2) = \{ x_1 +
  \lambda(x_2-x_1) \mid \lambda \in \setR\}$ is contained in $F$ but
  is not contained in $a^Tx\leq\beta$. This shows that $F$ is not
  contained in $P$ which is a contradiction.
\end{proof}

Exercise~\ref{item:19} asks for a proof of the following corollary. 
\begin{corollary}
  \label{co:8}
  Let $F_1$ and $F_2$ be two inclusion-wise minimal faces of
  $P=\{x\in\setR^n\colon Ax \leq b\}$, then $\dim(F_1) = \dim(F_2)$.
\end{corollary}

We say that a polyhedron contains a line $\ell(x_1,x_2)$ with $x_1
\neq x_2 \in P$ if $\ell(x_1,x_2) = \{ x_1 + \lambda(x_2-x_1) \mid
\lambda \in \setR \}\subseteq P$. A \emph{vertex} of $P$ is a
$0$-dimensional face of $P$. An \emph{edge} of $P$ is a
$1$-dimensional face of $P$.






\begin{example}
Consider a linear program $\min\{c^Tx \colon Ax = b,\, x\geq0\}$. A basic
feasible solution  defined by the basis $B\subseteq\{1,\ldots,n\}$ is a vertex of
the polyhedron $P = \{x \in \setR^n \colon Ax = b,\, x\geq0\}$. This can be seen
as follows. The inequality $a^Tx\geq0$ is valid for $P$, where $a_B =
\mathbf{0}$ and $a_{\overline{B}} = \mathbf{1}$. The inequality is
satisfied with equality by a point $x^* \in P$ if and only if
$x^*_{\overline{B}} = \mathbf{0}$. Since the columns of $A_B$ are
linearly independent, as $B$ is a basis, the unique point which
satisfies $a^Tx\geq0$ with equality is the basic feasible solution  
\end{example}

In exercise~\ref{item:co18} you are asked to show that the simplex
method can be geometrically interpreted as a walk on the graph $G =
(V,E)$, where $V$ is the set of basic feasible solutions and $uv\in  E$
if and only if $\conv\{u,v\}$ is a $1$-dimensional face of the
polyhedron defined by the linear program. 






\section*{Exercises}

\begin{enumerate}[1)]
\item Let $\{C_i\}_{i\in I}$ be a family of convex subsets of $\setR^n$.
  Show that the intersection $\bigcap_{i\in I} C_i$ is convex.
\item Show that the set of feasible solutions  of a linear program  is
  convex. \label{conv:item:1}
\item Prove Carath\'eodory's Theorem for convex hulls,
  Corollary~\ref{conv:co:1}. 
\item Let $A \in \setR^{n\times n}$ be a non-singular matrix and let
  $a_1,\ldots,a_n\in \setR^n$ be the columns of $A$.  Show that
  $\cone(\{a_1,\ldots,a_n\})$ is the polyhedron $P = \{ y \in \setR^n \colon
  A^{-1} y\geq0\}$. \label{conv:item:3} Show that $\cone(\{a_1,\ldots,a_k\})$ for
  $k\leq n$ is the set $P_k = \{y \in \setR^n \colon
  a_i^{-1} x\geq0, i=1,\ldots,k, \, a_i^{-1}x = 0, i=k+1,\ldots,n\}$, where
  $a_i^{-1}$ denotes the $i$-th row of $A^{-1}$. 
\item Prove that for a finite set $X\subseteq\setR^n$ the conic hull $\cone(X)$ is closed
  and convex. \label{conv:item:2} 
  
  \emph{Hint: Use Carath\'eodory's theorem and exercise~\ref{conv:item:3}.}
%{
%If $\cone(X)$ is not closed, then the exist a $y \in
%    \setR^n$ such that each $\varepsilon$-environment $U_\varepsilon(y) = \{x \in \setR^n \colon
%    \|x-y\|<\varepsilon\}$ contains an element of  $\cone(X)$. With
%    Carath\'eodory's theorem, there exist a sequence of points
%    $x_i\in\cone(\wt{X})$  which converges to $y$, where $\wt{X}$ is
%    linearly independent. Use now exercise~\ref{conv:item:3}.}
\item Find a countably infinite set $X\subset\setR^2$ such that $\cone(X)$ is not closed.
  Are there any cones that are open?
\item Prove Theorem~\ref{conv:thr:8}. 
\item Prove Theorem~\ref{conv:thr:9}. 
\item Let $f: \setR^n \to \setR^d$ be a linear map. 
  \begin{enumerate}[a)]
  \item Show that $f(K) = \{ f(x) \colon x \in K\}$  is convex if $K$ is
    convex.  Is the reverse also true? 
  \item  For $X\subseteq\setR^n$ arbitrary, prove that $\conv(f(X)) = f(\conv(X))$.
  \end{enumerate}
\item
  \label{conv:ex:10}
  Using Theorem~\ref{conv:thr:12}, prove the following variant of Farkas' lemma:
  Let $A\in\setR^{m\times n}$ be a matrix and $b\in\setR^m$ be a vector.
  The system $Ax \leq b$, $x\in\setR^n$ has a solution if and only if
  for all $\lambda\in\setR^m_{\geq0}$ with $\lambda^T A = 0$ one has $\lambda^T b \geq 0$.
\item 
  \label{ex:16}
  Provide an example of a convex and closed  set $K\subseteq\setR^2$ and a
  linear objective function $c^Tx$ such that $\inf\{c^Tx \colon
  x\in K\}>-\infty$ but there does not exist an $x^* \in K$ with $c^Tx^* \leq
  c^Tx$ for all $x \in K$. 
\item Consider the vectors
$$ x_1 = \left(\begin{array}{c} 3 \\ 1 \\ 2\end{array}\right), x_2 = \left(\begin{array}{c} 1 \\ 2 \\ 5 \end{array}\right), x_3 = \left(\begin{array}{c} 2 \\ 0 \\ 1 \end{array}\right), x_4 = \left(\begin{array}{c}  2 \\ 4 \\ 3 \end{array}\right), x_5 = \left(\begin{array}{c}  1 \\ 1 \\ 1 \end{array}\right). $$

Let $A=\{x_1, \ldots, x_5\}$.
Find two disjoint subsets $A_1,~A_2\subseteq A$ such that 
$$ \conv(A_1)\cap\conv(A_2)\neq \emptyset.$$ 
\emph{Hint: Recall the proof of Radon's lemma}
\item Consider the vectors
$$ x_1 = \left(\begin{array}{c} 3 \\ 1 \\ 2\end{array}\right), x_2 = \left(\begin{array}{c} 1 \\ 2 \\ 5 \end{array}\right), x_3 = \left(\begin{array}{c} 2 \\ 0 \\ 1 \end{array}\right), x_4 = \left(\begin{array}{c}  2 \\ 4 \\ 3 \end{array}\right), x_5 = \left(\begin{array}{c}  1 \\ 1 \\ 1 \end{array}\right). $$

The vector 
$$ v= x_1 + 3 x_2 + 2 x_3 + x_4 + 3 x_5 =  \left(\begin{array}{c}  15\\ 14 \\ 25 \end{array}\right)$$
is a conic combination of the $x_i$.

Write $v$ as a conic combination using only three vectors of the $x_i$.

\emph{Hint: Recall the proof of Carath\'eodory's theorem}

\item 
  Prove that each nonempty  polyhedron $P\subseteq\setR^n$ can be represented as $ P = L + Q$,
  where  $L\subseteq\setR^n$ is a linear space and $Q\subseteq\setR^n$ is a pointed
  polyhedron.   \label{po:ex:1}
\item Let $P\subset\setR^n$ be a polytope and $f:\setR^n\to\setR^m$ a linear map.
  \begin{enumerate}[i)]
  \item Show that $f(P)$ is a polytope.
  \item Let $y\in\setR^m$ be a vertex of $f(P)$. Show that there is a vertex $x\in\setR^n$ of $P$
    such that $f(x) = y$.
  \end{enumerate}
\item 
  Let $A \in \setR^{m\times n}$ and $b \in \setR^m$ and consider the polyhedron
  $P = P(A,b)$. Show that $\dim(P) = n - \rank(A^=)$.   \label{po:ex:2}
\item 
  \begin{enumerate}[i)]
  \item 
  Show that the dimension of each minimal face of a polyhedron $P$ is
  equal to $n - \rank(A)$. 
  \item
  Show that a polyhedron has a vertex if and only if the polyhedron
  does not contain a line. 
\end{enumerate}   \label{po:ex:3}
\item Show that the affine  dimension of the minimal faces of a 
  polyhedron $P = \{x \in \setR^n \colon Ax\leq b\}$ is invariant. \label{item:19}
\item 
  In this exercise you can assume that a linear program $\max\{c^Tx \mid
  Ax\leq b \}$ can be solved in polynomial time. Suppose that $P(A,b)$
  has vertices and that the linear program is bounded. Show how to
  compute an optimal \emph{vertex} solution of the linear
  program in polynomial time.    \label{po:ex:4}
\item  Let $P = \{x \in \setR^n \colon Ax = b, \, x\geq0\}$ be a polyhedron,
  where $A \in \setR^{m\times n}$ has full row-rank. Let $B_1,B_2$ be two bases
  such that $|B_1\cap B_2| = m-1$ and suppose that the associated basic
  solutions $x^*_1$ and $x^*_2$ are feasible. Show that, if
  $x_1\neq x_2$, then   $\conv\{x_1^*,x_2^*\}$ is a $1$-dimensional face
  of $P$. \label{item:co18} 
\end{enumerate}








% \bibliographystyle{abbrv}

% \bibliography{mybib,papers,books,my_publications}


%%% Local Variables: 
%%% mode: latex
%%% TeX-master: "lecture"
%%% TeX-master: "lecture"
%%% End: 

\chapter{The simplex method}
\label{cha:simplex-method}


In this chapter we describe the simplex method. The task is to solve a linear program 
\begin{equation}
  \label{eq:s-1}
  \max\{c^Tx \colon x \in \R^n, \,  Ax \leq b\}.
\end{equation}
We make the following assumption.
\begin{quote}
  The matrix $A \in\R^{m\times n}$ is of full column-rank. In other words, the columns of $A$ are linearly independent. 
\end{quote}
%
This assumption is not a restriction, since we can solve  the following equivalent linear program instead, where each $x_i$ is represented as the difference of two positive values $x_i = x_i^+ - x_i^-$. 
\begin{equation}
  \label{eq:s-2}
    \max\{c^Tx^+-c^Tx^- \colon x^+ ,x^- \in \R^n, \,  Ax^+-Ax^- \leq b, x^+ \geq 0, \, x^- \geq 0\}.
\end{equation}
The constraint matrix of the linear program~\eqref{eq:s-2} in inequality standard form is 
\begin{displaymath}
  \begin{pmatrix}
    A & -A \\
    -I_n & \mathbf{0} \\
    \mathbf{0} & -I_n
  \end{pmatrix},
\end{displaymath}
where $I_n$ is the $n\times n$ identity matrix and $\mathbf{0}$ is the $n\times n$ all-zero matrix. Clearly this matrix has linearly independent columns. 


\section{Adjacent vertices}
\label{sec:adjacent-vertices}


Let $P = \{x \in \R^n \colon Ax \leq b\}$ be the polyhedron of feasible solutions of \eqref{eq:s-1}. 

\begin{definition}
  \label{def:s-1}
  Two extreme points $x_1\neq x_2$ of $P$ are \emph{adjacent}, if there exists a valid inequality $d^Tx \leq \delta$ of $P$ such that 
  \begin{displaymath}
    P \cap \{x \in \R^n \colon d^Tx = \delta\} = \{ \lambda x_1 + (1-\lambda)x_2 \colon \lambda \in [0,1]\}. 
  \end{displaymath}
\end{definition}
In other words, $x_1 \neq x_2$ are adjacent if there exists a valid inequality for $P$ such that the points of $P$ that satisfy this inequality with equality are exactly the line-segment spanned by $x_1$ and $x_2$. 

Similar to the characterization of extreme points in Theorem~\ref{thr:1}, we can state and prove the following theorem. 

\begin{theorem}
  \label{thr:s-3}
  Two distinct vertices $x_1$ and $x_2$ of $P$ are adjacent if and
  only if there exists a sub-system $A'x \leq b'$ of $Ax \leq b$ such
  that
  \begin{enumerate}[i)]
  \item $A' \in \R^{(n-1)\times n}$ and the rows of $A'$ are linearly
    independent. \label{item:s-4}
  \item $x_1$ and $x_2$ satisfy $A'x \leq b'$ with
    equality. \label{item:s-3}
  \end{enumerate}
\end{theorem}

\begin{proof}
  Suppose that $x_1 \neq x_2$ are adjacent and suppose that $d^Tx \leq
  \delta$ is a valid inequality that asserts this fact. Consider the sub-system $\wt{A} x \leq \wt{b}$ of inequalities  of $Ax \leq b$ that are satisfied by $1/2 (x_1 + x_2)$ with equality. These are the inequalities that are satisfied by all points on the line-segment with equality. Since  $\wt{A}(x_1-x_2) = 0$, one has $\rank(\wt{A}) \leq n-1$. If $\rank(\wt{A}) < n-1$, then there exists a $v \in \R^n$ that is linearly independent from $(x_1-x_2)$ that satisfies $\wt{A} v = 0$. Consequently there exist a $\varepsilon >0$ such that 
  \begin{equation}
    \label{s:3}    
    \{ \frac{1}{2}(x_1+x_2) + \mu_1 (x_1-x_2) + \mu_2 v \colon -\varepsilon \leq \mu_1,\mu_2 \leq \varepsilon \} \subseteq P. 
  \end{equation}
All points of the set~\eqref{s:3} satisfy $d^Tx \leq \delta$ with equality and they are not a subset of the line-segment spanned by $x_1$ and $x_2$. From this we conclude that $\rank(\wt{A}) = n-1$ which implies that there exists a sub-system $A'x \leq b'$ satisfying~\ref{item:s-4}) and \ref{item:s-3}). 

Suppose on the other hand that there exists a sub-system $A'x \leq b'$
that satisfies \ref{item:s-4}) and \ref{item:s-3}).  The line spanned
by $x_1$ and $x_2$ is the set of points of $\R^n$ that satisfies $A'x
= b'$ and the intersection of this line with $P$ is, since $x_1$
and $x_2$ are vertices, the line-segment spanned by these two points.

The inequality $\mathbf{1}^T A'x \leq \mathbf{1}^T b'$ is valid for $P$ and is satisfied by the line-segment spanned by $x_1$ and $x_2$ with equality. Let $y^* \in P$ be a point  that does not lie on the line segment. Then one of the inequalities of $A'x \leq b'$ is satisfied by $y^*$ with strict inequality and thus $y^*$ does not satisfy $d^Tx \leq \delta$ with equality. 
\end{proof}




\section{Bases, feasible bases and vertices}
\label{sec:non-degenerate-case}

We will frequently use the following notation. Let $B \subseteq \{1,\dots,m\}$ then $A_B \in \R^{|B| \times n}$ is the matrix consisting of the rows of $A$ that are indexed by   $B$ and $b_B \in \R^{|B|}$  is the vector whose components are the ones of $b$ indexed by $B$. For example, for 
$    A = 
    \begin{pmatrix}
      3 & 2  \\
      7 & 1\\
      8 & 4\\
    \end{pmatrix}$,  $b = 
    \begin{pmatrix}
      3 \\ 2\\ 6
    \end{pmatrix}$ 
and  $B = \{2,3\}$, one has $A_B =  \begin{pmatrix}
          7 & 1\\
      8 & 4\\
    \end{pmatrix}$,  $b_B = 
    \begin{pmatrix}
      2\\ 6
    \end{pmatrix}$. 

    \begin{definition}
      \label{def:s-1}
      An index set $B \subseteq \{1,\dots,m\}$ is a \emph{basis} if
      $|B| = n$ and $A_B$ is non-singular. If in addition $x^* =
      A_B^{-1} b_B$ is feasible, then $B$ is called a \emph{feasible
        basis}.
    \end{definition}

    Theorem~\ref{thr:1} implies that every vertex $x^*$ of $P = \{x
    \in \R^n \colon Ax \leq b\}$ is \emph{represented} by a basis $B$,
    i.e. $x^*$ is the unique solution of $A_Bx^* = b_B$. This
    representation however must not be unique, see
    Figure~\ref{fig:7}. We say that a linear program is \emph{degenerate}, if
    there exists a basic solution $x^* \in \R^n$ that satisfies $n+1$
    inequalities with equality. Otherwise the linear program is called
    \emph{non-degenerate}. If the linear program is non-degenerate,
    then each vertex is represented by \emph{exactly one} basis.

    \begin{figure}
      \centering
      \includegraphics[height=3cm]{figures/Degenerate.pdf}
      \caption{The vertex $x^*$ is represented by each choice of two of the three tight constraints. The linear program is degenerate. The red vector is the objective function vector and the red labels are the indices of the constraints.}
      \label{fig:7}
    \end{figure}



    \begin{definition}
      \label{def:s-2}
       A basis $B$ is called \emph{optimal} if it is feasible and the
       unique $\lambda \in \setR^m$ 
       with 
       \begin{equation}
         \label{eq:s-5}
         \lambda^T A =  c^T \, \, \text{ and } \lambda_i = 0, \, i \notin B 
       \end{equation}
       satisfies $\lambda\geq0$. 
     \end{definition}


The basis $\{1,2\}$ in Figure~\ref{fig:7} is not optimal whereas the bases $\{2,3\}$ and $\{1,3\}$ are optimal bases. 


\begin{theorem}
  \label{thr:s-4}
  If $B$ is an optimal basis, then $x^* = A_B^{-1} b_B$ is an optimal solution of the linear program~\eqref{eq:s-1}. 
\end{theorem}

\begin{proof}
  The inequality $\lambda^T A x \leq \lambda^T b$ is valid for $P = \{x \in \R^n \colon Ax \leq b\}$. But $\lambda^T A =c^T$ and $\lambda^T b = \lambda^T Ax^* = c^Tx^*$. Consequently $x^*$ is an optimal solution of the linear program \eqref{eq:s-1}. 
\end{proof}


\section{Moving to an improving vertex}
\label{sec:moving-an-improving}

Suppose now that $B$ is a feasible but not optimal basis. Then the
unique $\lambda$ satisfying~\eqref{eq:s-5} has a negative component
$\lambda_i<0$ for some $i \in B$. The idea is now to move from $x^*_B
= A_B^{-1}b_B$ by remaining tight at all constraints indexed by $B$
except for $i$.

There is only one way this can be achieved. Namely by moving in the unique direction $d$ with
\begin{displaymath}
  a_j^T d =
  \begin{cases}
    0 & \text{ for } j \in B \setminus \{i\} \\
    -1 & \text{ if } j =i.
  \end{cases}
\end{displaymath}

When we do this, we follow the ray $x^*_B + \varepsilon \cdot d$ with $\varepsilon
\geq 0$. What happens to the objective function, as $\varepsilon$ grows? Since
$c^Td = \lambda^TA d = -\lambda_i >0$, the objective function strictly
grows with growing $\varepsilon$. There are now two cases. 

At some point, we hit the boundary of a constraint and further increase of $\varepsilon$ results in an infeasible point. Let $K \subseteq \{ 1,\dots,m\}$ be the set of indices 
\begin{equation}
  \label{eq:s-6}
  K = \{ k \colon 1 \leq k \leq m, \,a_k^T d >0\}. 
\end{equation}
Those are the indices of constraints that, at some point, will be violated. We can increase $\varepsilon$ until 
\begin{equation}
  \label{eq:s-8}
  \varepsilon^* = \min_{k \in K} \{ (b_k - a_k^Tx^*) / a_k^Td\}. 
\end{equation}
Now pick any $k \in K$ for which this minimum is achieved and set $B' = B \setminus \{i\} \cup \{k\}$. This is a feasible basis, $d$ is orthogonal to all rows  indexed by $B \setminus \{i\}$ but not to $a_k$. 

In the case where there is no constraint that puts an upper bound on $\varepsilon$, then the linear program is \emph{unbounded}. 
We have described one iteration of the simplex algorithm. We iterate this procedure until an optimal solution is found. 



\begin{algorithm}[Simplex algorithm]
 \begin{tabbing}
      Start with feasible basis $B$ \\[1ex]
      {\tt while} \= $B$ is not optimal \\ [.7ex]
      \> Let $i \in B$ be index with $\lambda_i<0$ \\
      \> Compute  $d \in \setR^n$ with $a_j^T d = 0, \, j \in B \setminus\{i\}$
      and $a_i^T d = -1$ \\ 
      \> Determine $K = \{ k \colon 1 \leq k \leq m, \, a_k^Td >0\}$\\[.7ex]  
      \> {\tt if} \= $K = \emptyset$ \\   
      \> \> \emph{assert LP unbounded} \\
      \> {\tt else} \\
      \> \> Let $k \in K$ index where 
     $
        \displaystyle \min_{k \in K} (b_k - a_k^Tx^*)/a_k^Td
      $
      is attained \\

      \> \>\emph{update} $B := B \setminus\{i\} \cup \{k\}$             
    \end{tabbing}
    
  \end{algorithm}
  


  \begin{theorem}
    \label{thr:s-5}
    If the linear program~\eqref{eq:s-1} is non-degenerate, then the simplex algorithm terminates. 
  \end{theorem}

  \begin{proof}
    In the non-degenerate case, $\varepsilon^*>0$ and the simplex
    algorithm makes progress, i.e., the objective function value
    strictly increases after each iteration. Since there is only a
    finite number of vertices, the algorithm terminates.
  \end{proof}

  
\section{Termination in the degenerate case}
\label{sec:term-degen-case}


In the case where the linear program~\eqref{eq:s-1} is degenerate, we cannot argue that the objective function value increases each iteration and that the simplex algorithm terminates. However, the simplex algorithm leaves us some choice. Namely, there can be several indices $i \in B$ such that $\lambda_i<0$. Also, there could be several indices $k \in K$ attaining the minimum in \eqref{eq:s-8}. If one adheres to the  \emph{smallest index rule}, then one can prove termination of the simplex algorithm also in the degenerate case. One iteration of the simplex algorithm is now as follows. 



\begin{algorithm}[Simplex algorithm with the smallest index rule]
 \begin{tabbing}
      Start with feasible basis $B$ \\[1ex]
      {\tt while} \= $B$ is not optimal \\ [.7ex]
      \> Let $i^* \in B$ be the \emph{ smallest index} with $\lambda_i<0$ \\
      \> Compute  $d \in \setR^n$ with $a_j^T d = 0, \, j \in B \setminus\{i\}$
      and $a_i^T d = -1$ \\ 
      \> Determine $K = \{ k \colon 1 \leq k \leq m, \, a_k^Td >0\}$\\[.7ex]  
      \> {\tt if} \= $K = \emptyset$ \\   
      \> \> \emph{assert LP unbounded} \\
      \> {\tt else} \\
      \> \> Let $k^* \in K$ the \emph{smallest index} where 
     $
        \displaystyle \min_{k \in K} (b_k - a_k^Tx^*)/a_k^Td
      $
      is attained \\

      \> \>\emph{update} $B := B \setminus\{i^*\} \cup \{k^*\}$             
    \end{tabbing}
    
  \end{algorithm}
  



  \begin{theorem}
    \label{thr:3}
    The simplex algorithm with the smallest index rule terminates. 
  \end{theorem}

  \begin{proof}
    We suppose that the simplex algorithm does not terminate. This means that the simplex algorithm iterates through a sequence of bases 
    \begin{displaymath}
      B_0,B_1,\dots,B_k
    \end{displaymath}
    with $B_k = B_0$. Inspecting two succeeding bases $B_\ell$ and
    $B_{\ell+1}$ for $0\leq \ell \leq k-1$, one has $B_{\ell+1}
    =B_{\ell} \setminus \{i\} \cup \{j\}$, i.e., $i$ \emph{leaves}
    and $j$ \emph{enters} $B_{\ell}$. Now let $j$ be \emph{the largest
      index} that leaves on that sequence. Since $B_0 = B_k$, $j$ also
    enters again at some point. Let $p$ and $q$ be the indices of
    bases, $0 \leq p,q <k$ where $j$ leaves and enters respectively. 

    Let $\lambda^{(p)}$ and $d^{(q)}$ be the corresponding $\lambda$ and $d$
    vectors from the iteration $p$ and $q$ of the simplex algorithm
    respectively. Since ${\lambda^{(p)}}^T A = c^T$ and $c^T d^{(q)}>0$ we conclude 
    \begin{equation}
      \label{eq:s15-1}
      {\lambda^{(p)}}^T A d^{(q)} >0.
    \end{equation}
    Let $i \in B_p$ be an index with 
    \begin{equation}
      \label{eq:s-15-2}
      \lambda^{(p)}_i a_i d^{(q)} >0
    \end{equation}
    where $a_i$ denotes the $i$-th row of $A$. 


    We now distinguish three cases. Let us suppose that $i>j$. Then,
    since $j$ is the largest index that ever leaves or enters one has
    that $i$ is also an element of $B_q$ implying that $ a_i d^{(q)}$
    is $0$ or $-1$. It cannot be $-1$, since this would mean that $i$
    leaves $B_q$ contradictory to the choice of $j$.

    Suppose then that $i<j$. Then $\lambda_i^{(p)}>0$ since $j$ is the
    smallest index that can leave the basis $B_p$. But $a_id^{(q)} >0$
    is not possible, otherwise index $i$ is an index where the minimum
    ($\varepsilon^* = 0$) in \eqref{eq:s-8} is attained and $j>i$ is
    the smallest index where this minimum is attained. Thus this case
    can also be ruled out.

    Finally, if $i=j$, then $\lambda_i^{(p)} <0$ and $a_i d^{(q)}>0$
    which also contradicts~\eqref{eq:s-15-2}.
  \end{proof}


\section{Finding an initial basic feasible solution}
\label{sec:finding-an-initial}

The simplex algorithm starts with a feasible basis. How can such a feasible basis be determined? In fact, this can be done with an auxiliary linear program. 

Any linear program has an equivalent form 
\begin{equation}
  \label{eq:s-3}
  \max\{ c^Tx \colon Ax \leq b, \, x \geq 0\}. 
\end{equation}
  We want to find an initial feasible basis for this linear program. Suppose first that we re-write the constraints $Ax \leq b$ as $A_1x\leq b_1$ and $A_2x\leq b_2$ with $b_1 \geq 0$ and $b_2<0$. Now consider the following linear program
  \begin{equation}
    \label{eq:s-151}    
  \min\{ \mathbf{1}^Ty \colon A_1x \leq b_1, \, A_2x \leq b_2 +y, \, x,y \geq 0, y \leq -b_2\}. 
\end{equation}
An initial basic feasible solution is $x=0$ and $y= -b_2$ with the basis corresponding to the inequalities $x\geq 0$ and $y \leq -b_2$. The simplex algorithm applied to this linear program terminates. It finds an optimal solution with objective value $0$ if and only if the original linear program is feasible. Let $B$ be the optimal basis in this case. Then $B$ without the indices corresponding to the constraints $y\geq 0$ is a feasible basis for the linear program~\eqref{eq:s-3}. 


\section{Removing degeneracy by perturbation}
\label{sec:remov-degen-pert}

In the following, we derive an alternative pivoting rule that also ensures termination of the simplex algorithm. We begin by \emph{perturbing} the right-hand-sides of our constraints 
\begin{equation}
  \label{eq:pert-2}
  Ax \leq b
\end{equation}
by adding a vector $p_\eps$ 
\begin{equation}
  \label{eq:per-1}
  Ax \leq b + p_\eps 
\end{equation}
where $p_\eps$ is the vector 
\begin{displaymath}
p_\eps = 
  \begin{pmatrix}
    \eps\\ \eps^2\\ \vdots \\ \eps^m
  \end{pmatrix}.
\end{displaymath}

\begin{lemma}
  \label{lem:per-2}
  If $Ax \leq b$ is feasible, then $ Ax \leq b + p_\eps $ is feasible
  for each $\eps>0$. If $B$ is an infeasible basis of $Ax \leq b$,
  then $B$ is an infeasible basis of $Ax \leq b + p_\eps$ for $\eps>0$
  sufficiently small.
\end{lemma}

\begin{proof}
  Clearly, the feasible  solutions of \eqref{eq:per-1} contain the feasible solutions of~\eqref{eq:pert-2}. Suppose now that $B$ is infeasible, then there exists an index $i$ and some $\delta>0$ such that 
  \begin{displaymath}
    a_i A_B^{-1}b_B \geq b_i + \delta,
  \end{displaymath}
  where $a_i$ and $b_i$ are the $i$-th row of $A$ and the $i$-th
  component of $b$ respectively. Now 
  \begin{displaymath}
    a_i A_B^{-1}(b_B + {p_\eps}_B) \geq b_i + \delta + A_B^{-1}{p_\eps}_B > b_i
  \end{displaymath}
  if $\eps>0$ is sufficiently small. 
\end{proof}

Furthermore, we can show that the constraint system \eqref{eq:per-1} is non-degenerate for $\eps>0$ small enough. 

\begin{lemma}
  \label{lem:pert-3}
  If $\eps>0$ is small enough, then \eqref{eq:per-1} is non-degenerate. 
\end{lemma}

\begin{proof}
  If the set of  inequalities \eqref{eq:per-1} is degenerate, then there exists a basis $B$ and an index $i \notin B$ such that 
  \begin{equation}
    \label{eq:pert-3}
    a_ix^*_{B,\eps} = b_i+\eps^i
  \end{equation}
where $x^*_{B,\eps} = A_B^{-1}(b_B + (p_\eps)_B)$. Notice that 
  \begin{displaymath}
    a_ix^*_{B,\eps} - b_i - \eps^i
  \end{displaymath}
  is a nonzero polynomial in $\eps$. It is non-zero, since the coefficient of $\eps^i$ is $-1$ as $\eps^i$ is not a component of $(p_\eps)_B$. A nonzero polynomial has only a finite number of roots. Thus, if $\eps>0$ is small enough, no equation of the form \eqref{eq:pert-3} can hold. 
\end{proof}

Suppose now that we want to solve the linear program
\begin{equation}
  \label{eq:pert-5}
  \max\{c^Tx \colon x \in \R^n, Ax \leq b \}
\end{equation}
The idea is run the simplex algorithm  on the perturbed linear program 
\begin{equation}
  \label{eq:pert-4}
  \max\{c^Tx \colon x \in \R^n, Ax \leq b+p_\eps \}
\end{equation}
where $\eps>0$ is sufficiently small. This linear program is
non-degenerate. What is nice is that this perturbation does not have
to be computed explicitly. We can formulate a pivot rule for the
non-perturbed linear program \eqref{eq:pert-5} that is conform with the pivoting that is
performed on the perturbed linear program~\eqref{eq:pert-3}. For this, assume that $B$ is a feasible basis of the perturbed linear program~\eqref{eq:pert-3}. By  Lemma~\ref{lem:per-2}, $B$ is also feasible for the unperturbed linear program. Let us now consider the iteration of the simplex algorithm at the basis $B$. As before, we choose $i \in B$ with $\lambda_i<0$ arbitrary and compute $d \in \R^n$. Now we have to determine the unique index of the inequality of \eqref{eq:pert-4} that is hit first, when moving in the direction of $d$ in the perturbed linear program. 

This index can be determined as follows. As before, we determine  $K = \{ k \colon 1 \leq k \leq m, \, a_k^Td >0\}$ and we let $\wt{K}\subseteq K $ be those indices of $K$ where the minimum $\min_{k \in K} (b_k - a_k^Tx^*)/a_k^Td$ is attained. In the perturbed program we would have to determine the unique minimum 
\begin{equation}\label{eq:p2}
  b_k + \eps^k  - a_k^T A_B^{-1} (b_B + (p_\eps)_B)/a_k^Td, \,\,\, k \in \wt{K},
\end{equation}
where we imagine that $\eps$ tends to zero from above.  Each of the expressions in~\eqref{eq:p2} is a polynomial in $\eps$. The index that will leave in the perturbed linear program is the one that corresponds to the lexicographically minimal polynomial in~\eqref{eq:p2}. 





\begin{algorithm}[Simplex with largest index leaving rule]
  \label{alg:pert-1}
 \begin{tabbing}
      Start with feasible basis $B$ of the \emph{perturbed} linear program~\eqref{eq:pert-4} \\[1ex]
      {\tt while} \= $B$ is not optimal \\ [.7ex]
      \> Let $i \in B$ be index with $\lambda_i<0$ \\
      \> Compute  $d \in \setR^n$ with $a_j^T d = 0, \, j \in B \setminus\{i\}$
      and $a_i^T d = -1$ \\ 
      \> Determine $K = \{ k \colon 1 \leq k \leq m, \, a_k^Td >0\}$\\[.7ex]  
      \> {\tt if} \= $K = \emptyset$ \\   
      \> \> \emph{assert LP unbounded} \\
      \> {\tt else} \\
      \> \> Let $k \in K$ be the index corresponding to the lexicographically\\
      \> \> smallest polynomial of the form~\eqref{eq:p2} \\

      \> \>\emph{update} $B := B \setminus\{i\} \cup \{k\}$             
    \end{tabbing}
    
  \end{algorithm}
%
Exercise \ref{item:pert-3} explains how to convert a feasible basis of~\eqref{eq:pert-5} into a feasible basis of~\eqref{eq:pert-4}. 

\begin{theorem}
  \label{thr:6}
  The variant of the simplex method described in Algorithm~\ref{alg:pert-1} terminates. 
\end{theorem}


\section*{Exercises} 
\begin{enumerate}
\item For each of the following assertion, provide a proof or a counterexample. 
  \begin{enumerate}[i)]
  \item An index that has just left the basis $B$ in the simplex
    algorithm cannot enter in the very next iteration.
  \item An index that has just entered the basis $B$ in the simplex
    algorithm cannot leave again in the very next iteration. 
  \end{enumerate}
\item Consider the auxiliary linear program to find an initial feasible basis~\eqref{eq:s-151}. The constraint matrix of this linear program is of the form
  \begin{displaymath}
    \begin{pmatrix}
      A & 0 \\
      -I_n & 0\\ 
      0 & -I_{m_2}\\
      0 & I_{m_2}
    \end{pmatrix},
  \end{displaymath}
  where $m_2$ is the number of rows of $A_2$. This matrix has $m+n+2\cdot m_2$ rows. Describe an initial feasible basis that corresponds to the basic feasible solution $x = 0$ and $y=0$. 

Suppose that the optimal value of the auxiliary linear program is $0$ and let $B'$   be an optimal basis found by the simplex algorithm. Prove that $B' \setminus \{m+n+1,\dots,m+n+m_2\}$ is a feasible \emph{basis} of the linear program~\eqref{eq:s-3}. 

\item Suppose that the linear program $\max \{c^Tx \colon x \in \R^n, \, Ax \leq b\}$ is non-degenerate and $B$ is an optimal basis. Show that the linear program has a unique optimal solution if and only if $\lambda_B>0$. 

\item Let $P = \{ x \in \setR^n \colon Ax\leq b\}$ be a polyhedron. Show that the
  following are equivalent for a feasible $x^*$:
  \begin{enumerate}[i)] 
  \item $x^*$ is a vertex of $P$. 
  \item There exists a set $B\subseteq \{1,\ldots,m\}$ such that $|B| = n$, $A_B$  is
    invertible and $A_B x^* = b_B$. Here the matrix $A_B$ and the
    vector $b_B$ consists of the rows of $A$ indexed by $B$ and the
    components of $b$ indexed by $B$ respectively. 
  \item For every feasible $x_1, x_2 \neq x^* \in P$ one has $x^* \notin
    \conv\{x_1,x_2\}$. 
  \end{enumerate}  \label{i:item:1} 

\item A polyhedron $P = \{x \in \setR^n \colon Ax\leq b\}$ \emph{contains a
    line}, if there exists a nonzero $v \in \setR^n$ and an $x^* \in R
^n$ such that for all $\lambda \in \setR$, the point $x^* + \lambda\cdot v \in P$. Show
that a nonempty polyhedron $P$ contains a line if and only if $A$ does not have full
column-rank. \label{xitem:10}
\item \label{item:pert-3} Let $x^*$ be a basic feasible solution of the non-perturbed linear program~\eqref{eq:pert-5} and let $C\subseteq \{1,\dots,m\}$ be the indices of inequalities that are tight at $x^*$. Prove that the following \emph{greedy algorithm} produces a feasible basis $B\subseteq C$ of the perturbed linear program.
 \begin{tabbing}
    Initialize $B = \emptyset$ \\
    {\tt while} \= $B$ is not a basis \\ 
      \> Let $i\in C$ be the largest index such that the rows \\
      \>  indexed by  $B \cup\{i\}$ are linearly independent \\[.1ex]
      \> Update $B = B \cup\{i\}$. 
    \end{tabbing}
\end{enumerate}




%%% Local Variables: 
%%% mode: latex
%%% TeX-master: "lecture"
%%% End: 


\section{One iteration of the simplex algorithm}
\label{sec:one-iter-simpl}


In the following we will analyze the complexity of one iteration of the simplex algorithm.  We suppose that the 
input data  $A \in \setQ^{m\times n}$, $c \in \setQ^n$, $b \in \setQ^{m}$ is rational. 

Now if $A_B^{-1}$ has been computed, then $\lambda_B^T = c^T \cdot
A_B^{-1} $ and $d$, which is the negative of a column of $A_B^{-1}$,
can be computed with $O(n^2)$ operations.  To compute $K$ we have to
compute $A\cdot d$. This can be done with $O(m\cdot n)$
operations. The index of element entering the basis can be determined
by computing $x^* = A_B^{-1} b_B$, $b - Ax^*$ and $Ad$.  Thus, if
$A_B^{-1}$ is known, this amounts to a total of
  \begin{displaymath}
    O(m \cdot n) 
  \end{displaymath}
  arithmetic operations.
  

In order to argue that one iteration of the simplex algorithm runs in polynomial time, we have to show that  each of the numbers of  $A_B^{-1}$ has size that is polynomial in the size of the input and that 
$A_B^{-1}$ can be quickly computed. 


Let us first see how large the size of the numbers in $A_B^{-1}$ can be. 
Suppose that 
\begin{displaymath} 
  A =
  \begin{pmatrix}
    p_{11}/q_{11} & \cdots & p_{1n}/q_{1n} \\
             & \cdots &  \\
   p_{n1}/q_{n1} & \cdots & p_{nn}/q_{nn} \\
  \end{pmatrix} \in \setQ^{n\times n}.  
\end{displaymath}
is invertible. 
The {size}  of the product of denominators $\prod_{i=1,j=1}^n q_{ij}$
is  
clearly  linear in the size of the input. 
Now write $A = 1/Q \cdot A'$ where $Q$ is product of denominators and $A'\in
  \setZ^{n\times n}$  
 Since $A^{-1} = Q\cdot (A')^{-1} $ we only have to answer  this question for $A'$ instead of $A$. In other words, we can  assume that $A$ is integral. 



 For $A \in \setR^{m\times n}$ and $1\leq i\leq m$ and $1\leq j\leq
 n$, $A_{ij}$ denotes the matrix obtained from $A$ by deleting the
 $i$-th row and $j$-th column.
The following 
matrix inversion formula is known as Cramer's rule. 
\begin{displaymath}
  A^{-1} = \frac{1}{\det(A)}
  \begin{pmatrix}
    \det(A_{11}) & - \det(A_{21}) & \det(A_{31}) & \hdots  \\\
    -\det(A_{12}) &  \det(A_{22}) & - \det(A_{32}) & \hdots \\\
    \det(A_{13}) & - \det(A_{23}) & \det(A_{33}) & \hdots \\\
    \vdots          &     \vdots        & \vdots          & \hdots \\
    \vdots          &     \vdots        & \vdots          & \hdots \\
  \end{pmatrix}
\end{displaymath}

\begin{theorem}[Hadamard bound] 
  Let $A \in \R^{n \times n}$ be non-singular. Then 
  \begin{displaymath}
    |\det(A)| \leq \prod_{i=1}^n \|a_i\|_2 \leq n^{n/2} \cdot B^n, 
  \end{displaymath}
  where $B$ is upper bound on absolute values of entries of $A$.
\end{theorem}

\begin{proof}
  The Gram-Schmidt orthogonalization of $A$ yields a factorization 
  \begin{displaymath}
    A = Q \cdot R,
  \end{displaymath}
where $R$ is an upper triangular matrix with ones on the diagonal. The matrix $Q$ has orthogonal columns, where the length of the $i$-th column $q^{(i)}$ is upper bounded by the length of the $i$-th column of $A$. 
The assertion follows from 
\begin{displaymath}
  \det(A)^2 = \det(Q)^2 = \det(Q^T) \det(Q) = \prod_i \|q^{(i)}\|^2. 
\end{displaymath}
\end{proof}



\begin{corollary}
  If $A\in \setZ^{n\times n}$ is integral and each entry in absolute
  value is bounded by $B$, then $\size(\det(A)) = O(n \log n+ n \cdot
  \size(B))$.
\end{corollary}
% 
  
\begin{corollary}  
  Let $A \in \setQ^{n\times n}$  be an invertible matrix. The size of $A^{-1}$
  is polynomial in the size of $A$. 
\end{corollary}


Now we known that the size of  $A_B^{-1}$ polynomial in the size of the input
($A,b,c$)?  Now, how  expensive is it to compute $A_B^{-1}$? 
Suppose basis $B$  is preceded by  $B'$ 
with 
\begin{displaymath}    
  \begin{array}{lcrcl}
    B'&=& \{b_1,\ldots,b_{k-1},& b'_{k}& ,b_{k+1},\ldots,b_n\} \\
    B &=& \{b_1,\ldots,b_{k-1},& b_{k}& ,b_{k+1},\ldots,b_n\} \\
  \end{array}
\end{displaymath}
Then  each row of  $A_B \cdot A_{B'}^{-1}$, except for row $k$, is the corresponding row of the $n\times n$ identity matrix except. Let the $k$-th row be $(v_1,v_2,\dots,v_n)$. We now only have to perform the elementary column operations that turn this row into the $k$-th unit vector on $A_{B'}^{-1}$ to obtain $A_B^{-1}$. In other words, the following algorithm computes $A_B^{-1}$ given $A_{B'}^{-1}$. 


\begin{itemize}
\item Compute $a_{b_k}^T \cdot A_{B'}^{-1} = (v_1,\ldots,v_k,\ldots,v_n)$ 
\item For each column $i \neq k$: Subtract $v_i / v_k$ times column
  $k$ from column $i$ 
\item Divide column $k$ by $v_k$ 
\end{itemize}

This amounts to a total number of $O(n^2)$ arithmetic operations for the update. 
We can conclude with the following theorem. 
\begin{theorem}
  \label{thr-a-4}
  One iteration of the simplex algorithm requires a total number of
  $O(m\cdot n)$ operations on rational numbers whose size is polynomial
  in the input size. 
\end{theorem}



%%% Local Variables: 
%%% mode: latex
%%% TeX-master: "lecture"
%%% End: 

\chapter{Polyhedra and Integer Programming} 
\label{po:cha:polyhedra}


In this section we give definitions and fundamental facts about
polyhedra. An excellent reference for this topic is the book by
Schrijver~\cite{Schrijver86}.  A \emph{polyhedron} $P$ is a set of
vectors of the form $P=\{ x \in \setR^n \mid Ax \leq b\}$, for some matrix $A \in \setR^{m
  \times n}$ and some vector $b \in \setR^m$. We write $P(A,b)$.  The polyhedron
is \emph{rational} if both $A$ and $b$ can be chosen to be rational.

Recall that a  finite set $V\subseteq\setR^n$ is \emph{affinely independent} if for each $v\in
V$ one has $v\notin\affhull(V \textbackslash{} \{v\})$. This is equivalent to $(V - v) \setminus \{0\}$ being
linearly independent for each $v \in V$. The \emph{dimension} of $V$ is
 the size of the largest subset of $V$ which is affinely independent
 minus one:
\[ \dim(V) = \max\{ |U| - 1 \mid U \subseteq V \text{ is affinely independent} \} \]

\begin{example} 
\begin{itemize}
  \item $\dim(\setR^n) = n$
  \item $\dim(\{x\}) = 0$ for every $x\in\setR^n$
  \item $\dim(\emptyset) = -1$
\end{itemize}
\end{example}
Notice that $V\subseteq\setR^n$ is affinely independent if and only if $(V - v)
\setminus \{0\}$ is linearly independent for each $v \in V$.


\begin{definition}
\label{po:def:5}
An inequality $a^Tx\leq\beta$ is called an \emph{implicit equality} of
$Ax\leq b$ if each $x^* \in P(A,b)$ satisfies $a^Tx^* = \beta$. We denote the
subsystem  consisting of implicit equalities of $Ax\leq b$ by $A^=x\leq b^=$
and the subsystem consisting of the other inequalities by
$A^\leq x\leq b^\leq$. A constraint is \emph{redundant} if its removal from
$Ax\leq b$ does not change the set of feasible solution of $Ax\leq b$.  
\end{definition}

In the following, a vector $x$ satisfies $Ax < b$ if and only if
$a_i^T x < b_i$ for all $1\leq i\leq m$, where $a_1$,\ldots,$a_m$ are the rows of $A$.

\begin{lemma}
  \label{po:lem:3}
  Let $P(A,b)$ be a non-empty polyhedron.
  Then there exists an $x \in P(A,b)$ with $A^\leq x<b^\leq$. 
\end{lemma}
\begin{proof}
  Suppose that the inequalities in  $A^\leq x\leq b^\leq$ are $a_1^Tx\leq\beta_1
  ,\ldots,a_k^Tx\leq\beta_k$. For each $1\leq i\leq k$ there exists an $x_i \in P$ with
  $a_i^Tx_i<\beta_i$. Thus the point $x = 1/k (x_1+\cdots+x_k)$ is a point of
  $P(A,b)$ satisfying  $A^\leq x<b^\leq$. 
\end{proof}


\begin{lemma}
  \label{po:lem:2}
  Let $Ax\leq b$ be a system of inequalities. One has 
  \begin{displaymath}
    \affhull(P(A,b)) = \{ x \in \setR^n \mid A^=x = b^=\} = \{ x \in \setR^n \mid A^=x\leq b^=\}.
  \end{displaymath}
\end{lemma}
\begin{proof}
  Let $x_1,\ldots,x_t \in P(A,b)$ and suppose that $a^Tx\leq\beta$ is an
  implicit equality. Then since $a^Tx_i = \beta$ one has
  $a^T(\sum_{j=1}^t\lambda_ix_i) = \beta$. Therefore the inclusions $\subseteq$
  follow. 

  Suppose now that $x_0$ satisfies $A^=x\leq b^=$. Let $x_1 \in P(A,b)$
  with $A^\leq x_1<b^\leq$. If $x_0=x_1$ then $x_0 \in P(A,b) \subseteq
  \affhull(P(A,b))$.  Otherwise the line segment between $x_0$ and
  $x_1$ contains more than one point in $P$ and thus $x_0 \in
  \affhull(P)$. 
\end{proof}


\section{Decomposition theorem for polyhedra}





A nonempty set  $C\subseteq\setR^n$ is a \emph{cone} if $\lambda\, x + \mu\,y \in C$
for each $x,y\in C$ and $\lambda,\mu\in \setR_{\geq0}$. A cone  $C$ is \emph{polyhedral}
if $C = \{ x \in \setR^n \mid Ax\leq0\}$. A cone \emph{generated by} vectors
$x_1,\ldots,x_m \in \setR^n$ is a set of the form $C = \{ \sum_{i=1}^m \lambda_i x_i
\mid \lambda_i\in \setR_{\geq0}, \, i=1,\ldots,m\}$.   A point $x = \sum_{i=1}^m \lambda_i x_i$
with  $\lambda_i\in \setR_{\geq0}, \, i=1,\ldots,m$ is called a \emph{conic
  combination}   of the $x_1,\ldots,x_m$. The set of conic combinations of
$X$ is denoted by $\cone(X)$. 


\begin{theorem}[Farkas-Minkowsi-Weyl theorem]
  \label{po:thr:3}
  A convex cone is polyhedral if and only if it is finitely
  generated. 
\end{theorem}


\begin{proof}
  Suppose that $a_1,\ldots,a_m$ span $\setR^n$ and consider the cone $C = \{
  \sum_{i=1}^m \lambda_i a_i \mid \lambda_i\geq0, \, i=1,\ldots,m\}$.
  Let $b \notin C$.
  Then the system $A\lambda = b$, $\lambda \geq 0$ has no solution.
  By Theorem~\ref{conv:thr:12} (Farkas' lemma), this implies that there exists a $y\in\setR^n$
  such that $A^Ty \leq 0$ and $b^Ty > 0$.

  Suppose that the columns of $A$ which correspond to
  inequalities in $A^Ty\leq0$   that are satisfied  by $y$ with equality
  have rank $<n-1$.
  Denote these columns by $a_{i_1},\ldots,a_{i_k}$.  
  Then there exists a $v\neq0$ which is orthogonal to
  each of these columns and to $b$, i.e., $a_{i_j}^Tv = 0$ for each
  $j=1,\ldots,k$ and $b^Tv =0 $. 
  There also exists a column $a^*$ of $A$ which is not in the set
  $\{a_{i_1},\ldots,a_{i_k}\}$ such that $(a^*)^Tv>0$ since the columns of
  $A$ span $\setR^n$. Therefore there exists an $\epsilon>0$ such that 
  \begin{enumerate}[i)]
  \item $A^T(y + \epsilon \cdot v)\leq0$ 
  \item The subspace generated by the columns of $A$ which correspond
    to inequalities of $A^Tx\leq0$ which are satisfied by $y + \epsilon \cdot v$
    with equality strictly contains $\langle a_{i_1},\ldots,a_{i_k}\rangle$. 
  \end{enumerate}
  
  Notice that we have $b^Ty = b^T(y + \epsilon \cdot v)>0$. 

  Continuing this way, we obtain a solution of the form $y + u$ of
  $A^Tx\leq0$ such that one has $n-1$ linearly independent columns of $A$
  whose corresponding inequality in $A^Tx\leq0$ are satisfied with
  equality.   Thus we see that each $b$ which does
  not belong to $C$ can be separated from $C$ with an inequality of
  the form $c^Tx\leq0$  which
  is uniquely defined by $n-1$ linearly independent vectors from the set
  $a_1,\ldots,a_m$.  This shows that $C$ is polyhedral. 

  Suppose now that $a_1,\ldots,a_m$ do not span $\setR^n$. Then there  exist
  linearly independent vectors $d_1,\ldots,d_k$ such that each $d_i$ is
  orthogonal to each of the $a_1,\ldots,a_m$ and $a_1,\ldots,a_m,d_1,\ldots,d_k$
  spans $\setR^n$.   The cone generated by
  $a_1,\ldots,a_m,d_1,\ldots,d_k$ is polyhedral and thus of the form $Ax\leq0$
  with some matrix $A\in \setR^{m\times n}$. Suppose that  $\langle a_1,\ldots,a_m\rangle = \{x
  \in \setR^n \mid Ux = 0\}$.  Now $C = \{ x \in \setR^n \mid Ax\leq0, \, Ux = 0\}$
  and $C$ is polyhedral. 


  Now suppose that $C = \{ x \in \setR^n \mid a_1^Tx\leq0,\ldots,a_m^Tx\leq0\}$.
  The cone
  \[ C' := \cone(a_1,\ldots,a_m) = \{ \sum_{i=1}^m \lambda_i a_i \mid
  \lambda_i\geq0,\,i=1,\ldots,m\} \]
  is polyhedral and thus of the form $C' = \{ x
  \in \setR^n \mid b_1^Tx\leq0, \ldots,b_k^Tx\leq0\}$. Clearly,
  $\cone(b_1,\ldots,b_k)\subseteq C$  since $b_i^Ta_j\leq0$. Suppose now that $y \in
  C \textbackslash{} \cone(b_1,\ldots,b_k)$. Then, since $\cone(b_1,\ldots,b_k)$ is
  polyhedral, there exists a $w\in \setR^n$ with $w^Ty>0$ and $w^Tb_i\leq0$
  for each $i=1,\ldots,k$. From the latter we conclude that $w \in
  C'$. From $y \in C$ and $w \in C'$ we conclude $w^Ty\leq0$, which is a contradiction.
\end{proof}

A set of vectors $Q = \conv(X)$, where $X\subseteq\setR^n$ is finite is called a
\emph{polytope}. 

\begin{theorem}[Decomposition theorem for polyhedra]
  \label{po:thr:4}
  A set $P\subseteq\setR^n$ is a polyhedron if and only if $P = Q + C$ for some
  polytope $Q$ and a polyhedral cone $C$. 
\end{theorem}

\begin{proof}
  Suppose $P = \{ x \in \setR^n \mid Ax\leq b\}$ is a polyhedron. Consider the
  polyhedral cone 
  \begin{equation}
    \label{po:eq:6}
    \left\{\mat{x\\\lambda} \mid x \in \setR^n, \, \lambda \in \setR_{\geq0}; Ax -
      \lambda b\leq0\right\} 
  \end{equation}
  is generated by finitely many vectors $\mat{x_i\\\lambda_i}$,
  $i=1,\ldots,m$. By scaling with a positive number we may assume that
  each $\lambda_i\in \{0,1\}$.  Let $Q$ be the convex hull of the $x_i$ with
  $\lambda_i=1$ and let $C$ be the cone generated by the $x_i$ with
  $\lambda_i=0$. A point $x \in \setR^n$ is in $P$ if and only if $\mat{x\\1}$
  belongs to~\eqref{po:eq:6} and thus if and only if 
  \begin{displaymath}
    \mat{x\\1} \in
    \cone\left\{\mat{x_1\\\lambda_1},\ldots,\mat{x_m\\\lambda_m}\right\}. 
  \end{displaymath}
  Therefore $P = Q + C$. 


  Suppose now that $P = Q+C$ for some polytope $Q$ and a polyhedral
  cone $C$ with $Q = \conv(x_1,\ldots,x_m)$ and $C = \cone(y_1,\ldots,y_t)$. A
  vector $x_0$ is in $P$ if and only if 
  \begin{equation}
    \label{po:eq:7}
    \mat{x_0\\1} \in \cone\left\{\mat{x_1\\1},\ldots, \mat{x_m\\1},\mat{y_1\\0},\ldots,\mat{y_t\\0}     \right\}
  \end{equation}
By Theorem~\ref{po:thr:3} \eqref{po:eq:7} is equal to 
\begin{equation}
\label{po:eq:12}
  \left\{ \mat{x\\\lambda} \mid Ax - \lambda b \leq0\right\}
  \end{equation}
for some matrix $A$ and vector $b$.  Thus $x_0\in P$  if and only if
$Ax_0\leq b$ and thus $P$ is a polyhedron.

\end{proof}




\begin{figure}[htbp]
  \begin{center}
    %\resizebox{4cm}{4cm}{%
    \includegraphics{figures/PicPolyhedra1.pdf}
  %}
  \end{center}
  \caption{A polyhedron and its decomposition into $Q$ and $C$\label{po:fig:decomp}}
\end{figure}




Let $P = \{ x \in \setR^n \mid Ax\leq b\}$. The \emph{characteristic cone} is 
$\charcone(P) =\{ y \mid y+x \in P \text{ for all } x \in P\} = \{y \mid Ay
\leq0\}$. One has
\begin{enumerate}[i)]
\item $y \in \charcone(P)$ if and only if there exists an $x \in P$ such
  that $x + \lambda\,y \in P$ for all $\lambda\geq0$ 
\item $P + \charcone(P)=P$
\item $P$ is bounded if and only if $\charcone(P)=\{0\}$. 
\item If the decomposition of $P$ is $P = Q +C$, then $C = \charcone(P)$. 
\end{enumerate}



The \emph{lineality space} of $P$ is defined as $\charcone(P) \cap -
\charcone(P)$. A polyhedron is \emph{pointed}, if its lineality space is
$\{0\}$.






\section{Faces}
\label{po:sec:faces}

An inequality $c^Tx\leq\delta$    is called \emph{valid} for $P$ if each $x
\in P$ satisfies $c^Tx\leq\delta$. If in addition $(c^Tx = \delta) \cap P \neq\emptyset$,
then $c^Tx\leq\delta$ is a \emph{supporting inequality} and $c^Tx = \delta$ is a
supporting hyperplane. 

A set $F\subseteq\setR^n$ is called a \emph{face} of $P$ if there exists a
valid inequality $c^Tx\leq\delta$ for $P$ with $F = P \cap (c^Tx = \delta)$. 


\begin{lemma}
  \label{po:lem:4}
  A set $\emptyset \neq F \subseteq \setR^n$ is a face of $P$
  if and only if $F = \{ x \in P \mid A'x = b'\}$ for a subset $A'x\leq b'$ of $Ax\leq b$.
\end{lemma}

\begin{proof}
  Suppose that $F = \{ x \in P \mid A'x = b'\}$. Consider the vector $c =
  1^TA'$ and $\delta = 1^Tb'$. The inequality $c^Tx\leq \delta$ is valid for
  $P$. It is satisfied with equality by each $x \in F$. If $x' \in P\textbackslash{}
  F$, then there exists an inequality $a^Tx\leq\beta$ of $A'x\leq b'$ such
  that $ a^Tx' < \beta$ and consequently $c^Tx'<\delta$. 

  On the other hand, if $c^Tx\leq\delta$ defines the face $F$,
  then by the linear programming duality
  \[ \max\{ c^Tx \mid Ax\leq b \} = \min\{ b^T\lambda \mid A^T\lambda = c, \lambda \geq 0 \} \]
  there exists a $\lambda \in \setR^m_{\geq 0}$ such that $c=\lambda^TA$ and $\delta = \lambda^Tb$.
  Let $A'x\leq b'$ be the
  subsystem of $Ax\leq b$ which corresponds to strictly positive entries
  in $Ax\leq b$. One has $F = \{ x \in P \mid A'x = b'\}$. 
\end{proof}





A \emph{facet} of $P$ is an inclusion-wise maximal face $F$ of $P$
with $F\neq P$.  An inequality $a^Tx\leq\beta$ of $Ax\leq b$ is called
\emph{redundant} if $P(A,b) = P(A',b')$, where $A'x\leq b'$ is the system
stemming from $Ax\leq b$ by deleting $a^Tx\leq\beta$.  A system $Ax\leq b$ is
irredundant if $Ax\leq b$ does not contain a redundant inequality.

\begin{lemma}
  \label{po:lem:1}
  Let $Ax\leq b$ be an irredundant system. 
  Then a set $F\subseteq P$ is a facet if and only if it is
  of the form $F = \{ x \in P \mid a^Tx = \beta\}$ for an
  inequality $a^Tx\leq\beta$ of $A^\leq x\leq b^\leq$. 
\end{lemma}


\begin{proof}
  Let $F$ be a facet of $P$. Then $F = \{x \in P \mid c^Tx\leq\delta\}$ for a valid
  inequality $c^Tx\leq\delta$ of $P$. There exists a $\lambda \in \setR_{\geq0}^m$ with
  $c=\lambda^TA$ and $\delta=\lambda^Tb$.  There exists an inequality  $a^Tx\leq\beta$ of
  $A^\leq x\leq b^\leq$ whose corresponding entry in $\lambda$ is strictly
  positive. Clearly $F\subseteq\{x \in P \mid a^Tx=\beta\}\subset P$. Since $F$ is an
  inclusion-wise maximal face one has $F = \{x \in P \mid a^Tx=\beta\}$. 
  
  Let $F$ be of the form $F = \{ x \in P \mid a^Tx = \beta\}$ for an inequality
  $a^Tx\leq\beta$ of $A^\leq x\leq b^\leq$.  Clearly $F \neq \emptyset$ since the system $Ax\leq b$ is
  irredundant. If $F$ is not a facet, then $F\subseteq F'=\{ x \in P \mid a'^Tx = \beta'\}$
  with another inequality $a'^Tx\leq\beta'$ of $A^\leq x\leq b^\leq$. Let $x^*\in \setR^n$ be a point with
  $a^Tx^*>\beta$ and which satisfies all other inequalities of $Ax\leq b$. Such an $x^*$
  exists, since $Ax\leq b$ is irredundant.  Let $\wt{x}\in P$ with
  $A^\leq\wt{x}<b^\leq$.  There exists   a point $\wb{x}$ on the
  line-segment $\wb{\wt{x}x^*}$ with   $a^T\wb{x}=\beta$.  This point is then
  also in $F'$ and thus $a'^Tx = \beta'$ follows. This shows that
  $a'^Tx^*>\beta'$
  and thus $a^Tx\leq\beta$ can be
  removed from the system.  This is a contradiction to 
   $Ax\leq b$ being irredundant.  
\end{proof}




\begin{lemma}
  \label{po:lem:5}
  A face $F$ of $P(A,b)$ is inclusion-wise minimal if and only if it
  is of the form $F = \{ x \in \setR^n \mid A'x=b'\}$ for some subsystem
    $A'x\leq b'$ of $Ax\leq b$. 
\end{lemma}


\begin{proof}
  Let $F$ be a minimal face of $P$ and let $A'x\leq b'$ a the subsystem
  of inequalities of $Ax\leq b$ with $F = \{ x \in P \mid A'x=b'\}$. Suppose that
  $F \subset \{x\in\setR^n  \mid A'x=b'\}$ and let $x_1 \in \setR^n \textbackslash{} P$ satisfy
  $A'x_1=b'$ and  $x_2 \in F$.  There exists ``a first''  inequality $a^Tx\leq\beta$ of
  $Ax\leq b$ which is ``hit'' by the line-segment  $\wb{x_2x_1}$. Let $x^* =
  \wb{x_2x_1}\cap(a^Tx=\beta)$. Then $x^*\in F$ and thus $F \cap (a^Tx=\beta) \neq
  \emptyset$. But $F \supset F \cap (a^Tx=\beta)$ since $a^Tx \leq\beta$ is not an inequality of
  $A'x\leq b'$. This is a contradiction to the minimality of $F$. 

  
  Suppose that $F$ is a face with $F = \{x \in \setR^n \mid A'x= b'\}
  = \{x \in P \mid A'x=b'\}$ for a subsystem $A'x\leq b'$ of $Ax\leq
  b$. Suppose that there exists a face $\wt{F}$ of $P$ with $\emptyset
  \subset \wt{F}\subset F$. By Lemma~\ref{po:lem:4} $\wt{F} = \{ x \in
  P \mid A'x=b', A^*x=b^*\}$, where $A^*x\leq b^*$ is a sub-system of
  $Ax\leq b$ which contains an inequality $a^Tx\leq\beta$ such that
  there exists an $x_1,x_2 \in F$ with $a^Tx_1<\beta$ and
  $a^Tx_2\leq\beta$.  The line $\ell(x_1,x_2) = \{ x_1 +
  \lambda(x_2-x_1) \mid \lambda \in \setR\}$ is contained in $F$ but
  is not contained in $a^Tx\leq\beta$. This shows that $F$ is not
  contained in $P$ which is a contradiction.
\end{proof}

Exercise~\ref{item:19} asks for a proof of the following corollary. 
\begin{corollary}
  \label{co:8}
  Let $F_1$ and $F_2$ be two inclusion-wise minimal faces of
  $P=\{x\in\setR^n\colon Ax \leq b\}$, then $\dim(F_1) = \dim(F_2)$.
\end{corollary}

We say that a polyhedron contains a line $\ell(x_1,x_2)$ with $x_1
\neq x_2 \in P$ if $\ell(x_1,x_2) = \{ x_1 + \lambda(x_2-x_1) \mid
\lambda \in \setR \}\subseteq P$. A \emph{vertex} of $P$ is a
$0$-dimensional face of $P$. An \emph{edge} of $P$ is a
$1$-dimensional face of $P$.






\begin{example}
Consider a linear program $\min\{c^Tx \colon Ax = b,\, x\geq0\}$. A basic
feasible solution  defined by the basis $B\subseteq\{1,\ldots,n\}$ is a vertex of
the polyhedron $P = \{x \in \setR^n \colon Ax = b,\, x\geq0\}$. This can be seen
as follows. The inequality $a^Tx\geq0$ is valid for $P$, where $a_B =
\mathbf{0}$ and $a_{\overline{B}} = \mathbf{1}$. The inequality is
satisfied with equality by a point $x^* \in P$ if and only if
$x^*_{\overline{B}} = \mathbf{0}$. Since the columns of $A_B$ are
linearly independent, as $B$ is a basis, the unique point which
satisfies $a^Tx\geq0$ with equality is the basic feasible solution  
\end{example}

In exercise~\label{item:18} you are asked to show that the simplex
method can be geometrically interpreted as a walk on the graph $G =
(V,E)$, where $V$ is the set of basic feasible solutions and $uv\in  E$
if and only if $\conv\{u,v\}$ is a $1$-dimensional face of the
polyhedron defined by the linear program. 



\subsection{Integer Programming}
\label{sec:integer-programming}

An \emph{integer program} is a problem of the form 
\begin{displaymath}
  \begin{array}{c}
    \max c^Tx \\
    Ax\leq b \\
    x \in \setZ^n,
  \end{array}
\end{displaymath}
where $A \in \setR^{m\times n}$ and $b \in \setR^m$. 




\begin{figure}[htbp]
  \begin{center}
   \includegraphics{figures/PicPolyhedra2.pdf}
\label{fig:inthull}
  \end{center}
  \caption{This picture illustrates a polyhedron $P$ an objective
    function vector $c$ and optimal points $u,v$ of the integer
    program and the relaxation respectively. }
\end{figure}



The difference to linear
programming is the \emph{integrality constraint} $x \in \setZ^n$. This
powerful constraint  allows to model discrete choices but, at the same
time, makes an integer program much more difficult to solve than a
linear program. In fact one can show that integer programming is
NP-hard, which means that it is \emph{in theory} computationally
intractable. However, integer programming has nowadays become an
important tool to solve difficult industrial optimization problems
efficiently. In this chapter, we characterize some integer programs
which are easy to solve, since the \emph{linear programming
  relaxation} $\max\{c^Tx \colon Ax\leq b\}$ yields already an optimal
integer solution. The following observation is crucial. 

\begin{theorem}
  \label{thr:14}
  Suppose that the optimum solution $x^*$ of the linear programming
  relaxation $\max\{c^Tx \colon Ax\leq b\}$  is integral, i.e., $x^* \in
  \setZ^n$, then $x^*$ is also an optimal solution to the integer
  programming problem $\max\{c^Tx \colon Ax\leq b, \, x \in \setZ^n\}$
\end{theorem}

Before we present an example for the power of integer programming we
recall the definition of an undirected graph. 

\begin{definition}[Undirected graph, matching] 
  An \emph{undirected graph} is a tuple $G = (V,E)$ where $V$ is a
  finite set, called the \emph{vertices} and $E\subseteq\binom{V}{2}$ is the
  set of \emph{edges} of $G$.  A \emph{matching} of $G$ is a subset
  $M\subseteq E$ such that for all $e_1\neq e_2\in M$ one has $e_1\cap e_2 = \emptyset$. 
\end{definition}


We are interested in the solution of the following problem, which is
called \emph{maximum weight matching} problem. Given a graph $G =
(V,E)$ and a weight function $w:E\to\setR$, compute a matching with
maximum weight $w(M) = \sum_{e \in M} w(e)$. 

For a vertex $v \in V$, the set $\delta(v) = \{e \in E \colon v \in e\}$ denotes
the \emph{incident} edges to $v$. 

The maximum weight matching problem  can now be modeled as an integer
program as follows. 
\begin{displaymath}
  \begin{array}{c}
    \max \sum_{e \in E} w(e) x(e) \\
    v \in V: \, \sum_{e \in \delta(v)} x(e)\leq1 \\
    e \in E:\, 0\leq x(e) \\
    x \in \setZ^{|E|}.
  \end{array}
\end{displaymath}

Clearly, if an integer vector $x \in \setZ^n$ satisfies the constraints
above, then this vector is the \emph{incidence vector}  of a matching of
$G$. In other words, the integral solutions to the constraints above
are the vectors $\{\chi^M \colon M \text{ matching of }G\}$, where $\chi^M(e)
= 1$ if $e \in M$ and $\chi^M(e)=0$ otherwise. 

  
\subsection{Integral Polyhedra}


In this section we derive sufficient conditions on an integer program
to be solved easily by an algorithm for linear programming. A central
notion is the one of an integral polyhedron. A rational polyhedron $P$
is called \emph{integral} if each minimal face  of $P$ contains an
integer point.  

\begin{theorem}
  \label{po:thr:5}
  Let $P = \{x \in \setR^n \mid Ax\leq b\}$  be a rational nonempty  polyhedron with
  vertices. $P$ is integral if and only if for all integral vectors $c
  \in \setZ^n$ with $\max\{c^Tx \mid x \in P\}<\infty$ one has $\max\{c^Tx \mid x
  \in P\} \in \setZ$. 
\end{theorem}

\begin{proof}
  Let $P$ be integral and $c \in \setZ^n$ with
  $\max\{c^Tx\mid x\in P\}=\delta<\infty$. Since the face    $F=\{x\in P\mid c^Tx=\delta\}$
  contains an integer point it follows   that $\delta\in\setZ$.  

  On the other hand let $x^*$ be a vertex of $P$ and assume that $x^*(i) \notin
  \setZ$. There exists a subsystem $A'x\leq b'$ of $Ax\leq b$ with $A'\in\setR^{n\times n}$,
  $A'$ nonsingular and $A'x^*=b'$. Let $a_1,\ldots,a_n$ be the rows of
  $A'$. Since $A'$ is invertible, there exists an integer vector $c \in
  \cone(a_1,\ldots,a_n)\cap\setZ^n$ such that $c\pm e_i \in
  \cone(a_1,\ldots,a_n)$. The point $x^*$ maximizes both $c^Tx$ and
  $(c+e_i)^Tx$. Clearly not both numbers $c^Tx^*$ and
  $(c+e_i)^Tx^*$ can be integral, which is a contradiction. 
\end{proof}



\begin{lemma}
  \label{po:lem:6}
  Let $A\in \setZ^{n\times n}$ be an integral and invertible matrix. One has
  $A^{-1}b \in \setZ^n$ for each $b \in \setZ^n$ if and only if $\det(A)=\pm 1$.
\end{lemma}


\begin{proof}
  Recall Cramer's rule which says $A^{-1} = 1/\det(A) \wt{A}$, where
  $\wt{A}$ is the adjoint matrix of $A$. Clearly $\wt{A}$ is
  integral. If $\det(A) = \pm 1$, then $A^{-1}$ is an integer matrix. 

  If $A^{-1}b$ is integral for each $b \in \setZ^n$, then $A^{-1}$ is an
  integer matrix. We have $1=\det(A\cdot A^{-1})=\det(A)\cdot\det(A^{-1})$.
  Since $A$ and $A^{-1}$ are integral it follows that $\det(A)$ and
  $\det(A^{-1})$ are integers. The only divisors of one in the integers
  are $\pm 1$. 
\end{proof}



A matrix $A \in \setZ^{m\times n}$ with $m\leq n$ is called \emph{unimodular} if
each $m\times m$ sub-matrix has determinant $0,\pm1$. 

\begin{theorem}
  \label{po:thr:14}
  Let $A \in \setZ^{m\times n}$ be an integral matrix of full row-rank. The
  polyhedron defined by $Ax=b, x\geq0$ is integral for each $b \in \setZ^m$
  if and only if $A$ is unimodular. 
\end{theorem}


\begin{proof}
  Suppose that $A$ is unimodular and $b$ is integral. The polyhedron
  $P = \{ x \in \setR^n \mid Ax=b, \, x\geq0\}$ does not contain a line and
  thus has vertices. A vertex $x^*$ is of the form  $x^*_B= A_B^{-1}b$
  and $x^*_{\wb{B}}=0$, where $B\subseteq\{1,\ldots,n\}$ is a basis. Since $A_B$
  is unimodular one has $x^* \in \setZ^n$. 

  If $A$ is not unimodular, then there exists a basis $B$ with
  $\det(A_B) \neq \pm1$. By Lemma~\ref{po:lem:6} there exists an integral $b
  \in \setZ^n$ with $(A_B)^{-1}b \notin \setZ^m$. Let $\lambda$ be the maximal
  absolute value of a component of $A_B^{-1}b$. Then $b' =  \lceil\lambda\rceil A_B \mathbf{1}
  +b$ is an integral vector with $ A_B^{-1}b' = \lceil\lambda\rceil \mathbf{1}+
  A_B^{-1}b \geq0$ and $A_B^{-1}b' \notin \setZ^m$.  The
  polyhedron $P = \{ x \in \setR^n \mid Ax = b', \, x\geq0\}$ has thus a fractional
  (non-integer) vertex.   
\end{proof}




An integral matrix $A \in \{0 , \pm 1\}^{m\times n}$ is called \emph{totally
  unimodular} if each of its square sub-matrices has determinant
$0,\pm1$. 

\begin{theorem}[Hoffman-Kruskal Theorem]
  \label{po:thr:16}
  Let $A\in \setZ^{m\times n}$ be an integral matrix. The polyhedron $P = \{ x
  \in \setR^n \mid Ax\leq b, \, x\geq0\}$ is integral for each integral $b \in
  \setZ^m$ if and only if $A$ is totally unimodular. 
\end{theorem}


\begin{proof}
  The polyhedron $P = \{ x  \in \setR^n \mid Ax\leq b, \, x\geq0\}$ is integral
  if and only if the polyhedron $Q = \{ z \in \setR^{n+m} \mid (A|I)z = b,
  \, z\geq0\}$ is integral. The assertion thus follows from
  Theorem~\ref{po:thr:14}. 
\end{proof}


If an integral polyhedron has vertices, then an optimal vertex
solution of a  linear program over this polyhedron is integral.  

\section{Applications of total unimodularity}

\subsubsection{Bipartite matching} 

% An undirected graph $G=(V,E)$ is a tuple, where $V$ is a finite set
% and $E$ is a set of unordered pairs of $V$. The set $V$ is called
% \emph{nodes} and the set $E$ are the \emph{edges} of $G$. We write
% $uv$ in short for the edge $\{u,v\}\subseteq V$. 

A graph is
\emph{bipartite}, if $V$ has a partition into sets $A$ and $B$ such
that each  edge $uv$ satisfies $u\in A$ and $v \in B$.
Recall that $\delta(v)$ is the set of edges incident to the vertex $v\in V$,
that is $\delta(v) = \{ e\in E \mid v\in e \}$.

A \emph{matching} of $G$ is a subset $M\subseteq E$ such that $e_1\cap e_2 = \emptyset$
holds for each $e_1\neq e_2\in M$. Let $c:E\longrightarrow\setR$ be a weight function. The
weight of a matching is defined as $c(M) = \sum_{e \in M}c(e)$.  The
\emph{weighted matching problem } is  defined as follows. Given a
graph $G = (V,E)$ and edge-weights $c:E\longrightarrow\setR$, compute a matching $M$
of $G$ with $c(M)$ maximal.


%We have already seen how to compute a maximum weight matching of a
%bipartite graph in polynomial time with a polynomial minimum cost
%network flow algorithm. We now take a different perspective. 
We now define
an \emph{integer program} for this problem and show that, for
bipartite graphs, an optimal vertex of the corresponding linear
program is integral. 


The idea is as follows. We have decision variables $x(e)$ for each
edge $e \in E$. We want to model the characteristic vectors $\chi^M\in
\{0,1\}^E$  of matchings, where $\chi^M(e)=1$ if  $e\in M$  and $\chi^M(e)=0$
otherwise. This is achieved with the following set of constraints. 
\begin{equation}
  \label{po:eq:3}
  \begin{array}{rcll}
    \sum_{e \in \delta(v)} x(e) &\leq&1, & \forall v \in V \\
    x(e) &\geq&0, & \forall e \in E. 
  \end{array}
\end{equation}

Clearly, the set of vectors $x \in \setZ^E$ which satisfy the
system~\eqref{po:eq:3} are exactly the characteristic vectors of
matchings of $G$. The matrix $A \in \{0,1\}^{V\times E}$ which is defined as 
\begin{displaymath}
  \label{po:sec:bipartite-matching}
  A(v,e) = 
  \begin{cases}
    1 & \text{ if } v \in e,\\
    0 & \text{ otherwise}
  \end{cases}
\end{displaymath}
is called \emph{node-edge incidence matrix} of $G$. 

\begin{lemma}
  \label{po:lem:9}
  If $G$ is bipartite, the node-edge incidence matrix of $G$ is
  totally unimodular. 
\end{lemma}

Lemma~\ref{po:lem:9} implies that each vertex of the polytope $P$ defined
by the inequalities~\eqref{po:eq:3} is integral. Thus an optimal vertex
of the linear program $\max\{c^Tx \mid x \in P\}$ corresponds to a
maximum weight matching. 



\begin{proof}[Proof of Lemma~\ref{po:lem:9}]
  Let  $G = (V,E)$ be a bipartite graph with bi-partition $V=V_1\cup V_2$. 
  
  Let $A'$ be a $k\times k$ sub-matrix of $A$. We are interested in the
  determinant of $A$. Clearly, we can assume that $A$ does not contain
  a column which contains only one $1$, since we simply consider the
  sub-matrix $A''$ of $A'$, which emerges from developing the
  determinant of $A'$ along this column. The determinant of $A'$ would
  be $\pm1\cdot \det(A'')$.
  
  Thus we can assume that each column contains exactly two ones. Now
  we can order the rows of $A'$ such that the first rows correspond to
  vertices of $V_1$ and then follow the rows corresponding to vertices
  in $V_2$. This re-ordering only affects the sign of the
  determinant. By summing up the rows of $A'$ in $V_1$ we obtain
  exactly the same row-vector as we get by summing up the rows of $A'$
  corresponding to $V_2$. This shows that $\det(A')=0$. 
\end{proof}



\subsubsection{Flows}

Let $G=(V,A)$ be a directed graph, see
chapter~\ref{cha:short-paths-graphs}.
The \emph{node-edge incidence matrix of a directed graph}
is a matrix $A \in \{0,\pm1\}^{V\times E}$ with 
\begin{equation}
  \label{po:eq:8}
  A(v,a) = 
  \begin{cases}
    1 & \text{ if } v \text{ is the starting-node of } a, \\
    -1 & \text{ if } v \text{ is the end-node of } a, \\
    0  & \text{ otherwise.}
  \end{cases}
\end{equation}


A \emph{feasible flow} $f$  of $G$ with capacities $u$ and in-out-flow $b$ is then
a solution $f \in \setR^A$ to the system $A\,f=b, \, 0\leq f\leq u$.  

\begin{lemma}
  \label{po:lem:10}
  The node-edge incidence matrix $A$ of a directed graph  is totally
  unimodular.  
\end{lemma}


\begin{proof}
  Let $A'$ be a $k\times k$ sub-matrix of $A$. Again, we can assume that in
  each column we have exactly one $1$ and one $-1$. Otherwise, we
  develop the determinant along a column which does not have this
  property.  But then, the $A'$ is singular, since adding up all rows
  of $A'$ yields the $0$-vector. 
\end{proof}
  

A consequence is that, if the $b$-vector and the capacities $u$ are
integral and an optimal flow exists, then there exists an integer
optimal flow. 
%We have seen that this follows from the cycle-cancelling
%algorithm, but total unimodularity gives another simple and elegant
%proof of this fact.  




\subsection{Further applications of polyhedral theory}
\label{po:sec:furth-appl-polyh}


\subsubsection{Doubly stochastic matrices}


A matrix $A \in \setR^{n\times n}$ is \emph{doubly stochastic} if it satisfies
the following linear constraints 
\begin{equation}
  \label{po:eq:9}
  \begin{array}{rcll}
    \sum_{i=1}^n A(i,j) & = & 1, & \forall j=1,\ldots,n\\
    \sum_{j=1}^n A(i,j) & = & 1, & \forall i=1,\ldots,n\\
    A(i,j)       & \geq & 0, & \forall 1 \leq i,j\leq n.
  \end{array}
\end{equation}

A permutation matrix is a matrix which contains exactly one $1$ per
row and column, where the other entries are all $0$. 

\begin{theorem}
  \label{po:thr:17}
  A matrix $A \in \setR^{n\times n}$ is doubly stochastic if and only if $A$ is
  a  convex combination  of permutation matrices. 
\end{theorem}

\begin{proof}
  Since a permutation matrix satisfies the constraints~\eqref{po:eq:9},
  then so does a convex combination of these constraints. 


  On the other hand it is enough to show that each vertex of the
  polytope defined by the system~\eqref{po:eq:9} is integral and thus a
  permutation matrix. However, the matrix defining the
  system~\eqref{po:eq:9}  is the node-edge incidence matrix of the
  complete bipartite graph having $2n$ vertices. Since such a matrix
  is totally unimodular, the theorem follows. 
\end{proof}





\subsection{The matching polytope}
\label{po:sec:matching-polytope}


We now come to a deeper theorem concerning the convex hull of
matchings. We mentioned several times in the course that the maximum
weight matching problem can be solved in polynomial time. We are now
going to show a theorem of Edmonds~\cite{Edmonds65b} which provides a
complete description of the matching polytope and present the
proof by Lov\'asz~\cite{Lovasz79}. 

Before we proceed let us inspect the symmetric difference $M_1\Delta M_2$
of two matchings of a graph $G$. If a vertex is adjacent to two edges
of $M_1\cup M_2$, then one of the two edges belongs to
$M_1$ and one belongs to $M_2$. Also, a vertex can never be adjacent
to three edges in $M_1 \cup M_2$. Edges which are both in $M_1$ and $M_2$
do not appear in the symmetric difference. We therefore have the
following lemma. 

\begin{lemma}
  \label{po:lem:11}
  The symmetric difference $M_1\Delta M_2$ of two matchings decomposes
  into node-disjoint  paths and cycles, where the edges on these paths
  and cycles alternate between $M_1$ and $M_2$. 
\end{lemma}



The \emph{Matching polytope} $P(G)$ of an undirected graph $G = (V,E)$
is the convex hull of incidence vectors  $\chi^M$ of matchings $M$ of
$G$. 


  
\begin{figure}[htbp]
  \centering 
   \includegraphics{figures/PicPolyhedra3.pdf}
  \caption{Triangle}
  \label{po:fig:triangle}
\end{figure}
  

The incidence vectors of matchings are exactly the $0/1$-vectors that
satisfy the following system of equations. 

\begin{equation}
\label{po:eq:10}
  \begin{array}{rcll}
     \sum_{e \in \delta(v)} x(e) & \leq &  1 & \forall v \in V\\
        x(e)& \geq & 0 &    \forall e \in E. 
  \end{array}
\end{equation}

However the triangle (Figure~\ref{po:fig:triangle}) shows that  the
corresponding polytope is not integral. The objective function $\max
\mathrm{1}^Tx$ has value $1.5$. However, one can show that a maximum
weight matching of an undirected graph can be computed in polynomial
time which is a result of Edmonds~\cite{Edmonds65}. 


The following (Figure~\ref{po:fig:edmonds}) is an illustration of an
Edmonds inequality. Suppose that $U$ is an odd subset of the nodes $V$
of $G$ and let $M$ be a matching of $G$. The number of edges of $M$
with both endpoints in $U$ is bounded from above by $\lfloor|U|/2\rfloor$. 

Thus the following inequality is valid for the integer points of the
polyhedron defined by~\eqref{po:eq:10}. 

\begin{equation}
  \label{po:eq:11}
  \sum_{e \in E(U)} x(e) \leq\lfloor|U|/2\rfloor,\quad \quad \text{ for each } U\subseteq V,
  \quad |U| \equiv 1 \pmod{2}. 
\end{equation}


\begin{figure}[htbp]    
  \begin{center}
 \includegraphics{figures/PicPolyhedra4.pdf}
\end{center}
\caption{Edmonds inequality.}
  \label{po:fig:edmonds}
\end{figure}


The goal of this lecture is a proof of the following theorem. 

\begin{theorem}[Edmonds 65]
\label{po:thr:18}
  The matching polytope is described by the following inequalities:
  \begin{enumerate}[i)]
  \item $x(e) \geq0$ for each $e \in E$,
  \item $\sum_{e \in \delta(v)} x(e) \leq 1$ for each $v \in V$,
  \item $\sum_{e \in E(U)} x(e) \leq \lfloor |U|  /2 \rfloor$ for each $U\subseteq V$
    %\fromSlide*{2}{ Gomory Cut!} 
  \end{enumerate}
\end{theorem}


\begin{lemma}
\label{po:lem:11}
  Let $G=(V,E)$ be connected and 
  let $w:E\longrightarrow\setR_{>0}$ be a weight-function.  Denote the set of maximum
  weight matchings of $G$ w.r.t. $w$ by $\eM(w)$. Then one of the following statements
  must be true:
  \begin{enumerate}[i)]
  \item $\exists\,v \in V$ such that $\delta(v) \cap M \neq \emptyset$ for each $M \in \eM(w)$
  \item $|M| =   \lfloor|V| /2\rfloor$ for each $M \in \eM(w)$ and $|V|$ is odd.
  \end{enumerate}
\end{lemma}


\begin{proof}
Suppose both $i)$ and $ii)$ do not hold.
Then there exists ${ M}\in \eM(w)$ leaving two exposed nodes $u$ and
$v$.  Choose $ M$ such that the  minimum  distance between   two exposed nodes
  $u,v$ is  minimized. 


Now let $t$ be on shortest path from $u$ to $v$. The vertex $t$ cannot
be exposed. 
\begin{figure}[htbp]
  \centering
    \begin{center}    
    \includegraphics{figures/PicPolyhedra5.pdf}
  \end{center}
  \caption{Shortest path between $u$ and $v$. }
  \label{po:fig:short}
\end{figure}


 Let ${ M'} \in \eM(w)$ leave $t$ exposed. 
 Both $u$ and $v$ are covered by  ${ M'}$ because the distance to
 $u$ or $v$ from $t$ is smaller than the distance of $u$ to $v$. 
 
 Consider the symmetric difference ${ M} \triangle { M'}$ which  decomposes into
 node disjoint paths and   cycles. 
The  nodes $u, \, v$ and $t$ have degree one in ${M}\triangle{ M'}$. Let 
$P$ be a  path with endpoint $t$ in ${ M}\triangle{ M'}$


\begin{figure}
  \centering
    
\includegraphics{figures/PicPolyhedra6.pdf}
\caption{Swapping colors. }\label{fig:2}
\end{figure}




 If we swap colors on $P$, see Figure~\ref{fig:2}, we obtain matchings  ${\wt{M}}$ and
 ${\wt{M'}}$ with 
 $w({ M}) + w({ M'}) = w({ \wt{M}})+w({  \wt{M'}}) $ and thus
 ${ \wt{M}} \in \eM(w)$.  

 The node $t$ is exposed in ${\wt{M}}$ and $u$ or $v$ is exposed  in
  ${\wt{M}}$. This is a  
contradiction to $u$ and $v$ being shortest distance exposed
  vertices 


\end{proof}



\begin{proof}[Proof of Theorem~\ref{po:thr:18}]

Let $w^Tx\leq\beta$ be a \emph{facet} of $P(G)$, we need to show that this
facet it is of the form
  \begin{enumerate}[i)]
  \item $x(e)\geq0$ for some $e \in E$\label{po:item:1}
  \item $\sum_{e \in \delta(v)} x(e)\leq1$  for some $v \in V$\label{po:item:2}
  \item $\sum_{e \in E(U)} x(e)\leq \lfloor|U|/2\rfloor$ for some $U \in P_{odd}$\label{po:item:3}
  \end{enumerate}
  
  
  To do so, we use the following method: One of the inequalities
  \ref{po:item:1}), \ref{po:item:2}), \ref{po:item:3}) is satisfied with
  equality by each $\chi^M, \,M \in \eM(w)$. This establishes the claim
  since the matching polytope is full-dimensional and a facet is a
  maximal face. 

  




  If $w(e)<0$ for some $e \in E$, then each $M \in \eM(w)$
  satisfies $e \notin M$ and thus satisfies $x(e)\geq0$ with equality. 

  Thus we can assume that  $w\geq0$. 
  
  Let $G^*=(V^*,E^*)$ be the graph induced by edges $e$ with $w(e)>0$.  Each $M
  \in \eM(w)$ contains maximum weight matching $M^* = M \cap E^*$ of
  $G^*$ w.r.t.  $w^*$. 

  If $G^*$ is not \emph{connected }, suppose that
  $V^*=V_1\cup V_2$, where $V_1\cap V_2 = \emptyset$ and $V_1,V_2 \neq\emptyset$ and there
  is no edge connecting $V_1$ and $V_2$, then
  $w^Tx\leq\beta$ can be written as the sum of $w_1^Tx\leq\beta_1$ and
  $w_2^Tx\leq\beta_2$, where $\beta_i$ is the maximum weight of a matching in
  $V_i$ w.r.t. $w_i$, $i=1,2$, see Figure~\ref{blobs}. This would also contradict the fact
  that $w^Tx\leq\beta$ is a facet, since it would follow from the previous
  inequalities and thus would  be a redundant
  inequality.
    
  \begin{figure}
    \centering
  
        
    \includegraphics{figures/PicPolyhedra7.pdf}
    \caption{$G^*$ is connected. }\label{blobs}
  \end{figure}
   
      Now we can use Lemma~\ref{po:lem:11} for $G^*$. 
      
      \begin{enumerate}[i)]
      \item $\exists v$ such that $\delta(v) \cap M = \emptyset$ for each $M \in
        \eM(w)$. This means that each $M$ in $\eM(w)$ satisfies
        \begin{displaymath}
          \sum_{e \in          \delta(v)} x(e)\leq1 \quad \text{ {with equality}}
         \end{displaymath}
      \item $|M\cap E^*| =   \lfloor|V^*| /2\rfloor$ for each $M \in \eM(w)$ and $|V^*|$
        is odd. This means that each $M$ in $\eM(w)$ satisfies
        \begin{displaymath}
           \sum_{e \in E(V^*)          } x(e)\leq \lfloor|V^*|/2\rfloor \quad \text{
             {with equality}} 
        \end{displaymath}
      \end{enumerate}

    \end{proof}


\subsection*{Exercises}

\begin{enumerate}
\item 
  Each nonempty  polyhedron $P\subseteq\setR^n$ can be represented as $ P = L + Q$,
  where  $L\subseteq\setR^n$ is a linear space and $Q\subseteq\setR^n$ is a pointed
  polyhedron.   \label{po:ex:1}
\item Let $P\subset\setR^n$ be a polytope and $f:\setR^n\to\setR^m$ a linear map.
  \begin{enumerate}[i)]
  \item Show that $f(P)$ is a polytope.
  \item Let $y\in\setR^m$ be a vertex of $f(P)$. Show that there is a vertex $x\in\setR^n$ of $P$
    such that $f(x) = y$.
  \end{enumerate}
\item 
  Let $A \in \setR^{m\times n}$ and $b \in \setR^m$ and consider the polyhedron
  $P = P(A,b)$. Show that $\dim(P) = n - \rank(A^=)$.   \label{po:ex:2}
\item 
  \begin{enumerate}[i)]
  \item 
  Show that the dimension of each minimal face of a polyhedron $P$ is
  equal to $n - \rank(A)$. 
  \item
  Show that a polyhedron has a vertex if and only if the polyhedron
  does not contain a line. 
\end{enumerate}   \label{po:ex:3}
\item Show that the affine  dimension of the minimal faces of a 
  polyhedron $P = \{x \in \setR^n \colon Ax\leq b\}$ is invariant. \label{item:19}
\item 
  In this exercise you can assume that a linear program $\max\{c^Tx \mid
  Ax\leq b \}$ can be solved in polynomial time. Suppose that $P(A,b)$
  has vertices and that the linear program is bounded. Show how to
  compute an optimal \emph{vertex} solution of the linear
  program in polynomial time.    \label{po:ex:4}
\item  Let $P = \{x \in \setR^n \colon Ax = b, \, x\geq0\}$ be a polyhedron,
  where $A \in \setR^{m\times n}$ has full row-rank. Let $B_1,B_2$ be two bases
  such that $|B_1\cap B_2| = m-1$ and suppose that the associated basic
  solutions $x^*_1$ and $x^*_2$ are feasible. Show that, if
  $x_1\neq x_2$, then   $\conv\{x_1^*,x_2^*\}$ is a $1$-dimensional face
  of $P$. \label{item:18} 
\end{enumerate}





\bibliographystyle{abbrv}

\bibliography{mybib,papers,books,my_publications}



%%% Local Variables: 
%%% mode: latex
%%% TeX-master: "lecture"
%%% End: 

%\chapter{The simplex method}
\label{cha:introduction}


In this chapter, we describe the simplex method. The task is to solve
a linear program 
\begin{equation}
  \label{eq:28}
  \max\{c^Tx \colon x \in \setR^n, \, Ax\leq b\}
\end{equation}
where $A \in \setR^{m\times n}$, $b \in \setR^m$ and $c \in \setR^n$. We make the
following assumption.

\dproblem{Full-rank assumption}{The matrix $A \in \setR^{m\times n}$ has full
  column-rank. In other words, the columns of $A$ are linearly
  independent}  
%
We will see later that this assumption can be made without loss of
generality. 

\section{Roofs}
\label{sec:two-variable-linear}


\begin{figure}[htbp]
  \begin{center}{
%    \psset{unit=.8cm}
    \begin{pspicture}(-4,-3)(4,4)%\showgrid
        \pspolygon[fillcolor=vlg,fillstyle=solid](-1,-2)(-2,1)(-1,2)(2,1)(1,-1)
        \psline(-4,1)(0,4)
        \psline(-2,4)(4,0)
%        \psline(-4,2)(4,4)
        \psline(-4,-1)(1,4)
        \psline(-4,3)(4,0.333)
%      \psline{->}(0,-3)(0,3)
%      \psline{->}(-3,0)(3,0)
%      \rput(-.5,-.5){$(0,0)$}
     % \psdot(2,2)
     % \psdot(1,-1)
     % \psdot(-1,-2)
     % \psdot(-.5,-1.2)
     % \psdot(.5,1.5)
     % \psdot(-2,1)
     % \psdot(-1,2)
        \psline[linecolor=red]{->}(3.5,1)(3.5,3)
        \psdot[linecolor=blue](-1,2)
        \psdot[linecolor=blue](-.95,3.3)
        \psdot[linecolor=green](3,.65)
        \rput(3,2){\red{$c$}}
    \end{pspicture}
    }
    
  \end{center}
  \caption{A linear program; the objective function vector $c$ is
    pointing vertically upwards. The blue dots mark two roofs. Notice
    that the lowest roof is the optimum of the linear program. The
    green point marks a non-roof. The two constraints
    satisfy~\ref{xitem:5}) and~\ref{xitem:6}) but not~\ref{xitem:7}). }
\end{figure}


Roofs are linear programs originating from~\eqref{eq:28} by selecting
a subset of the inequalities only. A roof should provide an upper
bound on the optimal value of the linear program~\eqref{eq:28} and at
the same time consist of $n$ ``linearly independent
constraints''. Here is the definition of a roof. 



\begin{definition}
  Consider the linear program~\eqref{eq:28} and let $B\subseteq\{1,\ldots,m\}$ be
  a subset of the row-indices. This set $B$ is a \emph{roof} if 
  \begin{enumerate}[i)]
  \item $|B| = n$, \label{xitem:5}
  \item The rows $a_i, \, i \in B$ are linearly independent, and  \label{xitem:6}
  \item The linear program  \label{xitem:7}
    \begin{equation}
      \label{eq:29}
      \max\{c^Tx \colon a_i^Tx \leq b(i), \, i \in B\}
    \end{equation}
    is bounded. 
  \end{enumerate}

\end{definition}







% \begin{figure}[htbp]
%   \centering
  
%   \epsfig{file=figures/highest_point.eps,height=10cm}
%   \caption{Roofs  and the optimal point}
%   \label{fig:highest_point}
% \end{figure}








\begin{figure}[htbp]
  \begin{center}{
%    \psset{unit=.8cm}
    \begin{pspicture}(-4,1)(4,4)%\showgrid
        \pspolygon[fillcolor=vlg,linecolor=vlg,fillstyle=solid](-4,1)(-1,3.25)(2,1)
        \psline(-4,1)(0,4)
        \psline(-2,4)(2,1)

         \psline[linecolor=green](-1,3.25)(-2.8,1)
        
        \psline[linecolor=red]{->}(3.5,1)(3.5,3)
%        \psdot[linecolor=blue](-1,2)
        \psdot[linecolor=blue](-1,3.25)
        \psdot[linecolor=blue](-2,2)
       % \psdot[linecolor=green](3,.65)
        \rput(-.5,3.25){$x^*_B$}
        \rput(-1.5,2){$y^*$}
       
    \end{pspicture}
    }
    
  \end{center}
  \caption{An illustration for the proof of Lemma~\ref{lem:2}. The
    green ray illustrates the set $\{x^*_B + \lambda( y^* - x^*_B) \colon \lambda \in
  \setR_{\geq0}\}$. }
\end{figure}



\noindent 
What is the optimal solution of a linear program~\eqref{eq:29} defined
by a roof? This question is answered in the next lemma. 

\begin{lemma}
  \label{lem:2}
  Let $B\subseteq\{1,\ldots,m\}$ be a roof of the linear program~\eqref{eq:28}
  and let $x^*_B$ be the unique solution of the linear system 
  \begin{displaymath}
    a_i \,x = b(i), \, i \in B,
  \end{displaymath}
  then $x^*_B$ is an optimal solution of the roof-linear
  program~\eqref{eq:29}.
\end{lemma}


\begin{proof}
  Suppose that $y^*$ is a feasible solution of the roof-linear program
  with $c^Ty^*>c^Tx^*_B$. We now show that each $x_B^* + \lambda (y^* -x_B^*)$
  for $\lambda\geq0$ is feasible. One has $a_i\, (x_B^* + \lambda (y^* -x_B^*)) =
  b(i) + \lambda\cdot ( a_i\,y^* -  b(i) ) \leq b(i)$ for each $i \in B$ and thus
  $x_B^* + \lambda (y^* -x_B^*)$ is feasible for each $\lambda\geq0$. 
  
  The objective function value of such a point  is $c^Tx_B^* +
  \lambda(c^Ty^*-c^Tx^*_B)$ which, for $\lambda \to \infty$ tends to infinity. This
  is a contradiction to $B$ being a roof
  (condition~\ref{xitem:7}). Thus $x^*_B$ must be  an optimal 
  solution to the roof-linear program~\eqref{eq:29}.   \qed
\end{proof}


Now that we know that a roof-linear program has an optimal solution,
we can define the value of a roof $B$.

\begin{definition}
  The \emph{value}  of a roof $B$ is the optimum value $c^Tx^*_B$ of
  the roof-linear program
  \begin{displaymath}
    \max\{c^Tx \colon a_i\,x\leq b(i), \, i \in B\}. 
  \end{displaymath}
\end{definition}



The next theorem is very simple, but in fact very important. It states
that the value of a roof is an upper bound on the optimum value of a
linear program $\max\{c^Tx \colon x \in \setR^n, \, Ax\leq b\}$. 

\begin{theorem}[Weak duality]
  \label{thr:6}
  The value of a roof is an upper bound on the objective function
  value of any feasible point of the linear program.   
\end{theorem}


\begin{proof}
  Let $B$ be a roof of the linear program $\max\{c^Tx \colon x \in \setR^n, \,
  Ax\leq b \}$. Any feasible point $x^*$ of this linear program is also
  a feasible point of the 
  roof-linear program $\max\{c^Tx \colon a_i\,x\leq b(i)\}$. Therefore
  $c^Tx^*_B \geq c^Tx^*$ and the claim follows. \qed 
\end{proof}


When is an index-set $B\subseteq\{1,\ldots,m\}$ satisfying  \ref{xitem:5}) and
\ref{xitem:6}) a roof? Consider the  example in see
Figure~\ref{fig:2}. The objective is 
to maximize $2x_1+x_2$ and the two roof-constraints are $x_1+x_2\leq5$
and $x_1\leq6$. From the picture, it is clear that the objective
function vector is in the cone of the two constraint vectors. In fact,
this is the characterization that holds in any dimension as we now
show. 






\begin{figure}[htbp]
  \begin{center}{
%    \psset{unit=.8cm}
    \begin{pspicture}(-1,-2)(5,4)
      \pspolygon[fillcolor=vlg,linecolor=vlg,fillstyle=solid](-1,4)(2,1)(2,-2)
      \psline(2,-2)(2,4)
      \psline(-1,4)(4,-1)
      \psline[linecolor=red]{->}(2.5,1.5)(4.5,2.5)
      \psline[linecolor=blue]{->}(0,3)(1,4)
      \psline[linecolor=blue]{->}(2,-1)(3,-1)
%      \showgrid
    \end{pspicture}

    \vspace{1cm}
    \begin{pspicture}(0,0)(4,4)
%      \rput(-.2,-.2){$0$}
      \pspolygon[fillcolor=vlg,linecolor=vlg,fillstyle=solid](0,0)(4,0)(4,4)
       \psline[linecolor=red]{->}(0,0)(2,1)
      \psline[linecolor=blue]{->}(0,0)(1,0)
      \psline[linecolor=blue]{->}(0,0)(1,1)
%      \showgrid
    \end{pspicture}
    }
    
  \end{center}

  \caption{A roof that is defined by the two constraints $x_1\leq2$ and
    $x_1+x_2\leq3$. The objective function  vector is $(2,1)^T$. Indeed
    $(1,1)^T = 1 \cdot (1,1)+ 1\cdot (1,0)$ which shows that it is in the
    cone generated by the constraint-defining vectors.  
    }
  \label{fig:2}
\end{figure}






\begin{lemma}
  \label{lem:4}  
  Let $B\subseteq\{1,\ldots,m\}$ satisfy \ref{xitem:5}) and
  \ref{xitem:6}). Then $B$ is a roof, if and only if $c \in \cone\{a_i\colon i \in B\}$. 
\end{lemma}

\begin{proof}
  Suppose that $c \in \cone\{a_i\colon i \in B\}$. Thus there exist
  $\lambda_i\geq0, \, i\in B$ with $c = \sum_{i \in B}\lambda_i\cdot a_i$. 
  The  unique solution $x^*_B$  to the
  system 
  \begin{equation}
    \label{eq:0}
    a_i\,x=b(i), \, i\in B
  \end{equation}
  is an optimal solution to
  $\max\{c^Tx \colon a_i\,x\leq b(i)\}$. Because if $\wt{x}$ is another feasible solution,
  then $c^T\wt{x} = \sum_{i\in B} \lambda_i \cdot a_i\, \wt{x}$. Since $\lambda \geq0$ and
  $a_i\,\wt{x}\leq b(i)$  we can write
  \begin{eqnarray}
    \label{eq:2}
    c^T\wt{x} & = &  \sum_{i\in B} \lambda_i\cdot a_i\,  \wt{x}\\
    & \leq & \sum_{i\in B} \lambda_i\cdot  b(i) \\
    & = &  \sum_{i\in B} \lambda_i\cdot  a_i\,x_B^*\\
    & = &  (\sum_{i\in B} \lambda_i\cdot  a_i )\, x_B^*\\
             & = & c^Tx_B^*.
  \end{eqnarray}
  Thus $B$ is a roof. 

  Suppose on the other hand that $B$ is a roof.  Then, since
  $a_i, i\in B$ is a basis of $\setR^n$, there exist $y_i \in  \setR, \, i \in B$ with 
  $c =  \sum_{i \in B}y_i \cdot a_i$. If all $y_i\geq0, \, i \in B$, then it
  follows that $c \in \cone\{a_i \colon i \in B\}$ and we are done. 
  Suppose therefore   that there exists an index  $j \in B$ with
  $y_{j}<0$. We will derive a contradiction. 

% Without loss of 
  %  generality assume  $y(1)<0$. 
  Consider the system of linear equations
  \begin{equation}
    \label{eq:1}
    a_{j}\,x = -1, a_i\,x=0, \, i \in B\setminus\{j\}. 
  \end{equation}
  This system~\eqref{eq:1} has a unique solution $0\neq v\in \setR^n$.  
  Let $x^*$ be a feasible solution to the roof-linear
  program. Clearly $x^* + \lambda \cdot v$ is also feasible for each 
  $\lambda>0$ (Exercise~\ref{xitem:9}). But $c^T (x^* + \lambda v) = c^Tx^* + \lambda \cdot \sum_{i=1}^n y_i  a_i\,  v = c^T x^*
  + \lambda \cdot y_{j} \cdot  a_{j}\,v = c^T x^*  - \lambda \cdot y_{j} $. This increases with $\lambda$ since $y_{j}<0$ and
  $a_{j}\,v<0$.  This contradicts the fact $B$ is a roof, since the
  roof-linear program is unbounded. \qed 
\end{proof}





\begin{definition}
  \label{def:5}
  Let $B$  be a roof of the linear program~\eqref{eq:28}. The unique
  solution 
  $x^*$  of  the
  system 
  \begin{equation}
    \label{eq:3}
 a_i^Tx = b(i), \, i \in B, 
  \end{equation}
  is the  \emph{vertex} of the roof. 
\end{definition}


% The proof of Lemma~\ref{lem:4} reveals the following fact. 

% \begin{proposition}
% \label{prop:1}
% Let $\ev$ be a v-shape of $\hpp(\eh,c)$. The vertex of $\ev$ is an
% optimal solution to $\hpp(\ev,c)$.
% \end{proposition}


Similarly one can prove the following fact. 

\begin{proposition}
  \label{prop:2}
  Let $B$ be a roof of the linear program~\eqref{eq:28}. 
    The vertex of a roof is the \emph{unique optimal solution} of the
    roof-linear program~\eqref{eq:29} if and only if   
    $c$ is a \emph{strictly positive} conic combination  of 
    the normal-vectors $a_i, i \in B$. 
\end{proposition}





\section{The simplex algorithm}
\label{sec:simplex-algorithm}


We now sketch one iteration of the simplex algorithm. Our task is to
solve a linear program~\eqref{eq:28} and we assume that we have a roof
$B$ to start with.


\begin{enumerate}[i)]
\item Compute the vertex $x^*_B$ of the roof $B$. \label{xitem:13}
\item Find an index $i \in  \{1,\ldots,m\}\setminus B$ with $a_i\,x^*_B>b(i)$. If
  there does not exist such an index, then $x^*_B$ is an optimal
  solution of the linear program~\eqref{eq:28}. \label{xitem:14} 
\item Determine an index $j \in B$ such that \label{xitem:17}
  \begin{enumerate}[a)]
  \item   $B' = B \cup \{i\} \setminus \{j\}$ is a roof, and \label{xitem:15}
  \item The vertex $x^*_{B'}$ of $B'$ is feasible for $B$. \label{xitem:16}
  \end{enumerate}
  If such an index does not exist, then the linear
  program~\eqref{eq:28} is infeasible. 
\end{enumerate}


The simplex algorithm iterates these steps until it has found an
optimal solution, or asserts that the linear program~\eqref{eq:28} is
infeasible. The big questions are how to determine an index $j$ such
that \ref{xitem:15}) and \ref{xitem:16}) hold in step \ref{xitem:17}) and
that the algorithm is correct. Furthermore, we want to understand
whether the simplex method eventually terminates.  







\subsection{Termination and degeneracy}




\begin{definition}[Degenerate roof and linear program]
  \label{def:4}
  A roof $B$ of a linear program~\eqref{eq:28} is \emph{degenerate} if
  the optimum solution of the roof-linear program~\eqref{eq:29} is not
  unique. A linear program is called degenerate, if it has degenerate
  roofs. 
\end{definition}


\begin{figure}[htbp]
  \centering
  \includegraphics[height=3cm]{figures/degnodeg.eps}
  \caption{A non-degenerate and a degenerate roof.}
  \label{fig:non-deg-deg}
\end{figure}



We now argue that the simplex algorithm  terminates if the linear
program is non-degenerate. 

\begin{theorem}
  \label{thr:1}
  If the linear program~\eqref{eq:28} is non-degenerate, then the
  simplex algorithm terminates. 
\end{theorem}



\begin{proof}
  The important observation is that the simplex method makes progress
  from iteration to iteration because of the non-degeneracy of the
  roofs. If $B'$ is the roof computed in step~\ref{xitem:17}), then,
  since $x^*_{B'}$ is contained in the feasible region of the roof
  $B$, and since $B$ is non-degenerate, we have
  $c^Tx^*_B>c^Tx^*_{B'}$. Since there is only a finite number of
  roofs, the algorithm thus terminates. \qed
\end{proof}


\subsection{Implementing step~\ref{xitem:17})}


% \begin{figure}[htbp]
%   \centering
%   \epsfig{file=figures/entering.eps,height=5cm}
%   \caption{On the left: A strictly lower roof. On the right: $\eh$
%     is infeasible}
%   \label{fig:two-poss}
% \end{figure}


% We now describe the procedure for this move. Suppose that $a_1,\ldots,a_n$
% are the normal-vectors of the initial v-shape $\ev$.

% \noindent 
% {\bf Step 1:}  Compute the vertex $x^*$ of $\ev$. 

% \begin{lemma}
%   \label{lem:5}
%   If there exists no halfspace $h \in \eh$ which is violated by $x^*$,
%   then $x^*$ is an optimal solution to the highest point problem. 
% \end{lemma}

% \begin{proof}
%   Easy!
% \end{proof}

% \noindent
% {\bf Step 2:} Find a halfspace $a\,x\leq\beta$ for which $x^*$ is not
% feasible. If such a halfspace does not exist, then $\ev$ is a lowest
% v-shape and we stop.   

The situation is as follows. We are having a roof $B$ and its vertex
$x^*_B$ and an index $i \in \{1,\ldots,m\}$ with $a_i\,x^*_B> b(i)$. We now
want to bring $i$ into the new roof and we have to determine a $j \in
B$ that is supposed to leave the roof. The idea is very similar now to
the proof of Carath\'eodory's theorem. 

Consider the systems of equations 
\begin{eqnarray}
  \sum_{k \in B} a_k z_k   & = &  c^T \label{eq:5}\\
  \sum_{k \in B}  a_k y_k   & = &  -a_i \label{eq:17}
\end{eqnarray}
with variables $z_k, \, k \in B$ and $y_k,\, k \in B$. 

Compute    solutions $z^* \in \setR^{n}$ of~\eqref{eq:5} and  $y^* \in \setR^{n}$
of~\eqref{eq:17}. 
Now we have for any $\lambda \in \setR$ 
\begin{equation}
  \label{eq:49}
   \sum_{k \in B} a_k (z^*_k + \lambda\, y^*_k) + \lambda a_i= c^T
\end{equation}
Notice that $z^*\geq0$. 
To bring the index $i$ into the roof, we want to increase $\lambda = 0$
until some other component of $ z^* + \lambda y^*$, component $j$ lets say, becomes
zero. So in virtue of finding an index which drops out of $B$, we
have to determine the largest $\lambda^* \in \setR_{\geq0}$ such that all components of
$x^* + \lambda\,y^*$ are nonnegative. This is done as follows. 

Compute   the index set $J  = \{ k \in B \colon y^*_k <0\}$.  Those are the 
indices we have to worry about, since only those components 
 can become negative with increasing $\lambda$. Still, how large can
$\lambda^*$ be? We have to ensure that
\begin{equation}
  \label{eq:8}
  z^*(k) +\lambda^* y^*(k) \geq0 \text{ for all } k \in J.
\end{equation}
In other words we have to ensure 
\begin{equation}
  \label{eq:4}
  \lambda^* \leq - \frac{z^*(k)}{y^*(k)} \text{ for all }  k\in J.
\end{equation}



{ If $J \neq \emptyset$, }
we pick 
 \begin{equation}
   \label{eq:9}
   \lambda^* = \min_{\substack{k\in J}}  - \frac{z^*(k)}{y^*(k)}. 
\end{equation}
We choose an index  $j \in J$ for which this minimum is achieved. 
This  index $j$ is the
one which leaves the roof. 

\begin{lemma}
  \label{lem:1}
  The index set $B'=B \backslash \{i\} \cup \{j\}$ is a roof and the new vertex
  $x^*_{B'}$ is contained in the feasible set of the roof $B$. 
\end{lemma}
\begin{proof}
  By construction,  $c$ is a nonnegative linear combination of
  the  vectors $a_k, \, k \in B'$.  Thus in order to conclude that $B'$
  is a roof, we  need to show that the 
  $a_k, \, k \in B' $ are linearly independent. The component 
  $y^*_j$ is nonzero. Since $y^*$ is a solution of
  equation~(\ref{eq:17}) it follows that $a_j$ is a linear combination
  of the normal-vectors of $B'$. Thus the $a_k, \, k \in B'$ are a
  basis of $\setR^n$ and since $|B'|=n$ they are linearly independent. 


  Let $x_B^*$ be the vertex of $B$ and let $w \in \setR^n$ be a solution to
  the system 
  \begin{equation}
    \label{eq:15}
    a_j\,w = -1, a_k\,w=0, \, k \in B\setminus\{j\}. 
  \end{equation}
  The half-line  $l(x_B^*,w) = \{ x^* + \lambda \, w \mid \lambda \in \setR_{\geq0}\}$  is feasible
  for $B$. We have the equation
  \begin{equation}
    \label{eq:16}
    a_i = -\sum_{k \in B} y^*_k a_k
  \end{equation}
  where $y^*_j < 0$.  Thus 
  \begin{eqnarray}
    a_i\, w  & = & - \sum_{k \in B} y^*_k a_k\, w \\
          & = &  y^*_k \\
          & < & 0.
  \end{eqnarray}
  Therefore the hal-fline  $l(x^*,w)$ enters at some point,  $x' \in \setR^n$ say, the
  halfspace $a_i\,x\leq b(i)$. Clearly, this $x'$  is the vertex $x^*_{B'}$ of $B'$. \qed
\end{proof}

 \begin{figure}[htbp]
    \begin{center}
      \includegraphics[height=5cm]{figures/example.eps}
      \caption{The initial roof  of Example~\ref{ex:1}.}
        \label{fig:ex:1.1}
    \end{center}    
  \end{figure}


\begin{example}
  \label{ex:1}
    
  Consider the linear program $\max\{x_2 \colon x \in \setR^n, \, (-1,1) x \leq
  1,\,  (2,1)x\leq1,\, (1,2) x \leq1 \}$. We start with the roof $B =
  \{1,2\}$ that consists of the first two inequalities see 
  Figure~\ref{fig:ex:1.1}. 
  
  \noindent 
  We compute first the  vertex $x^*_B$  which is the solution to the
  system 
  \begin{equation}
    \label{eq:19}
    \mat{-1 & 1 \\ 2 & 1}  x  = \mat{1\\1}.
  \end{equation}
  Thus the  vertex is the vector $x_B^* = \mat{0\\1}$.  
  
  \noindent 
  Next  we find that the halfspace $(1,2) x \leq1$ is not satisfied
  by $x_B^* = \mat{0\\1}$. We want to bring this index   into the new
  roof $B'$.
  
  \noindent 
  Step~3:
  Now we compute the solution  to the system 
  \begin{equation}
    \label{eq:20}
    \mat{-1 & 2 \\ 1 & 1} z = \mat{0\\1}
  \end{equation}
  and  find 
  $$z^* = \mat{2/3 \\ 1/3}.$$  
  
  Next we find a solution  to the system  
   \begin{equation}
     \label{eq:21}
    \mat{-1 & 2 \\ 1 & 1} y  = - \mat{1\\2}
  \end{equation} 
  and find  
  $$y^* = \mat{-1\\ - 1 }. $$
  
  \noindent
  The index set $J = \{1,2\}$ is not empty. The minimum~\eqref{eq:9} is
  achieved at  $j =2$. So the halfspace 
  $(2,1)x\leq1$ will leave the roof and $B' = \{1,3\}$. 
  This is also what we immediately see by
  looking at Figure~\ref{fig:ex:1.2}. 
   
\end{example}

  \begin{figure}[htbp]
    \begin{center}
      
      \includegraphics[height=5cm]{figures/example2.eps} 
      \caption{The new roof from  Example~\ref{ex:1}.} 
      \label{fig:ex:1.2}
    \end{center}
  \end{figure}


\noindent 

{ If $J = \emptyset$, } we assert that the linear program is infeasible based on
the following result. 

\begin{proposition}
  \label{thr:3}
  The half-spaces   $a_k\,x\leq b(k), \,k \in B $ and $a_i\,x\leq b(i)$
  define together an infeasible system 
  if and only   if $J = \emptyset$.
\end{proposition}

\begin{proof}
  
  If $J \neq \emptyset$, then   Lemma~\ref{lem:1} implies that the half-spaces
  $a_k\,x\leq b(k), \,k \in B $ and  $a_i\,x\leq b(i)$ define a feasible system. 

  
  The index set  $J$ is empty if and only if $y^*\geq0$.   
  We now show that, if $y^* \geq 0$, then the half-spaces $a_k\,x\leq b(k), \,k \in B $ and
  $a_i\,x\leq b(i)$ define an infeasible system. 
  Since  $\sum_{k \in B} y^*_k a_k + a_i = 0$  this assertion follows,
  once we have shown that  $\sum_{k \in B} y_k b(k) + b(i) < 0$. But 
  \begin{eqnarray*}
    \sum_{k \in B} y_k b(k) + b(i)  & = & \sum_{k \in B} y^*_k a_k\, x^*_B + b(i)  \\
               & <  & \sum_{k \in B} y^*_k a_k\, x^*_B + a_i\, x^*_B  \\ 
               & = &  \left(\sum_{k \in B} y^*_k a_k + a_i \right)\, x^*_B  \\ 
               & = &  0^T   x_B^* \\
               & = & 0. 
  \end{eqnarray*}
\end{proof}







% This algorithm which we have just described is called
% \emph{pivoting}. 
% Now we can describe the famous simplex algorithm which finds the
% lowest v-shape of a HPP. It simply keeps pivoting until one can
% assert  that  $\eh$ is infeasible or unbounded or until a lowest
% v-shape is found. 

% \begin{algorithm}[Simplex]~\\
%   \begin{tabular}{ll}
%     {\bf Input:}  & non-degenerate $\hpp(\eh,c)$ and a v-shape $\ev$\\
%     {\bf Output:} & lowest v-shape $\ev'$ \\
%     & or assertion that $\eh$ is infeasible 
%   \end{tabular}
  
%   \begin{enumerate}
%   \item $\wt{\ev} \gets \ev$
%   \item {\bf REPEAT}
%     \begin{enumerate}
%     \item Run {\bf Pivot} on $\wt{\ev}$
%     \item If this reveals that $\wt{\ev}$ is lowest v-shape
%       \begin{enumerate}
%       \item {\bf RETURN} $\wt{\ev}$
%       \end{enumerate}
%     \item If this reveals that $\eh$ is infeasible
%       \begin{enumerate}
%       \item {\bf RETURN} $\eh$ is infeasible
%       \end{enumerate}
%     \item If this returns v-shape $\ev'$ \label{xitem:2}
%       \begin{enumerate}
%       \item  $\wt{\ev} \gets \ev'$
%       \end{enumerate}
%     \end{enumerate}
%     \item {\bf UNTIL} 0
%   \end{enumerate}
% \end{algorithm}


% \begin{theorem}
%   \label{thr:4}
%   The simplex algorithm terminates after a finite number of steps on a
%   non-degenerate HPP. 
% \end{theorem}

% \begin{proof}
%   The only interesting step is~\ref{xitem:2}). The vertex of the new
%   v-shape is feasible for the old v-shape by
%   Proposition~\ref{prop:3}. Since the HPP is non-degenerate, this new
%   vertex has a strictly smaller objective function. Since there are
%   only a finite number of v-shapes, the simplex algorithm terminates. 
% \end{proof}

% This gives us also a proof of our strong duality theorem for
% non-degenerate HPP's. 






 \subsection{The degenerate case}
 \label{sec:degenerate-case}

 The termination argument for the non-degenerate case was that the
 value of the new roof  is strictly dropping and thus, that a
 roof can never be revisited. Since there are only a finite number
 of roofs, this implies that the simplex algorithm terminates.

 In the degenerate case, roofs could be revisited. This phenomenon
 is called \emph{cycling} and you are asked to construct such an
 example in the exercises. What can we do about it? The idea is  to change the
 objective vector $c \in \setR^n$ a little bit and turn it into a vector
 $c_\epsilon$ such that the following conditions hold. 

 \begin{enumerate}[1)]
 \item The linear program  
   \begin{equation}
     \label{eq:6}
      \max\{c_\varepsilon^Tx \colon  x \in  \setR^n, \, Ax \leq b\}
   \end{equation}
   has a roof.\label{xitem:19}
 \item Each non-roof of the linear program~\eqref{eq:28}  is a non-roof  of the
   linear program~\eqref{eq:6}. \label{xitem:18}  
 \item No roof of \eqref{eq:6} is degenerate. \label{xitem:20}
 \end{enumerate}


 Suppose we start with  an initial roof 
 $R$ at the beginning of the
 simplex algorithm and let ${A_R}\in \setR^{n\times n}$ be the matrix whose
 rows  are the vectors $a_i, \, i \in R$. 
 Notice that we implicitly assume that the set $R$ is ordered. We
 adhere to the following notation. For $i \in R$, the function $f_R(i)$
 denotes the position of $i$ in the ordered set $R$. For example, if
 $R = \{5,2,9\}$, then $f_R(5)=1$ and $f_R(9)=3$. 

 The system ${A_R}^T \,y =c$
 has a solution ${y^*}\geq0$, where some components of ${y^*}$ are zero
 if and only if $R$ is degenerate. 
  This is undesirable and we wish that ${y^*}$ is replaced by
 \begin{equation}
   \label{eq:7}
   {y^*} +  \mat{\epsilon \\ \epsilon^2 \\ \vdots \\\epsilon^n}
 \end{equation}
 for some $\epsilon>0$. Later it will
 become clear why we add the vector $(\epsilon,\ldots,\epsilon^n)^T$ instead of the vector
 $(\epsilon,\ldots,\epsilon)^T$. Now the vector~\eqref{eq:7} becomes a solution if we
 perturb $c$ and consider the vector 
 \begin{equation}
   \label{eq:18}
   c_\epsilon = c + {A_R}^T \mat{\epsilon \\ \epsilon^2 \\ \vdots \\\epsilon^n}
 \end{equation}
 instead, see figure~\ref{fig:perturb}. 
 \begin{figure}[h]
     \begin{center}{


    \psset{unit=.8cm}
    \begin{pspicture}(-1,-2)(5,4)
%      \showgrid
      \pspolygon[fillcolor=vlg,linecolor=vlg,fillstyle=solid](-1,4)(2,1)(2,-2)
      \rput(0,1){Roof $R$}
      \psline(2,-2)(2,1)
      \psline(-1,4)(2,1)
      \psline[linecolor=red,linewidth=1.5pt]{->}(2.5,1.5)(4.5,1.5)
      \rput(4.3,1.2){\red{$c$}}

      {
      \psline[linecolor=magenta,linewidth=1.5pt]{->}(2.5,1.5)(4.4,1.9)
      \rput(4.3,2.1){\magenta{$c_\varepsilon$}}
      }

      \psline[linecolor=blue]{->}(0,3)(1,4)
      \psline[linecolor=blue]{->}(2,-1)(3,-1)
%      \showgrid
    \end{pspicture}
\hfill 
    \begin{pspicture}(0,0)(4,4)
 %           \showgrid
%      \rput(-.2,-.2){$0$}
      \pspolygon[fillcolor=vlg,linecolor=vlg,fillstyle=solid](0,0)(4,0)(4,4)
      
      \rput(3,1.3){$\cone(A_R^T)$}
      
       \psline[linecolor=red]{->}(0,0)(2,0)
       
       {
          \psline[linecolor=magenta]{->}(0,0)(2.4,0.3)
       }
      \psline[linecolor=blue]{->}(0,0)(1,0)
      \psline[linecolor=blue]{->}(0,0)(1,1)
%      \showgrid
    \end{pspicture}
    }
    
  \end{center}

   \caption{An example of perturbation.}
   \label{fig:perturb}
 \end{figure}
If $\epsilon>0$, then $R$ is a non-degenerate roof  of
 the linear program~\eqref{eq:6}. Thus condition~\ref{xitem:19}) holds
 for any $\epsilon>0$. In 
 the sequel, we make $\epsilon$ smaller and smaller, such that also the
 conditions~\ref{xitem:18}) and~\ref{xitem:20}) will be satisfied. 



 Let us  first deal with condition~\ref{xitem:18}). Let 
 $B$ be a set of linear indices such that the vectors $a_i, \, i \in
 B$ are a basis of $\setR^n$ and suppose that $B$ is not a roof. We have
 to guarantee that $B$ is not a roof of the perturbed linear program. 

 Let $A_{B} \in  \setR^{n\times n}$ be the sub-matrix of $A$ that is defined
 by the rows of $A$ indexed by $B$. 
 Since $B$ is not a roof, the vector ${A_{B}}^{-T}\,c $ has a strictly
 negative component. Suppose that this component is the $i$-th
 component $(A_{B}^{-T}\,c)(i)<0$.  By choosing $\epsilon>0$ sufficiently small,
 we guarantee that   
 \begin{equation}
   (A_{B}^{-T}\,c_\epsilon)(i)  =  \left(A_{B}^{-T}\, c + A_{B}^{-T}A_R^T     \mat{\epsilon \\ \epsilon^2 \\ \vdots \\\epsilon^n}\right)(i)\label{eq:23}    <  0.               
 \end{equation}
               Thus condition~\ref{xitem:18}) is satisfied by choosing $\epsilon>0$
 sufficiently small. 

 For  condition~\ref{xitem:20}) we have to work only a little harder, and
 in fact, this is why we add the vector $(\epsilon,\ldots,\epsilon^n)$ instead of
 $(\epsilon,\ldots,\epsilon)$. Let $B$ be a roof  of \eqref{eq:6} 
 and let $A_{B} \in \setR^{n\times n}$ be again the above described
 sub-matrix. 
 We now argue that, if  $\epsilon$ is sufficiently small, then $A_{B}^{-T}c_\epsilon$
 does not have any component equal to  zero.  We are done then in
 showing that, for $\epsilon$ sufficiently small, the linear program~\eqref{eq:6} is
 non-degenerate. Because any  roof of~\eqref{eq:6}  will then be
 non-degenerate. 

 So let us inspect the vector 
 \begin{equation}
   \label{eq:24}
   A_{B}^{-T} c_\epsilon = A_{B}^{-T}c + A_{B}^{-T}\cdot {A_R}^T \mat{\epsilon \\ \vdots \\ \epsilon^n}. 
 \end{equation}
 Each component of \eqref{eq:24} is a nonzero polynomial with
 variable~$\epsilon$.  Nonzero, because $A_{B}$ and ${A_R}$ are non-singular
 matrices. It is well known, that a nonzero polynomial of degree~$n$
 has at most $n$ roots. Thus, for $\epsilon>0$, sufficiently  small, no
 component of~\eqref{eq:24} will be zero. 



 So the conditions~\ref{xitem:19}), \ref{xitem:18}) and
 \ref{xitem:20}) hold for $\epsilon>0$ sufficiently small. Thus we can
 modify a degenerate linear program into an equivalent non-degenerate
 linear program and apply the simplex algorithm to it.


 \begin{theorem}
   \label{thr:2}
   Let 
   \begin{equation}
     \label{eq:11}
        \max\{c^Tx \colon x \in \setR^n, \, Ax\leq b\}
   \end{equation}
   be a linear program. The simplex method
   terminates on the perturbed linear program~\eqref{eq:6}. It either
   returns  a roof
   $B$ of~\eqref{eq:11} and \eqref{eq:6} whose vertex $x^*_{B}$ is
   an optimal solution of~\eqref{eq:11} or it asserts
   that~\eqref{eq:11} is infeasible. 
 \end{theorem}
 

 \begin{proof}
   The simplex method terminates on~\eqref{eq:6} since this linear
   program is non-degenerate. If it asserts that~\eqref{eq:18} is
   infeasible, then it also follows that \eqref{eq:11} is infeasible,
   since the perturbation only changes the objective-function vector.
   
   If it returns a roof $B$ of \eqref{eq:6}, then this is also a roof
   of~\eqref{eq:11} by condition~\ref{xitem:18}). Furthermore, the
   vertex $x^*_{B}$ is feasible for~\eqref{eq:11}. It follows from
   weak  duality   (Theorem~\ref{thr:6}) that $x^*_{B}$ is an optimal
   solution of~\eqref{eq:11}.  \qed 
 \end{proof}





   \subsection{The lexicographic pivot rule}
   \label{sec:lexic-pivot-rule}


   We now show that we do not have to physically perform the
   perturbation which sends $c$ to $c_\epsilon$, but instead describe a rule to
   choose the leaving index from the possible candidates for which the
   minimum in~\eqref{eq:9} is attained. Rules for entering and exiting
   indices are called \emph{pivoting rules}.  Recallthat for  $i \in
   B$, the function $f_B(i)$ 
   denotes the position of $i$ in the ordered set $B$. For example, if
   $B = \{5,2,9\}$, then $f_B(5)=1$ and $f_B(9)=3$. We need to define
   the lexicographic order. A vector $u \in \setR^n$ is lexicographically
   smaller than a vector $v \in \setR^n$ if $u\neq v$ and the first nonzero
   component of $v-u$ is positive.  We write $u<_{lex}v$ and if
   $u<_{lex}v$ of $u=v$ we write $u\leq_{lex}v$. 
   
   

   \begin{algorithm}{Simplex algorithm with lexicographic pivoting
       rule} \label{alg:3}
     
     \begin{tabbing}
       {\bf Input:} $A \in \setR^{m\times n}$, $c \in \setR^n$, $b
       \in \setR^m$
       $R\subseteq\{1,\ldots,n\}$  initial roof \\
      {\bf Output:}  $B$ optimal roof  or assert that LP is infeasible\\
      \\ 
      $B := R$  \\
       {\bf while } \= $\exists i \in\{1,\ldots,n\}$ with $a_ix^*_B >b_i$ \\
                   \> compute solution  \=  $y^*$ of \\
\> \>    $\Sigma_{k \in  B}  a_k y_k   =  -a_i $ \\
\> Compute  $J  = \{ k \in B \colon y^*_k <0\}$ \\

\> {\bf if }   $J = \emptyset$ assert {\bf LP not feasible} \\
\> {\bf else}\=      \\
\> \> Choose \emph{unique}  $j  \in J$ for which the vector \\
\> \> $\frac{(A_B^{-T} [c \mid {A_R}^T])_{f_B(j)}}{  - y^*_j}$ is \emph{ lexicographically
  minimal} \\
\> \> $B := B \textbackslash{} \{j\} \cup\{i\}$       
    \end{tabbing}
  \end{algorithm}
  % 
\marginpar{Change $T$ to $R$}
  A close inspection reveals that this algorithm only mimics the simplex
  algorithm on the perturbed problem~\eqref{eq:6} if $\varepsilon$ is
  arbitrarily close to $0$. This is because the first component of the
  vector $(A_B^{-T} [c \mid {A_R}^T])_{f_B(j)}$ is $z^*_j$, where $z^*$ is
  the solution of \eqref{eq:5} and if one replaces $c$ in
  \eqref{eq:5} with $c_\varepsilon$, the corresponding $z^*_j$ is simply 
  \begin{equation}
    \label{eq:50}
    (A_B^{-T} [c \mid {A_R}^T])_{f_B(j)} \mat{1\\\varepsilon\\\varepsilon^2 \\ \vdots\\\varepsilon^n}. 
  \end{equation}
  
%    So let us inspect step~3 for the perturbed problem
%    $\hpp(\eh,c_\epsilon)$. The solution $v^*$ to~(\ref{eq:17}) is the same for
%    the perturbed and non-perturbed problem. However, the solution
%    to~(\ref{eq:5})  is equal to 
%    \begin{equation}
%      \label{eq:25}
%      y^* + A^{-T} \cdot \wt{A} \mat{\epsilon \\ \vdots \\ \epsilon^n},
%    \end{equation}
%    where we neglect the last zero component. 
%    Now again, let $I \subseteq\{1,\ldots,n\}$ be the in $v^*(i)<0$. In the perturbed
%    problem, we determine the index $\ell \in I$ such that 
%    \begin{equation}
%      \label{eq:26}
%      - (y^*(\ell) + A^{-T} \cdot \wt{A} \mat{\epsilon \\ \vdots \\ \epsilon^n})(\ell) / v^*(\ell) 
%    \end{equation}
%    is minimal. In the unperturbed problem, there could be a set $L\subseteq I$ of
%    indices, for which the minimum is attained. 

%    \begin{definition}[Lexicographic order, $\geq_{\lex}$]
%      \label{def:2}
%      Let $u$ and $v$ be vectors of $\setR^n$. We say that $u\geq_{\lex}v$ if the
%      first nonzero entry of $u-v$ is positive. We say that $u$ is
%      \emph{lexicographically larger} than $v$. If $u \neq v$ then we write $u>_{\lex}v$. 
%    \end{definition}
%    %
%    In the following we write $r_j(\epsilon)$ for the $j$-th polynomial  
%    \begin{equation}
%      \label{eq:27}
%      -(y^*(j) + A^{-T} \cdot \wt{A} \mat{\epsilon \\ \vdots \\ \epsilon^n})(j) / v^*(j), \, j \in L.
%    \end{equation}
%    %
%    This polynomial is defined by its coefficient vector $r_j\in
%    \setR^{n+1}$.  The vector $r_j$ is the $j$-th row of the matrix 
%    $- \left( y^*, A^{-T} \cdot \wt{A}\right)$ divided by $v^*(j)$. 
%    We imitate the perturbed choice by choosing the index $\ell$ such that 
%    $r_j$ is lexicographically smallest. 
%    Let us spend a little more time in explaining why.  Let $r_\ell$
%    be this lexicographically smallest vector. Since $ A^{-T} \cdot \wt{A}$ is
%    non-singular, each row $r_j$ is strictly larger than $r_\ell$. 
%    This means that $r_j - r_\ell$ is a vector whose first nonzero component
%    is strictly positive. Thus $(r_j - r_\ell)(\epsilon) >0$ for $\epsilon$ sufficiently
%    small. 


%    Thus the following lexicographic  pivoting rule guarantees the
%    termination of the simplex method. 

%    \begin{xitemize}
%    \item   Determine the set $L\subseteq I$ for which the minimum 
%      \begin{displaymath}
%        \min_{i \in I} -  y^*(i)/v^*(i)
%      \end{displaymath}
%      is attained.  
%      \item Choose the index $\ell \in L$  such that 
%      \begin{displaymath}
%        \frac{\ell\mbox{-th row} - \left( y^*, A^{-T}\cdot\wt{A}\right)}{v^*(\ell)} \leq_\lex 
%        \frac{j\mbox{-th row} - \left( y^*, A^{-T} \cdot \wt{A}\right)}{v^*(j)} \mbox{ for each } j \in L.
%      \end{displaymath}
%      \item The $\ell$-th halfspace leaves the v-shape. 
%    \end{xitemize}



  \section{Phase~I, finding an initial roof}
  \label{sec:phase-i-finding}


  So far, we always started with an initial roof. Where do we get 
  it  from?  This is where Phase~I of the simplex method is put to
  work. Above, we described Phase~II.   First we prove a little lemma. 
  
  \begin{lemma}
    \label{lem:3}
    If the linear program~\eqref{eq:28} is feasible and bounded, then
    it has an optimal roof. In particular, a feasible and bounded
    linear program has an optimal solution.
  \end{lemma}
  
  \begin{proof}
    We change the linear program~\eqref{eq:28} by adding the
    additional constraint $c^Tx \leq M+1$, where $c^Tx\leq M$ is valid for
    all feasible solutions. We then have a (degenerate) roof by
    choosing this inequality together with any subset of   $n-1$
    constraints whose normal-vectors together with $c$ form a basis of
    $\setR^n$. The     simplex algorithm will find an optimal roof. \qed 
  \end{proof}


  We now form an auxiliary linear program. Observe that the linear
  program 
  \begin{equation}
    \label{eq:22}
    \max\{c^Tx \colon x \in \setR^n, \, Ax\leq b\}
  \end{equation}
  has a roof if and only  
  if the linear program $\max\{c^Tx \colon x \in \setR^n, \,Ax\leq0\}$ has a
  roof.  The latter linear program is feasible since $0$ is a feasible
  solution. Furthermore, $0$ is an optimal solution of this linear
  program if and only if it has a roof. This is what we check with the
  simplex algorithm on the auxiliary program 
  \begin{equation}
    \label{eq:14}
    \max\{c^T x \colon x \in \setR^n, \, Ax\leq0, \, c^Tx \leq 1\}
  \end{equation}
  and we start with a (possibly) degenerate roof involving the
  inequality $c^Tx\leq1$ and $n-1$ of the constraints of $Ax\leq0$ whose
  normal-vectors, together with $c$ form a basis of $\setR^n$. The
  simplex algorithm terminates with an optimal roof. If the roof has
  vertex $0$, then we have found a roof of the original linear
  program~\eqref{eq:22} and can start the simplex algorithm for it. If
  the roof of Phase~I still contains the inequality $c^Tx\leq1$, then
  the original linear program~\eqref{eq:22} does not have any
  roof. This either means that the program is infeasible or
  unbounded. 
  
%   We now form an auxiliary linear program. Let $C^+$ and $C^-$ be the sets of indices
%   $C^+ = \{i\in\{1,\ldots,n\}\colon c(i) \geq0 \}$
%   $C^-=\{i\in\{1,\ldots,n\}\colon c(i)<0\}$.  The auxiliary linear program is 
%   \begin{equation}
%     \label{eq:12}
%     \max\{c^Tx \colon x \in \setR^n,\,  Ax\leq0, x(i) \leq1, \, i \in C^+, \, -x(i) \leq1, \, i \in C^-\}. 
%   \end{equation}

%   Notice that this linear program  is feasible and has a roof, namely
%   the inequalities
%   \begin{equation}
%     \label{eq:13}
%        x(i) \leq1, \, i \in C^+, \, -x(i) \leq1, \, i \in C^-.
%   \end{equation}
% %
%   The vertex $x^*$ of each roof that 
%   contains one of the inequalities~\eqref{eq:13} has a positive
%   objective value $c^Tx^*$. On the other hand, a roof  of
%   \eqref{eq:28}  yields a roof of~\eqref{eq:12}  with vertex $x^*=0$.
%   We start the simplex algorithm on this auxiliary linear program. If
%   we obtain a 
%   roof $B$ with vertex $x^*$ such that $c^Tx^*>0$, then the original
%   linear program  does not have roofs. Otherwise, we have identified
%   a subset  $B' \subseteq B$ of linear independent halfspaces  such that $c$ is a
%   conic combination of the normal-vectors of these half-spaces. If
%   $|B'|<n$,  we complete 
%   $B'$ to a degenerate roof $\wt{B}$ of the original linear program
%   and start the simplex algorithm on $\wt{B}$. 


%   \section{The running time of the simplex method}
%   \label{sec:running-time-simplex}


%   Starting at a v-shape $\ev$, it may happen, that one basically visits
%   all v-shapes of $\hpp(\eh,c)$ before one arrives at the lowest
%   v-shape. The number of v-shapes is bounded by $\binom{m}{n} =
%   \bigO(m^n)$. This is {\bf exponential} in the dimension. Thus the
%   simplex algorithm runs in polynomial time if the dimension is
%   fixed. We will see later, that there exists an algorithm, which solves
%   the highest point problem in polynomial time, even if the dimension
%   varies. 

%   However, the simplex method performs very well in practice. 

%   We close this section with a larger example.
%   \begin{example}
%     We consider a larger example with several iterations of the simplex
%     algorithm. Our goal is to solve a linear program~\eqref{eq:28} with 
%     \begin{displaymath}
%       A = \begin{pmatrix} 2& 4& 3 \\ 3& 1& 4 \\ 1& 2& 2 \\ 1& 0& 0 \\ 0&
%           1& 0 \\ 0& 0& 1 \end{pmatrix},\,    
%         b = \begin{pmatrix}-5\\ -3\\ -2 \\0 \\0 \\0 \end{pmatrix}\, ,
%         c = \begin{pmatrix} 5 \\ 9 \\ 10          \end{pmatrix}. 
%     \end{displaymath}
%   %
%     The set of row-indices is $\{1,2,3,4,5,6\}$. We start with the roof
%     $\{4,5,6\}$. 
%     \begin{xitemize}
%     \item[$\{4,5,6\}$]  
%     \end{xitemize}

  


  
%   \end{example}



\section{The full column-rank assumption}
\label{sec:full-column-rank}

There is one thing that we still have to deal with. The simplex
algorithm is based on the assumption that the columns of $A$ are
linearly independent. %  We use the following lemma from linear algebra
% to make this assumption without loss of generality. 
We now argue that this assumption can be made without loss of
generality. 

Suppose that $A$ can be written as $[A_1 \mid A_2]$ with $A_1 \in \setR^{m\times k}$
and $A_2 \in \setR^{m\times(n-k)}$ where the columns of $A_1$ are linearly
independent and each column of $A_2$ is a  linear combinations of the
columns of $A_1$.  We also write $c = \smat{c_1\\c_2}$ with $c_1 \in
\setR^k$ and $c_2 \in \setR^{n-k}$ and consider the linear program 
\begin{equation}
  \label{eq:32}
  \max\{c_1^Tx_1 \colon x_1 \in \setR^{k} , \, A_1x_1\leq b\}
\end{equation}


Since each column of $A_2$ is a linear combination of the columns of
$A_1$, there exists a uniquely determined  matrix $U \in \setR^{k\times(n-k)}$
with $A_2 = A_1\cdot U$.  


\begin{lemma}
  \label{lem:7}
  The linear program~\eqref{eq:32} is feasible if and only if the
  linear program~\eqref{eq:28} is feasible. 
\end{lemma}


\begin{proof}
  Suppose  that $x^* =\smat{x_1^*\\x_2^*}$ is a feasible
  solution of \eqref{eq:28}, i.e.,  $A_1x_1^* + A_2 x_2^*\leq b$. But
  $A_1x_1^* + A_2x_2^* = A_1 (x_1^* + U \cdot x_2^*)$ which yields a
  feasible solution of~\eqref{eq:32}. Likewise we see that any
  solution $x_1^*$ of $A_1x_1\leq b_1$ can be extended to a feasible
  solution $\smat{x^*_1 \\ 0}$ of~\eqref{eq:28}. \qed 
\end{proof}


\begin{lemma}
  \label{lem:8}
  If~\eqref{eq:32} is feasible and if $c_2^T \neq c_1^T\cdot U$, then
  \eqref{eq:28} is unbounded. 
\end{lemma}


\begin{proof}
  Since~\eqref{eq:32} is feasible, then also~\eqref{eq:28} is
  feasible. Let $x^*$ thus be a feasible solution
  of~\eqref{eq:28}. Then $x^* + \mu \smat{-U v \\ v}$ is feasible for
  any $v \in \setR^{n-k}$ and $\mu \in \setR$. Let $v$ satisfy $c_2^Tv \neq 0$,
  then $c_1^T (-U) v + c_2^T v \neq 0$ which implies that \eqref{eq:28}
  is unbounded. \qed 
\end{proof}


The idea is thus to re-order the columns of $A$ such that the first
$k$ columns are linearly independent and the last $n-k$ columns are
linear combinations of the first $k$ columns. This also yields a
matrix $U$ and we now solve the linear program~\eqref{eq:32} with the
simplex algorithm. If it is infeasible or unbounded, then so is the
linear program~\eqref{eq:28}. Otherwise the simplex algorithm finds an
optimal solution $x_1^*$  of~\eqref{eq:32}. If $c_2^T = c_1^T \cdot U$
then this optimal solution $\smat{x_1^*\\0}$ is also an optimal
solution of the linear program~\eqref{eq:32}. This is because a
feasible solution $\smat{x_1^*\\x_2^*}$ of \eqref{eq:28} yields a
feasible solution $x_1^*+U\cdot x_2^*$ of \eqref{eq:32} with same objective
value $c_1^T(x_1^* + U x_2^*) = c_1^Tx_1^* + c_2^Tx_2^*$. 

% \begin{lemma}
%   \label{lem:6}
%   Let $A \in \setR^{m\times n}$ where the first $k\leq n$ columns are linearly
%   independent and the last $n-k$ columns of $A$ are obtained as linear
%   . There exists an invertible
%   matrix $U \in \setR^{n\times n}$ with $A\cdot U = [B | 0]$, where $B$ has full
%   column-rank.
% \end{lemma}
% Here $0$ stands for the all zero matrix with $m$ rows and $n-k$
% columns. 
% This lemma is proved in a course on linear algebra. The algorithm with
% which one can compute the matrix $U$ is Gaussian elimination. We use
% this to compute from  the linear program
% \begin{equation}
%   \label{eq:30}
%   \max\{c^Tx \colon  x \in \setR^n, \,Ax \leq b\}
% \end{equation}
% the linear program
% \begin{equation}
%   \label{eq:31}
%   \max\{c_1^Ty \colon  y \in \setR^{n-k}, \, By \leq b\},
% \end{equation}
% where $c^T = (c_1^T,c_2^T)$ with $c_1 \in \setR^{n-k}$. 
% We observe that this linear program is feasible if and only
% if~\eqref{eq:30} is feasible. 


\section*{Exercises}

\begin{enumerate}[1)]
\item \label{xitem:8}
  Show that the linear program~\eqref{eq:29} that is defined by a roof
  is always feasible. 
\item  Let $B$ be a roof of the linear program~\eqref{eq:28} and 
  consider the system of linear equations
  \begin{displaymath}
    a_{j}^Tx = -1, a_i^Tx=0, \, i \in B\setminus\{j\}
  \end{displaymath}
  for an index $j \in B$. 
  Let $x^*$ be a feasible solution to the roof-linear
  program. Show that  $x^* + \lambda \cdot v$ is also feasible for each 
  $\lambda>0$.  \label{xitem:9} 


\item Prove Proposition~\ref{prop:2}. \label{xitem:11}


\item 
  \label{sec:high-point-probl}
  Let $B$ be a roof with vertex $x^*_B$. Show that the 
   set of  feasible points $\{x \in \setR^n \colon a_i^Tx\leq b(i), \, i\in B\}$
   of the roof is of the form 
  $x_B^* + \cone\{r_1,\ldots,r_n\}$ for some suitable vectors $r_i \in
  \setR^n$, $i=1,\ldots,n$. 

\item A cone $C\subseteq \setR^n$ is \emph{pointed} if it does not contain a line: There are no vectors $x\in C$, $v\in \setR^n$ such that $x+\lambda v\in C$ for all $\lambda\in \setR$.

Prove the following variant of Carath\'eodory's theorem. 
Given some set  $X\subseteq \setR^n$, $|X|>n$ such that $\operatorname{cone}(X)$ is pointed.
For any $x \in \operatorname{cone}(X)$, there exist at least two
  different subsets $X_1,X_2\subseteq X$ with  $|X_1| = |X_2| = n$ such that $x\in \operatorname{cone}(X_1)\cap \operatorname{cone}(X_2)$.
 \label{xitem:12}
\item 
Consider the problem
\begin{eqnarray}
\label{row0}    \max & z & \\
\label{row1}    \text{s.t.} &  x+2y      & \leq -3 \\
\label{row2}        &  -2x - 3y  & \leq 5 \\
\label{row3}        & -2x -y  +2z& \leq -1 \\
\label{row4}	    &  3x+y  & \leq 2 \\
\label{row5}	    &  x  & \leq 0 \\
\label{row6}	    &  y  & \leq 0 \\
\label{row7}	    &  z  & \leq 0.
\end{eqnarray}
Assume that we perform the simplex method, and at some point have the roof given by the rows \eqref{row1}, \eqref{row6} and \eqref{row7}.
Figure \ref{fig:IllustrDegeneracyExercise} shows the situation in the 2-dimensional subspace given by the hyperplane $z=0$.

Show that the simplex algorithm might not terminate, by giving a cycling sequence of roofs that might be selected by the simplex method.
Explain why your sequence is valid (it is sufficient to give drawings here, you do not need to compute the roof vertices explicitly).

\emph{Hint: Never let \eqref{row7} leave the roof. Then it is sufficient to consider the subspace as in the illustration.}

\begin{figure}[hbt]

  \begin{center}
\begin{pspicture}(-4cm,-4cm)(4cm,4cm)
\psset{unit=0.4cm}
\pspolygon[fillcolor=vlg,linecolor=vlg,fillstyle=solid](-10,3.4)(-10,2.9)(10,-7.1)(10,-6.6)
\pspolygon[fillcolor=vlg,linecolor=vlg,fillstyle=solid](-10,5.1)(-10,5.6)(10,-7.73)(10,-8.23)
\pspolygon[fillcolor=vlg,linecolor=vlg,fillstyle=solid](-4.4,10)(-3.9,10)(6.1,-10)(5.6,-10)
\pspolygon[fillcolor=vlg,linecolor=vlg,fillstyle=solid](-2.76,10)(-3.26,10)(3.5,-10)(3.9,-10)
\pspolygon[fillcolor=vlg,linecolor=vlg,fillstyle=solid](-0.1,10)(-0.6,10)(-0.6,-10)(-0.1,-10)
\pspolygon[fillcolor=vlg,linecolor=vlg,fillstyle=solid](-10,-0.1)(-10,-0.6)(10,-0.6)(10,-0.1)


      \psline(-10,3.5)(10,-6.5) %first
      \psline(-10,5)(10,-8.33) %second
      \psline(-4.5,10)(5.5,-10) % third
      \psline(-2.66,10)(4,-10) % fourth
      \psline(0,10)(0,-10) % fifth
      \psline(-10,0)(10,0) % sixt

\rput(-9.5,2.2){\eqref{row1}}
\rput(-9.5,6){\eqref{row2}}
\rput(-5.2,9.5){\eqref{row3}}
\rput(-1.86,9.5){\eqref{row4}}
\rput(0.6,9.5){\eqref{row5}}
\rput(-9.5,-1){\eqref{row6}}
      
      \psdot[linecolor=red](-3,0)
      %\psline[linecolor=blue]{->}(0,3)(1,4)
      %\psline[linecolor=blue]{->}(2,-1)(3,-1)
%      \showgrid
\end{pspicture}
\end{center}
\caption{The halfspaces defined by system \eqref{row0} in the subspace
  $\{(x,y,0)~:~x,y\in \setR\}$.}
\label{fig:IllustrDegeneracyExercise}
\end{figure}


\end{enumerate}




%%% Local Variables: 
%%% mode: latex
%%% TeX-master: "lecture"
%%% End: 

%%\chapter{The development of the simplex method}
\label{cha:introduction}



In this chapter, we describe an algorithm which solves linear
programming problems $\min\{c^Tx \colon Ax = b, \, x\geq0\}$ in
equation standard form, where $A \in \setR^{m\times n}$ has full row-rank, i.e.,
the rows of $A$ are linearly independent. Our presentation follows
closely the one in~\cite{BertsimasTsitsiklis97}. 




\section{Equation standard form}
\label{sec:equat-stand-form}

We have seen in Chapter~\ref{cha:introduction-1} that each linear
program can be described in inequality standard form, i.e., as
$\max\{c^Tx \colon Ax\leq b\}$. The simplex algorithm is conveniently
described for linear programs in \emph{equation standard form} $\min\{c^Tx \colon
Ax = b, \, x\geq0\}$.  
%, where one can assume that the rows of $A$ are linearly independent. 

The next example shows how to transform a linear program in inequality
standard form in equation standard form.  Consider the linear program 
\begin{equation}
\label{eq:28}
  \begin{array}{lrcl}
    \text{maximize}   & 3 x_1 - 2x_2 \\
    \text{subject to} & 2x_1 - x_2 & \leq & 4 \\
                      & x_1 + 3x_2 & \leq & 5 \\
                      & -x_1 & \leq & 0.     
  \end{array}
\end{equation}
Our goal is to transform this linear program into one of the form
$\min\{c^Tx \colon Ax = b, \, x\geq0\}$.  The objective function can be
re-written as $\min -3x_1 + 2x_2$.  Furthermore, we observe that the
last inequality of~\eqref{eq:28} can be re-written as $x_1\geq0$. The
variable $x_2$ is not bounded from below. The trick is to split $x_2$
into $x_2 = x_2^+ - x_2^-$, where $x_2^+,x_2^- \geq0$. The linear
program then becomes 
\begin{equation}
  \label{eq:29}
  \begin{array}{llcl}
    \text{minimize}   & -3 x_1 + 2x_2^+ -2 x_2^- \\
    \text{subject to} & 2x_1 - x_2^+ + x_2^- & \leq & 4 \\
                      & x_1 + 3x_2^+ -3 x_2^- & \leq & 5 \\
                      & x_1, x_2^+,x_2^- &\geq &0.
  \end{array}
\end{equation}
Next we introduce two new \emph{slack variables} $z_1,z_2\geq0$, which
model $z_1 = 4- (2x_1 - x_2^+ + x_2^-)$ and $z_2 = 5 - (x_1 + 3x_2^+
-3 x_2^-)$ and obtain a linear program in equation
standard form which is equivalent to~\eqref{eq:28} 
\begin{displaymath}
  \begin{array}{llcl}
    \text{minimize }   & -3 x_1 + 2x_2^+ -2 x_2^- \\
    \text{subject to }\, & 2x_1 - x_2^+ + x_2^- + z_1 & = & 4 \\
                      & x_1 + 3x_2^+ -3 x_2^- + z_2  & = & 5 \\
                      & x_1, x_2^+,x_2^-,z_1,z_2 & \geq &0.
  \end{array}
\end{displaymath}
%
In matrix notation, we can describe the transformation procedure from
above as follows. The input is a linear program in inequality standard
form $\max\{c^Tx \colon x \in \setR^n, \,Ax\leq b \}$ with $A \in \setR^{m\times n}$. From this, one creates the
equivalent linear program 
\begin{displaymath}
  \begin{array}{ll}
    \min & -c^T x^+ + c^T x^-  \\
    & A x^+ - A x^- + I_m z =  b\\
     &x^+, x^-, z  \geq 0  \\
    & x^+,x^- \in \setR^n, z \in \setR^m,     
  \end{array}
\end{displaymath}
where $I_m \in \setR^{m\times m}$ is the $m\times m$ identity matrix, i.e. for
$1\leq i,j\leq m$ one has 
\begin{displaymath}
  I_m(i,j) = 
  \begin{cases}
    1, & \text{if } i=j\\
    0 & \text{otherwise.}
  \end{cases}
\end{displaymath}


\subsubsection*{The full row-rank assumption}

We recall some definitions and facts from linear algebra.   For $A \in \setR^{m\times n}$
the \emph{row-rank} of $A$ is the maximum number of linearly independent
rows. The \emph{column-rank} of $A$ is the maximum number of linearly
independent columns. The row-rank and column rank of $A$ are
equal. This number is the \emph{rank} of $A$ and is denoted by
$\rank(A)$. 

\begin{definition}
  \label{def:8}
  A matrix $A \in \setR^{m\times n}$ has \emph{full row-rank}, if $\rank(A) =
  m$. Similarly, $A$ has \emph{full column-rank}, if $\rank(A) =n$. 
\end{definition}

We now show that  we furthermore can assume that the matrix $A$ in the
linear  program 
\begin{equation}
  \label{eq:30}
  \min\{c^Tx \colon x \in \setR^n, \, Ax =b,\, x\geq0\}
\end{equation}
has full row-rank, i.e., the rows of $A$ are linearly independent. 
To see this, suppose that $A\in \setR^{m\times n}$ does not have full row-rank. Then
there is a row $a_j$ of $A$ which is in the span of the other rows,
i.e., one has
\begin{displaymath}
  a_j = \sum_{\substack{i=1\\ i\neq j}}^m \lambda_i a_i \text{ with suitable numbers }
  \lambda_i \in \setR, i\in \{1,\ldots,m\} \setminus\{j\}. 
\end{displaymath}
If $\sum_{\substack{i=1\\ i\neq j}}^m \lambda_i  b(i) = b(j)$, then one has for all  $x \in \setR^n$
with 
$a_i^Tx = b(i), i=1,\ldots,m, i\neq j$ also  $a_j^Tx =  b(j)$ which means that the
$j$-th  equation in  $Ax = b$ can be removed. 
If $\sum_{\substack{i=1\\ i\neq j}}^m \lambda_i  b(i) \neq b(j)$, then there does not exist an $x \in
\setR^n$ with $Ax = b$ and the  LP~(\ref{eq:30}) is infeasible. 

\begin{example}
  \label{sec:full-row-rank}
  Consider the linear program 
  \begin{equation}
    \label{eq:46}
    \begin{array}{rl}
      \min  x_1 + 2 \,x_2 + x_3 + 4 x_4 \\
      x_1 + x_2 + 2 x_3 + 3 x_4 & = 5 \\
      x_1 + 2 x_2 + x_3 + 4 x_4 & =7 \\
      3x_1 +4x_2 + 5 x_3 + 10 x_4& = 17\\
      x_1,x_2,x_3,x_4 \geq0. 
    \end{array}
  \end{equation}
  The third equation is can be obtained by multiplying the first
  equation with $2$ and adding it to the second equation. This linear
  program is not of full row-rank and it is equivalent to the linear
  program 
  \begin{equation}
    \label{eq:47}
    \begin{array}{rl}
      \min  x_1 + 2 \,x_2 + x_3 + 4 x_4 \\
      x_1 + x_2 + 2 x_3 + 3 x_4 & = 5 \\
      x_1 + 2 x_2 + x_3 + 4 x_4 & =7 \\
      x_1,x_2,x_3,x_4 \geq0. 
    \end{array}
  \end{equation}
  which is of full row-rank. The linear program 
 \begin{displaymath}
    \begin{array}{rl}
      \min  x_1 + 2 \,x_2 + x_3 + 4 x_4 \\
      x_1 + x_2 + 2 x_3 + 3 x_4 & = 5 \\
      x_1 + 2 x_2 + x_3 + 4 x_4 & =7 \\
      3x_1 +4x_2 + 5 x_3 + 10 x_4& = 15\\
      x_1,x_2,x_3,x_4 \geq0. 
    \end{array}
  \end{displaymath}
  is not of full row-rank but is infeasible, since, as before, the
  third row of the constraint matrix is $2$ times the first row plus
  the second row, however $2\cdot 5 + 7 =17 \neq 15$. 
\end{example}



\begin{assumption} 
\label{assumption1}
We assume in the following that our linear program which
we want to solve is in equation standard form~\eqref{eq:30}, or in
\emph{standard   form} for short, with $A\in \setR^{m\times n}$ having full
row-rank. 
\end{assumption}



\section{Bases and basic solutions}
\label{sec:bases-basic-solut}


We use the following notation.  For an index set
$I=\{i_1,\ldots,i_k\}\subseteq\{1,\ldots,n\}$, with $i_1<i_2<\cdots<i_k$, $A_I$ denotes
the $m\times k$ matrix whose $j$-th column is $i_j$-th column $a^{i_j}$ of
$A$ for $j=1,\ldots,k$.
Similarly $c_I\in \setR^k$ denotes the vector whose $j$-th component is
$c(i_j)$, $j=1,\ldots,k$.   
\begin{example}
  \label{ex:-2}
  Let $A =\smat{2 & 3 & 1 \\ 4 & 5 & 6}$ and $I = \{2,3\}$, then $A_I
  = \smat{3&1\\5&6}$.
\end{example}
The proof of the next lemma is very similar to
the proof of Carath\'eodory's theorem, Theorem~\ref{conv:thr:7}.
It will help us to determine an optimal solution to~\eqref{eq:30}.



  \begin{lemma}
    \label{lem:8}    
    Let $x^*$ be a feasible solution of the linear
    program~\eqref{eq:30} and let $J = \{i \colon x^*(i)>0\}$ be the index
    set corresponding to those components of $x^*$ which are strictly
    positive. There exists a procedure that either asserts that the
    linear program is unbounded or computes a solution $\wt{x}$
    of~\eqref{eq:30} such that
    \begin{enumerate}[i)]
    \item     the index set $\wt{J}= \{ i \colon \wt{x}(i)>0\}$ is
       contained in $J$, $\wt{J}\subseteq J$,
    \item $A_{\wt{J}}$ has full column-rank,
    \item $c^Tx^* \geq c^T\wt{x}$. 
    \end{enumerate}

   % If the linear program~\eqref{eq:30}  is feasible and  bounded,
%    then there  exists an optimal solution  $x^* \in \setR^n$
%    of~\eqref{eq:30} such that    $A_J$    has full column rank, where
%    $J = \{ j \mid x^*(j)>0\}$.  
  \end{lemma}


  \begin{proof}    
    Let $x^*$ be a feasible solution and suppose that the columns of
    $A_J$ are linearly dependent.  The idea is now to compute a new
    solution $x'$ with $J'=\{ j \colon  x'(j)>0\} \subset J$ and $c^T
    x' \leq c^T x^*$ or to assert that the linear program is
    unbounded.  After at most $n$  repetitions of this procedure,
    one has found a solution $\wt{x}$  which satisfies the condition of the
    theorem or one has asserted that the linear program is unbounded. 
             

    Let $\wb{J}$ be the index set $\wb{J}= \{1,\ldots,n\} \setminus J$.  Since
    $A_J$ does not have full column rank there exists a $d \in \setR^n$
    with $d\neq0$, $A\, d =0$ and $d(j)=0$ for each $j \in \wb{J}$.  Consider
    the points 
    \begin{equation}
      \label{eq:32}
        x^* \pm \lambda\, d  \text{ for } \lambda \in \setR.
    \end{equation}
    Since $A (x^* \pm \lambda d)
    =b$ for each $\lambda \in \setR$, the only way to make  a point~\eqref{eq:32}
    infeasible  is to violate the lower bounds $x^* \pm \lambda d \geq0$. Since
    $x^*_J>0$ and $d_{\wb{J}} =0$ there exists 
    a sufficiently small $\varepsilon^*>0$ such that for each $\lambda \in \setR$ with
    $0 \leq \lambda \leq \varepsilon^*$ the point $x^* \pm \lambda d$ is feasible. 
    
    Suppose that $c^Td \neq0$. By multiplying $d$ with $-1$
    we can assume that $c^Td <0$. If $d\geq0$, then $x^*+\lambda\,d\geq0$ for
    each $\lambda\geq0$ and since $c^T(x^*+\lambda\,d) = c^Tx^* + \lambda c^Td$, we can
    make the objective value of $x^*+\lambda\,d$ as small as we wish. Thus we
    assert that the linear program is unbounded. Otherwise, let $K =
    \{i \colon d(i) <0\}$ be the index set corresponding to the negative
    components of $d$. How large can we choose $\lambda>0$ such $x^* + \lambda
    d$ is still feasible?  For $\lambda>0$, the point  $x^* + \lambda \,d$  is
    feasible if and only if 
    \begin{equation}
      \label{eq:33}
      \lambda \leq \min\{- x^*(k) / d(k) \colon  k \in K\}. 
    \end{equation}
    Let $\lambda_{\max} = \min\{  - x^*(i) / d(i) \colon i \in K\}$ and let $k^*\in
    K$
    be an index, where the minimum is attained. The point
    $x' = x^*+ \lambda_{\max} d$ is feasible, $c^Tx' < c^Tx^*$ and
    since $x'(k^*)=0$ one has with $J' = \{ i \colon x'(i) >0\}$ also $J'\subset J$. 
    
    Suppose now that $c^Td = 0$. By multiplying $d$ with $-1$ if
    necessary, we can assume that $d$ has strictly negative
    components. We proceed as above with the index set $K$ and
    $\lambda_{\max}$ to     construct a feasible  $x' = x^*+ \lambda_{\max} d$ with
    $c^Tx' \leq c^Tx^*$ and  $J' = \{ i \colon x'(i) >0\}\subset J$.     
  \qed
\end{proof}

\begin{example}
  \label{exa:ll}
  Consider the linear program~\eqref{eq:47} and the feasible solution
  $x^* = (1,1,0,1)$. The set $J$ from the proof above is $J =
  \{1,2,4\}$.  The vector $d^T = (-2, -1, 0, 1)$ satisfies $Ad=0$ and
  $d_{\wb{J}} =0$. A largest $\lambda>0$ such that $(1,1,0,1) + \lambda (-2, -1, 0, 1)
  \geq0$ is $\lambda_{\max} = 1/2$ and the new solution is $(1, 1, 0, 1)
  + (1/2)\cdot  (-2, -1, 0, 1) = (0,0.5,0,1.5)$. This is a basic feasible
  solution and the procedure described in the proof of
  Lemma~\ref{lem:8} ends. 
\end{example}



\begin{definition}[Basis, basic solution, basic feasible solution]
  \label{def:13}
   A set $B\subseteq\{1,\ldots,n\}$ with $|B| = m$ and $A_B$ non-singular is
   called \emph{basis}. A vector $x^* \in \setR^n$ is a \emph{basic
     solution}, if there exists a basis $B$ with $x^*_B = A_B^{-1}b$
   and $x^*_{\wb{B}}=0$, where $\wb{B} = \{1,\ldots,n\} \setminus B$. We say that
   $x^*$ is associated to the basis $B$. If
   additionally $x^*\geq0$ holds, then $x^*$ is a \emph{basic feasible
     solution}. 
\end{definition}


  \begin{lemma}
    \label{lem:3} Let  $x^* \in \setR^n$ be a feasible solution and let $J
    = \{ i \colon x^*(i) \neq 
    0\}$. If the columns of $A_J$ are linearly independent, then $x^*$
    is a basic feasible solution.
  \end{lemma}
  \begin{proof}
    Augment $J$ to index set $B\supseteq J$ such that $A_B$ is non-singular.
    One has $A_Bx^*_B=b$ and $x^*(j)=0$ for all $j \notin B$, thus $x^*$
    is a basic solution.  \qed
  \end{proof}
  

In exercise~\ref{item:5} you are to prove that there could be several
bases which are associated to a basic feasible solution $x^*$. 


Exercise~\ref{ex:16} shows that there exist optimization problems
$\min\{c^Tx  \colon x \in K\}$ where $K\subseteq\setR^n$ is closed and bounded which
are bounded but do not have optimum solutions. A corollary of
Lemma~\ref{lem:8} is that this cannot happen for linear programs. 

\begin{corollary}
  \label{co:2}
  A bounded and feasible linear program has an optimal basic feasible
  solution.  
\end{corollary}

\begin{proof}
  Lemma~\ref{lem:8} together with Lemma~\ref{lem:3} shows that for
  each feasible $x^*$, there exists a basic feasible solution with an
  objective function value which is at most the objective function
  value of $x^*$. 
  Since the number of basic feasible solutions is finite, 
  a bounded linear program has an optimal basic feasible solution.

\end{proof}



\section{A naive algorithm for linear programming} 
 \label{sec:naive-algor-line}
 
 Suppose that we have a bounded linear program $\min\{c^Tx \mid x \in
 \setR^n, \,Ax = b,\, x\geq0\}$ in standard form. We can already describe a
 finite algorithm which computes an optimal solution, if the linear
 program is feasible. Lemma~\ref{lem:8} tells us that, if there exists
 an optimal solution, then there exists also an optimal basic feasible
 solution. 

 
The algorithm maintains a value which reflects the \emph{best
  objective value} observed so far. In the beginning this is $\infty$. 

We now enumerate all index sets $B \in \binom{n}{m}$ and test, whether
$A_B$ is non-singular with Gaussian elimination. If this is the case,
we compute $x^*_B = A_B^{-1}b$ and check whether $x^*_B\geq0$. If this
is also the case, we check whether $c_B^Tx^*_B$ is smaller than the
previously best observed objective value and update this best
objective function value if this is again the case.

This algorithm is however very inefficient. The enumeration of all
index sets $B \in \binom{2m}{m}$ requires already at least $2^m$
steps. 
If $m$ is large this is enormous, see exercise~\ref{item:6}. 
The simplex algorithm is based on a smarter principle.
It \emph{improves} a basic feasible solution, by replacing it by
another basic feasible solution which is in the neighborhood of the
previous one. We will describe this next.
 



\section{Simplex method: First steps}
\label{sec:simplex-method}

 Consider the linear
programming problem 
\begin{equation}
  \label{eq:40}
  \begin{array}{r}
  \begin{matrix}               
    \min \quad \quad  & 2 x_1 & + & 3 x_2 & +&  4 x_3 & +&  4 x_4  & & \\
                      & x_1   & + & 2 x_2 & + & 3 x_3 & + & 4 x_4 & =&  3 \\
                      & x_1   & + & x_2  & + & x_3 & + & x_4 & = & 2                      \\                          
  \end{matrix} \\
  x_1,x_2,x_3,x_4 \geq0. 
\end{array}
\end{equation}
The first two columns of the constraint matrix $A$ are linearly
independent. Let us now compute $x^* \in \setR^n$ such that $x^*_B =
A_B^{-1} \cdot \smat{3\\2}$. For this, we subtract the second row of the 
system 
\begin{displaymath}
  \begin{matrix}    
    x_1   & + & 2 x_2 & + & 3 x_3 & + & 4 x_4 & =&  3 \\
    x_1   & + & x_2  & + & x_3 & + & x_4 & = & 2
    \\                          
  \end{matrix}
\end{displaymath}
from the first, to obtain 
\begin{displaymath}
  \begin{matrix}    
         &   &  x_2 & + & 2 x_3 & + & 3 x_4 & =&  1 \\
    x_1   & + & x_2  & + & x_3 & + & x_4 & = & 2
    \\                          
  \end{matrix}
\end{displaymath}
Then we subtract the first row from the second, to obtain 
\begin{displaymath}
  \begin{matrix}    
         &  &  x_2 & + & 2 x_3 & + & 3 x_4 & =&  1 \\
    x_1   & + &     & + & -1 x_3 & + & -2x_4 & = & 1
    \\                          
  \end{matrix}
\end{displaymath}
Finally we swap the first and second row, and obtain the system 
\begin{displaymath}
  \begin{matrix}       
    x_1   & + &     & + & -1 x_3 & + & -2x_4 & = & 1    
    \\                          
    &   &  x_2 & + & 2 x_3 & + & 3 x_4 & =&  1 
  \end{matrix}
\end{displaymath}

\begin{definition}[Elementary row-operations]
  \label{def:14}
  Let $A\in \setR^{m\times n}$ be a matrix. An \emph{elementary row operation}
  on $A$ is the addition of a multiple of one row, to another row. 
\end{definition}

The following lemma is known from linear algebra. 
\begin{lemma}
  \label{lem:12}
  Let $A,A' \in \setR^{m\times n}$ and $b , b' \in \setR^m$. If the matrix
  $\left[A'|b'\right]$ results from  $\left[A|b\right]$ by a finite
  series of elementary row operations, then  $\{x \in  \setR^n \colon Ax=b\} =
  \{x \in  \setR^n \colon A'x=b'\}$. 
\end{lemma}

\noindent 
In other words, we have re-written our linear program from above as 
\begin{displaymath}
 \begin{array}{r}
  \begin{matrix}               
    \min \quad  & 2 x_1 & + & 3 x_2 & +&  4 x_3 & +&  4 x_4  & & \\
 &      x_1   & + &     & + & -1 x_3 & + & -2x_4 & = & 1    
    \\                          
&    &   &  x_2 & + & 2 x_3 & + & 3 x_4 & =&  1 
  \end{matrix} \\
  x_1,x_2,x_3,x_4 \geq0. 
\end{array}
\end{displaymath}  
From this representation, we can easily read off the basic feasible
solution $x_1=1, x_2=1, x_3 = 0, x_4 = 0$ associated to $B =
\{1,2\}$. We now consider the system
\begin{displaymath}
   \begin{matrix}               
    2 x_1 & + & 3 x_2 & +&  4 x_3 & +&  4 x_4  & =  & 0 \\
    \hline
      x_1   & + &     & + & -1 x_3 & + & -2x_4 & = & 1    
    \\                          
    &   &  x_2 & + & 2 x_3 & + & 3 x_4 & =&  1 
  \end{matrix}. 
\end{displaymath}
If we subtract 2-times the second row and 3-times the third row from
the first row, we obtain the system 
\begin{displaymath}
  \begin{matrix}    
 & + &       & + &     0 x_3 & + &    -1x_4 & = &
-5\\ \hline 
1  x_1 & + &    0 x_2 & + &   -1x_3 & + &    -2 x_4 & = &    1\\
     0x_1 & + &     1x_2 & + &     2x_3& + &     3x_4 & = &     1
   \end{matrix}                             
\end{displaymath}
We have eliminated the basic variables in the first row. The entry
behind the equality sign of the first row is the negative of the
objective value of the basic feasible solution $(1,1,0,0)$. We write
this in a more compact form as follows 
\begin{equation}
  \label{eq:36}
  \begin{tabularx}{.3 \textwidth}{XXXX|X}
    0   &      0 &     0  &    -1   &   -5\\ \hline 
    1   &    0  &    -1  &    -2 &      1\\
    0 &     1 &     2&     3  &     1
  \end{tabularx}
\end{equation}
%
The general interpretation of this approach is as follows. We
start with 
\begin{displaymath}
   \begin{array}{c|c}
    c^T  & 0 \\ \hline 
   A        & b
  \end{array}
\end{displaymath}
and, given a basis $B$ of $A$ compute 
\begin{equation}
  \label{eq:37}
  \begin{array}{c|c}
    c^T - c_B^TA_B^{-1}A & -c_B^Tx^*_B \\ \hline 
    A_B^{-1}A        & x^*_B.
  \end{array}
\end{equation}
\begin{definition}[Tableau, reduced costs]
  The matrix \eqref{eq:37} is called the \emph{tableau} of the basis
  $B$. The tableau is \emph{feasible} if $x^*_B\geq0$. The first row of
  the tableau is the $0$-th row. The remaining $m$ rows constitute the
  \emph{system-matrix} of the tableau. 

  The vector $\overline{c}^T_B = c^T- c_B^TA_B^{-1}A$ is called the \emph{reduced cost vector} of $B$. For
  $i \in \{1,\ldots,n\}$ $\overline{c}_B(i)$ is the reduced cost of the
  variable $x_i$. 
\end{definition}
Now back to our concrete example. Is the solution $(1,1,0,0)$ optimal?
The fact that the reduced cost of $x_4$ is negative suggests to make a
\emph{basis change}. We want to include $x_4$ into the basis, since
then the elimination of the variable $x_4$ in the $0$-th row will make
the element in the top-right corner larger and, remembering that this
entry is the negative of the objective function value, we would have
obtained  a
better solution. % We increase the  value of $x_4$ from $0$
%to $\lambda>0$. How do we have to adapt $x_1$ and $x_2$?  We have to
%maintain the following conditions  
%\begin{displaymath}
%  \begin{array}{c}
%    x_1 -2 \lambda =1 \\
%    x_2 + 3\lambda =1
%  \end{array}
%\end{displaymath}
%which implies $x_1 = 1 + 2\lambda$ and $x_2 = 1- 3 \lambda$. We want to increase
%$\lambda$ by the maximum amount. The only restriction that we have to
%respect is that $x_1$ and $x_2$ must remain nonnegative. We could only
%violate the constraint $x_2\geq0$ and this, only if $\lambda>1/3$.
If $x_4$ should enter the basis, then which variable should leave,
$x_1$ or $x_2$? If $x_1$ should leave the basis, then we need to apply
elementary row-operations to turn the gray column below into the first
unit vector and if $x_2$ should leave the basis, then this column should
become the second unit vector. 

\begin{equation}
\label{eq:45}
  \begin{tabularx}{0.5 \textwidth}{AAAB|X}
\rowcolor{white}    0   &      0 &     0  &    -1   &   -5\\ \hline 
    1   &    0  &    -1  &    -2 &      1\\
    0 &     1 &     2&     3  &     1
  \end{tabularx}
\end{equation}
If $x_1$ leaves the basis, then the tableau becomes infeasible. This
is because the first entry of the gray column is negative.  We decide
therefore to make the following basis change $B = \{1,2\} \to B' =\{ 1, 4\}$. To
compute the new tableau for $B'$ we multiply the last row by $1/3$ 
and  add $2$-times the last row to the second
and add the last row to the first row of the tableau, to obtain 
\begin{equation}
  \label{eq:39}
  \begin{array}{cccc|c}
       0&   1/3&   2/3&     0& -14/3 \\ \hline
       1&    2/3&   1/3&     0&   5/3\\
     0&   1/3&   2/3&     1&   1/3
   \end{array}
 \end{equation}
%
The next lemma shows that the above is indeed the tableau of $B = \{1,4\}$.
\begin{lemma}
  \label{lem:13}
  Let $B\subset\{1,\ldots,n\}$ be an index set %and let  linear program $\min\{c^Tx \colon x \in \setR^n,
%  \, Ax=b, \, x\geq0\}$ 
and    consider a matrix 
  \begin{equation}
    \label{eq:38}
    \begin{array}{c|c}
      d^T & \beta \\ \hline 
      Q & q 
    \end{array}
  \end{equation}
  where 
  \begin{enumerate}[i)]
  \item $[Q|q]$ can be obtained from $[A|b]$ via a series
      of elementary row operations 
  \item $[d^T|\beta] = [c^T | 0] + v^T$, where $v^T$ is in the row-span
    of $[A|b]$ 
  \item $d_B=0$ 
  \item $Q_B = I_m$,
  \end{enumerate}
  then $B$ is a basis and \eqref{eq:38} is the tableau of $B$.
\end{lemma}



%The general description of we have done so far is the following. We
%consider the system 
%\begin{displaymath}
%  \begin{array}{c}
%  c^Tx = 0 \\ \hline 
%  Ax = b
%\end{array}
%\end{displaymath}
%and a basis $B$. Then we multiply $A$ by $A_B^{-1}$. The result is
%that $(A_B^{-1}A)_B = I_m$. Let $B = \{j_1,\ldots,j_m\}$ be the basis with
%$j_1<\cdots<j_m$. The $i$-th row of $A_B^{-1}A$ has a $1$ in the $j_i$-th
%column.  Next we subtract for each $j_i \in B$,  $c(j_i)$ times the  $i$-th
%row of $(A_B^{-1}A)_B$ from the row above the line to obtain 
%\begin{equation}
%  \label{eq:34}
%  \begin{array}{c}
%    (c^T - c_B^T A_B^{-1}A) x = -c_B^T x^*_B \\ \hline 
%  A_B^{-1} A x = x^*_B   
%  \end{array}
%\end{equation}


\section{The dual linear program}
\label{sec:dual-linear-program}

Our guess is that the tableau~\eqref{eq:39} is optimal. To confirm
this suspicion, we introduce the concept of the dual
linear program. 

\begin{definition}[Dual of a linear program in standard form]
  \label{def:15}
  Let $\min\{c^Tx \colon x \in \setR^n, \, Ax=b, \, x\geq0\}$ be a linear
  program in standard form. The linear program $\max\{ b^Ty \colon y \in
  \setR^m, \,A^Ty \leq  c\}$ is  the \emph{dual linear program} of
  $\min\{c^Tx \colon x \in \setR^n, \, Ax=b, \, x\geq0\}$.  The linear program
   $\min\{c^Tx \colon x \in \setR^n, \, Ax=b, \, x\geq0\}$ is  the
  \emph{primal linear program}.  
\end{definition}

\begin{example}
  \label{sec:dual-linear-program-1}
  The dual linear program of \eqref{eq:40} is 
  \begin{equation}
    \label{eq:41}
    \begin{matrix}
      \max \quad &  3 y_1 & + & 2 y_2 & & \\
      4 y_1 & + & y_2 & \leq & 4 \\
      3 y_1 & + & y_2 & \leq & 4\\
      2 y_1 & + & y_2 & \leq & 3\\
      y_1 & +& y_2 & \leq & 2
    \end{matrix}
  \end{equation}
\end{example}

In our example of the simplex method above, we have arrived at a basis
$B = \{1,4\}$ whose reduced cost vector did not have a negative
entry. The next theorem reveals  that such a basis, if it is
feasible, also yields an optimal solution. 


\begin{theorem}[Weak duality]
  \label{thr:15}
  Let $x^*$ and $y^*$ be feasible solutions of the primal and dual
  respectively, then $c^Tx^* \geq b^Ty^*$. 
\end{theorem}

\begin{proof}
  We have 
  \begin{equation}
    \label{eq:42}
    b^T y^* = ({x^*}^TA^T) y^* = {x^*}^T (A^T y^*) \leq {x^*}^T c = c^Tx^*  % = c^T y^* 
  \end{equation}
 where the inequality in \eqref{eq:42} follows from the fact that
 $x^*\geq0$ and $A^T y^* \leq c$. \qed  
\end{proof}


\begin{lemma}
  \label{lem:14}
  Let $x^*$ be a basic solution associated to the basis $B$. 
  If the reduced cost vector of $B$ is non-negative and if
  $x^*_B\geq0$, then $x^*$ is an optimal basic feasible solution. 
\end{lemma}

\begin{proof}
  We show that $y^* = {A_B^{-1}}^T  c_B$ is a feasible solution of the
  dual with objective value $b^Ty^* = c^Tx^*$. The assertion then
  follows from weak duality. We have 
  \begin{eqnarray*}
    A^T y^* & = & A^T {A_B^{-1}}^T  c_B \\
         & \leq & c,
  \end{eqnarray*}
  since $c^T - c_B^T A_B^{-1}A \geq0$. This means that $y^*$ is a
  feasible dual solution. Its objective value 
  is
  \begin{eqnarray*}
    b^Ty^* & = & {y^*}^T A x^* \\
        & = & c_B^T A_B^{-1} A x^* \\
        & = & c_B^T x^*_B\\
        & = & c^Tx^*
  \end{eqnarray*}
  and the assertion follows. \qed 
\end{proof}

This shows that we have found in~\eqref{eq:39}  an optimal basic
feasible solution $(5/3,0,0,1/3)$.  
\begin{definition}[Optimal basis]
  \label{def:2}
  A basis $B$ with $\overline{c}^T = c^T - c_B^T A_B^{-1}A \geq0$
  and with $A_B^{-1}b \geq0$ is   called an \emph{optimal basis}. 
\end{definition}
%
We have seen above that an optimal basis $B$ yields an optimal basic
feasible solution $x^*$ with $x^*_B = A_B^{-1}b$ and
$x^*_{\overline{B}} = 0$. 


\section{A larger example}
\label{sec:larger-example}

Before we describe the simplex method in full generality, we inspect
another example. We consider the linear program 


\begin{equation}
\label{simpleq:44}
  \begin{array}{c}
  \begin{matrix}
    \min \quad &  -5x_1 & - &  3 x_2 & - & 2 x_3 & &    \\
        & 2x_1 &+& 3x_2 &+& x_3 &\leq&5 \\
        & 4x_1 &+& x_2 &+& 2x_3 &\leq& 9 \\
        & 3 x_1 &+& 4 x_2 &+& 2 x_3 &\leq& 10 \\

       \end{matrix}\\
       x_1,x_2,x_3\geq0
\end{array}
\end{equation}

We add variables $x_4,x_5$ and $x_6$ and transform this linear program
in standard form
\begin{equation}
\label{eq:44}
  \begin{array}{r}
  \begin{array}{cccccccccccccc}
    \min \quad &   -5x_1 & - &  3 x_2 & - & 2 x_3  &   &    & & & & & &      \\
               & 2x_1    & + &   3x_2 & + & x_3    & + & x_4 & & & & &
               = & 5  \\
        & 4x_1 &+& x_2 &+& 2x_3 & & & + & x_5  & &   &=& 9 \\
      & 3 x_1 &+& 4 x_2 &+& 2 x_3 & & & & & & + & x_6= & 10 
       \end{array}\\
       x_1,x_2,x_3,x_4,x_5,x_6\geq0
\end{array}
\end{equation}
%
We start with the feasible basis $B = \{4,5,6\}$ and the corresponding
tableau, where the red-columns correspond to basis elements 
\begin{displaymath}
  \begin{tabularx}{0.5\textwidth}{AAARRR|X}   
\rowcolor{white}    -5&-4&-3&0&0&0&0\\\hline 
   2&3&1&1&0&0&5\\
   4&1&2&0&1&0&9\\
   3&4&2&0&0&1&10 
 \end{tabularx}
\end{displaymath}
The reduced cost of $x_1$ is negative, we decide to let $x_1$ enter
the basis. We now want to transform the gray column into a unit
vector such that the tableau is still feasible. The red columns
are the basis-columns. 
\begin{equation}
\label{eq:1}
  \begin{tabularx}{0.5\textwidth}{BAARRR|X}   
\rowcolor{white}    -5&-4&-3&0&0&0&0\\\hline 
   2&3&1&1&0&0&5\\
   4&1&2&0&1&0&9\\
   3&4&2&0&0&1&10 
 \end{tabularx}
\end{equation}
We now need to decide, which element  should leave $B = \{4,5,6\}$. If
it should be $4$, then the gray column should be transformed into
$(1,0,0)^T$. Let us see what happens if we apply elementary
row-operations on the system matrix of the tableau such that the first
column becomes $(1,0,0)^T$. 
\begin{displaymath}
  \begin{tabularx}{0.5\textwidth}{RAAARR|X}   
\rowcolor{white}    0&-4&-3&0&0&0&25/2\\\hline 
   1&3/2&1/2&1/2&0&0&5/2\\
   0&-5&0&-2&1&0&-1\\
   0&2&1/2&-3/2&0&1&1/2 
 \end{tabularx}
\end{displaymath}
The result is the tableau of $\{1,5,6\}$, the objective function value
has improved but the basic solution of $\{1,5,6\}$ is not feasible,
since $x^*(5) = -1$.


What went wrong? Let us abstract from the specific situation and
suppose that we have the following situation. We have a feasible basis
$B = \{B_1,\ldots,B_m\}$ and a tableau  of the basis $B$. The reduced cost
of $j$ are negative and we want to add $j$ to $B$. The picture below
represents this situation, where the $j$-th column is the gray
column. This column is also called the \emph{pivot column}. 

\begin{center}
  \begin{tabularx}{0.5\textwidth}{ABAAA|X}   
       \rowcolor{white} & $\overline{c}(j)$ & &  & & $-c^Tx^*$ \\\hline
       &$ u(1)$  & & & & $x^*(B_1)$  \\
       &$\vdots$  & & & & $\vdots$  \\
       &$ u(i)$  & & & & $x^*(B_i)$  \\
       &$\vdots$  & & & & $\vdots$  \\
       &$ u(m)$  & & & & $x^*(B_m)$  \\
 \end{tabularx}
\end{center}
%
If we now decide to let $B_i$ leave the basis, then we transform the
gray column into the $i$-th unit vector with elementary row-operations
and obtain 

\begin{center}
  \begin{tabular}{cGc|c}   
     \rowcolor{white}   \hspace{2cm} &  $\overline{c}(j)$
     &\hspace{2cm}  & 

$-c^Tx^*$ \\ \hline
&       $ 0 $  &&   $x^*(B_1) -  u(1) \cdot \frac{x^*(B_i)}{u(i)} $  \\
&       $ \vdots $ && $\vdots$  \\
&       $ 0 $  && $x^*(B_{i-1})- u(i-1) \cdot \frac{x^*(B_i)}{ u(i)}    $  \\
&       $ 1$  && $x^*(B_i) / u(i)$  \\
&       $ 0 $  &&  $x^*(B_{i+1})-   u(i+1)\cdot  \frac{x^*(B_i)}{ u(i)}  $  \\
&       $ $  && $\vdots$  \\
&       $ 0 $  &&   $x^*(B_m) - u(m) \cdot \frac{x^*(B_i)}{u(i)}$  
 \end{tabular}
\end{center}

Then we eliminate the $\overline{c}(j)$ by adding $-\overline{c}(j)$
times the $i$-th row of the system matrix to the zeroth row to obtain 



\begin{center}
  \begin{tabular}{cGc|c}   
      \rowcolor{white}   \hspace{2cm} &  $0$
        &\hspace{2cm}  &   $-c^Tx^* - \overline{c}(j)\cdot
       x^*(B_i) / u(i)$ \\ \hline 
   &       $ 0 $  &&   $x^*(B_1) -  u(1) \cdot \frac{x^*(B_i)}{u(i)} $  \\
   &     $ \vdots $ && $\vdots$  \\
   &    $ 0 $  &&  $x^*(B_{i-1})- u(i-1) \cdot \frac{x^*(B_i)}{ u(i)}    $  \\
   &    $ 1$  && $x^*(B_i) / u(i)$  \\
   &    $ 0 $  && $x^*(B_{i+1})-   u(i+1)\cdot  \frac{x^*(B_i)}{ u(i)}  $  \\
   &    $ \vdots$  && $\vdots$  \\
   &    $ 0 $  &&   $x^*(B_m) - u(m) \cdot \frac{x^*(B_i)}{u(i)}$  
 \end{tabular}
\end{center}

Let us deduce a set of conditions, under which these sequence of
operations are allowed and such that the obtained tableau is feasible. 


\begin{enumerate}[a)]
\item $u(i) >0$,\\ since if $u(i) \leq0$, then we either cannot divide by $u(i)$
  or the basis is not feasible since $x^*(B_i) / u(i) <0$.  \label{item:1}
\item $x^*(B_i) / u(i) = \min\{ x^*(B_j) / u(j) \colon B_j \in B, \, u(j)
  >0\}$, \\
  since the condition $x^*(B_j) - u(j) \cdot \frac{x^*(B_i)}{u(i)} \geq0$ is
  satisfied for each $B_j \in B$ if and only if the condition holds. \label{item:2}
\end{enumerate}


Looking back at our example~\eqref{eq:1} we see that
condition~\ref{item:2}) was violated. With $B_1=4,B_2=5,B_3=6$ we see
that $x^*(B_1)/u(1) = 5/2$ and this is larger than $x^*(B_2)/u(2) =
9/4$, which is the minimum of the ratios. We therefore let $B_2=5$
leave the basis and obtain 

\begin{displaymath}
  \begin{tabularx}{0.6\textwidth}{RAARAR|X}   
    \rowcolor{white}   
     0& -11/4&  -1/2&     0&   5/4&     0&  45/4 \\ \hline 
     0&   5/2&     0&     1&  -1/2&     0&   1/2 \\
     1&   1/4&   1/2&     0&   1/4&     0&   9/4 \\
     0&  13/4&   1/2&     0&  -3/4&     1&  13/4
  \end{tabularx}
\end{displaymath}
%
To arrive then at the feasible tableau of the basis $B' = \{1,4,6\}$
we still need to swap the first and second row of the system matrix: 
\begin{displaymath}
  \begin{tabularx}{0.6\textwidth}{RAARAR|X}   
    \rowcolor{white}   
     0& -11/4&  -1/2&     0&   5/4&     0&  45/4 \\ \hline 
     1&   1/4&   1/2&     0&   1/4&     0&   9/4 \\
     0&   5/2&     0&     1&  -1/2&     0&   1/2 \\
     0&  13/4&   1/2&     0&  -3/4&     1&  13/4
  \end{tabularx}
\end{displaymath}

\subsection{One iteration of the simplex method} 
\label{sec:one-iter-simpl}

 We start with a feasible basis  $B \subseteq\{1,\ldots,n\}$ and an associated tableau 
\begin{equation}
  \label{eq:2}
  \begin{array}{c|c}
    \overline{c}^T  & -c_B^Tx^*_B \\ \hline 
    A_B^{-1}A        & x^*_B.
  \end{array}
\end{equation}



\noindent 
{\bfseries The simplex algorithm} starts with a feasible tableau and
keeps iterating the procedure below.



\begin{enumerate}
\item If $\overline{c}\geq0$, then {\bf output optimal basis $B$}
  \label{item:3}
\item Otherwise, let $j$ be an index with $\overline{c}(j)<0$ and let
  $u = A_B^{-1}A^j$ be the $j$-th column of the system matrix of the
  tableau. If $u\leq0$, then {\bf output $-\infty$}, the linear program is
  unbounded. \label{item:7}
\item Otherwise compute the ratios $x^*(B_i)/ u(i)$ for all $i =
  1,\ldots,m$ with $u(i)>0$. Let $i^*$ be an index where this minimum is
  attained. The index $B_{i^*}$ leaves the basis and $j$ enters the
  basis, i.e., the new basis is $B' = B \setminus\{B_{i^*}\} \cup \{j\}$. 
\item Update the tableau to 
  \begin{displaymath}    
  \begin{array}{c|c}
    c^T - c_{B'}^TA_{B'}^{-1}A  & -c_{B'}^Tx^*_{B'} \\ \hline 
    A_{B'}^{-1}A        & x^*_{B'}.
  \end{array}
\end{displaymath}
\end{enumerate}
%
We need to justify that the we correctly assert in step~\ref{item:7} that the linear
program is unbounded if $u\leq0$. 
\begin{theorem}
  \label{thr:1}
  If $\overline{c}(j)<0$  and $u\leq0$, then the linear program is
  unbounded. 
\end{theorem}

\begin{proof}
  Consider the vector $d$ with $d_B = -u$ and $d(j) = 1$, which is
  zero at all other components. Since $A d=0$ and since $d\geq0$, we
  have that for each feasible $x^*$ and for each $\lambda\geq0$, also the
  point $x^* + \lambda d$ is feasible. What is the objective function value
  of $d$? 
  \begin{eqnarray*}
    c^Td & = & c(j) - c_B^T A_B^{-1} A \\
        & = & \overline{c}(j) \\
        & < & 0. 
  \end{eqnarray*}
  This means that  the objective function value of $x^* +
  \lambda d$ which is $c^Tx^* + \lambda c^Td$ can be made as small as we wish by
  increasing $\lambda$, without becoming infeasible.  The linear program is
  unbounded. 
  \qed 
\end{proof}



Let us continue our example. We decide to bring $2$ into the
basis, since $\overline{c}(2) = -11/4 <0$. 

\begin{displaymath}
  \begin{tabularx}{0.8\textwidth}{RBARAR|X}   
    \rowcolor{white}   
     0& -11/4&  -1/2&     0&   5/4&     0&  45/4 \\ \hline 
     1&   1/4&   1/2&     0&   1/4&     0&   9/4 \\
     0&   5/2&     0&     1&  -1/2&     0&   1/2 \\
     0&  13/4&   1/2&     0&  -3/4&     1&  13/4
  \end{tabularx}
\end{displaymath}

We compute the ratios: $9, 1/5, 1$ and decide that the second
 element of $\{1,4, 6\}$ should leave the basis, i.e., we want to have
 the second unit vector in the pivot column. 


\begin{displaymath}
  \begin{tabularx}{0.8\textwidth}{RRBAAR|X}   
    \rowcolor{white}       
    0&      0&   -1/2&  11/10&   7/10&      0&   59/5\\ \hline 
    1&      0&    1/2&  -1/10&   3/10&      0&   11/5\\
    0&      1&      0&    2/5&   -1/5&      0&    1/5\\
    0&      0&    1/2& -13/10&  -1/10&      1&   13/5\\
  \end{tabularx}
\end{displaymath}


After one more iteration we obtain an obtimal feasible tableau

\begin{displaymath}
  \begin{tabularx}{0.8\textwidth}{ARRAAR|X}   
    \rowcolor{white}       
        1&    0&    0&    1&    1&    0&   14\\ \hline 
        0&    1&    0&  2/5& -1/5&    0&  1/5\\
        2&    0&    1& -1/5&  3/5&    0& 22/5\\
       -1&    0&    0& -6/5& -2/5&    1&  2/5\\
  \end{tabularx}
\end{displaymath}





\subsection{What  next?}
\label{sec:further-questions}
 The two main questions that we
are next going to deal with   are the following.  
\begin{enumerate}[I)]
\item How do we find an initially feasible basis? \label{item:8}
\item If we start from a feasible basis, does the simplex method
  terminate? \label{item:9} 
\end{enumerate}




\section{Exercises}

\begin{enumerate}
\item Transform the following linear program into equation standard
  form \label{item:4}
  \begin{displaymath}
    \begin{array}{c}
      \begin{matrix}
        \max \quad &  3x_1& +& 4x_2& +& 2 x_3& & \\
             &  &  & x_2  & + & 2x_3 &\geq & 4 \\ 
             &  2x_1 &+  & x_2 &+  &3x_3 &\leq & 10\\
             &      &   & x_2   &   &    & \leq & 0 
      \end{matrix}
    \end{array}
  \end{displaymath}
\item  \label{item:5}    Provide an example of  a basic solution that can be
  associated to two   different bases. %, thereby showing that the
%  converse of Lemma~\ref{lem:3} is false. 
 
\item This exercise should give you a sense of exponential growth as
  one experiences it with the running time of the naive linear
  programming algorithm. Have a look at this program. What is its
  running time?  Experiment with different inputs and determine the
  largest exponent for which the program needs less running time than
  a day on your computer. Make an estimation on how many rows and
  variables a linear program could have such that the naive algorithm
  for linear programming finishes within a day on your computer.
  \label{item:6} Estimate how many rows and variables a linear needs
  to have such that the naive algorithm does not terminate on your
  computer before our sun will stop to shine. 
\begin{verbatim}
#include <iostream>
#include <cmath>

using namespace std;  

int main()
{
  
  double n; 
  
  cout << "\n Enter the exponent" ;
  cin >> n ;
  long i,j;
  j=0;
  for (i=1; i<= pow(2.0,n); i++)
    j = j+i;
  
  cout << j;

}


\end{verbatim}


\item Prove Lemma~\ref{lem:13}. 
\end{enumerate}




\bibliographystyle{abbrv}

\bibliography{mybib,papers,books,my_publications}

%%% Local Variables: 
%%% mode: latex
%%% TeX-master: "lecture"
%%% End: 

%%\chapter{Termination, Cycling, and Degeneracy}
\label{cha:cycling-degeneracy}


We now deal first  with the question, whether the simplex method
terminates. The quick answer is no, if it is implemented in a careless
way. Notice that we have left it open to choose the entering variable
(there could be several columns whose reduced cost are less than zero)
and we also did not specify how to choose the exiting variable (there
could be several basis elements $B_i$ with $u(B_i)>0$ and
$x^*(B_i)/u(B_i)$ minimal to choose from. \emph{Pivoting rules}
specify the alternative to choose from and there are pivoting rules,
which prevent the simplex algorithm from \emph{cycling}. 



The following is an example of cycling, see also
\cite[p.~104]{BertsimasTsitsiklis97}. 



\begin{displaymath}
  \begin{tabularx}{0.7\textwidth}{AAAARRR|X}   
    \rowcolor{white}  
    -3/4&   20& -1/2&    6&    0&    0&    0&    3\\ \hline
    1/4&   -8&   -1&    9&    1&    0&    0&    0\\
    1/2&  -12& -1/2&    3&    0&    1&    0&    0\\
    0&    0&    1&    0&    0&    0&    1&    1\\
  \end{tabularx}
  %\quad \quad \quad   B = \{5,6,7\} 
\end{displaymath}

         

\begin{displaymath}
  \begin{tabularx}{0.7\textwidth}{RAAAARR|X}   
    \rowcolor{white}  
    0&   -4& -7/2&   33&    3&    0&    0&    3\\ \hline
    1&  -32&   -4&   36&    4&    0&    0&    0\\
    0&    4&  3/2&  -15&   -2&    1&    0&    0\\
    0&    0&    1&    0&    0&    0&    1&    1\\
\end{tabularx}
%\quad \quad \quad   B = \{1,6,7\} 
\end{displaymath}


\begin{displaymath}
  \begin{tabularx}{0.7\textwidth}{RRAAAAR|X}   
    \rowcolor{white}  
     0&     0&    -2&    18&     1&     1&     0&     3\\ \hline
     1&     0&     8&   -84&   -12&     8&     0&     0\\
     0&     1&   3/8& -15/4&  -1/2&   1/4&     0&     0\\
     0&     0&     1&     0&     0&     0&     1&     1\\
\end{tabularx}
%\quad \quad \quad   B = \{1,2,7\} 
\end{displaymath}


\begin{displaymath}
  \begin{tabularx}{0.7\textwidth}{ARRAAAR|X}   
    \rowcolor{white}  
   1/4&     0&     0&    -3&    -2&     3&     0&     3\\ \hline
   -3/64&     1&     0&  3/16&  1/16&  -1/8&     0&     0\\
   1/8&     0&     1& -21/2&  -3/2&     1&     0&     0\\
  -1/8&     0&     0&  21/2&   3/2&    -1&     1&     1\\
\end{tabularx}
%\quad \quad \quad   B = \{2,3,7\} 
\end{displaymath}

\begin{displaymath}
  \begin{tabularx}{0.7\textwidth}{AARAAAR|X}   
    \rowcolor{white}  
 -1/2&   16&    0&    0&   -1&    1&    0&    3\\ \hline
 -5/2&   56&    1&    0&    2&   -6&    0&    0\\
 -1/4& 16/3&    0&    1&  1/3& -2/3&    0&    0\\
  5/2&  -56&    0&    0&   -2&    6&    1&    1\\
\end{tabularx}
%\quad \quad \quad   B = \{3,4,7\} 
\end{displaymath}


\begin{displaymath}
  \begin{tabularx}{0.7\textwidth}{AAARRAR|X}   
    \rowcolor{white}  
 -7/4&   44&  1/2&    0&    0&   -2&    0&    3\\ \hline
  1/6&   -4& -1/6&    1&    0&  1/3&    0&    0\\
  -5/4&   28&  1/2&    0&    1&   -3&    0&    0\\
    0&    0&    1&    0&    0&    0&    1&    1\\
\end{tabularx}
%\quad \quad \quad   B = \{4,5,7\} 
\end{displaymath}


\begin{displaymath}
  \begin{tabularx}{0.7\textwidth}{AAAARRR|X}   
    \rowcolor{white}  
 -3/4&   20& -1/2&    6&    0&    0&    0&    3\\ \hline
  1/4&   -8&   -1&    9&    1&    0&    0&    0\\
  1/2&  -12& -1/2&    3&    0&    1&    0&    0\\
    0&    0&    1&    0&    0&    0&    1&    1\\
  \end{tabularx}
%\quad \quad \quad   B = \{5,6,7\} 
\end{displaymath}

In fact we have not made any progress in the objective function value
during all these pivots. Inspecting the pivot steps, we see that the
problem lies in the fact that we have basic feasible solutions which
are zero at some of the basic components. If this is never the case,
then the simplex method terminates, as we show now. 


Recall that a basis $B$ with reduced costs $\overline{c}\geq0$ is
optimal, see Lemma~\ref{lem:14}. There is a weak converse to this
lemma. Before we state it, we define the notion of degeneracy. 

\begin{definition}[Degenerate basic solution]
  \label{def:1}
  A basic solution $x^*$ associated to the basis $B$ is \emph{degenerate}, if
  there exists an index $i \in B$ with $x^*(i) =0$. If $x^*(i) \neq0$ for
  all $i \in B$, then $x^*$ is called \emph{non-degenerate}. 
\end{definition}


\begin{lemma}
  \label{lem:1}
  Let $x^*$ be a non-degenerate basic feasible solution associated to
  the basis $B= \{B_1,\ldots,B_m\}$. If there exists an index $i \in \{1,\ldots,n\}$ with
  $\overline{c}(i)<0$, then $x^*$ is not optimal. 
\end{lemma}

\begin{proof}
  Suppose that the reduced cost of variable $j$ is negative, 
  $\overline{c}(j)<0$ and inspect the tableau of $B$ 
  \begin{displaymath}
    \begin{array}{c|c}
      c^T - c_B^TA_B^{-1}A & -c_B^Tx^*_B \\ \hline 
      A_B^{-1}A        & x^*_B.
    \end{array}
  \end{displaymath}
Let $u(1),\ldots,u(m)$ be the components of the $j$-th column of the
system matrix of the tableau. If $u(i)\leq0$ for all $i$, then the linear program is
unbounded and the assertion follows. Otherwise let $K = \{ i \colon u(i) >
0\}$ and let $i^*\in K$ be an index where the minimum 
$\min \{x^*(B_i)/u(i) \colon i \in K\}$ is attained. The new objective
function value after the pivot step is 
\begin{equation}
  \label{eq:3}
  c^Tx^* + \overline{c}(j) x^*(B_{i^*})/u(i^*)
\end{equation}
and since $\overline{c}(j)<0$, this is smaller than $c^Tx^*$. \qed
  
\end{proof}



This Lemma also shows that the simplex method terminates if the basic
feasible solutions are all non-degenerate. 

\begin{theorem}
  \label{thr:2}
  If all basic feasible solutions are non-degenerate, then the simplex
  method terminates, regardless of the choice among the alternatives
  among entering and leaving variables. 
\end{theorem}


\begin{proof}
  The proof of the previous lemma, in particular equation~\eqref{eq:3}
  shows that  the objective function value strictly improves each
  iteration of the simplex method. Each tableau is uniquely defined by
  a basis $B$. We never re-visit a tableau, since the objective
  function is strictly increasing. The number of bases is bounded by
  $\binom{n}{m}$. The assertion follows. \qed
\end{proof}




\subsubsection*{Perturbation}

The termination argument for the simplex algorithm for linear programs,
who do not have degenerate basic feasible solutions, is very simple and
nice and in fact one can twist it into a pivot-rule for the
simplex-algorithm which makes the simplex algorithm terminate.  

We start with the  linear program 
\begin{equation}
  \label{eq:4}
  \min\{c^Tx \colon x \in \setR^n, \, Ax=b, \, x\geq0\}
\end{equation}
and assume that $S\subseteq\{1,\ldots,n\}$ is a feasible basis with $A_S
=I_m$. This assumption can be made, since we can replace $Ax = b$ with
$A_S^{-1}Ax = A_S^{-1}b$.

The idea is to transform \eqref{eq:4} into 
into a linear program 
\begin{equation}
  \label{eq:5}
  \min\{c^Tx \colon x \in \setR^n, \, Ax=b', \, x\geq0\}
\end{equation}
such that the following conditions hold. 
\begin{enumerate}[i)]
\item The LP~\eqref{eq:5} does not have degenerate basic solutions. \label{item:11}
\item If $B$ is an infeasible basis of \eqref{eq:4}, then $B$ is also an
  infeasible basis of~\eqref{eq:5}. \label{item:12}
\item The LP~\eqref{eq:5} is feasible. \label{item:13}
\end{enumerate}



\begin{lemma}
  \label{lem:2}
  The simplex algorithm terminates on the linear
  program~\eqref{eq:5}. If it 
  asserts that the linear program~\eqref{eq:5} is unbounded, then the
  linear program~\eqref{eq:4} is also unbounded. If it 
  outputs an optimal basis of~\eqref{eq:5} then this basis is also
  optimal for~\eqref{eq:4}. 
\end{lemma}

\begin{proof}
  
  
  It follows from Theorem~\ref{thr:2} and the property~\ref{item:11})
  that the simplex method terminates on the linear
  program~\eqref{eq:5}. If it asserts that~\eqref{eq:5} is unmounded,
  then there exists a $d\geq0$ with $Ad=0$ and $c^Td<0$. This means that
  the linear program~\eqref{eq:4} is also unbounded. 
  
  Suppose therefore that the simplex method on~\eqref{eq:5} terminates
  by outputting a 
  basis $\wt{B}$. This basis 
  $\wt{B}$ is a feasible basis of~\eqref{eq:4} by
  property~\ref{item:12}). The reduced costs of this basis are the
  same for the linear programs~\eqref{eq:4}  and~\eqref{eq:5} which
  implies that $\wt{B}$ is also an optimal basis of \eqref{eq:5}.
  \qed 
\end{proof}


Now we describe the construction of \eqref{eq:5}, more precisely the
construction of $b'$, which is 
\begin{equation}
  \label{eq:6}
  b' = b+ \mat{\varepsilon\\\varepsilon^2 \\\vdots \\\varepsilon^m}
\end{equation}
for some small $\varepsilon>0$. Since $\varepsilon>0$ the basic solution $x^*$ associated to
$S$ with 
\begin{displaymath}
  x^*_S=b+\mat{\varepsilon \\ \vdots\\ \varepsilon^m}\text{ and }x^*_{\overline{S}}=0 
\end{displaymath}
  is  a basic
feasible solution of~\eqref{eq:5} and thus~\ref{item:13}) follows. We show that we can
choose an $\varepsilon>0$ such that \ref{item:12}) holds. Suppose that $B$ is an
infeasible basis. This means that there exists an entry $i$ of $B$
such that $(A_B^{-1} \, b)(i) < 0$. Clearly, we can choose $\varepsilon_B>0$
small enough such that 
\begin{equation}
  \label{eq:7}
  \left[A_B^{-1}\,\left(b+\mat{\varepsilon\\\vdots\\\varepsilon^m}\right)\right](i)<0
\end{equation}
and the basis $B$ is still infeasible for the linear
program~\eqref{eq:5}. The most interesting property to infer is
property~\ref{item:11}). Let $B$ be any basis and consider the vector
\begin{equation}
  \label{eq:8}
  A_B^{-1} \left(b + \mat{\varepsilon\\\vdots\\\varepsilon^m}\right) 
\end{equation}
which determines the components of $x^*_B$ of the basic solution. If
we treat $\varepsilon$ as a variable, then the $i$-th component of this
vector~\eqref{eq:8} is a polynomial in $\varepsilon$ which we denote by
$p_{iB}(\varepsilon)$. Each of the polynomials $p_{iB}(\varepsilon)$ is unequal to the
zero polynomial,
see exercise~\ref{item:14} which implies that there are only a finite
(at most $m$) real-roots of $p_{iB}$. Consequently, there exists an
$\varepsilon_{iB}>0$ for each $p_{iB}$ such that $p_{iB}(\varepsilon) \neq 0$ for each $\varepsilon
\in ]0, \varepsilon_{iB}[$. 
Thus if we choose $\varepsilon>0$ such that it is smaller than all $\varepsilon_B$ and
all $\varepsilon_{iB}$ from above, then the perturbed problem with 
\begin{displaymath}
  b' = b + \mat{\varepsilon\\\vdots\\\varepsilon^m}
\end{displaymath}
satisfies the properties~\ref{item:11}-\ref{item:13}) from above. 


\subsection{Symbolic computation with $\varepsilon$} 
\label{sec:symb-comp-with}

The good news is that we do not have to calculate with a concrete
$\varepsilon>0$ at all. Let us understand which variable would leave the basis
if we would calculate with a very small $\varepsilon$. 

Recall that we consider the index set $K = \{ i \colon u(i)>0\}$, where
$u$ is the pivot column of the tableau 
\begin{displaymath}
  \label{c:eq:37}
  \begin{array}{c|c}
    c^T - c_B^TA_B^{-1}A & -c_B^Tx^*_B \\ \hline 
    A_B^{-1}A        & x^*_B.
  \end{array}
\end{displaymath}
The leaving variable is  $B_i$, where $i \in K$ minimizes the expression
$x^*(B_i) / u(i), \, i \in K$. We have 
\begin{displaymath}
  x^*(B_i) / u(i) = [(A_B^{-1} b)(i) + (A_B^{-1} \mat{\varepsilon\\\vdots\\\varepsilon^m}) (i)] / u(i). 
\end{displaymath}

\begin{definition}[Lexicographic order]
  Let $u,v \in \setR^n$ be vectors. We have $u<_{lex}v$, if $u=v$ or if
  the first nonzero component of $u-v$ is strictly negative. 
\end{definition}

Since $\varepsilon>0$ is arbitrarily  small we see that we need to find $i$
such that the vector $ [ (A_B^{-1} b)(i) | (A_B^{-1})_i] / u(i)$ is lexicographically
minimal. This is  the {\bf  lexicographic
pivoting rule}: 


\begin{enumerate}
\item If $\overline{c}\geq0$, then {\bf output optimal basis $B$}
  \label{c:item:3}
\item Otherwise, let $j$ be an index with $\overline{c}(j)<0$ and let
  $u = A_B^{-1}A^j$ be the $j$-th column of the system matrix of the
  tableau. If $u\leq0$, then {\bf output $-\infty$}, the linear program is
  unbounded. \label{c:item:7}
\item Otherwise compute the 
  vectors $[ (A_B^{-1} b) (i) | (A_B^{-1})_i] / u(i)$ 
  for all $i =
  1,\ldots,m$ with $u(i)>0$. Let $i^*$ be an index for which this vector
  is lexicographically smallest. 
  The index $B_{i^*}$ leaves the basis and $j$ enters the  basis, i.e., the new basis is $B' = B \setminus\{B_{i^*}\} \cup \{j\}$. 
\item Update the tableau to 
  \begin{displaymath}    
  \begin{array}{c|c}
    c^T - c_{B'}^TA_{B'}^{-1}A  & -c_{B'}^Tx^*_{B'} \\ \hline 
    A_{B'}^{-1}A        & x^*_{B'}.
  \end{array}
\end{displaymath}
\end{enumerate}


\begin{theorem}
  \label{thr:3}
  The simplex method terminates, if the leaving variable is chosen
  according to the lexicographic pivot rule above. 
\end{theorem}


\subsection{Bland's rule}

The following is an alternative pivoting rule which avoids cycling,
see, e.g.~\cite{Chvatal83}. 

\begin{quote}
  During an iteration of the simplex method
  (Chapter~\ref{sec:one-iter-simpl}), find a smallest $j$ such  
  that $\overline{c}(j)$ is less than $0$. Let this $j$ enter the
  basis. 
  
  Out of all the indices $i$ that may leave the basis, choose the
  smallest one. 
\end{quote}

%

% Let us apply this rule for our example. We start with the tableau 
% \begin{displaymath}
%   \begin{tabularx}{0.7\textwidth}{AAAARRR|X}   
%     \rowcolor{white}  
%     -3/4&   20& -1/2&    6&    0&    0&    0&    3\\ \hline
%     1/4&   -8&   -1&    9&    1&    0&    0&    0\\
%     1/2&  -12& -1/2&    3&    0&    1&    0&    0\\
%     0&    0&    1&    0&    0&    0&    1&    1\\
%   \end{tabularx}
%   %\quad \quad \quad   B = \{5,6,7\} 
% \end{displaymath}
% and let $1$ enter the basis. The ratios $x(i)/u(i)$ for $u(i)>0$ are
% minimized at $i=1,2$. 





\subsection*{Exercises}

\begin{enumerate}
\item Provide an example of a degenerate basic feasible solution $x^*$
  and an associated basis $B$ which is optimal but whose reduced cost
  vector $\overline{c}^T =c^T - c_B^T A_B^{-1}A$ does not satisfy
  $\overline{c}\geq0$.\label{item:10} 
\item Show that  $p_{iB}(\varepsilon)\neq 0$ holds for each $i\in \{1,\ldots,m\}$ and
  each basis $B$. \label{item:14}
\item Provide an example of a valid perturbation, as described above,
  such that a basis $B$ which is feasible for the original
  LP~\eqref{eq:4} is infeasible for the problem~\eqref{eq:5}. 
  \label{item:15} 
\end{enumerate}




\bibliographystyle{abbrv}

\bibliography{mybib,papers,books,my_publications}


%%% Local Variables: 
%%% mode: latex
%%% TeX-master: "lecture"
%%% End: 

\chapter{Duality}
\label{cha:duality}



%\section{Weak and strong duality}


Via the termination argument for the simplex algorithm, we can now
prove the duality theorem. We are given a linear program
\begin{equation}
  \label{eq:91}
  \max\{c^T x \colon x \in \setR^n, \, Ax \leq b\}, 
\end{equation}
called the \emph{primal} and its \emph{dual}
\begin{equation}
  \label{d:eq:9}
  \min\{b^Ty \colon y \in \setR^m, \, A^T y = c,\, y\geq0\}.
\end{equation}
%
We again formulate the theorem of weak duality in this setting. 
\begin{theorem}[Weak duality]
  If $x^*$ and $y^*$ are primal and dual feasible solutions respectively, then
  $c^Tx^* \leq b^Ty^*$. 
\end{theorem}

\begin{proof}
  We have $c^Tx^* = {y^*}^T A x^* \leq  {y^*}^T b$. \qed 
\end{proof}

The
strong duality theorem tells us that if there exist feasible primal and
dual solutions, then there exist feasible primal and dual solutions
which have the same objective value. We can prove it with the simplex
algorithm. 


\begin{theorem}
  \label{thr:4}
  If the primal linear program is feasible and bounded, then so is 
  the dual linear program. Furthermore in this case, both linear
  programs have an optimal solution  and the optimal  values coincide. 
\end{theorem}


\begin{proof}
  Suppose first that $A$ has full column rank.  The simplex method
  finds an optimal basis $B$ of~\eqref{eq:91} with $x^*_B$ being an
  optimal feasible solution. At the same time, we have a $\lambda \in
  \R_{\geq 0}^m$ with $\lambda_i = 0$ if $i \notin B$ and $\lambda^TA
  = c^T$. Notice that $\lambda$ is  a feasible solution of the dual linear program~\eqref{d:eq:9}. One has 
  \begin{displaymath}
    c^T x^*_B = \lambda^T A x_B^* = \lambda^T b,
  \end{displaymath}
  which shows the theorem in this case. 

  If $A$ does not have full column rank, then we re-write the linear
  program~\eqref{eq:91} as
  \begin{equation}
    \label{eq:d-1}
    \max\{ c^T (x_1 - x_2) \colon A(x_1-x_2) \leq b, x_1\geq 0, x_2 \geq 0\}.  
  \end{equation}
  There is a dual solution that we partition into three parts
  $\lambda_1,\lambda_2,\lambda_3 \geq 0$. Its dual objective function
  value is $\lambda_1^Tb$. Furthermore 
  \begin{displaymath}
    \lambda_1^TA - \lambda_2 = c^T, \, - \lambda_1^TA - \lambda_3 = -c^T,
  \end{displaymath}
  which together with $\lambda_2,\lambda_3 \geq 0$ implies $
  \lambda_1^TA = c^T$ and $\lambda_2=\lambda_3=0$. This means that $\lambda_1$ is an optimal dual solution. 

\end{proof}


 We can formulate dual linear programs also if the linear program is
 not in inequality standard form. % Consider for example a linear program
%  \begin{equation}
%    \label{eq:92}
%    \max\{c^Tx \colon x \in \setR^n, \, Ax \leq b, \, x\geq0\}. 
%  \end{equation}
%  We transform this into standard form via slack variables $z\geq0$ 
%  \begin{displaymath}  
%    \min\{-c^Tx + 0^Tz  \colon x\in \setR^n,z \in \setR^m, \, Ax +z =b, \, x,z\geq0\}. 
%  \end{displaymath}
%  The dual of this linear program in equation standard form is 
%  \begin{displaymath}
%    \max\{ b^Ty \colon [A \mid I_m]^T y \leq \smat{-c \\ 0}\}
%  \end{displaymath}
%  This can be re-formulated as 
%  \begin{displaymath}
%    \max\{ b^Ty \colon A^T y \leq-c, \, y\leq0\}. 
%  \end{displaymath}
%  Again, this is the same as 
%  \begin{displaymath}
%    \min\{ b^T(-y) \colon A^T (-y) \geq c, \, -y\geq0\} 
%  \end{displaymath}
%  and this finally is equivalent to 
%  \begin{equation}
%    \label{eq:93}
%    \min\{ b^Ty \colon A^T y \geq c, \, y\geq0\}. 
%  \end{equation}
%
 The procedure above can be described as follows. We transform a linear
 program into a linear program in inequality  standard form and construct
 its dual linear program. This dual is then transformed into
 an equivalent  linear program again which is conveniently described. 

 Let us perform such operations  on the dual linear
 program 
 \begin{displaymath}
   \min\{ b^Ty \colon y \in \setR^m, \, A^T y = c, \, y\geq0\}
 \end{displaymath}
 of the primal $\max\{c^T x \colon x \in \setR^n, \, Ax \leq b\}$. 
 We transform it into inequality  standard form 
 \begin{displaymath}
   \begin{array}{rcl}
     \max - b^T y \\
     A^T y & \leq & c \\
     -A^Ty &\leq &-c\\
     -I y &\leq &0. 
   \end{array}
 \end{displaymath}
 The dual linear program of this is 
 \begin{displaymath}
   \begin{array}{rcl}
     \min c^T x_1 - c^T x_2  \\
     A x_1 - A x_2 - x_3 & = & -b \\
    x_1,x_2,x_3\geq0 &&
   \end{array}
 \end{displaymath}
 This is equivalent to 
 \begin{displaymath}
   \begin{array}{rcl}
     \max c^T (x_2-x_1) \\
     A (x_2-x_1) + x_3 & = & b \\
     x_1,x_2,x_3 & \geq & 0
   \end{array}
 \end{displaymath}
 which is equivalent to the primal linear program
 \begin{displaymath}
   \begin{array}{rcl}
     \max c^T x \\
     Ax & \leq & b. \\
   \end{array}
 \end{displaymath}
 Loosely formulated one could say that ``The dual of the dual is the
 primal''. But this, of course, is not to be understood as  a
 mathematical statement.  In any case we can state the following
 corollary. 

 \begin{corollary}
   \label{thr:5}
   If the dual linear program has an optimal solution, then so does the
   primal linear program and the objective values coincide.
 \end{corollary}



We present another example of duality that we will need later on. 
Consider a linear program 
\begin{equation}
  \label{eq:40}
  \begin{array}{rcl}
    \max &c^Tx \\
    Bx &  =&  b\\
    Cx & \leq & d. 
  \end{array}
\end{equation}
After
re-formulation, we obtain 
\begin{displaymath}
  \begin{array}{rcl}
    \max & c^Tx & \\ Bx & \leq &  b \\ -Bx & \leq & -b \\ Cx & \leq &  d
  \end{array}
\end{displaymath}
We can form the dual of the latter problem
and obtain 
\begin{displaymath}
  \begin{array}{l}
    \min\, b^Ty_1 - b^Ty_2 + d^T y_3 \\
  \begin{array}{rcl}
   B^Ty_1 - B^T y_2 + C^T  y_3 & = &  c \\
  y_1,y_2,y_3& \geq & 0.
\end{array}
\end{array}
\end{displaymath}
%where $m_1$ and $m_2$ are the number of
%rows of $B$ and $C$ respectively. 
But this linear program is
equivalent to the linear program 
\begin{equation}
  \label{eq:41}
  \begin{array}{c}
    \min \, b^Ty_1 + d^T y_2 \\
    \begin{array}{rcl}
      B^Ty_1 + C^T  y_2 & = &  c \\  y_2& \geq & 0. 
    \end{array}
  \end{array}
\end{equation}
%
This justifies to say that~\eqref{eq:41} is the dual of~\eqref{eq:40}.


\section{Zero sum games} 
\label{sec:zero-sum-games}

Consider the following two-player game defined by a matrix $A \in
\setR^{m\times n}$. The \emph{row-player} chooses a row $i \in \{1,\ldots,m\}$ and the
column-player chooses a column $j \in \{1,\ldots,n\}$. Both players make
this choice at the same time.  The \emph{payoff} for the
row-player is then the matrix-element $A(i,j)$ whereas $A(i,j)$ also
determines the \emph{loss} of the column player. In other words, the
column player pays $A(i,j)$ to the row-player. If this number is
negative, then the row-player actually pays the absolute value of
$A(i,j)$ to the column player. 

Consider for example the matrix 

\begin{equation}
  \label{eq:34}
  A =
  \begin{pmatrix}
    5 & 1 & 3 \\
    3 & 2 & 4 \\
    -3 & 0 &1 
  \end{pmatrix}. 
\end{equation}

If the row-player chooses the second row and the column player chooses
the second-column, then the payoff for the row-player is $2$, whereas
this is the loss of the column player.  

The row-player is now interested in finding a strategy that maximizes
his \emph{guaranteed} payoff. 
For example, if he chooses row~$1$, then the best choice of the column
player  would be column 2, since the second element of the first row
is the smallest element of that row. Thus the strategy that maximizes
the minimal possible payoff would be to choose row $2$. In other words
\begin{displaymath}
  \max_i \min_j A(i,j) = 2. 
\end{displaymath}
%
What would be the column-player's best hedging strategy? He wants to
choose a column such that the largest element in this column is
minimized. This column would be the second one. In other words
\begin{displaymath}
  \min_j \max_i A(i,j) = 2. 
\end{displaymath}
Is it always the case that $\max_i \min_j A(i,j) = \min_j \max_i
A(i,j)$? The next example shows that the answer is no: 
\begin{equation}
  \label{eq:36}
  \begin{pmatrix}
    -1 & 1 \\ 
    1 & -1
  \end{pmatrix}. 
\end{equation}
Here we have $ \max_i \min_j A(i,j) = -1$ and  $ \min_j \max_i A(i,j)
= 1$. This can be interpreted as follows. If the column player knows
beforehand, the row to be chosen by the row-player, then he would
choose a column that results in a gain for him.  Similarly, if the
row-player knows beforehand the column to be chosen by the
column-player, then he can guarantee him a gain of one. 

The idea is thus not to stick with a \emph{pure} strategy, but to play
with a \emph{random} or \emph{mixed} strategy. If the row-player
chooses each of the two rows above uniformly at random, then his
expected payoff is zero. Similarly, if the column player chooses each
of his two columns with probability $1/2$, then his expected payoff is
zero as well. 

\begin{definition}[Mixed strategy]
  \label{def:m1}
  Let $A \in \setR^{m\times n}$ define a two-player matrix game. A mixed
  strategy for the row-player is a vector $x \in \setR_{\geq0}^m$ with
  $\sum_{i=1}^m x_i = 1$. A mixed strategy for the column player is a
  vector $y \in  \setR_{\geq0}^n$ with $\sum_{j=1}^n y_i = 1$. 
\end{definition}
%
Such mixed strategies define a probability distribution on the row and
column indices respectively. If the row-player and column-player
choose a row and column according to this distribution respectively,
then the \emph{expected payoff} for the row-player is
\begin{equation}
  \label{eq:37}
  E[\text{Payoff}]=x^TAy. 
\end{equation}
%
For the game defined by~\eqref{eq:36} and $x^T=(1/2,1/2)$ and
$y^T=(1/2,1/2)$ the expected payoff is $0$. 

\begin{lemma}
  \label{d:lem:10}
  Let $A \in \setR^{m\times n}$, then 
   \begin{displaymath}
   \max_{x \in X} \min_{y\in Y} x^TA y \leq \min_{y\in Y} \max_{x\in X}
   x^TAy,       
  \end{displaymath}
  where $X$ and $Y$ denote the set of mixed row and column-strategies
  respectively. 
\end{lemma}

\begin{proof}
  Let $x'$ and $y'$  be some fixed mixed strategies of the row and
  column-player  respectively. Clearly 
  \begin{displaymath}
    \min_y {x'}^T A y \leq {x'}^T A y' \leq \max_x {x}^T A y',
  \end{displaymath}
  which implies the assertion. \qed 
\end{proof}


The next theorem is one of the best-known results in the field of
\emph{game theory}. It states that there are mixed strategies $x'$ and
$y'$ from above such that equality holds. It is proved with the
theorem of strong duality. 

\begin{theorem}[Minimax-Theorem]
  \label{d:thr:7}
  \begin{displaymath}
   \max_{x \in X} \min_{y\in Y} x^TA y = \min_{y\in Y} \max_{x\in X}
   x^TAy,       
  \end{displaymath}
  where $X$ and $Y$ denote the set of mixed row and column-strategies
  respectively. 
\end{theorem}
\begin{proof}
  Let us inspect the value $\max_{x \in X} \min_{y\in Y} x^TA y $. This
  can be understood as to maximize the function 
  \begin{displaymath}
    f(x) = \min\{(x^TA) \cdot y \colon \sum_{j=1}^n y_j = 1, \, y\geq0\} . 
  \end{displaymath}
  Thus the value $f(x)$ is the optimal solution of a bounded and
  feasible linear program. The dual of this linear program (for fixed
  $x$) has only one variable $x_0$ and reads 
  \begin{displaymath}
    \max\{ x_0 \colon  x_0 \in \setR, \,  \mathbf{1} x_0 \leq A^T x\}. 
  \end{displaymath}
  But this shows that the maximum value of $f(x)$, where $x$ ranges
  over all mixed row-strategies is the linear program 
  \begin{equation}
    \label{eq:38}
    \begin{array}{rcl}
      \max x_0 \\
      \mathbf{1} x_0 - A^Tx &\leq&0 \\
      \sum_{i=1}^m x_i & = & 1 \\
      x & \geq & 0.
    \end{array}
  \end{equation}
  Let us now inspect the value $\min_{y \in Y} \max_{x\in X} x^TA y
  $. Again, by applying duality this can be computed with the linear
  program 
  \begin{equation}
    \label{eq:39}
    \begin{array}{rcl}
      \min y_0 \\
      \mathbf{1} y_0 - Ay &\geq&0 \\
      \sum_{j=1}^n y_j & = & 1 \\
      y & \geq & 0.
    \end{array}
  \end{equation} 
It follows from the duality of \eqref{eq:41} and \eqref{eq:40} that
the linear programs~\eqref{eq:38} and 
\eqref{eq:39} are duals of each other. This proves the
Minimax-Theorem. \qed 



\end{proof}



% \section{The classical   simplex algorithm} 
% \label{sec:class-expl-simpl}

% I have chosen to explain the simplex algorithm as a method to obtain
% primal feasibility of the linear program 
% \begin{equation}
%   \label{eq:43}
%   \max\{c^Tx \colon x \in \setR^n , \, Ax\leq b\}
% \end{equation}
% via improving the roofs. By duality, our description
% can be re-interpreted as solving the linear program
% \begin{equation}
%   \label{eq:42}
%   \min\{ b^T y \colon y \in \setR^m, \, A^T y =c , \, y \geq0\}
% \end{equation}
% where $A^T \in \setR^{n\times m}$ has now full row-rank. Recall that a roof is
% set of row-indices $B$ of $A$ such that the corresponding rows are a basis
% of $\setR^n$ and $c$ is a conic combination of these rows. The
% translation of a roof to the classical setting~\eqref{eq:42} is a set
% of $n$ linearly independent columns such that $c$ is a conic
% combination of these columns. The corresponding \emph{basic feasible
%   solution} is the vector $x^*$, where $x^*_B = {A_B^T}^{-1} c$ and
% $x^*_{\overline{B}} = 0$. The simplex method can now be interpreted as
% a method that improves the basic-feasible solution by letting a new
% index enter the basis while another index leaves the basis. More on
% this classical simplex versus the roof-interpretation can be found in 
% exercise~\ref{d:item:3}


\section{A proof of the duality theorem via Farkas' lemma}
\label{sec:proof-dual-theor}

Remember Farkas' lemma (Theorem~\ref{conv:thr:12}) which states that
 $Ax =b, x\geq0$ has a solution if and only if for all $\lambda \in \setR^m$
with $\lambda^T A \geq0$ one also has $\lambda^Tb \geq0$. In fact the duality
theorem follows  from this. First,  we derive  another variant of
Farkas' lemma.  

\begin{theorem}[Second variant of Farkas' lemma]
\label{dual:thr:2farkas}
    Let $A \in \setR^{m\times n}$ and $b \in \setR^m$. The system $Ax\leq b$ has a
    solution if and only if for all $\lambda\geq0$ with $\lambda^TA =0$ one has
    $\lambda^Tb\geq0$. 
\end{theorem}

\begin{proof}
  Necessity is clear: If $x^*$ is a feasible solution, $\lambda\geq0$ and
  $\lambda^TA=0$, then $\lambda^T A x^* \leq\lambda^Tb$  implies $0 \leq \lambda^Tb$. 

  On the other hand, $Ax\leq b$ has a solution if and only if 
  \begin{equation}
    \label{eq:10}
    Ax^+ - Ax^-   + z = b, \, x^+,x^-,z\geq0
  \end{equation}
  has a solution. So, if $Ax\leq b$ does not have a solution, then also
  \eqref{eq:10} does not have a solution. 
  By  Farkas' lemma, 
   there exists a 
  $\lambda \in \setR^m$ with $\lambda^T \smat{A \ -A \ I_m}\geq0$ and $\lambda^Tb <0$. 
  For this $\lambda$ one also  has $\lambda^TA = 0$ and  $\lambda\geq0$.
\qed
\end{proof}

We are now ready to prove the theorem of strong duality via the second
variant of Farkas' lemma. %In fact we prove Corollary~\ref{thr:5},
%which serves our purpose, since, by the discussion preceding
%Corollary~\ref{thr:5}, Theorem~\ref{thr:4} and Corollary~\ref{thr:5}
%are equivalent.  

\begin{proof}[of strong duality  via Farkas' lemma]
  Let $\delta$ be the
  objective function value of an optimal solution of the dual
  $\max\{b^T y \colon y \in \setR^m, \, A^Ty \leq c\}$. 
  For all $\varepsilon>0$, the system $A^T y \leq c, -b^Ty \leq - \delta -\varepsilon$ does not have
  a solution. By the second variant of Farkas' lemma,
  there exists a $\lambda\geq0$ with $\lambda^T \smat{-b^T \\ 
    A^T} = 0$ and $\lambda^T \smat{-\delta -\varepsilon \\ c}<0$. 
  Write $\lambda$ as
  $\lambda = \smat{\lambda_1 \\ \lambda'}$ with $\lambda' \in \setR^n$.  
  If $\lambda_1$ were zero, we could apply the second variant of Farkas' lemma
  to the system $A^Ty \leq c$ and $\lambda'$, since we know that $A^Ty \leq c$ has a solution.
  Therefore, we can conclude $\lambda_1 > 0$.
  Furthermore, by scaling, we can assume $\lambda_1 =
  1$. One has $\lambda'^TA^T = b^T$ and $\lambda'^Tc < \delta+\varepsilon$. The first
  equation implies that $\lambda'$ is a feasible solution of the
  primal (recall $\lambda'\geq0$). The second equation shows that the objective function value
  of $\lambda'$ is less than $\delta+\varepsilon$. This means that 
  the optimum value of the  primal linear program is also $\delta$,
  since the primal has an optimal solution 
  and $\varepsilon$ can be chosen arbitrarily small. 
  \qed
\end{proof}




\subsection*{Exercises}

\begin{enumerate}
\item Formulate the dual linear program of 
  \begin{displaymath}
    \begin{array}{r c l}
      \max \ 2 x_1 + 3 x_2 -7 x_3 & & \\
      x_1 + 3 x_2 + 2 x_3 & = & 4\\
      x_1 + x_2 & \leq & 8 \\
      x_1 - x_3 & \geq  &-15\\
      x_1,x_2 & \geq &0
    \end{array}
  \end{displaymath}
\item Consider the following linear program 
  \begin{displaymath}
    \begin{array}{rcl}
      \max \ x_1 + x_2 \\
      2 x_1 + x_2 & \leq  & 6 \\
      x_1 + 2x_2  & \leq & 8 \\
      3x_1 + 4x_2 & \leq & 22 \\
      x_1 + 5 x_2 & \leq & 23
    \end{array}   
  \end{displaymath}
  Show that $(4/3, 10/3)$ is an optimal solution by providing a
  suitable feasible dual solution. 

\item Show that for $A \in \setR^{m\times n}$, one has 
  \begin{displaymath}
    \max_i \min_j A(i,j) \leq \min_j \max_i A(i,j). 
  \end{displaymath}\label{d:item:2} 

\item  \label{d:item:3} 
 In the lecture you have seen the simplex algorithm for linear programs of the form 
 $$ \max \{ c^Tx~:~Ax\leq b\}. $$
 We will now derive a simplex algorithm for linear programs of the form
 \begin{align} \label{eq:LPEqStandard} \min \{ c^Tx~:~Ax = b,~x\geq 0\} \end{align}
with $c\in \R^n$ and $A\in \R^{m\times n}$, $b\in \R^m$.
Throughout the exercise we assume that \eqref{eq:LPEqStandard} is feasible and bounded, and that $A$ has full row rank.

For $i\in \{1, \ldots, n\}$ we define $A_i$ as the $i$-th column of $A$. Moreover, for some subset $B\subseteq \{1, \ldots, n\}$, $A_B$ is the matrix $A$ restricted to the columns corresponding to elements of $B$.

A subset $B\subseteq \{1, \ldots, n\}$ with $|B|=m$ such that $A_B$ has full rank is called a \emph{basis}.
The vector $x\in \R^n$ defined as $x_i := 0$ for all $i\notin B$ and $x_B:=A_B^{-1}b$ is called the \emph{basic solution} associated to $B$.
Note that $x$ is a feasible solution to \eqref{eq:LPEqStandard} if and only if $x\geq 0$.



Given a basis $B$ and let $j\in \{1, \ldots, n\}$, $j\notin B$. The vector $d\in \R^n$ defined as $d_j = 1$, $d_i=0$ for all $i\notin B$ and
$d_B := - A_B^{-1} A_j$ is called the \emph{$j$-th basic direction}.

Assume that the solution $x$ associated to $B$ is feasible. Moreover assume that $x_B>0$.
\begin{enumerate}
 \item Show that there is a $\theta >0$ such that $x+ \theta d$ is a feasible solution. Give a formula to compute the largest $\theta$ such that $x+ \theta d$ is feasible.
 \item \label{it:basechange} Let $\theta^*$ be maximal. Show that there is a basis $B'$ such that $x+ \theta^* d$ is the basic solution associated to $B'$.
 \item \label{it:improve} Let $x'=x+\theta d$. Show that the objective value of $x'$ changes by $\theta \left(c_j - c_B^TA_B^{-1} A_j\right)$.
 \item \label{it:optimality} Consider a basis $B$ with basic feasible solution $x$. Show that if $c - c_B^T A_B^{-1} A \geq 0 $, then $x$ is an optimal solution to \eqref{eq:LPEqStandard}.
\end{enumerate}
This suggests the following algorithm: Start with some basis $B$ whose associated basic solution is feasible. Compute $\bar c := c - c_B^T A_B^{-1} A$. If $\bar c \geq 0$, we have an optimal solution (see \ref{it:optimality}). Otherwise, let $j$ be such that $\bar c_j < 0$. Part \ref{it:basechange} and \ref{it:improve} show that if we change the basis, 
we find a feasible solution with an improved objective value. We repeat these steps until the vector $\bar c$ is nonnegative.

This is the way the simplex algorithm usually is introduced in the literature. This algorithm is exactly the same as the one you learned in the lecture.
To get an intuition why this is true, show the following:
\begin{enumerate}
\setcounter{enumi}{4}
\item Given a basis $B$, show that its associated basic solution is feasible if and only if $B$ is a \emph{basis} of the LP dual to \eqref{eq:LPEqStandard}.
 \item Consider a basis $B$ and its associated feasible basic solution $x$. As seen before, $B$ is also a basis in the dual LP. 
  Let $y$ be the vertex of that basis. 

Show that for any $j\in \{1, \ldots, n\}$ we have $\bar c_j < 0$ if and
only if  $A_j^T y > c_j$.
\end{enumerate}
\end{enumerate}








%%% Local Variables: 
%%% mode: latex
%%% TeX-master: "lecture"
%%% End: 

\chapter{Integer Programming}
\label{sec:integer-programming}

An \emph{integer program} is a problem of the form 
\begin{displaymath}
  \begin{array}{c}
    \max c^Tx \\
    Ax\leq b \\
    x \in \setZ^n,
  \end{array}
\end{displaymath}
where $A \in \setR^{m\times n}$ and $b \in \setR^m$. 


 

\begin{figure}[htbp]
  \begin{center}
   \includegraphics{./figures/IntProg1.pdf}
\label{fig:inthull}
  \end{center}
  \caption{This picture illustrates a polyhedron $P$, an objective
    function vector $c$ and optimal points $u,v$ of the integer
    program and the relaxation respectively. }
\end{figure}



The difference to linear programming is the \emph{integrality
  constraint} $x \in \setZ^n$. This powerful constraint allows to
model discrete choices but, at the same time, makes an integer program
much more difficult to solve than a linear program. In fact one can
show that integer programming is NP-hard, which means that it is
\emph{in theory} computationally intractable. However, integer
programming has nowadays become an important tool to solve difficult
industrial optimization problems efficiently. In this chapter, we
characterize some integer programs which are easy to solve, since the
\emph{linear programming relaxation} $\max\{c^Tx \colon Ax\leq b\}$
yields already an optimal integer solution. The following observation
is crucial.

\begin{theorem}
  \label{thr:14}
  Suppose that   $x^*$ is an integral  optimal solution of the  
  linear programming 
  relaxation $\max\{c^Tx \colon Ax\leq b\}$, i.e., $x^* \in
  \setZ^n$, then $x^*$ is also an optimal solution of the integer
  programming problem $\max\{c^Tx \colon Ax\leq b, \, x \in \setZ^n\}$
\end{theorem}

Before we present an example for the power of integer programming we
recall the definition of an undirected graph. 

\begin{definition}[Undirected graph, matching] 
  An \emph{undirected graph} is a tuple $G = (V,E)$ where $V$ is a
  finite set of elements, called the \emph{vertices} or the \emph{nodes}, and $E\subseteq\binom{V}{2}$ is the
  set of \emph{edges} of $G$.  A \emph{matching} of $G$ is a subset
  $M\subseteq E$ such that for all $e_1\neq e_2\in M$ one has $e_1\cap e_2 = \emptyset$. 
\end{definition}



% \begin{tikzpicture}[every node/.style={draw,circle}]
% \draw[help lines] (0,0) grid (2,5);
% \begin{scope}[node distance=5mm]
% \node (a) at (1,1) {a};\node [left=of a] {1}; \node [right=of a] {2};
% \node [above=of a] {3}; \node [below=of a] {4};
% \node [above left=of a] {5}; \node [above right=of a] {6};
% \node [below left=of a] {7}; \node [below right=of a] {8};
% \end{scope}

% \begin{scope}[node distance=5mm and 5mm]
% \node (b) at (1,4) {b};
% \node [left=of b] {1}; \node [right=of b] {2};
% \node [above=of b] {3}; \node [below=of b] {4};
% \node [above left=of b] {5}; \node [above right=of b] {6};
% \node [below left=of b] {7}; \node [below right=of b] {8};
% \end{scope}
% \end{tikzpicture}




\begin{figure}
\begin{center}
\begin{tikzpicture}[style=thick]
%\draw[help lines] (0,0) grid (2,5);
\tikzstyle{vertex}=[draw,shape=circle]
\tikzstyle{edge} = [draw,thick,-]
\tikzstyle{weight} = [font=\small];

\draw (0,0) node[vertex] (s) {$s$};
\draw (3,2) node[vertex] (u) {$u$};
\draw (3,-2) node[vertex] (v) {$v$};
\draw (6,0) node[vertex] (t) {$t$};

\draw[red]  (s) --  (u);
\draw        (s) --  (v);
\draw  (u) -- (v);
\draw (u) -- (t);
\draw [red] (v) -- (t);
\end{tikzpicture}
\end{center}
  
  \caption{A graph with 4 nodes $V = \{s,u,v,t\}$ and 5 edges $E = \{
    \{s,u\}, \{s,v\}, \{u,v\}, \{u,t\},\{v,t\} \}$. The red edges
    are a matching of the graph}
\end{figure}



We are interested in the solution of the following problem, which is
called \emph{maximum weight matching} problem. Given a graph $G =
(V,E)$ and a weight function $w:E\to\setR$, compute a matching with
maximum weight $w(M) = \sum_{e \in M} w(e)$. 

For a vertex $v \in V$, the set $\delta(v) = \{e \in E \colon v \in e\}$ denotes
the \emph{incident} edges to $v$. 
The maximum weight matching problem  can now be modeled as an integer
program as follows. 
\begin{equation}
  \label{eq:45}
  \begin{array}{c}
    \max \sum_{e \in E} w(e) x(e) \\
    v \in V: \, \sum_{e \in \delta(v)} x(e)\leq1 \\
    e \in E:\,  x(e)  \geq0 \\
    x \in \setZ^{|E|}.
  \end{array}
\end{equation}
% 
Clearly, if an integer vector $x \in \setZ^n$ satisfies the constraints
above, then this vector is the \emph{incidence vector}  of a matching of
$G$. In other words, the integral solutions to the constraints above
are the vectors $\{\chi^M \colon M \text{ matching of }G\}$, where $\chi^M_e
= 1$ if $e \in M$ and $\chi^M_e=0$ otherwise. 

  
\section{Integral Polyhedra}


In this section we derive sufficient conditions on an integer program
to be solved easily by an algorithm for linear programming. A central
notion is the one of an integral polyhedron.  


\begin{definition}[Valid inequality, face, vertex] 
  \label{def:8}
  Let $P = \{ x \in \setR^n \colon Ax\leq b\}$ be a polyhedron. An
  inequality $c^Tx\leq\beta$ is \emph{valid} for $P$ if $c^Tx^* \leq \beta$ for
  all $x^* \in P$.  A \emph{face} of $P$ is a set of the form $P \cap \{x
  \in \setR^n \colon c^Tx = \beta\}$ for a valid inequality $c^Tx\leq\beta$ of $P$. If
  a face consist of one point, then it is called a \emph{vertex} of
  $P$. 
\end{definition}




\begin{figure}[htbp]
  \begin{center}{
%    \psset{unit=.8cm}
   \includegraphics{./figures/IntProg2.pdf}
    }    
  \end{center}
  \caption{A polyhedron with a valid inequality defining a vertex. }
  \label{fig:3}
\end{figure}




\begin{definition}
  \label{def:12}  
  A rational polyhedron   is called \emph{integral} 
  if  each nonempty face of $P$ contains an integer vector. 
\end{definition}



\begin{lemma}
  \label{lem:9}
  Let $P = \{ x \in \setR^n \colon Ax\leq b\}$ be an integral polyhedron with $A
  \in \setR^{m\times n}$ full-column rank. If the linear program 
  \begin{equation}
    \label{eq:44}
    \max\{c^Tx \colon x \in \setR^n, \, Ax\leq b\}
  \end{equation}
  is feasible and bounded, then the simplex method computes an optimal
  integral solution to the linear program. 
\end{lemma}


\begin{proof}
  Recall that, if the linear program~\eqref{eq:44} is bounded, the simplex method finds an optimal basis $B
  \subseteq\{1,\ldots,m\}$ of~\eqref{eq:44} and the vertex of the basis $x^*_B$ is
  an optimal solution to~\eqref{eq:44}. We have to
  show that $x^*_B$ is integral. This will follow from the fact that
  $\{x^*_B\}$ is a face of $P$.

  Theorem~\ref{thr:1} implies that $x^*_B$ is the unique optimum
  solution of the linear program $\max\{\wt{c}^T x \colon x \in 
  \setR^n, \,
  a_i^Tx \leq b_i, \, i \in B\}$, where $\wt{c} = \sum_{i \in B} a_i$.
  Consequently $x^*_B$ is the unique solution of the linear program
  \begin{displaymath}
    \max\{\wt{c}^Tx \colon x \in P\}
  \end{displaymath}
  which implies that $\{x^*_B\}$ is a face defined by the valid
  inequality $\wt{c}^Tx \leq \wt{c}^Tx^*_B$.  \qed 
\end{proof}


% \begin{theorem}
%   \label{po:thr:5}
%   Let $P = \{x \in \setR^n \mid Ax\leq b\}$  be a rational nonempty  polyhedron with
%   vertices. $P$ is integral if and only if for all integral vectors $c
%   \in \setZ^n$ with $\max\{c^Tx \mid x \in P\}<\infty$ one has $\max\{c^Tx \mid x
%   \in P\} \in \setZ$. 
% \end{theorem}

% \begin{proof}
%   Let $P$ be integral and $c \in \setZ^n$ with
%   $\max\{c^Tx\mid x\in P\}=\delta<\infty$. Since the face    $F=\{x\in P\mid c^Tx=\delta\}$
%   contains an integer point it follows   that $\delta\in\setZ$.  

%   On the other hand let $x^*$ be a vertex of $P$ and assume that $x^*(i) \notin
%   \setZ$. There exists a subsystem $A'x\leq b'$ of $Ax\leq b$ with $A'\in\setR^{n\times n}$,
%   $A'$ nonsingular and $A'x^*=b'$. Let $a_1,\ldots,a_n$ be the rows of
%   $A'$. Since $A'$ is invertible, there exists an integer vector $c \in
%   \cone(a_1,\ldots,a_n)\cap\setZ^n$ such that $c\pm e_i \in
%   \cone(a_1,\ldots,a_n)$. The point $x^*$ maximizes both $c^Tx$ and
%   $(c+e_i)^Tx$. Clearly not both numbers $c^Tx^*$ and
%   $(c+e_i)^Tx^*$ can be integral, which is a contradiction. 
% \end{proof}



\begin{lemma}
  \label{po:lem:6}
  Let $A\in \setZ^{n\times n}$ be an integral and invertible matrix. One has
  $A^{-1}b \in \setZ^n$ for each $b \in \setZ^n$ if and only if $\det(A)=\pm 1$.
\end{lemma}


\begin{proof}
  Recall Cramer's rule which says $A^{-1} = \wt{A}/\det(A) $, where
  $\wt{A}$ is the adjoint matrix of $A$. Clearly $\wt{A}$ is
  integral. If $\det(A) = \pm 1$, then $A^{-1}$ is an integer matrix. 

  If $A^{-1}b$ is integral for each $b \in \setZ^n$, then $A^{-1}$ is an
  integer matrix. We have $1=\det(A\cdot A^{-1})=\det(A)\cdot\det(A^{-1})$.
  Since $A$ and $A^{-1}$ are integral it follows that $\det(A)$ and
  $\det(A^{-1})$ are integers. The only divisors of one in the integers
  are $\pm 1$.  \qed 
\end{proof}



% A matrix $A \in \setZ^{m\times n}$ with $m\leq n$ is called \emph{unimodular} if
% each $m\times m$ sub-matrix has determinant $0,\pm1$. 

% \begin{theorem}
%   \label{po:thr:14}
%   Let $A \in \setZ^{m\times n}$ be an integral matrix of full row-rank. The
%   polyhedron defined by $Ax=b, x\geq0$ is integral for each $b \in \setZ^m$
%   if and only if $A$ is unimodular. 
% \end{theorem}


% \begin{proof}
%   Suppose that $A$ is unimodular and $b$ is integral. The polyhedron
%   $P = \{ x \in \setR^n \mid Ax=b, \, x\geq0\}$ does not contain a line and
%   thus has vertices. A vertex $x^*$ is of the form  $x^*_B= A_B^{-1}b$
%   and $x^*_{\wb{B}}=0$, where $B\subseteq\{1,\ldots,n\}$ is a basis. Since $A_B$
%   is unimodular one has $x^* \in \setZ^n$. 

%   If $A$ is not unimodular, then there exists a basis $B$ with
%   $\det(A_B) \neq \pm1$. By Lemma~\ref{po:lem:6} there exists an integral $b
%   \in \setZ^n$ with $(A_B)^{-1}b \notin \setZ^m$. Let $\lambda$ be the maximal
%   absolute value of a component of $A_B^{-1}b$. Then $b' =  \lceil\lambda\rceil A_B \mathbf{1}
%   +b$ is an integral vector with $ A_B^{-1}b' = \lceil\lambda\rceil \mathbf{1}+
%   A_B^{-1}b \geq0$ and $A_B^{-1}b' \notin \setZ^m$.  The
%   polyhedron $P = \{ x \in \setR^n \mid Ax = b', \, x\geq0\}$ has thus a fractional
%   (non-integer) vertex.   
% \end{proof}



\begin{definition}[Total unimodularity]
  \label{def:skip1}  
  An integral matrix $A \in \{0 , \pm 1\}^{m\times n}$ is
  called \emph{totally
  unimodular} if each of its square sub-matrices has
  determinant
  $0,\pm1$.
\end{definition}
   

\begin{theorem}[Hoffman-Kruskal Theorem]
  \label{po:thr:16}
  Let $A\in \setZ^{m\times n}$ be an integral matrix. The polyhedron $P = \{ x
  \in \setR^n \mid Ax\leq b, \, x\geq0\}$ is integral for each integral $b \in
  \setZ^m$ if and only if $A$ is totally unimodular. 
\end{theorem}


\begin{proof}
  Let $A\in \setZ^{m\times n}$ be totally unimodular  and $b \in \setZ^m$. 
  Let $x^*$ be vertex of $P$ and suppose that this vertex is defined
  by the valid inequality $c^Tx \leq \delta$. Notice that the matrix $\smat{A \\ -I}$
  has full column-rank. If one applies the simplex algorithm to the
  problem 
  \begin{displaymath}
    \max\{c^Tx \colon x \in \setR^n, \, \smat{A\\-I} x\leq \smat{b\\0} \},
  \end{displaymath}
  it finds an optimal basis $B\subseteq\{ 1,\ldots,m+n\}$ with $x^*_B = x^*$. If
  $A_B$ denotes the matrix whose rows are those rows of $\smat{A \\
    -I}$ indexed by $B$ and if $b_B$ denotes the vector whose
  components are those of $\smat{b\\0}$ indexed by $B$, then $x^* =
  A_B^{-1} b_B$. We are done, once we conclude that $\det(A_B) = \pm 1$,
  since then $A_B^{-1}$ is an integer matrix and since $b_B$ is an
  integer vector $x^* =  A_B^{-1} b$ is integral as well. We can
  permute the  columns of $A_B$ in such a way that one obtains a
  matrix of the form 
  \begin{displaymath}
    \mat{ \wb{A} & \wt{A} \\
      0      & - I_k}
  \end{displaymath}
  where $\wb{A}$ is a $(n-k)\times (n-k)$ sub-matrix of $A$ and $I_k$ is
  the $k\times k$ identity matrix. Here $k = | B \cap \{
  m+1,\ldots,m+n\} |$. Clearly $ 0 \neq \det(A_B) = \pm \det(\wb{A}) = \pm 1$. 

  
  For the converse, suppose that $A$ is not totally unimodular. Then
  there exists an index set $B \subseteq\{1,\ldots,m+n\}$ with $|B| = n$ such that
  the matrix $A_B$ defined as above satisfies $|\det(A_B)| \geq  2$. We can suppose w.l.o.g. that $B = \{1, \ldots, n\}.$ By
  Lemma~\ref{po:lem:6} there exists choices for the components of
  $b_B$ making $A_B^{-1} b_B$ non-integral. In fact, if we split $B$
  into components $L \subseteq B$  corresponding to lines of $A$ and $C$
  corresponding to lines of $-I$ we can choose those components of
  $b_B$ corresponding to $L$ being equal to zero. Now let $v$ be the
  vector with $v_i = 1$ for all $i\in C$ and $v_i = 0 $ for all $i \in
  L$. 
  By choosing $\gamma \in \setN$ large enough the point $x^*_B = A_B^{-1} (b_B
  + \gamma A_B v)$ is non-integral and positive. Notice that
  starting from now we will consider a new vector $\tilde{b}$
  instead of 
  $b$, where $\tilde{b}_B = b_B$. In the next lines we will 
  say $\tilde{b}_{\{1, \ldots, m+n\} \textbackslash B}$ has to
  be to finish the proof.
  The set $B$ is a
  basis of the linear program 
  \begin{displaymath}
       \max\{\wb{c}^Tx \colon x \in \setR^n, \, \smat{A\\-I} x\leq \smat{\tilde{b}\\0} \},
  \end{displaymath}
  where $\wb{c} = \sum_{i \in B} a_i$ and $a_i$ denotes the $i$-th row of
  $\smat{A\\-I}$. If we define for $j \in \{1,\ldots,m\} \setminus B$,  $\tilde{b}_j =
  \lceil a_j^Tx^*_B\rceil$, then $x^*_B$ is feasible and thus a vertex of $P$
  that is non-integral. 
  \qed
  
\end{proof}

A direct consequence of theorem~\ref{po:thr:16} is the following 
corollary.

\begin{corollary}
   If $A \in \setZ^{m \times n}$ is totally unimodular, 
   $b \in \setZ^m$ and if 
   $max\{c^Tx: x \in \setR^n, Ax \leq b, x \geq 0\}$ is bounded, 
   then 
   \begin{displaymath}
      \max\{c^Tx: x \in \setR^n, Ax \leq b, x \geq 0\} = \max\{c^Tx: x \in \setZ^n, Ax \leq b, x \geq 0\}.
   \end{displaymath}
\end{corollary}

\section{Applications of total unimodularity}

\subsection{Bipartite matching} 

% An undirected graph $G=(V,E)$ is a tuple, where $V$ is a finite set
% and $E$ is a set of unordered pairs of $V$. The set $V$ is called
% \emph{nodes} and the set $E$ are the \emph{edges} of $G$. We write
% $uv$ in short for the edge $\{u,v\}\subseteq V$. 

A graph is
\emph{bipartite}, if $V$ has a partition into sets $A$ and $B$ such
that each  edge $uv$ satisfies $u\in A$ and $v \in B$.
Recall that $\delta(v)$ is the set of edges incident to the vertex $v\in V$,
that is $\delta(v) = \{ e\in E \mid v\in e \}$.

% A \emph{matching} of $G$ is a subset $M\subseteq E$ such that $e_1\cap e_2 = \emptyset$
% holds for each $e_1\neq e_2\in M$. Let $c:E\longrightarrow\setR$ be a weight function. The
% weight of a matching is defined as $c(M) = \sum_{e \in M}c(e)$.  The
% \emph{weighted matching problem } is  defined as follows. Given a
% graph $G = (V,E)$ and edge-weights $c:E\longrightarrow\setR$, compute a matching $M$
% of $G$ with $c(M)$ maximal.


% %We have already seen how to compute a maximum weight matching of a
% %bipartite graph in polynomial time with a polynomial minimum cost
% %network flow algorithm. We now take a different perspective. 
% We now define
% an \emph{integer program} for this problem and show that, for
% bipartite graphs, an optimal vertex of the corresponding linear
% program is integral. 


% The idea is as follows. We have decision variables $x(e)$ for each
% edge $e \in E$. We want to model the characteristic vectors $\chi^M\in
% \{0,1\}^E$  of matchings, where $\chi^M(e)=1$ if  $e\in M$  and $\chi^M(e)=0$
% otherwise. This is achieved with the following set of constraints. 
% \begin{equation}
%   \label{po:eq:3}
%   \begin{array}{rcll}
%     \sum_{e \in \delta(v)} x(e) &\leq&1, & \forall v \in V \\
%     x(e) &\geq&0, & \forall e \in E. 
%   \end{array}
% \end{equation}

% Clearly, the set of vectors $x \in \setZ^E$ which satisfy the
% system~\eqref{po:eq:3} are exactly the characteristic vectors of
% matchings of $G$. The matrix $A \in \{0,1\}^{V\times E}$ which is defined as 
% \begin{displaymath}
%   \label{po:sec:bipartite-matching}
%   A(v,e) = 
%   \begin{cases}
%     1 & \text{ if } v \in e,\\
%     0 & \text{ otherwise}
%   \end{cases}
% \end{displaymath}
% is called \emph{node-edge incidence matrix} of $G$. 

The \emph{node-edge} incidence matrix of a graph $G = (V,E)$ is the 
matrix $A \in \{0,1\}^{|V|\times|E|}$ with 
\begin{displaymath}
  A(v,e) = \
  \begin{cases}
    1, & \text{if } v \in e, \\
    0 & \text{otherwise.}
  \end{cases}
\end{displaymath}


The integer program~\eqref{eq:45} can thus be formulated as 
\begin{equation}
  \label{eq:46}
  \max\{w^Tx \colon Ax\leq1, \, x\geq0, \, x \in \setZ^E\}. 
\end{equation}
The next lemma implies that the simplex algorithm can be used to
compute a maximum-weight matching of a bipartite graph. 

\begin{lemma}
  \label{po:lem:9}
  If $G$ is bipartite, the node-edge incidence matrix of $G$ is
  totally unimodular. 
\end{lemma}


\begin{proof}[By induction]
  Let  $G = (V,E)$ be a bipartite graph with bi-partition 
  $V=V_1\cup V_2$. 
  
  The case where $A'$ is a $1 \times 1$ sub-matrix of $A$ is 
  trivial.
  Suppose the lemma is proven for $k-1 \geq 1$ and 
  let $A'$ be a $k\times k$ sub-matrix of $A$. We are interested 
  in the determinant of $A$. Clearly, we can assume that $A$ 
  does not contain a column which contains no $1$ or only one 
  $1$, since we simply consider the $(k-1) \times (k-1)$ 
  sub-matrix $A''$ of $A'$, which emerges from developing the 
  determinant of $A'$ along this column. By the induction 
  hypothesis the determinant of $A'$ would be zero or 
  $\pm1\cdot \det(A'')$.
  
  Thus we can assume that each column contains exactly two ones. Now
  we can order the rows of $A'$ such that the first rows correspond to
  vertices of $V_1$ and then follow the rows corresponding to vertices
  in $V_2$. This re-ordering only affects the sign of the
  determinant. By summing up the rows of $A'$ in $V_1$ we obtain
  exactly the same row-vector as we get by summing up the rows of $A'$
  corresponding to $V_2$. This shows that $\det(A')=0$.  \qed 
\end{proof}


\subsection{Bipartite vertex cover} 


A \emph{vertex cover} of a graph $G = (V,E)$ is a subset $C\subseteq V$ of the
nodes such $e \cap C \neq \emptyset$ for each $e \in E$. Let us formulate an
integer program for the \emph{minimum-weight vertex-cover}
problem. Here, one is given a graph $G = (V,E)$ and weights $w \in
\setR^V$. The goal is to find a vertex cover $C$ with minimum weight
$w(C) = \sum_{v \in V} w(v)$. 

\begin{equation}
\label{eq:47}
  \begin{array}{c}
    \min \sum_{v \in V} w(v) x_v \\
    uv \in E: x_u + x_v \geq1 \\
    v \in V:\,  x_v  \geq0 \\
    x \in \setZ^{V}.
  \end{array}
\end{equation}
Clearly, this is the integer program
\begin{equation}
  \label{eq:48}
  \min\{w^Tx \colon A^Tx \geq 1, \, x\geq0, \, x \in \setZ^V\},
\end{equation}
where $A$ is the node-edge incidence matrix of $G$. 
A matrix $A$ is totally unimodular if and only if $A^T$ is totally
unimodular. Thus the simplex algorithm can be used to compute a
minimum weight vertex-cover of a bipartite graph. Furthermore we have
the following theorem. 

\begin{theorem}[K\"onig's theorem]
  \label{thr:16}
  In any bipartite graph, the number of edges in a maximum matching
  equals the number of vertices in a minimum vertex cover. 
\end{theorem}

\begin{proof}  
  Let $A$ be the node-edge incidence-matrix of the bipartite graph $G
  = (V,E)$. 
  The linear programs 
  $\max\{1^T x \colon Ax\leq1, \, x\geq0\}$ and $\min\{1^T x \colon Ax\geq1, \,
  x\geq0\}$ are duals of each other. Since $A$ is totally unimodular,
  the value of the linear programs are the cardinality of a maximum
  matching and minimum vertex-cover respectively. Thus the theorem
  follows from strong duality.
  \qed
\end{proof}

\subsection{Flows}

Let $G=(V,A)$ be a directed graph.
The \emph{node-edge incidence matrix of a directed graph}
is a matrix $A \in \{0,\pm1\}^{V\times E}$ with 
\begin{equation}
  \label{po:eq:8}
  A(v,a) = 
  \begin{cases}
    1 & \text{ if } v \text{ is the starting-node of } a, \\
    -1 & \text{ if } v \text{ is the end-node of } a, \\
    0  & \text{ otherwise.}
  \end{cases}
\end{equation}


A \emph{feasible flow} $f$  of $G$ with capacities $u$ and in-out-flow $b$ is then
a solution $f \in \setR^A$ to the system $A\,f=b, \, 0\leq f\leq u$.  

\begin{lemma}
  \label{po:lem:10}
  The node-edge incidence matrix $A$ of a directed graph  is totally
  unimodular.  
\end{lemma}


\begin{proof}[By induction]
  The case where $A'$ is a $1 \times 1$ sub-matrix of $A$ is 
  trivial. Let $A'$ be a $k\times k$ sub-matrix of $A$ and suppose 
  we have proven the lemma for every 
  $(k-1) \times (k-1)$ sub-matrix with $k-1 \geq 1$. 
  Again, we can assume 
  that in each column we have exactly one $1$ and one $-1$. 
  Otherwise, we develop the determinant along a column which does not 
  have this property.  But then, the matrix $A'$ is singular, 
  since adding 
  up all rows of $A'$ yields the $0$-vector. 
\end{proof}
  

A consequence is that, if the $b$-vector and the capacities $u$ are
integral and an optimal flow exists, then there exists an integer
optimal flow. 
%We have seen that this follows from the cycle-cancelling
%algorithm, but total unimodularity gives another simple and elegant
%proof of this fact.  





\subsection{Doubly stochastic matrices}


A matrix $A \in \setR^{n\times n}$ is \emph{doubly stochastic} if it satisfies
the following linear constraints 
\begin{equation}
  \label{po:eq:9}
  \begin{array}{rcll}
    \sum_{i=1}^n A(i,j) & = & 1, & \forall j=1,\ldots,n\\
    \sum_{j=1}^n A(i,j) & = & 1, & \forall i=1,\ldots,n\\
    A(i,j)       & \geq & 0, & \forall 1 \leq i,j\leq n.
  \end{array}
\end{equation}

A permutation matrix is a matrix which contains exactly one $1$ per
row and column, where the other entries are all $0$. 

\begin{theorem}
  \label{po:thr:17}
  A matrix $A \in \setR^{n\times n}$ is doubly stochastic if and only if $A$ is
  a  convex combination  of permutation matrices. 
\end{theorem}

\begin{proof}
  Since a permutation matrix satisfies the constraints~\eqref{po:eq:9},
  then so does a convex combination of these constraints. 


  On the other hand it is enough to show that each vertex of the
  polytope defined by the system~\eqref{po:eq:9} is integral and thus a
  permutation matrix. However, the matrix defining the
  system~\eqref{po:eq:9}  is the node-edge incidence matrix of the
  complete bipartite graph having $2n$ vertices. Since such a matrix
  is totally unimodular, the theorem follows. 
\end{proof}





\section{The matching polytope}
\label{po:sec:matching-polytope}


We now come to a deeper theorem concerning the convex hull of
matchings. We mentioned several times in the course that the maximum
weight matching problem can be solved in polynomial time. We are now
going to show a theorem of Edmonds~\cite{Edmonds65b} which provides a
complete description of the matching polytope and present the
proof by Lov\'asz~\cite{Lovasz79}. 

Before we proceed let us inspect the symmetric difference $M_1\Delta M_2$
of two matchings of a graph $G$. If a vertex is adjacent to two edges
of $M_1\cup M_2$, then one of the two edges belongs to
$M_1$ and one belongs to $M_2$. Also, a vertex can never be adjacent
to three edges in $M_1 \cup M_2$. Edges which are both in $M_1$ and $M_2$
do not appear in the symmetric difference. We therefore have the
following lemma. 

\begin{lemma}
  \label{pox:lem:11}
  The symmetric difference $M_1\Delta M_2$ of two matchings decomposes
  into node-disjoint  paths and cycles, where the edges on these paths
  and cycles alternate between $M_1$ and $M_2$. 
\end{lemma}



The \emph{Matching polytope} $P(G)$ of an undirected graph $G = (V,E)$
is the convex hull of incidence vectors  $\chi^M$ of matchings $M$ of
$G$. 


  
\begin{figure}[htbp]
  \centering 
   \includegraphics{figures/IntProg3.pdf}
  \caption{Triangle}
  \label{po:fig:triangle}
\end{figure}
  

The incidence vectors of matchings are exactly the $0/1$-vectors that
satisfy the following system of equations. 

\begin{equation}
\label{po:eq:10}
  \begin{array}{rcll}
     \sum_{e \in \delta(v)} x_e & \leq &  1 & \forall v \in V\\
        x_e& \geq & 0 &    \forall e \in E. 
  \end{array}
\end{equation}

However the triangle (Figure~\ref{po:fig:triangle}) shows that  the
corresponding polytope is not integral. The objective function $\max
\mathrm{1}^Tx$ has value $1.5$. However, one can show that a maximum
weight matching of an undirected graph can be computed in polynomial
time which is a result of Edmonds~\cite{Edmonds65}. 


The following (Figure~\ref{po:fig:edmonds}) is an illustration of an
Edmonds inequality. Suppose that $U$ is an odd subset of the nodes $V$
of $G$ and let $M$ be a matching of $G$. The number of edges of $M$
with both endpoints in $U$ is bounded from above by $\lfloor|U|/2\rfloor$. 

Thus the following inequality is valid for the integer points of the
polyhedron defined by~\eqref{po:eq:10}. 

\begin{equation}
  \label{po:eq:11}
  \sum_{e \in E(U)} x_e \leq\lfloor|U|/2\rfloor,\quad \quad \text{ for each } U\subseteq V,
  \quad |U| \equiv 1 \pmod{2}. 
\end{equation}


\begin{figure}[htbp]    
  \begin{center}
   \includegraphics{figures/IntProg4.pdf}
\end{center}
\caption{Edmonds inequality.}
  \label{po:fig:edmonds}
\end{figure}


The goal of this lecture is a proof of the following theorem. 

\begin{theorem}[Edmonds 65]
\label{po:thr:18}
  The matching polytope is described by the following inequalities:
  \begin{enumerate}[i)]
  \item $x_e \geq0$ for each $e \in E$,
  \item $\sum_{e \in \delta(v)} x_e \leq 1$ for each $v \in V$,
  \item $\sum_{e \in E(U)} x_e \leq \lfloor |U|  /2 \rfloor$ for each $U\subseteq V$
    %\fromSlide*{2}{\red Gomory Cut!} 
  \end{enumerate}
\end{theorem}


\begin{lemma}
\label{po:lem:11}
  Let $G=(V,E)$ be connected and 
  let $w:E\longrightarrow\setR_{>0}$ be a weight-function.  Denote the set of maximum
  weight matchings of $G$ w.r.t. $w$ by $\eM(w)$. Then one of the following statements
  must be true:
  \begin{enumerate}[i)]
  \item $\exists\,v \in V$ such that $\delta(v) \cap M \neq \emptyset$ for each $M \in \eM(w)$
  \item $|M| =   \lfloor|V| /2\rfloor$ for each $M \in \eM(w)$ and $|V|$ is odd.
  \end{enumerate}
\end{lemma}


\begin{proof}
Suppose both $i)$ and $ii)$ do not hold.
Then there exists ${ M}\in \eM(w)$ leaving two exposed nodes $u$ and
$v$.  Choose $ M$ such that the  minimum  distance between   two exposed nodes
  $u,v$ is  minimized. 


Now let $t$ be on shortest path from $u$ to $v$. The vertex $t$ cannot
be exposed. 
\begin{figure}[htbp]
  \centering
    \begin{center}    
   \includegraphics{figures/IntProg5.pdf}
  \end{center}
  \caption{Shortest path between $u$ and $v$. }
  \label{po:fig:short}
\end{figure}


 Let ${\color{red} M'} \in \eM(w)$ leave $t$ exposed. 
 Both $u$ and $v$ are covered by  ${\color{red} M'}$ because the distance to
 $u$ or $v$ from $t$ is smaller than the distance of $u$ to $v$. 
 
 Consider the symmetric difference ${ M} \triangle {\color{red} M'}$ which  decomposes into
 node disjoint paths and   cycles. 
The  nodes $u, \, v$ and $t$ have degree one in ${M}\triangle{\color{red} M'}$. Let 
$P$ be a  path with endpoint $t$ in ${ M}\triangle{\color{red} M'}$


\begin{figure}
  \centering
    
 \includegraphics{figures/IntProg6.pdf}
\caption{Swapping colors. }\label{po:fig:2}
\end{figure}




 If we swap colors on $P$, see Figure~\ref{po:fig:2}, we obtain matchings  ${\wt{M}}$ and
 $\color{red}{\wt{M'}}$ with 
 $w({ M}) + w({\color{red} M'}) = w({ \wt{M}})+w({ \color{red} \wt{M'}}) $ and thus
 ${ \wt{M}} \in \eM(w)$.  

 The node $t$ is exposed in ${\wt{M}}$ and $u$ or $v$ is exposed  in
  ${\wt{M}}$. This is a  
contradiction to $u$ and $v$ being shortest distance exposed
  vertices 


\end{proof}



\begin{proof}[Proof of Theorem~\ref{po:thr:18}]

Let $w^Tx\leq\beta$ be a \emph{facet} of $P(G)$, we need to show 
that this facet is of the form
  \begin{enumerate}[i)]
  \item $x_e\geq0$ for some $e \in E$\label{po:item:1}
  \item $\sum_{e \in \delta(v)} x_e\leq1$  for some $v \in V$\label{po:item:2}
  \item $\sum_{e \in E(U)} x_e\leq \lfloor|U|/2\rfloor$ for some $U \in P_{odd}$\label{po:item:3}
  \end{enumerate}
  
  
  To do so, we use the following method: One of the inequalities
  \ref{po:item:1}), \ref{po:item:2}), \ref{po:item:3}) is satisfied with
  equality by each $\chi^M, \,M \in \eM(w)$. This establishes the claim
  since the matching polytope is full-dimensional and a facet is a
  maximal face. 

  




  If $w(e)<0$ for some $e \in E$, then each $M \in \eM(w)$
  satisfies $e \notin M$ and thus satisfies $x_e\geq0$ with equality. 

  Thus we can assume that  $w\geq0$. 
  
  Let $G^*=(V^*,E^*)$ be the graph induced by edges $e$ with $w(e)>0$.  Each $M
  \in \eM(w)$ contains maximum weight matching $M^* = M \cap E^*$ of
  $G^*$ w.r.t.  $w^*$. 

  If $G^*$ is not \emph{connected }, suppose that
  $V^*=V_1\cup V_2$, where $V_1\cap V_2 = \emptyset$ and $V_1,V_2 \neq\emptyset$ and there
  is no edge connecting $V_1$ and $V_2$, then
  $w^Tx\leq\beta$ can be written as the sum of $w_1^Tx\leq\beta_1$ and
  $w_2^Tx\leq\beta_2$, where $\beta_i$ is the maximum weight of a matching in
  $V_i$ w.r.t. $w_i$, $i=1,2$, see Figure~\ref{blobs}. This would also contradict the fact
  that $w^Tx\leq\beta$ is a facet, since it would follow from the previous
  inequalities and thus would  be a redundant
  inequality.
    
  \begin{figure}
    \centering
  
        
     \includegraphics{figures/IntProg7.pdf}
    \caption{$G^*$ is connected. }\label{blobs}
  \end{figure}
   
      Now we can use Lemma~\ref{po:lem:11} for $G^*$. 
      
      \begin{enumerate}[i)]
      \item $\exists v$ such that $\delta(v) \cap M = \emptyset$ for each $M \in
        \eM(w)$. This means that each $M$ in $\eM(w)$ satisfies
        \begin{displaymath}
          \sum_{e \in          \delta(v)} x_e\leq1 \quad \text{ {with equality}}
         \end{displaymath}
      \item $|M\cap E^*| =   \lfloor|V^*| /2\rfloor$ for each $M \in \eM(w)$ and $|V^*|$
        is odd. This means that each $M$ in $\eM(w)$ satisfies
        \begin{displaymath}
           \sum_{e \in E(V^*)          } x_e\leq \lfloor|V^*|/2\rfloor \quad \text{
             {with equality}} 
        \end{displaymath}
      \end{enumerate}

    \end{proof}


\subsection*{Exercises}

\begin{enumerate}
\item 
Let $M\in \setZ^{n\times m}$ be totally unimodular. Prove that the following matrices
are totally unimodular as well:
\begin{enumerate}[i)]
\item $M^T$
\item  $( M \quad I_n )$
\item $(M \quad -M)$
\item $M \cdot (I_n - 2 e_j e_j^T )$ for some $j$
\end{enumerate}\label{i:item:6}

$I_n$ is the $n\times n$ identity matrix, and
$e_j$ is the vector having a $1$ in the $j^{th}$ component, and $0$ in the other components.

% \begin{exercise}
% Recall the definition of a \emph{directed graph} from the last exercise sheet:
% A \emph{directed graph} $D=(V,A)$ is a tuple consisting of a set of \emph{vertices} $V$ and a set of \emph{arcs} $A\subseteq V\times V$.
% Given an arc $a=(u,v)\in A$, the vertex $u$ is called the \emph{tail} of $a$ and $v$ is called the \emph{head} of $a$.

% The \emph{node-arc incidence matrix} $A\in \{-1,0,1\}^{|V|\times|A|}$ of $D$ is defined as follows:
% $$ A(v,a) = \begin{cases} 1,~\text{if $v$ is the head of $a$} \\
%              1,~\text{if $v$ is the head of $a$} \\
% 	     -1,~\text{if $v$ is the tail of $a$}\\
% 	     0,~\text{else}
%             \end{cases}.$$
% Show that $A$ is totally unimodular.
% \end{exercise}

% \begin{solution}
% \end{solution}

\item 
A family $\mathcal{F}$ of subsets of a finite groundset $E$ is
\emph{laminar}, if for all  $C,D\in \mathcal{F}$, one of the following holds:
$$ (i)~C\cap D = \emptyset,~~(ii)~C \subseteq D,~~(iii)~D\subseteq C. $$

Let $\mathcal{F}_1$ and $\mathcal{F}_2$ be two laminar families of the same groundset $E$ and consider its union $\mathcal{F}_1\cup\mathcal{F}_2$.
Define the $|\mathcal{F}_1\cup\mathcal{F}_2|\times |E|$ adjacency matrix $A$ as follows:
For $F\in\mathcal{F}_1\cup\mathcal{F}_2$ and $e\in E$ we have
$A_{F,e}=1$, if $e\in F$ and $A_{F,e}=0$ otherwise.

Show that $A$ is totally unimodular.
\label{i:item:7}


\item 
  Consider the following scheduling problem: Given $n$ tasks with
  periods $p_1, \ldots, p_n\in \setN$, we want to find offsets $x_i\in \setN_0$,
  such that every task $i$ can be executed periodically at times $x_i
  + p_i \cdot k$ for all $k\in \setN_0$.  In other words, for all pairs $i,j$
  of tasks we require $x_i+k\cdot p_i \neq x_j +l\cdot p_j$ for all
  $k,l\in\setN_0$.

  Formulate the problem of finding these offsets as an integer program
  (with zero objective function).
\label{i:item:8}
% \item 
%   Each nonempty  polyhedron $P\subseteq\setR^n$ can be represented as $ P = L + Q$,
%   where  $L\subseteq\setR^n$ is a linear space and $Q\subseteq\setR^n$ is a pointed
%   polyhedron.   \label{po:ex:1}
% \item Let $P\subset\setR^n$ be a polytope and $f:\setR^n\to\setR^m$ a linear map.
%   \begin{enumerate}[i)]
%   \item Show that $f(P)$ is a polytope.
%   \item Let $y\in\setR^m$ be a vertex of $f(P)$. Show that there is a vertex $x\in\setR^n$ of $P$
%     such that $f(x) = y$.
%   \end{enumerate}
% \item 
%   Let $A \in \setR^{m\times n}$ and $b \in \setR^m$ and consider the polyhedron
%   $P = P(A,b)$. Show that $\dim(P) = n - \rank(A^=)$.   \label{po:ex:2}
% \item 
%   \begin{enumerate}[i)]
%   \item 
%   Show that the dimension of each minimal face of a polyhedron $P$ is
%   equal to $n - \rank(A)$. 
%   \item
%   Show that a polyhedron has a vertex if and only if the polyhedron
%   does not contain a line. 
% \end{enumerate}   \label{po:ex:3}
% \item Show that the affine  dimension of the minimal faces of a 
%   polyhedron $P = \{x \in \setR^n \colon Ax\leq b\}$ is invariant. \label{item:19}
%\item 
%. Suppose that $P(A,b) = \{x \in \setR^n \colon Ax\leq b\}$
%  has vertices and that the linear program is bounded. Show how to
%  compute an optimal \emph{vertex} solution of the linear
%  program in polynomial time.    \label{po:ex:4}
% \item  Let $P = \{x \in \setR^n \colon Ax = b, \, x\geq0\}$ be a polyhedron,
%   where $A \in \setR^{m\times n}$ has full row-rank. Let $B_1,B_2$ be two bases
%   such that $|B_1\cap B_2| = m-1$ and suppose that the associated basic
%   solutions $x^*_1$ and $x^*_2$ are feasible. Show that, if
%   $x_1\neq x_2$, then   $\conv\{x_1^*,x_2^*\}$ is a $1$-dimensional face
%   of $P$. \label{item:18} 


\item Show the following: A polyhedron $P \subseteq\setR^n$ with vertices is
  integral, if and only if each vertex is integral. \label{i:item:3}
\item Consider the polyhedron $P = \{ x \in \setR^3 \colon x_1 + 2\,x_2 + 4\,
  x_3 \leq 4, \, x\geq0\}$. Show that this polyhedron is
  integral. \label{i:item:2}
\item Which of these matrices is totally unimodular? Justify your
  answer. \label{i:item:4}
  \begin{displaymath}
    \begin{pmatrix}
      1 & 1  & 0 &  1 & 0 \\
      0 & 0 &  1 &  0 &  1 \\
      0 & 1 &  1 &  1 & 0 \\
      1 & 1 &  0 &  0 & 0 \\
      0 & 1 &  0 &  0 & 1
    \end{pmatrix}
    \quad \quad 
    \begin{pmatrix}
      1&  1& 1& 1& 1 \\
      1& 1& 0& 0& 0\\
      1& 1& 1& 0& 0\\
      1&  1& 1& 1& 0
    \end{pmatrix}
  \end{displaymath}
\item Consider the complete graph $G_n$ with $3$ vertices, i.e., $G =
  (\{1,2,3\}, \binom{3}{2} )$. Is the polyhedron of the linear
  programming relaxation of the vertex-cover integer program integral?
  \label{i:item:5} 
\end{enumerate}




% \bibliographystyle{abbrv}

% \bibliography{mybib,papers,books,my_publications}



%%% Local Variables: 
%%% mode: latex
%%% TeX-master: "lecture"
%%% End: 

%%\include{dualsimplex}
\chapter{Paths, cycles and flows in graphs} 
\label{cha:short-paths-graphs}


Suppose you want to find a shortest path from a given starting point
to a given destination. This is a common scenario in driver assistance
systems (GPS) and can be modeled as one of the most basic
combinatorial optimization problems, the \emph{shortest path problem}.
In this chapter, we introduce directed graphs, shortest paths and
flows in networks. We focus in particular on the maximum-flow problem,
which is a linear program that we solve with direct methods, versus
the simplex method, and analyze the running time of these direct
methods.

% \section{Growth of functions}
% \label{sec:growth-functions}


% In the analysis of algorithms, it is more appropriate to  investigate
% the asymptotic running time 
% of an  algorithm depending on the input and not the
% precise running time itself. We review the $O,\Omega$ and $\Theta$-notation. 

% \begin{definition}[$O,\Omega,\Theta$-notation]
  
%   Let $T,f: \setN \to \setN$ be two functions 
%   \begin{itemize}
%   \item $T(n)$ is  in $O(f(n))$, if there exist positive  constants   $n_o\in \setN$ and  $c\in\setR_{>0}$
%     with $T(n) \leq c \cdot f(n)$ for all $n\geq n_0$. 
%   \item $T(n)$ is in $\Omega(f(n))$, if there exist constants   $n_o\in \setN$ and  $c\in\setR_{>0}$
%     with $T(n) \geq c \cdot f(n)$ for all $n\geq n_0$.  
%   \item $T(n)$ is in $\Theta(f(n))$  if $T(n)$ is both  in $O(f(n))$ and 
%     in $\Omega(f(n))$. 
%   \end{itemize}
% \end{definition}



% \begin{example}
%   The function $T(n)=2n^2 + 3n +1$ is in $O(n^2)$, since for all
%   $x\geq1$ one has $2n^2 + 3n + 1 \leq 6n^2$. Here   $n_0 = 1$ and $c =
%   6$.

%   Similarly $T(n) = \Omega(n^2)$, since for each $n\geq1$ one has $2n^2 + 3n
%   +1\geq n^2$. Thus $T(n)$ is in $\Theta(n^2)$. 
% \end{example}




\section{Graphs}
\label{sec:graphs}


\begin{definition}
  \label{def:3}
  A \emph{directed  graph} is a tuple $G = (V,A)$, where $V$ is a finite set of elements, called the
    \emph{vertices} of $G$ and  $A\subseteq(V\times V)$ is the set of
  \emph{arcs} of $G$. We denote an arc by its two defining nodes $(u,v) \in
  A$. The nodes $u$ and $v$ are called  \emph{tail}  and  \emph{head}
  of the arc $(u,v)$ respectively. 
\end{definition}



\begin{figure}
  \centering 
  \tikzstyle{vertex}=[circle,fill=black!25,minimum size=15pt,inner sep=0pt]
   \tikzstyle{selected vertex} = [circle,fill=black!25,minimum size=15pt,inner sep=0pt]
   \tikzstyle{edge} = [draw,thick,->]
   \tikzstyle{weight} = [font=\small]
   \tikzstyle{selected edge} = [draw,thick, ->,blue!50]
   \tikzstyle{ignored edge} = [draw,line width=5pt,-,black!20]
   
   \begin{tikzpicture}[scale=1.5, auto,swap]
     \foreach \pos/\name in {{(2,1)/u}, {(4,1)/v},
       {(3,0)/x}, {(2,-1)/y}, {(4,-1)/z}}
     \node[vertex] (\name) at \pos {$\name$};
     
     \foreach \source/ \dest  in { u/y, x/u, y/x, z/x, x/y}
     \path[edge] (\source) --   (\dest);

     
     \path  (u) edge [bend right, ->, thick] (v);
     
     \path  (v) edge [bend right, ->, thick] (u);
     
   \end{tikzpicture}
   
  \caption{Example of a directed graph with $5$ nodes and $7$ arcs.}
\label{ex:graph}
\end{figure}




\begin{definition}[Walk, path, distance]
  \label{f:def:6}
  A \emph{walk} is a sequence of the form
  $$P=(v_0,a_1,v_1,\ldots,v_{m-1},a_m,v_m),$$ where  $a_i =
  (v_{i-1},v_i)\in A$ for $i=1,\ldots,m$. If the nodes $v_0,\ldots,v_m$ are all
  different, then $P$ is a \emph{path}. The \emph{length } of $P$ is
  $m$. The \emph{distance} of two nodes  $u$ and $v$ is the length of
  a shortest path from $u$ to $v$. It is denoted by $d(u,v)$. 
\end{definition}


\begin{example}
  The following is a walk and a path of the graph in
  Figure~\ref{ex:graph}. 
  \begin{displaymath}
    \begin{array}{c}
     z,(z,x),x,(x,u),u,(u,v),v,(v,u),u,(u,y),y,(y,x),x\\
     y,(y,x),x,(x,u),u,(u,v),v\\
    \end{array}
  \end{displaymath}
\end{example}


% \begin{figure}
%   \centering
%  $\psmatrix[colsep=0.2cm,mnode=circle]
% v_0&v_1 & v_3 & v_4 & v_5 & v_6 &   & v_0&v_1 & v_3 & v_4 & v_5 & v_6 \\
% \ncline{->}{1,1}{1,2}
% \ncline{->}{1,2}{1,3}
% \ncarc[arcangle=30]{->}{1,3}{1,2}
% \ncarc[arcangle=30]{->}{1,5}{1,1}
% \ncline{->}{1,3}{1,4}  
% \ncline{->}{1,4}{1,5}  
% \ncline{->}{1,5}{1,6}  
% \ncline{->}{1,8}{1,9}
% \ncline{->}{1,9}{1,10}
% \ncline{->}{1,10}{1,11}  
% \ncline{->}{1,11}{1,12}
% \ncline{->}{1,12}{1,13}
% \endpsmatrix
% $
%  \caption{Example of a walk and a path. }
% \end{figure} Das Beispiel 



\section{Representing graphs and computing the distance of two nodes}

\label{sec:repr-graphs-comp}

We represent a graph with $n$ vertices $v_1,\ldots,v_n$  as an array $A[v_1,\ldots,v_n]$, 
where the entry $A[v_i]$ is a pointer to a linked list of vertices, 
the \emph{neighbours of $v_i$}. $N(v_i) = \{ u \in V \colon (v_i,u) \in A\}$.

\begin{figure}
  \centering
 \includegraphics{figures/adjacency.pdf}
  \caption{Adjacency list representation of the graph in
    Figure~\ref{ex:graph}. } 
\end{figure}

\subsection{Breadth-first search}
\label{subsec:BFs}

We next describe a very basic algorithm that  computes the distances 
from a designated node  $s \in V$ to all other nodes. 
The \emph{distance} from $s$ to $v$ is denoted by $d(s,v)$. It is the
smallest integer $i$ such that there exists a path from $s$ to $v$ of
length $i$.  If there does not exist such a path, then  $s$ and $v$
are \emph{not connected} and we define $d(s,v) = \infty$. For $ i \in
\setN_0$,  $V_i\subseteq V$ denotes  the set of vertices that have distance $i$
from $s$.  Notice that $V_0 = \{s\}$. 

\begin{lemma}
  \label{lem:15}
  For $i=1,\ldots,n-1$, the set   $V_{i}$ is equal to the set of 
  vertices $v \in V \backslash (V_0\cup\cdots\cup V_{i-1})$ such that there exists an arc 
  $(u,v) \in A$ with $u \in V_{i-1}$.
\end{lemma}


\begin{proof}
  Suppose that $v \notin V_0\cup\cdots\cup V_{i-1}$ and there exists an arc $uv\in A$
  with $u \in V_{i-1}$. Since $u \in V_{i-1}$, there exists a path
  $s,a_1,v_1,a_2,v_2,\ldots,a_{i-1},u$ of length $i-1$ from $s$ to $u$. The
  sequence  $s,a_1,v_1,a_2,v_2,\ldots,a_i,u,uv,v$ is a path of length
  $i$ from $s$ to $v$ and thus $v \in V_{i}$. 

  If, on the other hand, $v \in V_{i}$, then there exists a path 
  \begin{displaymath}
    s,a_1,v_1,\ldots,a_{i-1},u,a_{i},v
  \end{displaymath}
  of length $i$ from $s$ to $v$. We need to show that $u \in V_{i-1}$
  holds. Clearly, since there exists a path of length $i-1$ from $s$ to
  $u$, one has $u \in V_j$ with $j\leq i-1$. If $j<i-1$, then there exists
  a path $s,a_1',v_1',\ldots,a_j',u$ of length $j$ which can be extended
  to a path of length $j+1<i$ from $s$ to $v$
  \begin{displaymath}
    s,a_1',v_1',\ldots,a_j',u,a_{i},v
  \end{displaymath}
  which contradicts $v\in V_i$.  \qed

\end{proof}


The  \emph{breadth-first search algorithm} is an implementation of
Lemma~\ref{lem:15}. 
The algorithm maintains arrays 
\begin{displaymath}
  \begin{array}{c}
    D[v_1=s,v_2,\ldots,v_n]\\
    \pi[v_1=s,v_2,\ldots,v_n]\\
  \end{array}
\end{displaymath}
and a queue $Q$ that contains only $s$ in the beginning. The array
$D$ contains at termination of the algorithm the distances from $s$ to
all other nodes and is initialized with $[0,\infty,\ldots,\infty]$.  
The array $\pi$ contains predecessor information for
shortest paths, in other words, when the algorithm terminates, $\pi[v]
= u$, where $uv$ is an arc and $D[u]+1 = D[v]$. The array $\pi$ is
initialized with $[0,\ldots,0]$.


After this initialization, the algorithm proceeds as follows. 

\begin{tabbing}
  {\bf while} \= $Q \neq \emptyset$  \\
              \> $u := head(Q)$ \\
              \> {\bf for} \= each $v \in \delta^+(u)$ \\
              \>           \> {\bf if} \= ($D[v]=\infty$) \\
              \>           \>          \> $\pi[v]:=u$ \\ 
              \>           \>          \> $D[v]:=D[u]+1$ \\
              \>           \>          \> $enqueue(Q,v)$ \\
              \> $dequeue(Q)$ 
\end{tabbing}

Here the function $head(Q)$ returns the next element in the queue and
$dequeue(Q)$ removes the first element of $Q$, while $enqueue(Q,v)$
adds $v$ to the queue $Q$ as last element.



\begin{lemma}
  \label{lem:17}
  The breadth-first search algorithm assigns \emph{distance labels} $D$
  correctly. 
\end{lemma}


\begin{proof}  
  We show the following claim by induction on $i \in \{0,\ldots,n-1\}$. 
  \begin{quote}
    For each $i \in \{1,\ldots,n-1\}$ there exists a point in time where:
    \begin{enumerate}[i)]
    \item     $Q$ contains precisely  the elements  of $V_i$
      \label{item:f1}
    \item for each $v \in V_i$,  $D[v] = d(s,v)$ \label{item:f2}
    \item for each $v \in V_i$  one has   $\pi[v]v$
      is an arc and $\pi[v] \in V_{i-1}$. 
    \end{enumerate}
  \end{quote}
Once this claim is shown, the lemma follows, because the labels $D[v]$
and $\pi[v]$ are only changed once, if at all, from $\infty$ or $0$ 
to an integer
or a vertex respectively. 

Since $V_0 = \{s\}$ and since  $Q = [s]$ and $D[s] = 0$ after the initialization, the
claim holds for $ i = 0$. Suppose $i>0$. By the induction hypothesis,
there is a point in time, where  $Q$ contains precisely $V_{i-1}$. By
Lemma~\ref{lem:15}, after the last element of $V_{i-1}$ is dequeued
$Q$ contains precisely the elements in $V_i$. Also, since $D[u] =
d(s,u) = i-1$ for all $u \in V_{i-1}$, we have for each $v \in V_i$ that  $D[v]
= d(s,v)=i$. Also $\pi[v]v$ is an arc, by virtue of the algorithm, and
$\pi[v] \in V_{i-1}$.  \qed 
\end{proof}

% \begin{proof}
%   Let $v \in V$. We show by induction on $d(s,v)$ that the labels are
%   correctly assigned. 

%   If $d(s,v)=0$, then $s=v$ and $D[v]=0$.  If $d(s,v)=1$, then $v$ is
%   a neighbor of $s$ and $D[v]=1$ is set correctly in the first
%   iteration of the {\bf while} loop. 
  
%   Let $d(s,v)>1$. Then there exists a $u\neq s,v$ with $d(s,u)=d(s,v)-1$
%   and $uv \in A$. By induction, the label $D[u]=d(s,u)$ is set
%   correctly by the breadth-first-search algorithm. Also, since the
%   breadth-first-search algorithm computes for $v$ a path of length
%   $D[v]$ from $s$ to $v$, the node $v$ receives a label which is
%   greater than or equal to $d(s,v)$. If we consider the sequence (over
%   time) of assigned labels, that breadth-first-search is assigning,
%   then it is easy to see that this sequence is monotonously
%   increasing, see exercise~\ref{item:16}.  The node $v$ is thus
%   explored at the latest, when $u$ is dequeued. This shows that the
%   label of $v$, $D[v]$ is assigned correctly.  \qed
% \end{proof}


\begin{definition}[Directed tree]
  \label{def:t6}
  A \emph{directed  tree} is a directed graph $T = (V,A)$ with 
  $|A| = |V| -1$ and containing a node $r \in V$ such that there 
  exists a path from $r$ to all other nodes of $T$. 
\end{definition}



\begin{lemma}
  \label{lem:18}
  Consider the arrays $D$ and $\pi$ after the termination of the
  breadth-first-search algorithm. 
  The graph $T = (V',A')$ with $V' = \{v \in V \colon D[v]<\infty\}$ and $A' =
  \{ \pi(v) v \colon 1\leq D[v]<\infty\}$ is a tree. 
\end{lemma}


\begin{proof}
  Clearly, $|A'| = |V'| -1$. For any $i \in \{1,\ldots,n-1\}$, by
  backtracking the $\pi$-labels from any $v \in V_i$, we will eventually
  reach $s$.
\end{proof}

\begin{definition}
  \label{def:t7}
  The tree $T$ from lemma~\ref{lem:18} is the \emph{shortest-path-tree}
   of the (unweighted) directed graph $G = (V,A)$. 
\end{definition}


% %% Adjacency matrix of graph
% %% \  a  b  c  d  e  f  g
% %% a  x  7     5
% %% b  7  x  8  9  7
% %% c     8  x     5
% %% d  5  9     x 15  6
% %% e     7  5 15  x  8  9
% %% f           6  8  x 11
% %% g              9  11 x

  \tikzstyle{vertex}=[circle,fill=black!25,minimum size=15pt,inner sep=0pt]
  \tikzstyle{selected vertex} = [vertex, fill=red!24]
  \tikzstyle{edge} = [draw,thick,->]
  \tikzstyle{weight} = [font=\small]
  \tikzstyle{selected edge} = [draw,thick, ->,blue!50]
  \tikzstyle{ignored edge} = [draw,line width=5pt,-,black!20]

 \begin{figure}
   \begin{subfigure}[t]{0.45\textwidth}
    \begin{tikzpicture}[scale=1, auto,swap]
     % Draw a 7,11 network
     % First we draw the vertices
     \foreach \pos/\name in {{(2,1)/a}, {(4,1)/d},
       {(3,0)/b}, {(2,-1)/c}, {(4,-1)/e}}
     \node[vertex] (\name) at \pos {$\name$};
     
     \foreach \pos/\name in {{(0,0)/s}}
     \node[selected vertex] (\name) at \pos {$\name$};
     % Connect vertices with edges and draw weights
   
     \foreach \source/ \dest  in { s/a, s/c,
       b/s, b/d,d/e,
       b/e,e/c, a/b}
     \path[edge] (\source) --   (\dest);
     % Start animating the vertex and edge selection. 
   %   \foreach \vertex / \fr in {d/1,a/2,f/3,b/4,e/5,c/6,g/7}
 %         \path<\fr-> node[selected vertex] at (\vertex) {$\vertex$};
     % For convenience we use a background layer to highlight edges
     % This way we don't have to worry about the highlighting covering
     % weight labels. 
    %  \begin{pgfonlayer}{background}
 %         \pause
 %         \foreach \source / \dest in {d/a,d/f,a/b,b/e,e/c,e/g}
 %             \path<+->[selected edge] (\source.center) -- (\dest.center);
 %         \foreach \source / \dest / \fr in {d/b/4,d/e/5,e/f/5,b/c/6,f/g/7}
 %             \path<\fr->[ignored edge] (\source.center) -- (\dest.center);
 %     \end{pgfonlayer}
   \end{tikzpicture}
 \caption{The breadth-first search algorithm starts with the queue
   $Q=[s]$.  The distance labels for $[s,a,b,c,d,e]$ are 
   $[0,\infty,\infty,\infty,\infty,\infty]$ respectively.}
\end{subfigure}
\hfill 
\begin{subfigure}[t]{0.45\textwidth}
 \begin{tikzpicture}[scale=1, auto,swap]
     % Draw a 7,11 network
     % First we draw the vertices
     \foreach \pos/\name in { {(4,1)/d},
       {(3,0)/b},  {(4,-1)/e}}
     \node[vertex] (\name) at \pos {$\name$};
     
     \foreach \pos/\name in {{(2,1)/a},{(2,-1)/c}, {(0,0)/s}}
     \node[selected vertex] (\name) at \pos {$\name$};
     % Connect vertices with edges and draw weights
   
     \foreach \source/ \dest  in { s/a, s/c,
       b/s, b/d,d/e,
       b/e,e/c, a/b}
     \path[edge] (\source) --   (\dest);
 \end{tikzpicture}
\caption{After the first iteration of the {\bf while} loop the queue
  is $Q = [a,c]$ and the distance labels are $[0,1,\infty,1,\infty,\infty]$ respectively. }
\end{subfigure}

\begin{subfigure}[t]{0.45\textwidth}
  \begin{tikzpicture}[scale=1, auto,swap]
     \foreach \pos/\name in { {(4,1)/d},
        {(4,-1)/e}}
     \node[vertex] (\name) at \pos {$\name$};
     
     \foreach \pos/\name in {{(2,1)/a}, {(3,0)/b}, {(2,-1)/c}, {(0,0)/s}}
     \node[selected vertex] (\name) at \pos {$\name$};
     % Connect vertices with edges and draw weights
   
     \foreach \source/ \dest  in { s/a, s/c,
       b/s, b/d,d/e,
       b/e,e/c, a/b}
     \path[edge] (\source) --   (\dest);
  \end{tikzpicture} 
\caption{After the second iteration of the {\bf while} loop the queue
 is $Q = [c,b]$ and the distance labels are $[0,1,2,1,\infty,\infty]$ respectively. }   
\end{subfigure}
\hfill
\begin{subfigure}[t]{0.45\textwidth}
  \begin{tikzpicture}[scale=1, auto,swap]
     % Draw a 7,11 network
     % First we draw the vertices
     \foreach \pos/\name in { {(4,1)/d},
        {(4,-1)/e}}
     \node[vertex] (\name) at \pos {$\name$};
     
     \foreach \pos/\name in {{(2,1)/a}, {(3,0)/b}, {(2,-1)/c}, {(0,0)/s}}
     \node[selected vertex] (\name) at \pos {$\name$};
     % Connect vertices with edges and draw weights
   
     \foreach \source/ \dest  in { s/a, s/c,
       b/s, b/d,d/e,
       b/e,e/c, a/b}
     \path[edge] (\source) --   (\dest);
  \end{tikzpicture}    
\caption{After the third iteration of the {\bf while} loop the queue
 is $Q = [b]$ and the distance labels are unchanged, since $c$ does
 not have any neighbors.  }
\end{subfigure}

\begin{subfigure}[t]{0.45\textwidth}
  \begin{tikzpicture}[scale=1, auto,swap]
       
   \foreach \pos/\name in {{(4,1)/d},
     {(4,-1)/e}, {(2,1)/a}, {(3,0)/b}, {(2,-1)/c}, {(0,0)/s}}
   \node[selected vertex] (\name) at \pos {$\name$};
   % Connect vertices with edges and draw weights
   
   
     \foreach \source/ \dest  in { s/a, s/c,
       b/s, b/d,d/e,
       b/e,e/c, a/b}
     \path[edge] (\source) --   (\dest);
  \end{tikzpicture}    
\caption{{After the fourth iteration of the {\bf while} loop the queue
 is $Q = [d,e]$ and the distance labels are  $[0,1,2,1,3,3]$ respectively. }}
\end{subfigure}
\hfill 
\begin{subfigure}[t]{0.45\textwidth}
  \begin{tikzpicture}[scale=1, auto,swap]
    
     \foreach \pos/\name in {{(4,1)/d},
        {(4,-1)/e}, {(2,1)/a}, {(3,0)/b}, {(2,-1)/c}, {(0,0)/s}}
     \node[selected vertex] (\name) at \pos {$\name$};
     % Connect vertices with edges and draw weights
   
     \foreach \source/ \dest  in {b/s, d/e, e/c}
     \path[edge] (\source) --   (\dest);

     
     \foreach \source/ \dest  in { s/a, s/c, b/d,  b/e, a/b}
     \path[selected edge] (\source) --   (\dest);

  \end{tikzpicture}    
\caption{After the fifth and sixth iteration of the {\bf while} loop 
 the queue is empty $Q = []$ and the distance labels remain unchanged. 
 The blue edges denote the shortest path tree.}
\end{subfigure}
  


\caption{An example-run of breadth-first search}\label{g:fig:4}
\end{figure}





\begin{theorem}
  \label{f:thr:26}
  The breath-first-search algorithm runs in time $O( |V| + |A|)$. 
\end{theorem}

\begin{proof}
  Each vertex is queued and dequeued at most once. These queuing
  operations take constant time each. Thus queuing and dequeuing costs
  $O( |V| )$ in total. 

  When a vertex $u$ is dequeued, its neighbors are
  inspected and the operations in the {\bf if} statement cost constant
  time each. Thus one has an additional cost of $O( |A | )$, since
  these constant-time  operations are carried out for each arc $a \in
  A$. 
  \qed
\end{proof}

\section{Shortest Paths}
\label{sec:shortest-paths}

\begin{definition}[Cycle]
  A walk in which starting node and end-node agree is called a
  \emph{cycle}. 
\end{definition}

Suppose we are given a directed graph $D=(V,A)$ and a length function
$c:~A\longrightarrow\setR$. The \emph{length} of a walk $W$ is defined as 
\begin{displaymath}
  c(W) = \sum_{\substack{a \in A\\ a\in W}}  c(a). 
\end{displaymath}
%
We now study how to determine a shortest path in a weighted directed
graph
$G$ efficiently, in case of the absence of cycles of negative length (such cycles are called \emph{negative cycles}). 

\begin{theorem}
  \label{f:thr:1}
  Suppose that each cycle in $D$ has non-negative length and suppose
  there exists an $s-t$-walk in $D$. Then there exists a path
  connecting $s$ with $t$ which has minimum length among all walks
  connecting $s$ and $t$. 
\end{theorem}

\begin{proof}
  If there exists an $s-t$-walk, then there exists an $s-t$-path. Since
  the number of arcs in a path is at most $|V| - 1 $, there must exist a
  shortest \emph{path}  $P$  connecting $s$ and $t$. We claim that
  $c(P)\leq c(W)$ for all $s-t$-walks $W$. Suppose that there exists
  an $s-t$-walk $W$ with $c(W)<c(P)$. Then let $W$ be such a walk with
  a minimum number of arcs. Clearly $W$ contains a cycle $C$.  Since the
  cycle has non-negative length, then it can be removed from $W$ to
  obtain a walk whose length is at most $c(W)$ and whose number of
  arcs is strictly less than $|W|$. This is a contradiction to the 
  minimality of the number of arcs in $W$.
  \qed
\end{proof}

We use the notation $|W|,|C|,|P|$ to denote the number of arcs in a
walk $W$, a cycle $C$ or a path $P$. 


As a conclusion we can note here: 
\begin{quote}
  If there do not exist negative cycles in $D$, and $s$ and $t$ are
  connected, then there exists a shortest walk traversing at most 
  $|V|  - 1$ arcs.
\end{quote}


\subsection*{The Bellman-Ford algorithm}

Let $n=|V|$. We calculate functions $f_0,f_1,\ldots,f_n:V\longrightarrow\setR\cup\{\infty\}$
successively by the following rule. 

\begin{enumerate}[i)]
\item $f_0(s) = 0$, $f_0(v) = \infty$ for all $v \neq s$ 
\item For $k<n$ if $f_k$ has been found, compute 
  \begin{displaymath}
    \displaystyle f_{k+1}(v) = \min\{f_k(v), \min_{(u,v)\in A}\{f_k(u)+c(u,v)\} \}  
  \end{displaymath}
  for all $v \in V$. 
\end{enumerate}





% %% Adjacency matrix of graph
% %% \  a  b  c  d  e  f  g
% %% a  x  7     5
% %% b  7  x  8  9  7
% %% c     8  x     5
% %% d  5  9     x 15  6
% %% e     7  5 15  x  8  9
% %% f           6  8  x 11
% %% g              9  11 x

  \tikzstyle{vertex}=[circle,fill=black!25,minimum size=15pt,inner sep=0pt]
  \tikzstyle{selected vertex} = [vertex, fill=red!24]
  \tikzstyle{edge} = [draw,thick,->]
  \tikzstyle{weight} = [font=\small]
  \tikzstyle{selected edge} = [draw,thick, ->,blue!50]
  \tikzstyle{ignored edge} = [draw,line width=5pt,-,black!20]

 \begin{figure}
 
      \begin{subfigure}[t]{0.45\textwidth}
     \begin{tikzpicture}[scale=1, auto,swap]
     % Draw a 7,11 network
     % First we draw the vertices
     \foreach \pos/\name in {{(2,1)/a}, {(4,1)/d},
       {(3,0)/b}, {(2,-1)/c}, {(4,-1)/e}}
     \node[vertex] (\name) at \pos {$\name$};
     
     \foreach \pos/\name in {{(0,0)/s}}
     \node[selected vertex] (\name) at \pos {$\name$};
     % Connect vertices with edges and draw weights
   
     \foreach \source/ \dest/ \weight   in { s/a/3, s/c/4,
       b/c/-2, b/d/3 ,d/e/-3,
       b/e/2,e/c/2, a/b/1}
     \path[edge] (\source) -- node[weight] {$\weight$}   (\dest);
     % Start animating the vertex and edge selection. 
   %   \foreach \vertex / \fr in {d/1,a/2,f/3,b/4,e/5,c/6,g/7}
 %         \path<\fr-> node[selected vertex] at (\vertex) {$\vertex$};
     % For convenience we use a background layer to highlight edges
     % This way we don't have to worry about the highlighting covering
     % weight labels. 
    %  \begin{pgfonlayer}{background}
 %         \pause
 %         \foreach \source / \dest in {d/a,d/f,a/b,b/e,e/c,e/g}
 %             \path<+->[selected edge] (\source.center) -- (\dest.center);
 %         \foreach \source / \dest / \fr in {d/b/4,d/e/5,e/f/5,b/c/6,f/g/7}
 %             \path<\fr->[ignored edge] (\source.center) -- (\dest.center);
 %     \end{pgfonlayer}
   \end{tikzpicture}
     
     \caption{
   The algorithm is initialized with distance labels for
   $s,a,b,c,d,e$ being $[0,\infty,\infty,\infty,\infty,\infty]$ respectively 
}
   \end{subfigure}\hfill 
      \begin{subfigure}[t]{0.45\textwidth}
      \begin{tikzpicture}[scale=1, auto,swap]
     % Draw a 7,11 network
     % First we draw the vertices
     \foreach \pos/\name in {{(2,1)/a}, {(4,1)/d},
       {(3,0)/b}, {(2,-1)/c}, {(4,-1)/e}}
     \node[vertex] (\name) at \pos {$\name$};
     
     \foreach \pos/\name in {{(0,0)/s}}
     \node[selected vertex] (\name) at \pos {$\name$};
     % Connect vertices with edges and draw weights
   
     \foreach \source/ \dest/ \weight   in {b/c/-2, b/d/3 ,d/e/-3,b/e/2,e/c/2, a/b/1}
     \path[edge] (\source) -- node[weight] {$\weight$}   (\dest);
     
     
     \foreach \source/ \dest/ \weight   in { s/a/3, s/c/4      }
     \path[selected edge] (\source) -- node[weight] {$\weight$}   (\dest);

     % Start animating the vertex and edge selection. 
   %   \foreach \vertex / \fr in {d/1,a/2,f/3,b/4,e/5,c/6,g/7}
 %         \path<\fr-> node[selected vertex] at (\vertex) {$\vertex$};
     % For convenience we use a background layer to highlight edges
     % This way we don't have to worry about the highlighting covering
     % weight labels. 
    %  \begin{pgfonlayer}{background}
 %         \pause
 %         \foreach \source / \dest in {d/a,d/f,a/b,b/e,e/c,e/g}
 %             \path<+->[selected edge] (\source.center) -- (\dest.center);
 %         \foreach \source / \dest / \fr in {d/b/4,d/e/5,e/f/5,b/c/6,f/g/7}
 %             \path<\fr->[ignored edge] (\source.center) -- (\dest.center);
 %     \end{pgfonlayer}
   \end{tikzpicture}
     
     \caption{
   After the first iteration the labels are $[0,3,\infty,4,\infty,\infty]$
}
   \end{subfigure}




      \begin{subfigure}[t]{0.45\textwidth}
      \begin{tikzpicture}[scale=1, auto,swap]
     % Draw a 7,11 network
     % First we draw the vertices
     \foreach \pos/\name in {{(2,1)/a}, {(4,1)/d},
       {(3,0)/b}, {(2,-1)/c}, {(4,-1)/e}}
     \node[vertex] (\name) at \pos {$\name$};
     
     \foreach \pos/\name in {{(0,0)/s}}
     \node[selected vertex] (\name) at \pos {$\name$};
     % Connect vertices with edges and draw weights
   
     \foreach \source/ \dest/ \weight   in {b/c/-2, b/d/3 ,d/e/-3,b/e/2,e/c/2}
     \path[edge] (\source) -- node[weight] {$\weight$}   (\dest);
     
     
     \foreach \source/ \dest/ \weight   in { s/a/3, s/c/4, a/b/1      }
     \path[selected edge] (\source) -- node[weight] {$\weight$}   (\dest);

     % Start animating the vertex and edge selection. 
   %   \foreach \vertex / \fr in {d/1,a/2,f/3,b/4,e/5,c/6,g/7}
 %         \path<\fr-> node[selected vertex] at (\vertex) {$\vertex$};
     % For convenience we use a background layer to highlight edges
     % This way we don't have to worry about the highlighting covering
     % weight labels. 
    %  \begin{pgfonlayer}{background}
 %         \pause
 %         \foreach \source / \dest in {d/a,d/f,a/b,b/e,e/c,e/g}
 %             \path<+->[selected edge] (\source.center) -- (\dest.center);
 %         \foreach \source / \dest / \fr in {d/b/4,d/e/5,e/f/5,b/c/6,f/g/7}
 %             \path<\fr->[ignored edge] (\source.center) -- (\dest.center);
 %     \end{pgfonlayer}
   \end{tikzpicture}

     
     \caption{
   After the second iteration the labels are $[0,3,4,4,\infty,\infty]$
}
   \end{subfigure}\hfill 
      \begin{subfigure}[t]{0.45\textwidth}
      \begin{tikzpicture}[scale=1, auto,swap]
     % Draw a 7,11 network
     % First we draw the vertices
     \foreach \pos/\name in {{(2,1)/a}, {(4,1)/d},
       {(3,0)/b}, {(2,-1)/c}, {(4,-1)/e}}
     \node[vertex] (\name) at \pos {$\name$};
     
     \foreach \pos/\name in {{(0,0)/s}}
     \node[selected vertex] (\name) at \pos {$\name$};
     % Connect vertices with edges and draw weights
   
     \foreach \source/ \dest/ \weight   in {d/e/-3,e/c/2, s/c/4}
     \path[edge] (\source) -- node[weight] {$\weight$}   (\dest);
     
     
     \foreach \source/ \dest/ \weight   in {b/e/2,b/c/-2, b/d/3 , s/a/3, a/b/1      }
     \path[selected edge] (\source) -- node[weight] {$\weight$}   (\dest);

     % Start animating the vertex and edge selection. 
   %   \foreach \vertex / \fr in {d/1,a/2,f/3,b/4,e/5,c/6,g/7}
 %         \path<\fr-> node[selected vertex] at (\vertex) {$\vertex$};
     % For convenience we use a background layer to highlight edges
     % This way we don't have to worry about the highlighting covering
     % weight labels. 
    %  \begin{pgfonlayer}{background}
 %         \pause
 %         \foreach \source / \dest in {d/a,d/f,a/b,b/e,e/c,e/g}
 %             \path<+->[selected edge] (\source.center) -- (\dest.center);
 %         \foreach \source / \dest / \fr in {d/b/4,d/e/5,e/f/5,b/c/6,f/g/7}
 %             \path<\fr->[ignored edge] (\source.center) -- (\dest.center);
 %     \end{pgfonlayer}
   \end{tikzpicture}
     
     \caption{
   After the third iteration the labels are $[0,3,4,2,7,6]$
}
   \end{subfigure}





      \begin{subfigure}[t]{0.45\textwidth}
      \begin{tikzpicture}[scale=1, auto,swap]
     % Draw a 7,11 network
     % First we draw the vertices
     \foreach \pos/\name in {{(2,1)/a}, {(4,1)/d},
       {(3,0)/b}, {(2,-1)/c}, {(4,-1)/e}}
     \node[vertex] (\name) at \pos {$\name$};
     
     \foreach \pos/\name in {{(0,0)/s}}
     \node[selected vertex] (\name) at \pos {$\name$};
     % Connect vertices with edges and draw weights
   
     \foreach \source/ \dest/ \weight   in { b/e/2, e/c/2, s/c/4}
     \path[edge] (\source) -- node[weight] {$\weight$}   (\dest);
     
     
     \foreach \source/ \dest/ \weight   in {d/e/-3,b/c/-2, b/d/3 , s/a/3, a/b/1      }
     \path[selected edge] (\source) -- node[weight] {$\weight$}   (\dest);

     % Start animating the vertex and edge selection. 
   %   \foreach \vertex / \fr in {d/1,a/2,f/3,b/4,e/5,c/6,g/7}
 %         \path<\fr-> node[selected vertex] at (\vertex) {$\vertex$};
     % For convenience we use a background layer to highlight edges
     % This way we don't have to worry about the highlighting covering
     % weight labels. 
    %  \begin{pgfonlayer}{background}
 %         \pause
 %         \foreach \source / \dest in {d/a,d/f,a/b,b/e,e/c,e/g}
 %             \path<+->[selected edge] (\source.center) -- (\dest.center);
 %         \foreach \source / \dest / \fr in {d/b/4,d/e/5,e/f/5,b/c/6,f/g/7}
 %             \path<\fr->[ignored edge] (\source.center) -- (\dest.center);
 %     \end{pgfonlayer}
   \end{tikzpicture}

     
     \caption{
   After the fourth iteration the labels are $[0,3,4,2,7,4]$
}
   \end{subfigure}\hfill 
      \begin{subfigure}[t]{0.45\textwidth}
      \begin{tikzpicture}[scale=1, auto,swap]
     % Draw a 7,11 network
     % First we draw the vertices
     \foreach \pos/\name in {{(2,1)/a}, {(4,1)/d},
       {(3,0)/b}, {(2,-1)/c}, {(4,-1)/e}}
     \node[vertex] (\name) at \pos {$\name$};
     
     \foreach \pos/\name in {{(0,0)/s}}
     \node[selected vertex] (\name) at \pos {$\name$};
     % Connect vertices with edges and draw weights
   
     \foreach \source/ \dest/ \weight   in { b/e/2, e/c/2, s/c/4}
     \path[edge] (\source) -- node[weight] {$\weight$}   (\dest);
     
     
     \foreach \source/ \dest/ \weight   in {d/e/-3,b/c/-2, b/d/3 , s/a/3, a/b/1      }
     \path[selected edge] (\source) -- node[weight] {$\weight$}   (\dest);

     % Start animating the vertex and edge selection. 
   %   \foreach \vertex / \fr in {d/1,a/2,f/3,b/4,e/5,c/6,g/7}
 %         \path<\fr-> node[selected vertex] at (\vertex) {$\vertex$};
     % For convenience we use a background layer to highlight edges
     % This way we don't have to worry about the highlighting covering
     % weight labels. 
    %  \begin{pgfonlayer}{background}
 %         \pause
 %         \foreach \source / \dest in {d/a,d/f,a/b,b/e,e/c,e/g}
 %             \path<+->[selected edge] (\source.center) -- (\dest.center);
 %         \foreach \source / \dest / \fr in {d/b/4,d/e/5,e/f/5,b/c/6,f/g/7}
 %             \path<\fr->[ignored edge] (\source.center) -- (\dest.center);
 %     \end{pgfonlayer}
   \end{tikzpicture}

     
     \caption{
   After the fifth  iteration the labels are unchanged. The shortest
   path distances have been computed. 
}
   \end{subfigure}
   
 \caption{An example-run of the Bellman-Ford algorithm. The blue
     edges represent the tree whose paths have the corresponding lengths.} \label{f:fig:5}

 \end{figure}




The following theorem shows us that the Bellman-Ford algorithm is a 
method to compute minimum length walks.


\begin{theorem}
  \label{thr:1f2}
  For each $k=0,\ldots,n$ and for each $v \in V$ 
  \begin{displaymath}
    f_k(v)  = \min\{ c(P) \colon  P \text{ is an } s-v\text{-walk traversing
    at most } k \text{ arcs}\}.
  \end{displaymath}
\end{theorem}

\begin{theorem}
   \label{bobo:f:thr:1}
   Given a directed graph $D=(V,A)$, $s \in V$ and a length 
   function $c:A \to \setR$, one has $f_n = f_{n-1}$ if and only 
   if $D$ does not have a cycle 
   of negative length that is reachable from $s$.
\end{theorem}

\begin{proof}
   $(\Leftarrow)$ Suppose there exists $t \in V$ such that 
   $f_n(t) < f_{n-1}(t) < \infty$. This implies that the shortest 
   $s-t$-walk traversing at most $n$ arcs (call it $W$) traverses 
   exactly $n$ arcs and thus contains a cycle (call it $C$). Consider 
   the $s-t$-walk $W'$ obtained by eliminating $C$ from $W$. 
   $W'$ traverses at most $n-1$ arcs and thus we have $c(W') > c(W)$.
   Since $c(W) = c(W') + c(C)$, we have that $c(C) < 0$.
   
   $(\Rightarrow)$ Let $C = v_0,v_1,\cdots,v_k,v_0$ be a cycle 
   reachable from $s$. Notice that 
   $f_{n-1}(v_i)<\infty$ $\forall 0 \leq i \leq k$. 
   We can show that $C$ is a non-negative cycle by these simple 
   calculations (notice that the node indices $i$ and $i+1$ are 
   considered modulo $k$ in the following sums).
   \begin{displaymath}
      0 = \sum_{i = 0}^k f_n(v_{i+1}) - f_n(v_i) \leq
      \sum_{i = 0}^k c(v_i,v_{i+1}) = c(C)
   \end{displaymath}
\end{proof}

Theorem~\ref{bobo:f:thr:1} can be generalized for every $n \geq |V|$.
This allows us to obtain the following corollary.

\begin{corollary}
  \label{co:7}
  If $D = (V,A) $ does not contain negative cycles w.r.t. $c$, then
  $f_n(v)$ is equal to the length of a shortest $s-v$-path. The
  numbers $f_n(v)$ can be computed in time $O( |V| \cdot |A| )$. 
\end{corollary}

Notice that by using theorem~\ref{bobo:f:thr:1} 
we can see if a directed graph $D(V,A)$ has negative cycles in the 
following way. We obtain a new graph $D'=(V',A')$ add a node $s$ to 
$D$ and we connect $s$ to all other vertices with an outgoing arc 
of weight $0$. Then, we apply Bellman-Ford algorithm to $D'$ with 
starting node $s$. Notice that $D$ contains a negative cycle, if 
and only if $D'$ contains a negative cycle (since every cycle in $D$ 
is reachable from $s$ in $D'$). Thus $D$ contains a 
negative cycle if and only if there exists $t \in V$ such that 
$f_n(t) < f_{n-1}(t)$ in $D'$.
This gives us the following corollary.

\begin{corollary}
  \label{co:9}
  In time $O(|V| \cdot |A| )$  one can test whether $D = (V,A)$ has a
  negative cycle w.r.t. $c$ and eventually return one. 
\end{corollary}

\begin{proof}
   By our previous words we know that we simply have to apply 
   Bellman-Ford algorithm to the graph $D'$. The corollary follows 
   by the fact that $|V'| = |V|+1$ and $|A'| = |A| + |V|$ and by 
   using corollary~\ref{co:7}.
\end{proof}

\section{Maximum $s-t$-flows}
\label{sec:maximum-s-t}

We now turn our attention to a linear programming problem which we
will solve by direct methods, motivated by the nature of the 
problem. We often use the following notation. If $f:A\longrightarrow B$ denotes a
function and if $U\subseteq A$, then $f(U)$ is defined as $f(U)= \sum_{a\in U}
f(a)$.  




\begin{definition}[Network, $s-t$-flow]
  A network with capacities  consists of a directed simple graph
  $D=(V,A)$ and a   \emph{capacity function} $u:A\to\setR_{\geq0}$. 
  We also require that if there is an arc $uv \in A$, then there is 
  no reverse arc $vu$, and we disallow self-loops. 
  A function $f:A\to\setR_{\geq0}$ is   called an \emph{$s-t$-flow}, if 
  \begin{equation}
    \sum_{e \in \delta^{out}(v)} f(e) = \sum_{e \in \delta^{in}(v)} f(e), \, \mbox{ for all } v \in V - \{s,t\},
  \end{equation}
  where $s,t\in V$. These two vertices are called \emph{source} and 
  \emph{sink} respectively. 
  The flow is \emph{feasible}, if $f(e) \leq u(e)$ for all
  $e \in A$. 
  The \emph{value} of $f$ is defined as 
  $value(f) =  \sum_{e \in \delta^{out}(s)} f(e) - \sum_{e \in \delta^{in}(s)} f(e)$.
  The \emph{maximum $s-t$-flow problem} is the problem of determining
  a maximum feasible $s-t$-flow. 
\end{definition}


Here, for $U \subseteq V$, $\delta^{in}(U)$ denotes the arcs which are entering $U$
and $\delta^{out}(U)$ denotes the arcs which are leaving $U$.   Arc sets
of the form $\delta^{out}(U)$   are called a \emph{cut} of $D$. The
\emph{capacity of a cut } $u(\delta^{out}(U))$ is the sum of the capacities of
its arcs. 


Thus the maximum  flow problem is a linear program of the form 

\begin{eqnarray}
  \max \sum_{e \in \delta^{out}(s)} x(e) & - &   \sum_{e \in \delta^{in}(s)} x(e) \\
   \sum_{e \in \delta^{out}(v)} x(e) & = &  \sum_{e \in \delta^{in}(v)} x(e), \, \mbox{ for all
   } v \in V-\{s,t\} \\
   x(e) & \leq & u (e) , \, \mbox{ for all
   } e \in A \\
   x(e) & \geq & 0 , \, \mbox{ for all
   } e \in A
\end{eqnarray}


\begin{definition}[excess function]
  For any $f:A \to \setR$, the excess function is the function 
  $excess_f:  2^V \to \setR$  defined by $excess_f(U) = \sum_{e \in \delta^{in}(U)}
  f(e) - \sum_{e \in \delta^{out}(U)} f(e)$. 
\end{definition}


\begin{theorem}
\label{f:thr:6}
Let $D = (V,A)$ be a digraph, let $f:A\to\setR$ and let $U\subseteq V$, then
\begin{equation}
  \label{f:eq:28}
  excess_f(U) = \sum_{v \in U} excess_f(v).
\end{equation}  
\end{theorem}


\begin{proof}
  An arc which has both endpoints in $U$ is counted twice with
  different parities on the right, and thus cancels out. An arc which
  has its tail in $U$ is subtracted once on the right and once on the
  left.  An arc which has its head in $U$ is added once on the right
  and once on the left. \qed
\end{proof}

A cut $\delta^{out}(U)$ with $s \in U$  and $t \notin U$ is called an $s-t$-cut.  

\begin{theorem}[Weak duality]
  \label{f:thr:7}
  Let $f$ be a feasible $s-t$-flow and let $\delta^{out}(U)$ be an
  $s-t$-cut, then $value(f) \leq u(\delta^{out}(U))$. 
\end{theorem}

\begin{proof}
  $value(f) = -excess_f(s) = -excess_f(U) = f(\delta^{out}(U)) -
  f(\delta^{in}(U)) \leq f(\delta^{out}(U)) \leq u(\delta^{out}(U))$. \qed
\end{proof}

For an arc $a = (u,v) \in A$ the arc $a^{-1}$ denotes the arc $(v,u)$. 

\begin{definition}[Residual graph]
  Let $f:A \to \setR$,  and $u: A \to \setR$ where $0 \leq f \leq u$. Consider
  the sets of arcs
  \begin{equation}
    \label{f:eq:30}
    A_f = \{ a \mid a \in A, \, f(a) < u(a)\} \cup \{ a^{-1} \mid a \in A, \, f(a) >    0\}. 
  \end{equation}
  The digraph $D(f) = (V,A_f)$ is called the \emph{residual graph} of
  $f$ (for capacities $u$).   
\end{definition}



\begin{corollary}
  \label{co:6}
  Let $f$ be a feasible $s-t$-flow and suppose that $D(f)$ has no path
  from $s$ to $t$, then $f$ has maximum value.
\end{corollary}


\begin{proof}
  Let $U$ be the set of nodes which are reachable in $D(f)$ from
  $s$. Clearly $\delta^{out}(U)$ is an $s-t$-cut. Now $value(f) = f(\delta^{out}(U)) -
  f(\delta^{in}(U)$.  Each arc leaving $U$ is not an arc of $D(f)$ and thus
  $f(\delta^{out}(U)) = u (\delta^{out}(U))$. Each arc entering $U$ does not
  carry any flow   and thus $f(\delta^{in}(U)) = 0$. It follows that  
  $value(f) = u (\delta^{out}(U))$ and $f$ is optimal by
  Theorem~\ref{f:thr:7}.   \qed
\end{proof}

\begin{definition}[undirected walk]
\label{f:def:7}
  An \emph{undirected walk} is a  sequence of the form
  $P=(v_0,a_1,v_1,\ldots,v_{m-1},a_m,v_m)$, where  $a_i \in A$ for
  $i=1,\ldots,m$ and $a_i = (v_{i-1},v_i)$ or $a_i = (v_{i},v_{i-1})$. If
  the nodes $v_0,\ldots,v_m$ are all 
  different, then $P$ is an \emph{undirected path}. 
% The \emph{length } of $P$ is
%   $m$. The \emph{distance} of two nodes  $u$ and $v$ is the length of
%   a shortest path from $u$ to $v$. 
\end{definition}


Any directed path $P$ in $D(f)$ yields an undirected path  in
$D$. Define for such a path $P$ the vector $\chi^P \in \{0,\pm1\}^A$ as
\begin{equation}
  \label{f:eq:29}
  \chi^P(a) = 
  \begin{cases}
    1 & \mbox{ if $P$ traverses $a$},\\
    -1&  \mbox{ if $P$ traverses $a^{-1}$},\\
    0 & \mbox{ if $P$ traverses neither $a$ or $a^{-1}$}. 
  \end{cases}
\end{equation}


\begin{theorem}[max-flow min-cut theorem, strong duality]
  \label{f:thr:8}
  The maximum value of a feasible $s-t$-flow is equal to the minimum
  capacity of an $s-t$ cut.
\end{theorem}

\begin{proof}
  Let $f$ be a maximum $s-t$-flow. Consider the residual graph
  $D(f)$. If this residual graph contains an 
  $s-t$-path $P$, then we can route flow along this path. More
  precisely, there exists an $\epsilon>0$ such that $f + \epsilon \, \chi^P$ is feasible.
  We have $value(f + \epsilon \, \chi^P) = value(f) + \epsilon $. This contradicts the
  maximality of $f$ thus there exists no $s-t$-path in $D(f)$. 
  
  Let $U$ be the nodes reachable from $s$ in $D(f)$. Then 
  $value(f)=u(\delta^{out}(U))$ and $\delta^{out}(U)$ is an
  $s-t$-cut of minimum capacity by the weak duality theorem.
\end{proof}
 


This suggests the algorithm of Ford and Fulkerson to find a maximum
flow. Start with $f = 0$. Next iteratively apply the following
\emph{flow augmentation algorithm}. 


Let $P$ be a directed $s-t$-path in $D(f)$. Set $f \gets f + \epsilon\chi^P$, where
$\epsilon$ is as large as possible to maintain $0\leq f\leq u$. 

\begin{exercise}
   Define a \emph{residual capacity} for $D(f)$. Then determine 
   the maximum $\epsilon$ such that $0\leq f\leq u$.
\end{exercise}

\begin{theorem}
  \label{f:thr:9}
  If all capacities are rational, this algorithm terminates. 
\end{theorem}

\begin{proof}
   By multiplying all the capacities by $10^x$ with $x$ 
   sufficiently large, we can suppose the capacities to be 
   integral. Since the value of the flow is augmented at 
   every step at least by $1$ and the capacities are finite, our 
   algorithm finds a maximum flow in finitely many iterations.
\end{proof}


\begin{center}
\begin{tikzpicture}[style=thick]
\tikzstyle{vertex}=[draw,shape=circle]
\tikzstyle{edge} = [draw,thick,-]
\tikzstyle{weight} = [font=\small];

\draw (0,0) node[vertex] (s) {$s$};
\draw (3,2) node[vertex] (u) {$u$};
\draw (3,-2) node[vertex] (v) {$v$};
\draw (6,0) node[vertex] (t) {$t$};

\path[edge,->] (s) -- node[weight,xshift=-3ex] {$M$} (u);
\path[edge,->] (s) -- node[weight,xshift=-3ex] {$M$} (v);
\path[edge,->] (u) -- node[weight,xshift=1ex] {$1$} (v);
\path[edge,->] (u) -- node[weight,xshift=3ex] {$M$} (t);
\path[edge,->] (v) -- node[weight,xshift=3ex] {$M$} (t);
\end{tikzpicture}
\end{center}
The example above shows that, if the augmenting paths are chosen in a
disadvantageous way, then the Ford-Fulkerson algorithm may take $\Omega(M)$
iterations, where $M$ is the largest capacity in the network. This
happens if all augmenting paths use the arc $uv$ or $vu$ respectively
in the residual network. 




\begin{corollary}[integrity theorem]
  If $u(a) \in \setN$ for each $a \in A$, then there exists an integer maximum
  flow ($f(a) \in \setN$ for all $a \in A$). 
\end{corollary}


\begin{proof}
This follows from the fact that the residual capacities remain
integral and thus the augmented flow is always integral. \qed 
\end{proof}



\begin{theorem}
\label{f:thr:10}
If we choose in each iteration a shortest $s-t$-path in $D(f)$ as a
flow-augmenting path, the number of iterations is at most $|V| \cdot|A|$. 
  
\end{theorem}

\begin{definition}
  Let $D = (V,A)$ be a digraph, $s,t \in V$  and let  $\mu(D)$ denote the
  length of a  shortest path from $s$ to $t$. Let $\alpha(D)$ denote the set
  of arcs contained in at least one shortest $s-t$ path.
\end{definition}

\begin{theorem}
  \label{f:thr:11}
  Let  $D = (V,A)$ be a digraph and  $s,t \in V$. Define $D' = (V,A \cup
  \alpha(D)^{-1})$. Then $\mu(D) = \mu(D')$ and $\alpha(D) = \alpha(D')$.
\end{theorem}


\begin{proof}
  It suffices to show that $\mu(D)$ and $\alpha(D)$ are invariant if we add
  $a^{-1}$ to $D$ for one arc $a \in \alpha(D)$. Suppose not, then there is a
  directed $s-t$-path $P_1$ traversing $a^{-1}$ of length at most
  $\mu(D)$. As $a \in \alpha(D)$ there is a path $P_2$ traversing $a$ of length
  $\mu(D)$. If we follow $P_2$  until the tail of $a$ is reached  and
  from thereon follow $P_1$, we obtain another $s-t$ path $P_3$ in $D$.
  Similarly if we follow $P_1$ until the head of $a$ is reached and then follow $P_2$,
  we obtain a fourth $s-t$ path $P_4$ in $D$. However $P_3$ or $P_4$ has length less
  than $\mu(D)$. This is a contradiction. \qed
\end{proof}


\begin{proof}[of Theorem~\ref{f:thr:10}]
  Let us augment flow $f$ along a shortest $s-t$-path $P$ in $D(f)$
  obtaining flow $f'$. The residual graph $D_{f'}$ is a subgraph of
  $D' = (V, A_f \cup\alpha(D(f))^{-1})$. Hence $\mu(D_{f'}) \geq \mu(D')=\mu(D(f))$. If
  $\mu(D_{f'})=\mu(D(f))$, then $\alpha(D_{f'})\subseteq\alpha(D') = \alpha(D(f))$. At least one
  arc of $P$ does not belong to $D_{f'}$, (the arc of minimum residual
  capacity!) thus the inclusion is
  strict. Since $\mu(D(f))$ increases at most $|V|$~times and, as long as
  $\mu(D(f))$ does not change, $|\alpha(D(f))|$ decreases at most $2\, |A|$~times, we
  have the theorem. \qed
\end{proof}

In the following let $m = |A|$ and $n = |V|$. 

\begin{corollary}
  \label{co:2}
  A maximum flow can be found in time $\bigO(n \, m^2)$.
\end{corollary}




\section{Minimum cost network flows, MCNFP}
\label{sec:minimum-cots-flows}

In contrast to the maximum $s-t$-flow problem, the goal here is to
route a flow, which comes from several sources and sinks through a
network with capacities and \emph{costs} in such a way, that the total
cost is minimized. 

\begin{example}
  Suppose you are given a directed graph $D=(V,A)$ with arc weights  
  $c: A \to \setR_{\geq0}$ and your task is to compute a shortest 
  path from a particular node $s$ to all other nodes in the graph 
  and assume that such paths exist. Then one can model this as a 
  MCNFP (minimum cost network flow problem) by sending a flow of 
  value $|V|-1$ into the source node 
  and by letting a flow of value $1$ leave each node. The costs on 
  the arcs are defined by $c$. The arcs have infinite capacities. 
  We will see later, that this MCNFP 
  has an integral solution which corresponds to the shortest
  paths from $s$ to all other nodes. 
\end{example}


\begin{figure}
  \centering
  \includegraphics{figures/flows1.pdf}   \includegraphics{figures/flows2.pdf} 
  \caption{A Network with in/out-flow, costs and capacities and a
    feasible flow of cost $13$. }
  \label{ex:net:1}
\end{figure}

 


Here is a formal definition of a MCNFP. In
this notation, vertices are indexed with the letters $i,j,k$ and arcs 
are denoted by their tail and head respectively, for example
$(i,j)$ denotes the arc from $i$ to $j$. 
  
% For a node $i\in V$, the subset $I(i)\subseteq V$ denotes the incoming nodes of
% $i$, i.e., the set $I(i) = \{ j \mid (j,i) \in A\}\subseteq V$. Similarly $O(i) = \{ j
% \mid (i,j) \in A\}\subseteq V$ denotes the set of outgoing nodes.

 A \emph{network} is now a
directed graph $D = (V,A)$ together with a capacity function $u: A \to
\setQ_{\geq0}$, a cost function $c: A \to\setQ$ and an external flow $b: V \to \setQ$.
The value of $b_i$ denotes the amount of flow which comes from the
exterior. If $b_i>0$, then there is flow from the outside, entering
the network through node $i$. If $b_i<0$, there is flow which leaves
the network through $i$.

In the following we often use the notation $f(i,j)$ for the flow-value
on the arc $(i,j)$ (instead of $f((i,j))$). Similarly we write
$c(i,j)$ and $u(i,j)$. 

A \emph{feasible flow} is a function $f:A \to \setQ_{\geq0}$ which
satisfies the following constraints. 
\begin{displaymath}  
  \begin{array}{cl}
    \sum_{e \in \delta^{out}(i)} f{(e)} - 
    \sum_{e \in \delta^{in}(i)} f(e) =b_i & 
    \text{ for all } i \in V, \\
    0 \leq f{(e)} \leq u{(e)}              & \text{ for all } e \in A.
  \end{array}
\end{displaymath}


The goal is to find a feasible flow with minimum cost: 


  \begin{displaymath}
    \begin{array}{rcl}
      \text{minimize}     & \sum_{e \in A} c(e) f(e) &     \\
      \text{ subject to}  & \sum_{e \in \delta^{out}(i)} f(e) - \sum_{e\in \delta^{in}(i)} f(e) =b_i & \text{ for all } i \in V,\\
                          & 0 \leq f(e) \leq u(e)              &
                          \text{ for all } e \in A 
    \end{array}
  \end{displaymath}
  

  \begin{example}
    Imagine you are a pilot and fly a passenger airplane in
    hops from airport $1$ to airport $2$ to airport $3$ and so on, 
    until airport $n$.  At airport $i$ there are $b_{ij}$ passengers 
    that want to travel to airport $j$, where $j>i$. You may decide 
    how many of the $b_{ij}$ passengers you will take on board. Each 
    of the passengers will pay $c_{ij}$ dollars for the trip. The 
    airplane can accommodate $p$ people. 

    You are a greedy pilot and think of a plan to pick up and deliver
    passengers on your hop from $1$ to $n$ which maximizes your
    revenue. 

    Finding this plan can be modeled as a MCNFP. 
    Your network has nodes $1,\ldots,n$ and arcs
    $(i,i+1), i=1,\ldots,n-1$ with capacities $p$ and  without
    costs. These nodes do not have  in/out-flow from   the outside. 
    You furthermore have nodes $i\to j$ for $i<j$ and 
    $i,j \in \{1,\ldots,n\}$ which are excess nodes with in-flow 
    $b_{ij}$ from the outside. Each node $i\to j$ is connected to
    $i$ and to $j$ with a directed arc. The capacities on these arcs
    are infinite. The cost of the arc $(i\to j,i)$ is $-c_{ij}$. 
    The cost of the arc $(i\to j,j)$ is zero. The outflow on the 
    node $j$ is the total number of passengers that want to fly to 
    node $j$. An integral optimal flow to this problem is an optimal 
    plan for you. 
  \end{example}


  
  Throughout this chapter we make the following assumptions.
  
  \begin{enumerate}
  \item All data (cost, supply, demand and capacity) are integral.
  \item The network contains an incapacitated directed path between
    every pair of nodes. 
  \item The supplies/demands at the nodes satisfy the condition 
    $\sum_{i \in V} b_i=0$ and the MCNFP has a feasible solution. 
  \item All arc costs are nonnegative. 
  \item The graph does not contain a pair of reverse arcs. 
  \end{enumerate}

\begin{exercise} 
   Show how to transform a MCNFP on a digraph with pairs of 
   reverse arcs into a MCNFP on a digraph with no pairs of 
   reverse arcs.  The number of arcs and nodes should 
   asymptotically remain the same.
\end{exercise}
  
  An \emph{arc-flow} of $D$ is a flow vector, that satisfies the
  nonnegativity and capacity constraints. 
 
  \begin{eqnarray*}
        \sum_{e\in \delta^{in}(i)} f(e) -  \sum_{e \in \delta^{out}(i)} f(e)  =  g(i) & &   \text{ for all } i \in V,\\
        0 \leq f(e) \leq u(e)                 & &   \text{ for all }
        e \in A. 
  \end{eqnarray*}


  \begin{itemize}
  \item  If $g(i) >0$, then $i$ is an \emph{excess node} (more inflow than
    outflow).
  \item  If $g(i) <0$, then $i$ is a \emph{deficit node} (more outflow than
    inflow).
  \item  If $g(i) = 0$ then $i$ is a \emph{balanced node}.
  \end{itemize}
  
  
\begin{exercise}
   Prove that $\sum_{i \in V} g(i) = 0$ holds and thus that a 
   feasible flow only exists if the sum of the $b(i)$ is equal 
   to zero. 
\end{exercise}

  
  Let $\paths$ be the collection of directed paths  of $D$ and
  let $\cycles$ be the collection of directed cycles of $D$. A
  path-flow is a function $\beta: \paths \cup \cycles \to \setR_{\geq0}$ which assigns
  flow values to paths and cycles. 

  For $(i,j)\in A$ and $P \in \paths$ let $\delta_{(i,j)}(P)$ be $1$ if $(i,j) \in
  P$ and $0$ otherwise. For $C \in \cycles$ let $\delta_{(i,j)}(C)$ be $1$ if
  $(i,j)\in C$ and $0$ otherwise. 

  A path-flow $\beta$ determines a unique   arc-flow
  \begin{displaymath}
    f(i,j) = \sum_{P \in \paths} \delta_{(i,j)}(P) \beta(P) + \sum_{C \in \cycles} \delta_{(i,j)}(C) \beta(C). 
  \end{displaymath}
    

  
  \begin{theorem}
    \label{f:thr:Decomp}
    Every path and cycle flow  has a unique representation
    as a  nonnegative arc-flow. Conversely, every nonnegative arc-flow
    $f$ can be represented as a path and cycle flow with the following
    properties:
    \begin{enumerate}
    \item Every directed path with positive flow connects a deficit
      node with an excess node.
    \item At most $n+m$ paths and cycles have nonzero flow and at most
      $m$ cycles have nonzero flow.
    \end{enumerate}
    If the arc-flow $f$ is integral, then so are the path and cycle
    flows into which it decomposes. 
  \end{theorem}

  \begin{proof}
   
  
  ``$\Rightarrow$'' See discussion above.
  
  ``$\Leftarrow$'' 
  
  Let $f$ be an arc-flow. Suppose $i_0$ is a deficit node. Then there
  exists an incident arc $(i_0,i_1)$ which carries a positive flow. If
  $i_1$ is an excess node, we have found a path from deficit to excess
  node. Otherwise, the flow balance constraint at $i_1$ implies that
  there exists an arc $(i_1,i_2)$ with positive flow. Repeating this
  procedure, we finally must arrive at an excess node or revisit a
  node. This means that we either have constructed a directed path $P$
  from deficit node to excess node or a directed cycle $C$, both
  involving only arcs with strictly positive flow.
  
  In the first case, let $P = i_0,\ldots,i_k$ be the directed path from
  deficit node $i_0$ to excess node $i_k$. We set $\beta(P) =
  \min\{-e_{i_0}, e_{i_k}, \min\{f{(i,j)} \mid (i,j) \in P\} \}$ and $f{(i,j)} =
  f{(i,j)} - \beta(P), \, (i,j) \in P$.  In the second case, set $\beta(C) =
  \min\{f{(i,j)} \mid (i,j) \in C$ and $f{(i,j)} = f{(i,j)} - \beta(C), \, (i,j)
  \in C$.  Repeat this procedure until all node imbalances are zero.
  
  Now find an arc with positive flow and construct a cycle $C$ by
  following only positive arcs from there. Set 
  $\beta(C) = \min\{f{(i,j)} \mid  (i,j) \in C\}$ and 
  $f{(i,j)} = f{(i,j)} - \beta(C),\, (i,j) \in C\}$. Repeat this process until
  there are no positive flow-arcs left. 

  Each time a path or a cycle is identified, the excess/deficit of
  some node is set to zero or some arc is set to zero. This implies
  that we decompose into at most $n+m$ paths and cycles. Since cycle
  detection sets an arc to zero we have at most $m$ cycles.  \qed
\end{proof}
  
  An arc flow $f$ with $g(i)=0$ for each $i \in V$ is called a
  \emph{circulation}. 
  
  \begin{corollary}
    A circulation can be decomposed into at most $m$ cycle-flows.
  \end{corollary}

  
  Let $D = (V,A)$ be a network with capacities 
  $u{(i,j)}, \,  (i,j) \in A$ and costs $c{(i,j)}, \, (i,j) \in A$ and let
  $f$ be a feasible flow of the network. The \emph{residual network} $D(f)$ is
  defined as follows.

  \begin{itemize}
  \item We replace each arc $(i,j) \in A$ with two arcs $(i,j)$ and
    $(j,i)$.
  \item The arc $(i,j)$ has cost $c{(i,j)}$ and \emph{residual capacity}
    $r{(i,j)} = u{(i,j)} - f{(i,j)}$.  
  \item The arc $(j,i)$    has cost $-c{(i,j)}$ and
    residual capacity $r{(j,i)}=f{(i,j)}$. 
  \item      Delete all arcs which do not have strictly positive residual
    capacity. 
  \end{itemize}

  

  \begin{figure}
  \centering
\includegraphics{figures/flows6.pdf}
\caption{The residual network of the flow in Figure~\ref{ex:net:1} and
a negative cycle marked by the red edges.} \label{ex:net:2}
\end{figure}



  A directed cycle (or path) in $D(f)$ is called an 
  \emph{augmenting cycle (or path)} of $f$.  

 
  \begin{lemma}
    \label{lem:6}
    Suppose that $f$ and $f^\circ$ are feasible flows, then 
    $f  - f^\circ$  is a circulation in $D(f^\circ)$.  Here $f  -
    f^\circ$ is the flow  
    \begin{displaymath}
      (f-f^\circ)(e) = 
      \begin{cases}
        \max\{0, f(e) - f^\circ(e)\}, & \text{ if } e \in  A(D)\\
        \max\{0, f^\circ(e) - f(e)\}, & \text{ if } e^{-1} \in  A(D)\\
        0, & \text{ otherwise.}
      \end{cases}
    \end{displaymath}
  \end{lemma}
   
  


  \begin{proof}
    It is very easy to see that the flow $f - f^\circ$ satisfies the
    capacity constraints. One also has for each $v \in V$
    \begin{displaymath}
      \sum_{e \in \delta^{out}(v)} (f (e) - f^\circ(e))  - \sum_{e \in \delta^{in}(v)}
      (f (e) - f^\circ(e)) = 0.
    \end{displaymath}
    If a term $(f (e) - f^\circ(e))$ is negative, it is replaced by its
    absolute value and charged as flow on the arc $e^{-1}$ in
    $D(f^\circ)$ which leaves its contribution to the sum above
    invariant. 
    \qed     
  \end{proof}




  \begin{figure}
  \centering
% \psset{unit=1.5cm}  
 \includegraphics{figures/flows4.pdf}
\caption{Two  arcs $e_1,e_2 \in A$ labeled with $f(e)/f^\circ(e)/u(e)$ and the
  corresponding flow on these arcs (or their reverse) in
  $D(f^\circ)$. Arcs in $D(f^\circ)$ are labeled with flow and capacity
  values respectively. }\label{fig:diff}
\end{figure}


  \begin{theorem}[Augmenting Cycle Theorem]
    \label{f:thr:augcyc}
    Let $f$ and $f^\circ$ be any two feasible flows of a network flow
    problem. Then $f$ equals $f^\circ$ plus the flow of at most $m$
    directed cycles in $D(f^\circ)$.  Furthermore the cost of $f$ equals
    the cost of $f^\circ$ plus the cost of flow on these augmenting
    cycles. 
  \end{theorem}
  
  \begin{proof}
    This can be seen by applying flow decomposition on the  flow 
    $f - f^\circ$  in $D(f^\circ)$. \qed
  \end{proof}
  

 
  
  \begin{theorem}[Negative Cycle Optimality Conditions]
    \label{f:thr:13}
    A feasible flow $f^*$ is an optimal solution of the MCNFP, 
    if and only if it satisfies the negative
    cycle optimality conditions: the residual network $D(f^*)$
    contains no directed cycle of negative cost.    
  \end{theorem}

  \begin{proof}
    

  ``$\Rightarrow$'' Suppose that $f$ is a feasible flow and that $D(f)$ contains
  a negative directed cycle. Then $f$ cannot be optimal, since we can
  augment positive flow along the corresponding cycle in the
  network. Therefore, if $f^*$ is an optimal flow, then $D(f^*)$
  cannot contain a negative directed cycle. 

  ``$\Leftarrow$'' Suppose now that $f^*$ is a feasible flow and suppose that
  $D(f^*)$ does not contain a negative cycle. Let $f^\circ$ be an optimal
  flow with $f^\circ \neq f^*$. The vector $f^\circ-f^*$ is a circulation in
  $D(f^\circ)$ with
  non-positive cost $c^T(f^\circ-f^*) \leq0$. It follows from
  Theorem~\ref{f:thr:augcyc} that the cost of $f^\circ$ equals the cost of
  $f^*$ plus  the cost of directed cycles in the residual network
  $D(f^*)$.  The cost of these cycles is nonnegative, and therefore
  $c(f^\circ) \geq c(f^*)$ which implies that $f^*$ is optimal. 
  \qed
\end{proof}




\begin{figure}
  \centering
  
  \includegraphics{figures/flows5.pdf} 
\caption{The result of augmenting a flow of one along the negative
  cycle in Figure~\ref{ex:net:2}. This flow has cost 12 but is not
  optimal, since the residual network still contains a negative
  cycle.}\label{ex:net:3}
\end{figure}





\begin{algorithm}[Cycle Canceling Algorithm]
  \label{alg:c1}
  ~\\
  \begin{enumerate}
  \item  establish a feasible flow $f$ in the network
  \item {\tt WHILE} $D(f)$ contains a negative cycle
    \begin{enumerate}
    \item  detect a negative cycle $C$ in $D(f)$
    \item let $\delta=\min\{r{(i,j)} \mid (i,j) \in C\}$
    \item augment $\delta$ units of flow along the cycle $C$
    \item update $D(f)$
    \end{enumerate}
  \item {\tt RETURN}  $f$
  \end{enumerate}
\end{algorithm}




  \begin{theorem}
    \label{f:thr:12}
    The cycle canceling algorithm terminates after a finite number of
    steps if the MCNFP has an optimal solution. 
  \end{theorem}
  \begin{proof}   
  The cycle canceling algorithm reduces the cost in each iteration.
  We have assumed that the input data is integral. Thus the cost
  decreases by at least one unit each iteration. 
  Therefore the number of iterations is finite.    
  \qed
\end{proof}


\begin{corollary}
  \label{co:3}
  If the capacities are  integral and if the MCNFP has a optimal flow, then
  it has an optimal flow with integer values only. 
\end{corollary}



%Consider the MCNFP

%\begin{displaymath}
%    \begin{array}{rcl}
%      \text{minimize}     & \sum_{(i,j) \in A} c{(i,j)} f{(i,j)} &     \\
%      \text{ subject to}  & \sum_{j \in O(i)} f{(i,j)} - \sum_{j\in I(i)} f{(j,i)} =b_i & \text{ for all } i \in V,\\
%                          & 0 \leq f{(i,j)} \leq u{(i,j)}              & \text{ for all } (i,j) \in A.
%    \end{array}
%  \end{displaymath}
    

%We transform this problem into standardform via slackvariables 
%$z{(i,j)}\geq0,\, (i,j) \in A$:

%\begin{displaymath}
%    \begin{array}{rrcll}
%      \text{minimize}     & \sum_{(i,j) \in A} c{(i,j)} f{(i,j)} &     \\
%      \text{ subject to}  & \sum_{j \in O(i)} f{(i,j)} - \sum_{j\in I(i)} f{(j,i)}
%      & = & b_i & \text{ for all } i \in V,\\
%                          & f{(i,j)} + z{(i,j)} & = &  u{(i,j)}
%                          & \text{ for all } (i,j) \in A, \\
%                    & f{(i,j)},z{(i,j)} & \geq & 0. &       \\
%    \end{array}
%  \end{displaymath}



  
%  The dual of this problem has variables $\pi_i, \, i \in V$ and 
%  $\alpha{(i,j)},  \,(i,j) \in A$ and is defined as follows:
  
  
%  \begin{displaymath}
%    \begin{array}{rrcll}
%      \text{maximize}     & \sum_{i \in V} b_i \pi_i  - \sum_{(i,j) \in A} u{(i,j)} \alpha{(i,j)}      \\
%      \text{ subject to}  & \pi_i - \pi_j - \alpha{(i,j)} 
%      & \leq & c{(i,j)} & \text{ for all } (i,j) \in A,\\
%                          & \alpha{(i,j)} & \geq & 0
%                          & \text{ for all } (i,j) \in A.
%  \end{array}
%  \end{displaymath}


Let $\pi : V  \to \setR$ be a function (\emph{node potential}). The
\emph{reduced cost} of an arc $(i,j)$ w.r.t. $\pi$ is 
$c_\pi((i,j))=c((i,j))+\pi(i) - \pi(j)$. The potential $\pi$ is called
\emph{feasible} if $c_\pi((i,j))\geq0$ for all arcs $(i,j)\in A$. 

\begin{lemma}
  \label{lem:7}
  Let $D = (V,A)$ be a digraph with arc weights $c:A\to\setR$. Then $D$
  does not have a negative cycle if and only if there exists a
  feasible node potential $\pi$ of $D$. 
\end{lemma}


\begin{proof}
  Consider a directed path $P = i_0,i_1,\ldots,i_k$. The cost of this
  path is 
  \begin{displaymath}
    c(P) = \sum_{j=1}^{k}c((i_{j-1},i_j)).
  \end{displaymath}
  The reduced cost of this path is equal to 
  \begin{displaymath}
    c_\pi(P) = \sum_{j=1}^{k}c((i_{j-1},i_j)) + \pi(i_0) - \pi(i_k).
  \end{displaymath}
  If $P$ is a cycle, then $i_0$ and $i_k$ are equal, which means that
  its cost and reduced cost coincide. Thus, if there exists a feasible
  node potential, then there does not exist a negative cycle. 

  
  On the other hand, suppose that $D$ does not contain a negative
  cycle w.r.t. $c$. Add a vertex $s$ to $D$ and the arcs $(s,i)$ for all 
  $i \in  V$. The weights (costs) of all these new arcs is $0$. Notice
  that in this way, no new cycles are created, thus still there does
  not exist a negative cycle. This means we can compute the shortest
  paths from $s$ to all other nodes $i \in V$. Let $\pi$ be the function
  which assigns these shortest paths lengths. Clearly $c_\pi((i,j)) =
  \pi(i) - \pi(j) + c((i,j))\geq0$, since the shortest-path length to $j$ is
  at most the shortest-path length to $i$ plus $c((i,j))$.  \qed
\end{proof}

  
This means that we have again a nice way to prove that a flow is
optimal. Simply equip  the residual network with a feasible node
potential.  



\begin{corollary}[Reduced Cost Optimality Condition]
  \label{co:1}
  A feasible flow $f^*$ is optimal if and only if there exists a node
  potential $\pi$ such that the reduced costs $c_\pi{(i,j)}$ of  each arch
  $(i,j)$ of $D(f)$ are nonnegative. 
\end{corollary}






  The cycle canceling algorithm is only pseudopolynomial. If we could
  always chose a minimum cycle (cycle with best improvement) as an
  augmenting cycle, we would have a polynomial number of
  iterations. Finding minimum cycles is $NP$-hard. Instead we augment
  along \emph{minimum mean cycles}. One can find minimum mean cycles
  in polynomial time.  

  The \emph{mean cost} of a cycle $C \in \cycles$ is the cost of $C$
  divided by the number of arcs in $C$:
  \begin{displaymath}
    \sum_{(i,j) \in C} c{(i,j)})  / |C|.
  \end{displaymath}
  

  \begin{algorithm}[Minimum Mean Cycle Canceling, MMCC]
    \label{alg:2}
    ~\\
    \begin{enumerate}
    \item establish a feasible flow $f$ in the network
    \item {\tt WHILE} $D(f)$ contains a negative cycle
      \begin{enumerate}
      \item  detect a minimum mean cycle $C$ in $D(f)$ 
      \item   $\delta=\min\{r{(i,j)} \mid (i,j) \in C\}$
      \item   augment $\delta$ units of flow along the cycle $C$
      \item   update $D(f)$ 
      \end{enumerate}
    \item {\tt RETURN} $f$
    \end{enumerate}
  \end{algorithm}



  We now analyze the MMCC-algorithm. Let $\mu(f)$ denote the minimum
  mean-weight of a cycle in $D(f)$. 
  
  \begin{lemma}[See Korte \& Vygen \cite{MR1897297}]
    \label{lem:8}
    Let $f_1,f_2,\ldots$ be a sequence of feasible flows such that
    $f_{i+1}$ results from $f_i$ by augmenting flow along $C_i$, where
    $C_i$ is a minimum mean cycle of $D(f_i)$, then
    \begin{enumerate}
    \item \label{item:5} $\mu(f_k)\leq\mu(f_{k+1})$ for all $k$.
    \item \label{item:6} $\mu(f_k) \leq \frac{n}{n-1} \mu(f_l)$, where $k<l$
      and $C_k \cup C_l$ contains a pair of reversed arcs.
    \end{enumerate}
  \end{lemma}



  \begin{proof}
    
    \ref{item:5}): Suppose $f_k$ and $f_{k+1}$ are two subsequent flows in this
    sequence. Consider the multi-graph $H$ which results from $C_k$
    and $C_{k+1}$ by deleting pairs of opposing arcs.  The arcs of
    $H$ are a subset of the arcs of $D(f_k)$, since an arc of
    $C_{k+1}$ which is not in $D(f_k)$ must be a reverse arc of $C_k$.

    Each node in $H$ has even degree.  Thus $H$ can be
    decomposed into cycles, each of mean weight at least $ \mu(f_k)$.  
    Thus we have $c(A(H)) \geq \mu(f_k) |A(H)|$. 

    Since the total weight of each reverse pair of arcs is zero we
    have 
    \begin{displaymath}
      c(A(H)) = c(C_k) + c(C_{k+1}) = \mu(f_k) |C_k|   + \mu(f_{k+1}) |C_{k+1}|.
    \end{displaymath}

    Since $|A(H)| \leq |C_k| + |C_{k+1}|$ we conclude 
    \begin{eqnarray*}
       \mu(f_k) ( |C_k| + |C_{k+1}|) & \leq &  \mu(f_k)|A(H)| \\
                              & \leq & c(A(H)) \\
                              & = & \mu(f_k) |C_k|   + \mu(f_{k+1}) |C_{k+1}|.
    \end{eqnarray*}

    Thus $\mu(f_k) \leq \mu(f_{k+1})$. 

    
    \ref{item:6}): By the first part of the theorem, it is enough to
    prove the statement for $k,l$ such that $C_i \cup C_l$ does not
    contain a pair of reverse arcs for each  $ i, \, k < i < l$.  

    Again, consider the graph $H$ resulting from  $C_k$
    and $C_{l}$ by deleting pairs of opposing arcs. $H$ is a subgraph
    of $D(f_k)$, since any arc of $C_l$ which does not belong to
    $D(f_k)$ must be a reverse arc of $C_k,C_{k+1},\ldots,C_{l-1}$. But
    only $C_k$ contains a reverse arc of $C_l$. So as above we have 
    \begin{displaymath}
      c(A(H)) = c(C_k) + c(C_{l}) = \mu(f_k) |C_k|   + \mu(f_{l})      |C_{l}|. 
    \end{displaymath}
    
    Since $|A(H)| \leq |C_k| + |C_{l}| -2$ we have $|A(H)| \leq
    \frac{n-1}{n}( |C_k| + |C_{l}|)$. Thus we get
    \begin{eqnarray*}
       \mu(f_k)\frac{n-1}{n} ( |C_k| + |C_{l}|) & \leq &  \mu(f_k)|A(H)| \\
                              & \leq & c(A(H)) \\
                              & = & \mu(f_k) |C_k|   + \mu(f_{l})|C_{l}|\\          
                              & \leq & \mu(f_{l}) (|C_k|   + |C_{l}|) \\                       
    \end{eqnarray*}
    This implies that $\mu(f_k) \leq \frac{n}{n-1}  \mu(f_l)$.  
    \qed
  \end{proof}






  \begin{corollary}
    \label{co:4}
    During the execution of the MMCC-algorithm, $|\mu(f)|$ decreases by
    a factor of $1/2$ every $n\cdot m$ iterations. 
  \end{corollary}

  \begin{proof}
    Let $C_1, C_2,\ldots$ be the sequence of augmenting cycles. Every
    $m^{th}$ iteration, there must be an arc of the cycle, which is
    reverse to one of the succeeding $m-1$ cycles, because every
    iteration, one arc of the residual network will be deleted. 
    Thus after $n \, m$ iterations, the absolute value of $\mu$ has
    dropped by $\left(\frac{n-1}{n} \right)^n \leq e^{-1} \leq 1/2$.  \qed
  \end{proof}

  \begin{corollary}
    \label{co:5}
    If all data are integral, then the
    MMCC-algorithm runs in polynomial time. 
  \end{corollary}
  
  \begin{proof}
    \begin{itemize}
    \item A lower bound on $\mu$ is the smallest cost $c_{min}$
    \item $|\mu|$ drops by $1/2$ every $m \, n$ iterations. 
    \item After $m n \log n |c_{min}|$ iterations, absolute value
    of minimum mean weight cycle   drops below $1/n$, thus is zero.
    \item {\bf We need to prove that a minimum mean cycle can be found        in polynomial time} 
    \end{itemize} \qed
  \end{proof}


This is a so-called \emph{weakly polynomial} bound, since the binary
encoding length of the numbers in the input (here the costs)
influences the running time. We now prove that the MMCC-algorithm is
\emph{strongly polynomial}. 






\begin{theorem}[See Korte \& Vygen~\cite{MR1897297}]
  The MMCC-algorithm requires  $\bigO(m^2 \, n \log n)$ iterations
  (mean weight cycle cancellations). 
\end{theorem}

\begin{proof}
  One shows that every $m\, n  (\lceil\log n\rceil+1)$ iterations, at least one
  arc is \emph{fixed}, which means that the flow through  this arc
  does not change anymore. 
  
  Let $f_1$ be some flow at some iteration and let $f_2$ be the flow
  $m\, n  (\lceil\log n\rceil+1)$ iterations later. 
  It follows from Corollary~\ref{co:4} that 
  \begin{equation}
    \label{f:eq:33}
    \mu(f_1)\leq 2 \, n \, \mu(f_2) 
  \end{equation}
  %
  holds. 
  
  Define the costs $c'(e) = c(e) - \mu(f_2)$ for the residual network
  $D(f_2)$. There exists no negative cycle in $D(f_2)$ w.r.t. this
  cost $c'$. (A cycle $C$ has weight $c'(C) = \sum_{e\in C} c(e) - |C|
  \mu(f_2)$ and thus $c'(C) / |C| = \sum_{e\in C} c(e) / |C| - \mu(f_2)\geq0$).
  By Lemma~\ref{lem:7} there exists a feasible node potential $\pi$ for
  these weights. One has $0 \leq c'_\pi(e) = c_\pi(e) - \mu(f_2)$ and thus
  \begin{equation}
    \label{f:eq:32}
     c_\pi(e) \geq \mu(f_2), \text{ for all } e \in A(D(f_2)). 
  \end{equation}

  Let $C$ be a minimum mean cycle of $D(f_1)$. One has
  \begin{equation}
    \label{f:eq:34}
    c_\pi(C) = c(C) = \mu(f_1) \, |C| \leq 2 \, n \, \mu(f_2) |C|.
  \end{equation}
  
  It follows that there exists an arc $e_0$ of $C$ such that
  \begin{equation}
    \label{f:eq:36}
    c_\pi(e_0)\leq 2 \, n \, \mu(f_2)
  \end{equation}
  %
  holds. The inequalities~\eqref{f:eq:32}
  imply that  $e_0 \notin A(D(f_2))$ 

  We now make the following
  claim:
  %
  \begin{quote}
    Let $f'$ be a feasible flow such that $e_0 \in D(f')$, then 
    $\mu(f') \leq    \mu(f_2)$. 
  \end{quote}
  
  If we have shown this claim, then it follows from Lemma~\ref{lem:8}
  that $e_0$ cannot be anymore in the residual network of a flow after
  $f_2$.  Thus the flow along the arc $e_0$ (or $e_0^{-1}$) is
  fixed. 

  Let $f'$ be a flow such that $e_0\in A(D(f'))$. Recall that $f' - f_2$
  is a circulation in $D(f_2)$ where $e_0\notin D(f_2), \,e_0^{-1}\in D(f_2)$  and this
  circulation sends flow over $e_0^{-1}$. This circulation can be
  decomposed into cycles and one of these cycles $C$ contains
  $e_0^{-1}$.  One has $c_\pi(e_0^{-1}) = -c_\pi(e_0) \geq -  2 \, n \,
  \mu(f_2)$ (eq.~\eqref{f:eq:36}). Using~\eqref{f:eq:32} one obtains
  \begin{eqnarray}
    \label{f:eq:37}
    c(C) & = &  \sum_{e \in C} c_\pi(e) \\
         & \geq &  -  2 \, n \,  \mu(f_2) + (n-1) \mu(f_2) \\
         & = & - (n +1)\, \mu(f_2) \\
         & > & -n \, \mu(f_2). 
  \end{eqnarray}
  %
  The reverse of $C$ is an augmenting cycle for $f'$ with total weight
  at most $n \, \mu(f_2)$ and thus with mean weight at most
  $\mu(f_2)$. Thus $\mu(f') \leq \mu(f_2)$.  \qed
\end{proof}






\section{Computing a minimum cost-to-profit ratio cycle}
\label{sec:comp-minim-cost}

Given a digraph $D = (V,A)$ with costs $c:A\to\setZ$ and profit $p:A \to
\setN_{>0}$, the task is to compute a cycle $C \in \cycles$ with minimum
ratio 
\begin{equation}
  \label{f:eq:31}
  \frac{c(C)}{p(C)}.
\end{equation}

Notice that this is the largest number $\beta\in \setQ$ which satisfies 
\begin{equation}
  \beta\leq  \frac{c(C)}{p(C)}, \, \text{ for all } C \in \cycles. 
\end{equation}

By rewriting this inequality, we understand this to be the largest
number $\beta\in \setQ$ such that 
\begin{equation}
  c(C) - \beta\, p(C) \geq0 \, \text{ for all } C \in \cycles. 
\end{equation}
In other words, given a digraph $D= (V,A)$ with costs 
$c_\beta:A \to \setQ$, where $c_\beta(e) = c(e) - \beta\,p(e)$,
we search the largest number $\beta \in \setQ$ such that the
\begin{equation}
   c_{\beta} \geq 0.
\end{equation}

We need a routine to check whether $D$ has a negative cycle for a
given weight function $c$. For this we assume w.l.o.g. that each 
vertex is reachable from the vertex $s$, if necessary by introducing a
new vertex $s$ from which there is an arc with cost and profit $0$ to
all other nodes. The minimum cost-to-profit ration cycle w.r.t. this
new graph is then the minimum cost-to-profit ratio cycle w.r.t. the
original graph, since $s$ is not a vertex of any cycle. 

Recall the following single-source shortest-path algorithm  of
Bellman-Ford which we now apply with weights $c_\beta$: 


{\small
  \begin{quote}
    Let $n=|V|$ and $m=|A|$. We calculate functions 
    $f_0,f_1,\ldots,f_n:V\longrightarrow\setR\cup\{\infty\}$
    successively by the following rule. 
    
    \begin{enumerate}[i)]
    \item $f_0(s) = 0$, $f_0(v) = \infty$ for all $v \neq s$ 
    \item For $k<n$ if $f_k$ has been found, compute 
      \begin{displaymath}
        \displaystyle f_{k+1}(v) = \min\{f_k(v), \min_{(u,v)\in A}\{f_k(u)+c_\beta(u,v)\}  
      \end{displaymath}
      for all $v \in V$. 
    \end{enumerate}
  \end{quote}
}

As seen when we talked about the Bellman-Ford algorithm, there 
exists a negative cycle w.r.t. $c_\beta$ if and only if 
$f_n(v) <f_k(v)$  for some $v \in V$ and $1\leq k<n$. Thus we can 
test in $O(m\cdot n)$ steps whether $D$ contains a 
negative cycle w.r.t. $c_\beta$. 

We now apply the following idea to  search for the correct value of
$\beta$. We keep an interval $I = [L,U]$ with the invariant that the
value  $\beta$ that we are searching lies in this interval $I$. As
starting values, we can chose $L = c_{min}$ and $U = c_{max}$, 
where $c_{min}$ and $c_{max}$ are the smallest and largest cost
respectively. In one iteration we compute $M = (L + U) /2$. We then
check whether $D$, together with $c_M$ contains a negative cycle. If
yes, we know that $\beta$ is at least $M$ and we set $L\gets M$. If not, then
$\beta$ is at most $M$ and we update the upper bound $U \gets M$. 

When can we stop this procedure? We can stop it, if we can assure that
only one valid cost-to-profit ratio cycle  lies in $[L,U]$. Suppose that $C_1$
and $C_2$ have different cost-to-profit ratios. Then 
\begin{eqnarray}
 | c(C_1) /  p(C_1) - c(C_2)/ p(C_2) | & = & \left| \frac{c(C_1)\, p(C_2) -
 c(C_2)p(C_1)}{  p(C_1) \, p(C_2)}\right| \\
                                & \geq & 1/ (n^2 p_{max}^2). 
\end{eqnarray}
%
Thus we can stop our process, if $U-L< 1/ (n^2 p_{max}^2)$, since we
know then that there can be only one cycle $c \in \cycles$ with
$c(C)/p(C) \in [L,U]$.

Suppose that $[L,U]$ is the final interval. 
We know then that 
\begin{displaymath}
  L \leq c(C) / p(C) \text{ for all } C \in \cycles 
\end{displaymath}
and 
\begin{displaymath}
  U> c(C)/p(C) \text{ holds for some } C \in \cycles. 
\end{displaymath}
Let $C$ be a minimum weight cycle w.r.t. the arc costs $c_L$. Clearly
$ U > c(C)/p(C) \geq L$ holds and thus $C$ is the minimum cost-to-profit
cycle we have been looking for. 

Let us analyze the number of required iterations. We need to halve the
starting interval-length $2 \, c$, where $c$ is the largest absolute
value of a cost, until the length is at most $1/ (n^2 p_{max}^2)$. 
We search the minimal $i\in \setN$ such that
\begin{equation}
  \label{f:eq:35}
  (1/2)^i c \leq 1/(n^2 p_{max}^2).
\end{equation}
This shows us that we need $\bigO(\log( c \,p_{max}^2 n^2))$
iterations which is $\bigO(\log n \log K)$, where $K$ is the largest
absolute value of a cost or a profit. 



\begin{theorem}[Lawler~\cite{MR0439106}]
  \label{f:thr:15}
  Let $D$ be a digraph  with costs $c:A\to\setZ$ and profit $p: A\to \setN_{>0}$
  and let $K \in \setN$   such that $|c(e)| + |p(e)| \leq K$ for all $e \in \setN$. A
  minimum cost-to-profit ratio cycle of $G$ can be computed in time
  $\bigO(m\,n \, \log n\, \log K)$.
\end{theorem}



But we knew a weakly polynomial algorithm for MCNFP from the
exercises. So you surely ask: Can we do better for minimum
cost-to-profit cycle computation? The answer is ``Yes''!



\subsection{Parametric search} 
\label{sec:parametric-search}




Let us first roughly describe  the  idea on  how to obtain a
strongly polynomial algorithm, see~\cite{Megiddo79}.  The Bellman-Ford
algorithm tells us whether our  current $\beta$ is too large or too
small, depending on whether $D$ with weights $c_\beta$ contains a
negative cycle or not. Recall that the B-F algorithm computes labels
$f_i(v)$ for $v \in V$ and $1\leq i\leq n$.  If these labels are computed
with costs $c_\beta$, then they are \emph{piecewise linear} functions in
$\beta$ and we denote them by $f_i(v)[\beta]$. 

Denote the optimal $\beta$ that we look for by $\beta^*$ and suppose that we
know an interval $I$ with such that $\beta^* \in I$ and each function
$f_i(v)[\beta]$ is linear if it is restricted to this domain $I$. Then we
can determine $\beta^*$ as follows.

Let $I = [L,U]$ be the interval and  remember that we are searching for the
largest value of $\beta \in I$ such that $f_{n}(v)[\beta]=f_{n-1}(v)[\beta]$
holds for each $v \in V$. Clearly this holds for $\beta = L$. Thus we only
need to check whether $\beta = U$ by computing the values $f_{n}(v)[U]$
and $f_{n-1}(v)[U]$ for each $v \in V$ and check whether one of these
pairs consists of different numbers. 


The idea is now to compute such an interval $I = [L,U]$ in strongly
polynomial time. 

Consider the function $f_1(v)[\beta]$. Clearly one has 
\begin{displaymath}
f_1(v)[\beta] = 
  \begin{cases}
    c(s,v) - \beta \cdot p(s,v) &  \text{ if } (s,v)\in A, \\
    \infty                  & \text{ otherwise}. 
  \end{cases}
\end{displaymath}

This shows that $f_1(v)[\beta]$ is a linear function in $\beta$ for each $v
\in V$. 

Now suppose that $i\geq1$ and that we have computed an interval
$I=[L,U]$ with  $\beta^*\in I$ and each function $f_i(v)[\beta]$ is a linear
function if $\beta$ is restricted to $I$. 

Now consider the function  $f_{i+1}(v)[\beta]$ for a particular $v \in
V$. Recall the formula 
\begin{equation}
  \label{f:eq:1}
  f_{i+1}(v)[\beta] = \min\{ f_i(v)[\beta], \min_{(u,v)\in A} \{ f_i(u)[\beta] +
  c(u,v) - \beta \cdot p(u,v)\} \}. 
\end{equation}


Each of the functions $ f_i(v)[\beta]$ and $ f_i(u)[\beta] +  c(u,v) - \beta
\cdot p(u,v)$  are linear on $I$.  The function $f_i(v)[\beta]$ can be retrieved by
computing a shortest path $P_i(v)$ from $s$ to $v$ with arc weights $c_\beta$ for
some $\beta$ in $(L,U)$ which uses at most $i$ arcs. If $\beta$ is then
allowed to vary, the line which is defined by $f_i(v)[\beta]$ on $I$ is
then the length of this path $P$   with parameter $\beta$.  Similarly we
can retrieve the functions (lines) $ f_i(u)[\beta] +  c(u,v) - \beta
\cdot p(u,v)$  for each $(u,v) \in A$. With the Bellman-Ford algorithm,
this amounts to a running time of $O(m \cdot n)$. 

We now have $n$ lines and can now compute the lower envelope of these
lines in time $O(n \log n)$ alternatively we can also compute all
intersection points of these lines and sort them w.r.t. increasing
$\beta$-coordinate. This would amount to $O(n^2 \log n)$.
Let $\beta_1,\ldots,\beta_k$ be the sorted list of  these  $\beta$-coordinates.
Now  $\beta_{trial}:= \beta_{\lfloor k/2\rfloor}$ and
check whether $\beta^*>\beta_{trial}$. If yes, we can replace $L$ by
$\beta_{trial}$ and we can delete the numbers
$\beta_1,\ldots,\beta_{\lfloor k/2\rfloor-1}$. Otherwise, we replace $U$ by $\beta_{trial}$ and
delete $\beta_{\lfloor k/2\rfloor+1},\ldots,\beta_k$. In
any case, 
we halved the number of possible $\beta$-coordinates and continue in this
way.  Such a check requires a negative cycle test in the graph
$D$ with arc weights $\beta_{trial}$ and costs $O(m \cdot n)$. In the end we have two consecutive $\beta$-coordinates and have an
interval $[L,U]$ on which  $f_{i+1}(v)[\beta]$ is linear.  To  find
an interval $I$ such that $f_{i+1}(v)[\beta]$ is linear on $I$ and $\beta^*
\in I$ costs thus $O(m \cdot n \log n)$ steps. 


We now continue to tighten \emph{this interval} such that all functions
$f_{i+1}(v)[\beta], v \in V$ are linear on $[L,U]$.  Thus in step $i+1$ this
amounts to a running time of 
\begin{displaymath}
  O\left(n \cdot ( m \cdot n \log n) \right). 
\end{displaymath}

The total running time  is thus 
\begin{displaymath}
  O(n^3 \cdot m \cdot \log n).  
\end{displaymath}

\begin{theorem}
  \label{f:thr:2}
  Let $D=(V,A)$ be a directed graph and let $c: A \longrightarrow\setR$ and
  $p:A\longrightarrow\setR_{>0}$ be functions. One can compute a cycle $C$ of $D$
  minimizing $c(C)/p(C)$ in time  $O(n^3 \cdot m \cdot \log n)$.  
\end{theorem}




\subsubsection{Exercises}

\begin{enumerate}[1)]
\item Show that there are no two different paths from $r$ to another
  node in a directed  tree $T = (V,A)$.  
\item Prove Lemma~\ref{lem:18}. 
\item  
    Why can we assume without loss of generality  that a minimum cost
    network has      a path from $i$ to $j$ for      all $i\neq j\in V$
    which is incapacitated? 
\item 
    Provide an example of a MCNFP for which the simple cycle-canceling
    algorithm from above can require an exponential number of
    cancels, if the cycles are chosen in a disadvantageous way. 
\item   Provide a proof of Theorem~\ref{f:thr:9}. 
\item Let $Q = <u_1,\ldots,u_k>$ be the queue before an iteration of the
  {\bf while} loop of the breadth-first-search algorithm. Show that
  $D[u_i]$ is monotonously increasing and that $D[u_1]+1\geq D[u_k]$. 
  Conclude that the sequence of assigned  labels (over time) is a
  monotonously increasing sequence. \label{item:16}
\item Extend the definitions and properties of flows and networks 
      in a multi-source and multi-sink network. Show that any flow 
      in a multi-source and multi-sink network corresponds to a 
      flow of identical value in the single-source and 
      single-sink network obtained by unifying the sources and 
      sinks in a single super-source and super-sink respectively.
\end{enumerate}



%%% Local Variables: 
%%% mode: latex
%%% TeX-master: "lecture"
%%% End: 

\chapter{The ellipsoid method}
\label{el:cha:ellipsoid-method}


It is not known whether the  simplex algorithm is  an algorithm that
runs in polynomial 
time. For many pivoting rules it was even proved to require an
exponential number of iterations~\cite{MR0332165}.  It was long open, whether there exists a polynomial time
algorithm for linear programming until Khachiyan~\cite{Khachiyan79}
showed that the ellipsoid method\cite{Shor77,MR736268} can solve
linear programs in polynomial time. The remarkable fact is that the
algorithm is polynomial in the binary encoding length of the linear
program. In other words, if the input consists of the problem
$\max\{c^Tx\colon x \in \setR^n, \, Ax \leq b\}$, where $A\in \setQ^{m\times n}$ and $b
\in \setQ^m$, then the algorithm runs in polynomial time in $m+n+s$, where
$s$ is the largest binary encoding length of a rational number
appearing in $A$ or $b$. The question, whether there exists an
algorithm which runs in time polynomial in $m+n$ and performs
arithmetic operations on numbers, whose binary encoding length remains
polynomial in $m+n+s$ is one of the most prominent open problems in
theoretical computer science and discrete optimization.

Initially, the ellipsoid method can be used to solve the following
problem. 
\begin{quote}
  Given a matrix $A \in \setZ^{m\times n}$ and a vector $b \in \setZ^m$, determine a
  feasible point $x^*$ in the polyhedron $P = \{ x \in \setR^n \mid Ax\leq b\}$ or
  assert that $P$ is \emph{not full-dimensional} or $P$ is unbounded.   
\end{quote}
%
After we understand how the ellipsoid method solves this problem in
polynomial time, we discuss why linear programming can be solved in
polynomial time. 

Clearly, we can assume that $A$ has full column rank. Otherwise, we
can find with Gaussian elimination an invertible matrix $U \in
\setR^{n\times n}$ with $A \cdot U = \smat{A' & 0}$ where $A'$ has full column rank. The
system $A'x\leq b$ is then feasible if and only if $Ax\leq b$ is feasible. 

\begin{exercise}
  \label{el:ex:6}
  Let $x'$ be a feasible solution of $A'x\leq b$ and suppose that $U$ from
  above is given. Show how to compute a feasible solution $\tilde{x}$ of
  $Ax\leq b$. Also vice versa, show how to compute $x'$, if 
  $\tilde{x}$ is
  given. 
\end{exercise}

The \emph{unit ball} is the set $B = \{ x \in \setR^n \mid \|x\|\leq1\}$ and an
\emph{ellipsoid} $E(A,b)$ is the image of the unit ball under a linear map
$t:\setR^n\to\setR^n$ with $t(x) = Ax +b$, where $A \in \setR^{n\times n}$ is an invertible
matrix and $b \in \setR^n$ is a vector.  Clearly 
\begin{equation}
  \label{el:eq:113}
  E(A,b) =  \{ x \in \setR^n \mid \|  A^{-1} x - A^{-1} b \| \leq1\}. 
\end{equation}
%



\begin{exercise}
  \label{el:ex:8}
  Consider the mapping $t(x) = \smat{1 & 3 \\ 2 & 5}
  \smat{x_1\\x_2}$. Draw the ellipsoid which is defined by $t$. What
  are the axes of the ellipsoid? 
\end{exercise}


The \emph{volume} of the unit ball is denoted by $V_n$, where $V_n \sim
\frac{1}{\pi\,n} \left( \frac{2\,e\,\pi}{n}\right)^{n/2}$. It follows that
the volume of the ellipsoid $E(A,b)$ is equal to $|\det(A)| \cdot V_n$.
The next lemma is the key to the development of the ellipsoid method. 

\begin{lemma}[Half-Ball Lemma]
  \label{el:lem:12}
  The half-ball $H=\{x \in \setR^n \mid \|x\|\leq1, x_1\geq0\}$ is contained in the
  ellipsoid 
  \begin{equation}
    \label{el:eq:13}
    E = \left\{ x \in \setR^n \mid \left(\frac{n+1}{n}\right)^2  \left( x_1 -
      \frac{1}{n+1}\right)^2 + \frac{n^2-1}{n^2} \sum_{i=2}^n x_i^2 \leq1 \right\}
  \end{equation}
\end{lemma}



\begin{figure}[htbp]
  \centering 
  {
    \includegraphics{figures/ellipsoid1.pdf}
  }
  \caption{Half-ball lemma.}
  \label{el:fig:half-ball}
\end{figure}


\begin{proof} 
  Let $x$ be contained in the unit ball, i.e., $\|x\|\leq1$ and suppose
  further that $0\leq x_1$ holds.  We need to show that 
  \begin{equation}
    \label{el:eq:16}
    \left(\frac{n+1}{n}\right)^2  \left( x_1 -
      \frac{1}{n+1}\right)^2 + \frac{n^2-1}{n^2} \sum_{i=2}^n x_i^2 \leq1 
  \end{equation}
  holds. 
  Since $\sum_{i=2}^nx_i^2\leq1 - x_1^2$ holds we have 
  \begin{equation}
    \label{el:eq:15}
    \begin{split}      
    \left(\frac{n+1}{n}\right)^2  \left( x_1 -
      \frac{1}{n+1}\right)^2 & + \frac{n^2-1}{n^2} \sum_{i=2}^n x_i^2 \\    
                            & \leq     \left(\frac{n+1}{n}\right)^2  \left( x_1 -
      \frac{1}{n+1}\right)^2 + \frac{n^2-1}{n^2} (1-x_1^2)     
  \end{split}
\end{equation}
This shows that~\eqref{el:eq:16} holds if $x$ is contained in the
half-ball and $x_1=0$ or $x_1=1$. Now consider the right-hand-side
of~\eqref{el:eq:15} as a function of $x_1$, i.e., consider 
\begin{equation}
  \label{el:eq:17}
  f(x_1) = \left(\frac{n+1}{n}\right)^2  \left( x_1 -
    \frac{1}{n+1}\right)^2 + \frac{n^2-1}{n^2} (1-x_1^2).    
\end{equation}
The first derivative is 
\begin{equation}
\label{el:eq:18}
  f'(x_1) = 2\cdot \left(\frac{n+1}{n}\right)^2  \left( x_1 -
    \frac{1}{n+1}\right) -  2\cdot \frac{n^2-1}{n^2} x_1.    
\end{equation}
We have $f'(0)<0$ and since both $f(0)=1$ and $f(1)=1$ (and since 
$f(x_1)$ is a 2-degree polynomial w.r.t. $x_1$),  we have
$f(x_1)\leq1$ for all $0\leq x_1\leq1$ and the assertion follows. 
\end{proof}



In terms of a matrix $A$ and a vector $b$, the ellipsoid $E$ is
described as $E = \{ x \in \setR^n \mid \| A^{-1} x - A^{-1}b\|\}$, where $A$
is the diagonal matrix with diagonal  entries
\begin{displaymath}
  \frac{n}{n+1},\sqrt{\frac{n^2}{n^2-1}},\ldots,\sqrt{\frac{n^2}{n^2-1}}
\end{displaymath}
and $b$ is the vector $b = (1/(n+1),0,\ldots,0)$. Our ellipsoid $E$ is thus
the image of the unit sphere under the linear transformation
$t(x)=Ax+b$. 
The determinant of $A$ 
is thus $\frac{n}{n+1}\left(\frac{n^2}{n^2-1}\right)^{(n-1)/2}$ which is
bounded by 
\begin{equation}
  \label{el:eq:19}
  e^{-1/(n+1)} e^{(n-1)/(2\cdot (n^2-1))} = e^{-\frac{1}{2(n+1)}}.  
\end{equation}
%
We can conclude the following theorem.
\begin{theorem}
  \label{el:thr:19}
  The half-ball $\{ x \in \setR^n \mid x_1\geq0, \, \|x\|\leq1\}$ is contained in an
  ellipsoid $E$, whose volume is bounded by $ e^{-\frac{1}{2(n+1)}} \cdot
  V_n$. 
\end{theorem}
%
% Consider now the unit-ball $B$ intersected with a half-space through
% the origin, thus the set 
% \begin{equation}
%   \label{el:eq:14}
%  H_c= B \cap (c^Tx\geq0) = \{ x \in \setR^n \mid \|x\|\leq1, \, c^Tx\geq0\}. 
% \end{equation}
% Clearly $H_c$ is the image of a \emph{rotation} $q(x) = Q\cdot x$, where
% $Q\in \setR^{n\times n}$ is an orthonormal matrix. The rotation  $q$ maps 
% $c^Tx\geq0$ to the half-space $x(1)\geq0$. Recall 
% that a matrix is called \emph{orthonormal} if $Q^TQ = I$. 
% \begin{exercise}
%   \label{el:ex:5}
%   Show that such a $Q$ exists. \emph{Hint: Recall the Gram-Schmidt
%     orthogonalization procedure and apply it to a basis of $\setR^n$ which
%     contains $c$}    
% \end{exercise}
% Since $H_c$ is the image of $H$ under the linear map $q(x)=Q\,x$ it
% follows from Theorem~\ref{el:thr:19} that $H_c$ is contained in the image
% of $E$ under the mapping $q$.  In fact, this is an ellipsoid, since it
% is the image of the unit sphere under the mapping $q\circ t$. 


Recall the following notion from linear algebra. A symmetric matrix $A
\in \setR^{n\times n}$ is called positive definite if all its eigenvalues are
positive.  Recall the following theorem. 
\begin{theorem}
  \label{el:thr:22}
  Let $A \in \setR^{n\times n}$ be a symmetric matrix. The following are
  equivalent. 
  \begin{enumerate}[i)]
  \item $A$ is positive definite. 
  \item $A = L^TL $, where $L\in \setR^{n\times n} $ is a uniquely determined upper
    triangular matrix. 
  \item $x^TAx>0$ for each $x \in \setR^n \setminus\{0\}$. 
  \item $A = Q^T \textrm{diag}(\lambda_1,\ldots,\lambda_n) Q$, where $Q \in \setR^{n\times n}$ is an
    orthogonal matrix and $\lambda_i \in \setR_{>0}$ for $i=1,\ldots,n$. 
  \end{enumerate}
\end{theorem}


It is now convenient to switch to a different representation of an
ellipsoid. An ellipsoid $\eE(A,a)$ is the set  $\eE(A,a) = \{ x \in \setR^n \mid
(x-a)^TA^{-1}(x-a)\leq 1\}$, where $A \in \setR^{n\times n}$ is a symmetric positive
definite matrix and $a\in \setR^n$ is a vector, called the 
\emph{center} of the ellipsoid.  Consider the half-ellipsoid
$\eE(A,a) \cap (c^Tx\leq c^Ta)$.

Our goal is a similar lemma as the half-ball-lemma for ellipsoids.
Geometrically it is clear that each half-ellipsoid $\eE(A,a) \cap
(c^Tx\leq c^Ta)$ must be contained in another ellipsoid $\eE(A',b')$
with $\vol(\eE(A',a'))/ \vol(\eE(A,a))\leq e^{-1/(2n)}$.  More precisely
this follows from the fact that the half-ellipsoid is the image of the
half-ball under a linear transformation. Therefore the image of the
ellipsoid $E$ under the same transformation contains the
half-ellipsoid.  Also, the volume-ratio of the two ellipsoids is
invariant under a linear transformation.


We now record the formula for the ellipsoid $\eE'(A',a')$. 
It is defined by 
\begin{eqnarray}
  \label{el:eq:20}
  a' &  = & a - \frac{1}{n+1}b \\
  A' & = & \frac{n^2}{n^2-1}\left( A - \frac{2}{n+1}b\, b^T   \right),
\end{eqnarray}
where $b$ is the vector $b = A\, c/ \sqrt{c^TA\,c} $.  
The proof of the correctness of this formula  can be found
in~\cite{GroetschelLovaszSchrijver88}.  
\begin{lemma}[Half-Ellipsoid-Theorem]
  \label{el:thr:20}
  The half-ellipsoid $\eE(A,b) \cap (c^Tx\leq c^Ta)$ is contained in the
  ellipsoid $\eE'(A',a')$ and one has $\vol(\eE')/
  \vol(\eE)\leq e^{-1/(2n)}$. 
\end{lemma}



%\begin{proof}
%  If $c^T$ is the vector $(-1,0,\ldots,0)$, and $a=0$, then 
%  \begin{eqnarray*}
%    b & = & (-1,0,\ldots,0)\\
%    a' &=& a - \frac{1}{n+1}b = (1/(n+1),0,\ldots,0)^T\\
%    A' &=& \diag\left(\frac{n^2}{(n+1)^2},\frac{n^2}{n^2-1}
%      ,\ldots,\frac{n^2}{n^2-1}\right) 
%  \end{eqnarray*}
%  and thus the assertion follows from Lemma~\ref{el:lem:12}.


%  We now denote the matrix
%  $\diag\left(\frac{n^2}{(n+1)^2},\frac{n^2}{n^2-1}
%    ,\ldots,\frac{n^2}{n^2-1}\right)$ by $D$ and the vector
%  $(1/(n+1),0,\ldots,0)^T$ by $d$.
% To this end, recall the notation from above, where $B$ denotes the
% unit ball and $E$ denotes the ellipsoid $\eE(D,d)$ which contains
% the half-ball $B \cap (x(1)\geq0)$.  

% Let $\eE(A,a)$ be any ellipsoid and let $\eE'(A',a')$ be the
% ellipsoid defined by~\eqref{el:eq:20} above. The matrix $A$ is
% positive definite and thus can be factored into $A = F^TF$
% Clearly there exists a \emph{rotation} or orthonormal matrix $Q$
% which maps $Fc$ to a positive multiple of $(-1,0,\ldots,0)^T$, i.e., with
%\begin{equation}
%  \label{el:eq:21}
%  (-1,0,\ldots,0)^T = Q F c/ \|Q F c\|. 
%\end{equation}
%The function
%\begin{equation}
%  \label{el:eq:22}
%  T(x) =  F Q^T x +a, 
%\end{equation}
%is a bijective affine transformation with
%$T^{-1}(x)=QF^{-1}(x-a)$.  
%One has
%\begin{eqnarray*}
%  T(B) & = & \{ T(y) \mid y^Ty\leq1\} \\
%       & = & \{ x \mid (T^{-1}(x))^T(T^{-1}(x))\leq1 \\
%       & = & \{ x \mid (x-a)^T (F^{-1})^TQQ^TF^{-1}\\
%       & = & \{ x \mid (x-a)^T A^{-1} (x-a)\leq1\}\\
%       & = & \eE.
%\end{eqnarray*}
%The image of $E$ under $T(x)$ is 
%\begin{eqnarray*}
% T(B) & = & \{ T(y) \mid (y-d)^TD^{-1} (y-d)\leq1\} \\
%       & = & \{ x \mid (T^{-1}(x)-d)^T D^{-1} (T^{-1}(y)-d)\leq1 \\
%       & = & \{ x \mid (x-d-a)^T (F^{-1})^TQDQ^TF^{-1}\\
%       & = & \{ x \mid (x-a)^T A^{-1} (x-a)\leq1\}\\
%       & = & \eE.
%\end{eqnarray*}
%\end{proof}

Before talking about the method, we give a useful definition.

\begin{definition}
   A polyhedron $P$ is full-dimensional if it has positive volume.
\end{definition}

\section{The method}
\label{el:sec:method}



Suppose  we know the following things of our polyhedron $P$. 
\begin{enumerate}[I)]
\item We have a number $L$ such that $\vol(P)\geq L$ if $P$ is
  full-dimensional.
\item We have an ellipsoid $\eE_{init}$ which contains $P$ if $P$ is
  bounded.
\end{enumerate}
The ellipsoid method is now easily described.


\begin{algorithm}[Ellipsoid method exact version]
\label{el:alg:3}
~

\begin{enumerate}[a)]
\item (Initialize): Set $\eE(A,a) := \eE_{init}$ \label{el:item:2}
\item If $a \in P$, then assert $P \neq \emptyset$ and stop  \label{el:item:1} 
\item If $\vol(\eE)<L$, then assert that $P$ is unbounded or $P$ is
  not full-dimensional \label{el:item:3}
\item Otherwise, compute an inequality $c^Tx\leq\beta$ which is valid for
  $P$ and satisfies $c^Ta>\beta$ and replace $\eE(A,a)$ by $\eE(A',a')$
  computed with formula~\eqref{el:eq:20} and goto
  step~\ref{el:item:1}).\label{el:item:4}  
\end{enumerate}
\end{algorithm}
%
\begin{theorem}
  \label{el:thr:21}
  The ellipsoid method computes a point in the polyhedron $P$ or
  asserts that $P$ is unbounded or not full-dimensional. The number of
  iterations is bounded by  $2\cdot n \ln (\vol(\eE_{init}) / L)$.
\end{theorem}


\begin{proof}
  Unless $P$ is unbounded, we start with an ellipsoid which contains
  $P$. This then holds for all the subsequently computed
  ellipsoids. After $i$ iterations one has
  \begin{equation}
    \label{el:eq:23}
    \vol(\eE)/ \vol(\eE_{init}) \leq e^{-\frac{i}{2n}}.
  \end{equation}
 Since we stop when $\vol(\eE)<L$, we stop at least after $2\cdot n
 \ln(\vol(\eE_{init})/ L)$ iterations. This shows the claim.   
\end{proof}



\section{Deciding feasibility}
\label{sec:deciding-feasibility}
Suppose that we have to decide the feasibility of $P = \{x ∈  ℝ^n : Ax ≤ b\}$
using the ellipsoid method described above. Here $A ∈ ℤ^{m ×n}$
and $b ∈ ℤ^{m}$.
In the following $B$
is an upper bound on the absolute values of the components of $A$
and $b$.

To apply the ellipsoid method we have to take care of the following items: 
\begin{enumerate}[i)]
\item We have to find a starting ball that contains a subset of the feasible region. \label{item:4}
\item We have to deal with the fact that $P$ might not be full dimensional even if it is feasible.  \label{item:7}
\item We have to bound the volume $P$ from below, if $P$ is full dimensional. \label{item:8}
\end{enumerate}


We first deal with the issue~\ref{item:4}). 
Suppose that there exists a feasible point $x^*$ and consider the index sets 
$I = \{ i : x^*_i ≥0 \}$ and $J = \{j : x^*_j ≥0\}$.    The polyhedron 
\begin{displaymath}
P' =   \{ x ∈ ℝ^n : Ax ≤b, x_i ≥ 0, i ∈ I, \, x_j ≤ 0, j ∈ J\}
\end{displaymath}
is contained in $P$ and it has vertices. Clearly $P'$ is described by a system of inequalities $Cx≤d$ where each component of $C$ and $d$ is bounded by $B$ in absolute value. 

A vertex $v^*$ of $P'$ is of the form $v^* = C_B^{-1} b_B$ for a basis $B$ of $C$. Since $C_B^{-1} = \widetilde{C_B} / \det(C_B)$ where $\widetilde{C_B} \in ℤ^{n ×n}$ is the adjoint of $C_B$. Each component of $\widetilde{C_B}$ is, by the Hadamard bound, bounded by $B^n ⋅ n^{n/2}$.  It follows that $\|v^*\|_∞$ is bounded by 
 $n^{n/2+1} ⋅ B^{n+1}$. We can thus conclude that $P$ is infeasible if and only if 
 \begin{displaymath}
   P ∩ \{ x ∈ ℝ^n : \|x\|_∞ ≤ n^{n/2+1} ⋅ B^{n+1} \}
 \end{displaymath}
is infeasible. 
For this polyhedron, we have a starting ellipsoid which is the ball of radius $n^{n/2+1} ⋅ B^{n+1} $ around $0$.  


Next, we deal with the issue~\ref{item:7}). 
\begin{exercise}
  \label{el:ex:11}
  Let $P = \{ x \in \setR^n \mid Ax\leq b\}$ be a polyhedron and $\varepsilon>0$ be a
  real number. Show that $P_\varepsilon = \{ x \in \setR^n \mid Ax\leq b+ \varepsilon\cdot
  \mathbf{1}\}$ is full-dimensional if $P \neq \emptyset$.
\end{exercise}


The above exercise raises the following question.  Is there an $\varepsilon>0$
such that $P_\varepsilon = \emptyset$ if and only if $P = \emptyset$ and furthermore is the
binary encoding length of this $\varepsilon$ polynomial in the binary encoding
length of $A$ and $b$? 

We recall the Farkas' lemma (Theorem~\ref{thr:7}). 
\begin{theorem}
  \label{el:thr:24}
  The system $Ax\leq b$ does not have a solution if and only if there
  exists a nonnegative vector $\lambda \in \setR^m_{\geq0}$ such that $\lambda^TA = 0$
  and $\lambda^Tb=-1$. 
\end{theorem}


%Again, suppose  $A\in \setZ^{m\times n}$ and $b\in \setZ^m$ and let $B$ be the largest absolute
%value of a coefficient of $A$ and $b$.  If $Ax\leq b$ is not feasible,
%then there exists a $\lambda\geq0$ such that $\lambda^T(A|b) = (\mathbf{0}|-1)$. 
If $Ax ≤b$ is infeasible, then the polyhedron 
\begin{displaymath}
  \{ λ ∈ ℝ^m : λ^TA = 0, \, λ^T b = -1 , \, λ ≥0 \} 
\end{displaymath}
has vertices. As we have argued before, there exists a vertex $λ^*$ of $P$ such that $\|λ^*\| ≤ n^{n/2+1} ⋅  B^{n+1}  $. 

Now let $\varepsilon =  (2 ⋅ m ⋅ n^{n/2+1} ⋅ B^{n+1})^{-1}$. Then $|{λ^*}^T \mathbf{1} \cdot \varepsilon |<1$ and thus 
\begin{equation}
  \label{el:eq:42}
  \lambda^T (b+ \varepsilon \cdot \mathbf{1}) <0. 
\end{equation}
%
Consequently the system $Ax\leq b+\varepsilon \cdot \mathbf{1}$ 
is infeasible if and only
if $Ax\leq b$ is infeasible. Furthermore, if $Ax\leq b+\varepsilon \cdot \mathbf{1}$  is feasible, then the polyhedron described by this set of inequalities is full-dimensional. 
Notice again that the encoding length of
$\varepsilon$ is polynomial in the encoding length of $Ax\leq b$.  % and we conclude
% with the main theorem of this section. 

% This discussion is summarized in the following lemma. 

% \begin{lemma}
%   \label{lem:2}
%   Let $A ∈ ℤ^{m ×n}$ and $b ∈ ℤ^m$ and let $B$ be an upper bound on the absolute value of the entries of $A$ and $b$ and let $K = (2 ⋅ m ⋅ n^{n/2+1} ⋅ B^{n+1})$. 
% Then $Ax ≤ b $ is feasible if and only of $ K ⋅ A x ≤ K ⋅ b + \mathbf{1}$ is feasible. The binary encoding length of $ K ⋅ A x ≤ K ⋅ b + \mathbf{1}$ is polynomial in the binary encoding length of $Ax ≤b$. 
% \end{lemma}

Next we deal with issue~\ref{item:8}.  We ask ourselves the following question. If $P = \{x ∈ ℝ^n : Ax ≤b\}$ is full-dimensional and bounded, what is a lower bound on $\vol(P)$? 

If $P$ is full-dimensional, then $P$ contains $n+1$ affinely independent vertices $v_1,\dots,v_{n+1}$, see Exercise~\ref{exercise:3}.  Each of these vertices is of the form $v_i = v_i'/q_i$, where $v_i' ∈ ℤ^n$ and $q_i$ is a $n×n$-sub-determinant of $A$. Consequently 
\begin{eqnarray*}
  \vol(P) & ≥ & (1/  ∏_{i=1}^{n+1} q_i) ⋅ \vol(\conv(v'_1,\dots,v'_{n+1})) \\
          & ≥ & (B^{n}⋅ n^{n/2} )^{-n}⋅ 1/n!,
\end{eqnarray*}where the last inequality follows from the Hadamard bound and Exercise~\ref{exercise:1} and Exercise~\ref{exercise:2}. 
A crude bound on $\vol(P)$ is thus $\vol(P) ≥ 1 / (B^{n^2} ⋅ n^{2⋅n^2})$. 


We can conclude with the following theorem. 

\begin{theorem}
  \label{thr:11}
  Let $Ax ≤ b$ be a system of linear inequalities with $A∈ℤ^{m×n}$ and $b∈ℤ^{m}$. Let $B$ be an upper bound on the absolute value of the entries of $A$ and $b$. The ellipsoid medthod can decide in time polynomial in $n,m$ and $\log B$ whether $Ax ≤ b$ is feasible. 
\end{theorem}


The next theorem follows from binary search. Its proof is an exercise. 



\begin{theorem}
\label{thr:12}
Let $\max\{c^Tx : x ∈ ℝ^n, \, Ax ≤b\}$ be a linear program with $A ∈ ℤ^{m×n}$, $b ∈ ℤ^{m}$, $c ∈ Ζ^{n}$.  Let $B$ be an upper bound on the absolute value of the entries of $A$ and $b$. The ellipsoid method finds an optimal solution of the linear program if one exists in time polynomial in $n,m$ and $\log B$. 
  Let $Ax ≤ b$ be a system of linear inequalities with $A∈ℤ^{m×n}$ and $b∈ℤ^{m}$. Let $B$ be an upper bound on the absolute value of the entries of $A$ and $b$. The ellipsoid can decide in time polynomial in $n,m$ and $\log B$ whether $Ax ≤ b$ is feasible. 
\end{theorem}


\subsection*{Exercises} 

\begin{enumerate}
\item Show the following. If $P⊆ℝ^n$ is a bounded and full-dimensional polyhedron, then there exist vertices $v_1,\dots,v_{n+1}$ of $P$ that are affinely independent, i.e., $v_2-v_1, v_3-v_1, \dots, v_{n+1}-v_1$ are linearly independent. \emph{Hint: If $a^Tx = β$ is some hyperplane, where $a ∈ ℝ^n \setminus \{0\}$, then there exists a vertex of $P$ that is not contained in that hyperplane.}\label{exercise:3}
\item Show that $\vol(\conv(0,e_1,\dots,e_n)) = 1/n!$ \label{exercise:1}
\item Let $a_1,\dots,a_n ∈ℤ^n$ be linearly independent. Show that
  \begin{displaymath}
    \vol(\conv(0,a_1,\dots,a_n)) = |\det(a_1,\dots,a_n)| /n!. 
  \end{displaymath} \label{exercise:2}
\end{enumerate}

\section{The separation problem}
\label{el:sec:separation-problem}

At this point we can already notice a very important fact. Inspect
step~\ref{el:item:4} of the algorithm. What is required here? An
inequality which is valid for $P$ but not for the center $a$ of
$\eE(A,a)$. Such an inequality is readily at hand if we have the
complete inequality description of $P$ in terms of a system
$Cx\leq d$. Just pick an inequality which is violated by $a$. Sometimes
however, it is not possible to describe the polyhedron of a
combinatorial optimization problem with an inequality system
efficiently, simply 
because the number of inequalities is too large.
An example of such a polyhedron is the matching polytope, see
Theorem~\ref{po:thr:18}. 

The great power of the ellipsoid method lies in the fact that we do
not have to \emph{write down} the polyhedron entirely. We only have to
solve the so-called separation problem for the polyhedron, which is
defined as follows. 

\fancybox[Separation Problem]{ Given a
  point $a \in \setR^n$ determine, whether $a\in P$ and if not, compute an
  inequality $c^Tx\leq\beta$ which is valid for $P$ with $c^Ta>\beta$. }



\begin{exercise}
  \label{el:ex:10}
  We are given an undirected graph $G=(V,E)$. A \emph{spanning tree}
  $T$ is a subset $T\subseteq E$ of the edges such that $T$ does not contain
  a cycle and $T$ connects all the vertices $V$. Consider the
  following \emph{spanning tree polytope} $P_{span}$
\begin{eqnarray}
 \sum_{e\in E} x(e)       & = & n-1  \label{el:eq:26}\\
 \sum_{e\in \delta(U)} x(e)   & \geq & 1  \quad \quad \forall \emptyset \subset U \subset V  \label{el:eq:25}\\
                 x(e)  & \leq & 1 \quad \quad  \forall e\in E \\
                 x(e)  & \geq & 0 \quad \quad  \forall e\in E. 
\end{eqnarray}

Let $x$ be an integral solution of $P_{span}$ and define $T=\{e\in E \mid
x(e)=1\}$. The inequality~\eqref{el:eq:26} ensures that exactly $n-1$ edges are
picked. The  inequalities~\eqref{el:eq:25} ensure that $T$ connects the
vertices 
of $G$.  Thus $T$ must be a spanning tree.  Clearly, there are
exponentially many inequalities of type~\eqref{el:eq:25}.  Nevertheless,
a fractional solution of this polytope can be computed using the
ellipsoid method.

Show that the separation problem for $P_{span}$ can be solved in
polynomial time. 

\emph{Hint:} To verify whether a vector $x\in \setR_{\geq0}^{|E|}$ fulfills
inequalities of type~\eqref{el:eq:25}, it is a good idea to recall the
MinCut or MaxFlow problem.

{\small 
\noindent 
Via binary search even an optimal solution can be computed in
polynomial time (in the input length) if we introduce edge costs (you
don't have to show that). In the next semester you will see that any
optimal basis solution is integral and hence defines an optimal
spanning tree w.r.t. the edge costs.
}
\end{exercise}

\begin{exercise}
  \label{el:ex:7}
  Consider the triangle defined by
  \begin{eqnarray*}
    -x_1-x_2 & \leq & -2 \\
    3 x_1     & \leq & 4\\
    -2x_1 + 2x_2 & \leq & 3.
  \end{eqnarray*}
  
  Draw the triangle and simulate the ellipsoid method with starting
  ellipsoid being the ball of radius $6$ around $0$. Draw each of the
  computed ellipsoids with your favorite program (pstrics, maple,\ldots ).
  How many iterations does the ellipsoid method take?

  \emph{Ignore the occurring rounding errors!} 
\end{exercise}

% \section{How to start and when to stop }
% \label{el:sec:how-start-when-1}

% In our description of the ellipsoid method, we did not explain yet what
% the initial ellipsoid  is and when we can stop with asserting
% that $P$ is either not full-dimensional or unbounded. 


% Suppose therefore that $P = \{ x \in \setR^n \mid Ax\leq b\}$ is
% full-dimensional and bounded with $A\in \setZ^{m\times n}$ and $b\in \setZ^m$. Let
% $B$ be the largest absolute value of a component of $A$ and $b$. In this 
% section we will show the following things.  
% \begin{enumerate}[i)]
% \item The vertices of $P$ are in the box 
% $\{x \in \setR^n \mid -n^{n/2} B^n \leq x_i \leq n^{n/2} B^n$ 
% $\forall 1 \leq i \leq n \}$. 
% Thus $P$ is contained in the ball around $0$
%   with radius $n^{n} B^n$ . Observe that the encoding length of this
%   radius is $\size(n^{n}B^n) = O(n \log n + n\size(B))$ which is
%   polynomial in the dimension $n$ and the largest encoding length of a
%   coefficient of $A$ and $b$. \label{el:item:7}
% \item The volume of $P$ is bounded from below by  $1 /
% (n\cdot B)^{3n^2}$. \label{el:item:8}
% \end{enumerate}

% The following lemma is proved in any linear algebra course. 
% \begin{lemma}[Inverse formula and Cramer's rule]
%   \label{el:lem:11}
%   Let $C \in \setR^{n\times n}$ be a nonsingular matrix. Then 
%   \begin{displaymath}
%     C^{-1}(j,i) = (-1)^{i+j} \det(C_{ij}) / \det(C),
%   \end{displaymath}
%   where $C_{ij}$ is the matrix arising from $C$ by the deletion of the
%   $i^{th}$ row and $j^{th}$ column. If $d\in \setR^n$ is a vector 
%   then the $j^{th}$ component of $C^{-1}d$ is given by
%   $\det(\wt{C})/\det(C)$, where $\wt{C}$ arises from $C$ be 
%   replacing the $j^{th}$ column with $d$. 
%  \end{lemma}

% We now define the size of a rational number $r = p/q$ with $p$ and $q$
% relatively prime integers, a vector $c \in \setQ^n$ and a matrix $A \in
% \setQ^{m\times n}$: 
% \begin{itemize}
% \item $\size(r) = 1 + \lceil\log ( |p|+1) \rceil + \lceil \log( |q|+1) \rceil$
% \item $\size(c) = n+ \sum_{i=1}^n \size(c_i)$
% \item $\size(A) = m\cdot n + \sum_{i=1}^n\sum_{j=1}^m \size(A(i,j))$
% \end{itemize}



% We recall the Hadamard inequality which states that for $A \in
% \setR^{n\times n}$ one has 
% \begin{equation}
%   \label{el:eq:10}
%   |\det(A)| \leq \prod_{i=1}^n \|a_i\|,
% \end{equation}
% where $a_i$ denotes the $i^{th}$ column of $A$.  In particular, if $B$
% is the largest absolute value of an entry in $A$, then 
% \begin{equation}
%   \label{el:eq:11}
%   |\det(A)| \leq n^{n/2} B^{n}. 
% \end{equation}


% %\begin{exercise}[Computing the determinant with modular arithmetic]
% %  \label{el:ex:9}
% %  In this exercise we show how to compute the determinant of a matrix
% %  $A \in \setZ{n\times n}$ in polynomial time. We are all aware of the fact that
% %  Gaussian elimination performs $O(n^3)$ 
% %  \emph{arithmetic operations} to compute an upper triangular matrix
% %  $R \in \setQ^{n\times n}$ from $A$ such that $\det(R) = \det(A)$ and $\det(A)
% %  = \prod_{i=1}^n r(i,i,)$. But could the number of \emph{bits} which are
% %  necessary to represent the intermediate steps in Gaussian
% %  elimination explode~? We do not want to answer this question in this
% %  exercise. Instead we circumvent it by using the Hadamard
% %  inequality.   
% %\end{exercise}

% Now let us inspect the vertices of a polyhedron $P = \{ x \in \setR^n \mid
% Ax\leq b\}$, where $A$ and $b$ are integral and the largest absolute
% value of any entry in $A$ and $b$ is bounded by $B$. A vertex is
% determined as the unique solution of a linear system $A'x=b'$, where
% $A'x\leq b'$ is a subsystem of $Ax\leq b$ and $A'$ is invertible. Using
% Cramer's rule and our observation~\eqref{el:eq:11} we see that the
% vertices of $P$ lie in the box $\{x \in \setR^n \mid -n^{n/2} B^n \leq x_i \leq n^{n/2} B^n,$ $i = 1,\cdots,n\}$.  This
% shows~\ref{el:item:7}).


% Now let us consider a lower bound on the volume of $P$. Since $P$ is
% full-dimensional, there exist $n+1$ affinely independent vertices
% $v_0,\ldots,v_n$ of $P$ which span a \emph{simplex} in $\setR^{n}$ 
% (i.e. the $n$-dimensional polytope equivalent to the convex-hull of 
% these $n+1$ vertices). The
% volume of this simplex is determined by the formula
% \begin{equation}
%   \label{el:eq:24}
%   \frac{1}{n!} \cdot \left|\det \mat{1 & \cdots & 1 \\ v_0& \ldots&v_n} \right|. 
% \end{equation}
% %
% By Cramer's rule and the Hadamard inequality, the common  denominator
% of each component of $v_i$ can be  bounded by $n^{n/2} B^n$. Thus
% \eqref{el:eq:24} is bounded by 
% \begin{equation}
%   \label{el:eq:44}
%   1 / \left(   n^n (n^{\frac{n}{2}} \cdot B^n)^{n+1}      \right) \geq  1 /
%   \left(   n^{3n^2}   B^{2n^2}      \right) \geq 1 / (n\cdot B)^{3\cdot n^2},
% \end{equation}
% which shows \ref{el:item:8}). 


% Now we plug these values into our analysis in
% Theorem~\ref{el:thr:21}. Our initial volume $\vol(\eE_{init})$  is bounded by   the
% volume of the box with side-lengths $2 (n\cdot B)^n$. Thus
% \begin{equation}
%   \label{el:eq:27}
%   \vol(\eE_{init}) \leq(2\cdot n\cdot B)^{n^2}.  
% \end{equation}
% %
% Above we have shown that 
% \begin{equation}
%   \label{el:eq:38}
%   L \geq 1 / (n\cdot B)^{3n^2}. 
% \end{equation}
% Clearly 
% \begin{equation}
%   \label{el:eq:39}
%    \vol(\eE_{init}) / L  \leq 2^{n^2} \cdot (n \cdot B)^{4 \cdot n^2} 
%    \leq (2n \cdot B)^{4\cdot n^2}. 
% \end{equation}
% %
% By Theorem~\ref{el:thr:21} the ellipsoid method performs 
% \begin{equation}
%   \label{el:eq:40}
%   O\left(2\cdot n \cdot \ln \left((2n \cdot B)^{4\cdot n^2}\right) \right)
% \end{equation}
% iterations. This is bounded by 
% \begin{equation}
%   \label{el:eq:41}
%   O(n^3 \cdot \ln (n \cdot B)).
% \end{equation}
% Now recall that $\log B$ is the number of \emph{bits}  which are
% needed to encode the coefficient with the largest absolute value of
% the constraint system $Ax\leq b$ and that $n$ is the number of variables
% of this system. Therefore the expression~\eqref{el:eq:41} is polynomial
% in the binary input encoding of the system $Ax\leq b$. We conclude the
% following theorem. 

% \begin{theorem}
%   \label{el:thr:23}
%   The ellipsoid method (exact version)  performs a polynomial number
%   of iterations.  
% \end{theorem}




% \section{The boundedness and full-dimensionality condition}
% \label{el:sec:bound-full-dimens}



% In this section we want to show how the ellipsoid method can be used
% to solve the following problem. 
% \begin{quote}
%   Given a matrix $A \in \setZ^{m\times n}$ and a vector $b \in \setZ^m$, determine a
%   feasible point $x^*$ in the polyhedron $P = \{ x \in \setR^n \mid Ax\leq b\}$ or
%   assert that $P = \emptyset$. 
% \end{quote}
% %
% \subsection{Boundedness}
% \label{el:sec:boundedness}

% We have argued that the matrix $A \in \setZ^{m\times n}$ can be assumed to have
% full column rank. So, if $P$ is not empty, then $P$ does have at least
% one vertex. The vertices are contained in the box $\{x \in \setR^n \mid
% -n^{n/2} B^n \leq x_i \leq n^{n/2} B^n,$ $i=1,\cdots,n\}$. Therefore, we can append the
% inequalities $-n^{n/2} B^n \leq x_i \leq n^{n/2} B^n$ 
% $(i=1,\cdots,n)$ to $Ax\leq b$ without
% changing the status of $P \neq \emptyset$ or $P = \emptyset$. Notice that the
% \emph{binary encoding length} of the new inequalities is polynomial in
% the binary encoding length of the old inequalities.



% \begin{theorem}
%   \label{el:thr:25}
%   The ellipsoid method can be used to decide whether a system of
%   inequalities $Ax\leq b$ contains a feasible point, where $A\in\setZ^{m\times n}$
%   and $b\in \setZ^m$. The number of iterations is bounded by a polynomial
%   in $n$ and $\log B$, where $B$ is the largest absolute value of a
%   coefficient of $A$ and $b$.
% \end{theorem}

\section{The ellipsoid method for optimization}
\label{el:sec:ellips-meth-optim}

Suppose that you want to solve a linear program 
\begin{equation}
  \label{el:eq:43}
  \max \{ c^Tx \mid x \in \setR^n, \, Ax\leq b\} 
\end{equation}
and recall that if~\eqref{el:eq:43} is bounded and feasible, then so is
its dual and the two objective values are equal. Thus, we can use the
ellipsoid method to find a point $(x,y)$  with $c^Tx = b^Ty$, $Ax\leq b$
and $A^Ty=c, y\geq0$. 


However, we mentioned that the strength of the ellipsoid method lies
in the fact that we do not need to write the system $Ax\leq b$ down
explicitly. The only thing which has to be solvable is the
\emph{separation problem}. This is to be exploited in the next
exercise. 


\begin{exercise}
  \label{el:ex:12}
  Show how to solve the optimization problem $\max\{c^Tx \mid Ax\leq b\}$
  with a polynomial number of calls to an algorithm which solves the
  separation problem for $Ax\leq b$. You may assume that $A$ has full
  column rank and the polynomial bound on the number of calls to the
  algorithm to    solve the separation problem can depend on $n$ and
  the largest size of a component of $A,b$ and $c$. 
\end{exercise}


\section{Numerical issues}
\label{el:sec:numerical-issues}

We did not discuss the numerical details on how to implement the
ellipsoid method such that it runs in polynomial time. One issue is
crucial. 

\fancybox[!]{We only want to compute with a precision which is  polynomial in the input encoding! }


 In the formula~\eqref{el:eq:20} the vector $b$ is defined by
 taking a square root. The question thus rises on how to round the
 numbers in the intermediate ellipsoids such that they can be handled
 on a machine.  Also one has to analyze the growth of the numbers in
 the course of the algorithm. All these issues  can be overcome but we
 do not discuss them in this course. I would like to refer you to the
 book of Alexander Schrijver~\cite{Schrijver86}  for further
 details. They are not difficult, but a little technical. 




%\bibliographystyle{abbrv}

%\bibliography{mybib,papers,books,my_publications}


%%% Local Variables: 
%%% mode: latex
%%% TeX-master: "lecture"
%%% End: 

\chapter{Primal-Dual algorithm}
\label{cha:pdalgo}


In this chapter we talk about matching problems, a very
important topic in combinatorial optimization.

\section{Graphs and Matchings}
\label{sec:gandm}

%Recall the following definitions, which will be very useful in this 
%chapter.

%\begin{definition}[Undirected Graph]
%   An \emph{undirected graph} $G=(V,E)$ consists of a set $V$
%   of  \emph{vertices} (also called \emph{nodes}) and a set 
%   $E$ of \emph{edges}. An edge $e=\{u,v\}$ is a 
%   2-elements subset of $V$, where $u$ and $v$ are called the 
%   \emph{endpoints} of $e$. 
%\end{definition} 

%\begin{definition}[Directed Graph]
%   A \emph{directed graph} (also called \emph{digraph)} $G=(V,A)$ 
%   consists of a set $V$ of  \emph{vertices} (also called 
%   \emph{nodes}) and a set $A$ of \emph{arcs}. An arc $a$ is an 
%   ordered pair $(u,v) \in V\times V$, where $u$ and $v$ are called 
%   the \emph{tail} and the \emph{head} of $a$ respectively. 
%   Moreover $u$ and $v$ are also called the \emph{endpoints} of $a$.
%\end{definition} 

We begin by giving some useful definitions.

\begin{definition}
   A graph $G=(V,E)$ (or $G=(V,A)$) is called \emph{bipartite} 
   if we can split $V$ into two disjoint subsets $V_1$ and $V_2$,
   such that for every $e \in E$ (or $a \in A$) one endpoint is
   in $V_1$ and the other one is in $V_2$.
\end{definition}


We recall an important result about bipartite graphs.

\begin{lemma}
   A graph is bipartite if and only if it does not contain
   an odd cycle (that is, a cycle of odd length).
\end{lemma}

\begin{definition}[Matching]
   Let $G=(V,E)$, or $G=(V,A)$, be a graph. A \emph{matching} 
   in $G$ is a subset $M \subseteq E$, or $M \subseteq A$, of
   pairwise disjoint edges, or arcs (that is, edges, or arcs,
   which do not share any vertex).
   A vertex $v$ that is an endpoint of an edge (or arc) in 
   $M$ is called a \emph{matched node}, otherwise it is 
   called an \emph{exposed node}.
   A matching is called \emph{perfect} if there are no
   exposed nodes.
\end{definition}

\begin{definition}
   Let $G=(V,E)$ be a graph and $M \in E$ be a matching in $G$.
   An \emph{alternating path} is a path that alternates between
   edges in $M$ and edges in $E \backslash M$. An alternating
   path that starts and ends at an exposed node is called an
   \emph{augmenting path}.
\end{definition}

\begin{definition}[Vertex cover]
   Let $G=(V,E)$ be a graph. A \emph{vertex cover} is a
   set $C \subseteq V$ such that every $e \in E$ has at least
   one endpoint in $C$.
\end{definition}

\begin{figure}[htbp]
  \begin{center}
   \includegraphics{figures/bipartite-graph.pdf}
  \end{center}
    \caption{This picture shows an example of an undirected
    bipartite graph $G=(V,E)$, where $V = \{a,b,c,d,w,x,y,z\}$,
    $E=\{\{a,w\},\{b,y\},\{b,z\},\{c,z\},\{d,w\},\{d,x\}\}$,
    $V_1=\{a,b,c,d\}$ and $V_2=\{w,x,y,z\}$. 
    The thicker edges form a matching, vertices $a$ and $x$ 
    (marked in grey) are exposed nodes and $a,w,d,x$ is an 
    alternating and augmenting path. $C=\{b,w,x,z\}$ is a vertex 
    cover.}
    \label{fig:bipgraph}
\end{figure}

\section{Matching problems}
\label{sec:matprob}

Now that we know some basic definitions, we can
concentrate on two of the most important problems about
matchings:
\begin{enumerate}
   \item \textbf{Maximum cardinality matching problem}:
    Find a matching $M$ of maximum size.
   \item \textbf{Maximum (or minimum) weight matching problem}:
   Given a graph $G=(V,E)$ and a weight function 
   $w:~E\longrightarrow\setR$, find a matching $M$ of maximum
   (or minimum) weight, where the weight of a matching is
   \begin{displaymath}
        w(M) = \sum_{\substack{e \in M}}  w(e). 
   \end{displaymath}
\end{enumerate}

From now we will focus on the case of an undirected bipartite 
graph $G=(V,E)$.

\subsection{The maximum cardinality matching problem}
\label{subsec:maxcarmat}

Before giving a method to find a maximum cardinality matching
in a graph, we want to show how we can prove its optimality. 
For this purpose we recall theorem~\ref{thr:16}, that gives
an upper bound on the size of any matching in a given bipartite
graph, and we prove another theoreme, which gives us a method
to see if a matching is of maximum size or not.

\begin{theorem}[K\"onig's theorem]
   In any bipartite graph, the number of edges in a 
   maximum cardinality matching equals the number of vertices in 
   a minimum vertex cover. 
\end{theorem}

\begin{theorem}
   A matching $M$ of a graph $G=(V,E)$ is of maximum 
   cardinality if and only if there are no augmenting paths
   with respect to $M$.
\end{theorem}

\begin{proof}
   $(\Rightarrow)$ Suppose there exists an augmenting path in
   $G$, call it $P$. Consider the set $M' = M \triangle P$.
   By definition of augmenting path, we have that $M'$ has
   exactly one edge more than $M$. Since $M$ is a matching
   and $P$ starts and ends in an exposed node, the only edges
   of $M$ with an endpoint in $P$ are the edges of $P$. This implies
   that $M'$ is a matching.
   
   $(\Leftarrow)$ Let $M$ be a maximum cardinality matching
   and let $\tilde{M}$ be a strictly smaller matching.
   Consider $X=M \triangle \tilde{M}$. $X$ is formed by
   cycles and alternating paths (with respect to $M$ and
   $\tilde{M}$). Since $ |M| \textgreater
   |\tilde{M}|$, $X$ contains at least one alternating 
   path $P$ with more edges from $M$ than from $\tilde{M}$.
   By definition of $X$, we see that $P$ is an 
   augmenting path. \qed
\end{proof}

We are now ready to show an algorithm that allows us to
find a maximum cardinality matching in any bipartite graph.

Let $G=(V,E)$ be a bipartite graph with bipartition 
$V = A \sqcup B$ and $M$ a matching.
Now turn $G$ into a directed graph $D=(V,A)$ by directing 
matching edges from $A$ to $B$ and non-matching edges from $B$ to $A$.
We are interested in a method to find augmenting paths or to
assert that there aren't any (and thus, that $M$ is of 
maximum size). For this purpose we illustrate the following
claim and theorem.

\begin{claim}
   Let $G$ and $D$ be as above.
   A path in $D$ between two exposed nodes that starts in
   an exposed node in $B$ (resp. in $A$), ends in an exposed 
   node in $A$ (resp. in $B$).   
\end{claim}

\begin{theorem}
   Let $G$ and $D$ be as above and $M$ be a matching.
   There exists an augmenting path in $G$ if and only if 
   there exists a path from an exposed node in $B$ to an
   exposed node in $A$ in the directed graph $D$.
\end{theorem}

\begin{proof}
   $(\Rightarrow)$ This is a direct consequence of the choice
   of the direction of the arcs in $D$ and of the previous 
   claim.
   
   $(\Leftarrow)$ Trivial.
\end{proof}

Let $G$ and $D$ be as illustrated before.
How can we use the previous theorem for our purpose?  
Add a vertex $s$ to our directed graph $D$ and connect it to all 
exposed nodes in $B$ by an arc whose tail is $s$. 
Now we have that finding a directed path, in $D$, from an exposed 
node $u$ in $B$ to an exposed node $v$ in $A$ is equivalent 
to finding a directed path from $s$ to $v$ passing by $u$. 
Thus, we deduce that there is an augmenting path in $G$ if and 
only if there is at least one exposed node in $A$ reachable from $s$. 
To find if there are such nodes reachable from $s$ and the 
corresponding augmenting paths, we can use the 
\emph{Breadth-First search} algorithm (see chapter~\ref{subsec:BFs}) 
applied to vertex $s$ in $D$.
If there is any exposed node $u$ in $A$ with finite distance 
from $s$, then we can obtain an augmenting path by taking
the shortest path $s,a_1,a_2,\cdots,u$ from $s$ to $u$
(which is given by Bread-First search) without $s$ (i.e., the augmenting path would be $a_1,a_2,\cdots,u$).

To resume, we obtain the following algorithm.

\begin{algorithm}
  \begin{tabbing}
     \\
     {\bf Initialise} \= $M = \emptyset$  \\
     {\bf while}      \=  $There$ $exists$ $M$-$augmenting$
                          $path$        \\  
                      \>  $Update$ $M$  \\
     {\bf return}     \= $M$ 
  \end{tabbing}
\end{algorithm}

Finally, we are interested in the running time of our algorithm.

\begin{theorem}
   A maximum cardinality matching in a bipartite graph 
   $G=(V,E)$ can be computed in time $O(|V| \cdot (|V|+|E|))$.
   If we assume that $G$ does not have any isolated vertex,
   then we can consider the running time to be 
   $O(|V| \cdot |E|)$.
\end{theorem}

\begin{proof}
   The \emph{while loop} runs at most $|V|/2$ times and its
   execution requires $O(|V|+|E|)$ 
   ($=O(|E|)$ if we have the assumption) operations.
\end{proof}

\begin{figure}[htbp]
  \begin{center}
   \includegraphics{figures/augpath.pdf}
  \end{center}
    \caption{This picture illustrate how transform a bipartite
    graph to apply the \emph{maximum cardinality matching 
    algorithm}, using the graph of figure~\ref{fig:bipgraph} as
    an example.}
    \label{fig:augpath}
\end{figure}

\subsection{The maximum weight matching problem}
\label{subsec:maxwmat}

In this section we consider a graph $G=(V,E)$ and a weight 
function $w:~E\longrightarrow\setR$. First of all, notice that,
by changing the sign of the weights of all edges, we have that 
finding a maximum weight matching or a minimum weight matching
are equivalent problems.
 
We can also prove 
that the problem of finding a minimum weight matching
can always be replaced by the problem of finding a minimum 
weight \emph{perfect} matching. To see that, it is 
sufficient to create a copy $G'$ of our graph $G$ and add
edges of weight $0$ connecting all vertices $v$ in $G$
to their copy $v'$ in $G'$. This will give us a new graph
$\tilde{G}$. Notice that $\tilde{G}$ has at least a perfect
matching, thus there is one, call it $\tilde{M}$, of minimum 
weight.
If we consider only the edges of $\tilde{M}$ in $G$, we obtain
a minimum weight matching in $G$.

Moreover, we have an important relation between the difficulty
of this two problems, given by the following theorem.

\begin{theorem}
   If there exists a polynomial time algorithm for the minimum
   weight perfect matching problem, then there exists a polynomial 
   time algorithm for the minimum weight matching problem.
\end{theorem}

\begin{proof}
   Let $G=(V,E)$ be a graph. Using the process described above,
   we create the graph $\tilde{G} = (\tilde{V}, \tilde{E})$.
   Notice that $|\tilde{V}| = 2|V|$ and 
   $|\tilde{E}| = 2|E|+|V|$.
   By hypothesis, we can find a minimum weight perfect 
   matching in $\tilde{G}$ (and thus a minimum weight matching
   in $G$) in time $O((\tilde{V} + \tilde{E})^k) =
                    O((2|V|+2|E|+|V|)^k) = O((|V|+|E|)^{k'})$,
   for some $k,k' \in \setN$ (that is, a polynomial time with 
   respect to the size of $G$).
\end{proof}

Since we concentrate on bipartite graphs, we show
how to reduce the problem of finding a minimum weight 
\emph{bipartite} matching to the problem of finding a minimum
weight \emph{bipartite} perfect matching. Let $G=(V,E)$ be a 
bipartite graph with bipartition $V = A \sqcup B$. We can
suppose w.l.o.g. that $|A| \geq |B|$. If $|A|>|B|$, add a
set $C$ of cardinality $|A|-|B|$ and add edges of weight 
zero connecting all nodes of $A$ to all nodes of $B \cup C$.
Call $G'$ the bipartite graph obtained by this process. Obviously,
$G'$ admits at least one bipartite perfect matching, and thus 
also one of minimum weight, call it $M'$. 
By taking out all edges of weight zero from $M'$, we obtain a
minimum weight bipartite matching in $G$.

Also in this case, we have an important relation between
the difficulty of these two matching problems.

\begin{theorem}
   If there exists a polynomial time algorithm for the minimum
   weight bipartite perfect matching problem, then there exists a 
   polynomial time algorithm for the minimum weight bipartite 
   matching problem.
\end{theorem}

\begin{proof}
   Let $G=(V,E)$ be a bipartite graph with bipartition 
   $V = A \sqcup B$. Construct $G'=(V',E')$ as shown above.
   Notice that $|V'| = 2|A| \leq 2|V|$ and 
   $|E'| = |E| + |A| \cdot (|A|-1) \leq |E| + |V|^2$.
   Hence, by hypothesis, finding a minimum weight bipartite 
   matching in $G$ (by finding a minimum weight bipartite perfect 
   matching in $G'$) can be done in time $O((2|V| + |E| + |V|^2)^k)
   = O((|V|+|E|)^{k'})$, for some $k,k' \in \setN$ (that is, a 
   polynomial time with respect to the size of $G$).
\end{proof}

Let $G=(V,E)$ be a bipartite graph. We have showed that the
problem of finding a maximum weight (bipartite) matching can
be reduced to the problem of finding a minimum weight 
(bipartite) perfect matching. Moreover, if there exists a
polynomial time algorithm to solve the first problem, then 
there exists a polynomial time algorithm to solve the second one.

The following theorem is crucial for the introduction of the
\emph{Primal-Dual algorithm}.

\begin{theorem}[Complementary Slackness]
\label{thr:compslack}
   Let 
   \begin{displaymath}
      \begin{array}{c}
         max \ c^T x \\
         Ax \leq b
      \end{array}
   \end{displaymath}
   be a primal LP and let 
   \begin{displaymath}
      \begin{array}{c}
         min \ b^T y \\
         A^T y = c \\
         y \geq 0
      \end{array}
   \end{displaymath}
   be his dual. Let $x^*$ and $y^*$ be primal and dual feasible
   solutions respectively. Then they are both optimal if and
   only if 
   \begin{displaymath}
      (b - A x^*)^T y^* = 0.
   \end{displaymath}
\end{theorem}

\begin{proof}
   \begin{displaymath}
      \begin{array}{c}
         (b - A x^*)^T y^* = 0 \\
         \Leftrightarrow \\
         y^*_i > 0 \Rightarrow a^T_i x^* = b_i \\
         \Leftrightarrow \\
         b^T y^* = (A x^*)^T y^* = (x^*)^T A^T y^* = c^T x
      \end{array}
   \end{displaymath}
   The result follows from strong duality theorem 
   (theorem \ref{thr:4})
\end{proof}

Let $G=(V,E)$ be a bipartite graph.
In order to prove the next theorem, we need to recall the 
\emph{LP-relaxation} of the minimum weight perfect matching 
problem (we will call it the \emph{dual LP}) and his dual (which 
we will call the \emph{primal LP}).

\begin{displaymath}
   \begin{array}{c c}
   PRIMAL                & DUAL \\
   max \sum_{v \in V} y_v & min \sum_{e \in E} w_e \cdot x_e \\
   \ uv \in E: y_u+y_v \leq w_{uv}  \ &
   \ v \in V: \sum_{e \in \delta(v)} x_e = 1 \ \\
   y \in \setR^{|V|}     & x \geq 0
   \end{array}
\end{displaymath}

\begin{theorem}
\label{thr:optmat}
   Suppose $x^*$ and $y^*$ are feasible solutions of the recalled 
   primal and dual linear programs respectively. 
   Then, they are optimal solutions if and only if
   \begin{displaymath}
      \forall uv \in E: 
      x^*_{uv} > 0 \Rightarrow y^*_u + y^*_v = w_{uv}
   \end{displaymath}
\end{theorem}

\begin{proof}
   The proof follows immediately from theorem~\ref{thr:compslack}.
\end{proof}

Another direct consequence of theorem~\ref{thr:compslack} is 
the following criterion for the optimality of a minimum
weight perfect matching.

Let $y^* \in \setR^{|V|}$ be feasible in the primal LP and 
let $G_{y^*}=(V,E_{y^*})$ be the graph with edge set
\begin{displaymath}
   E_{y^*} = {uv \in E: y^*_u + y^*_v = w_{uv}}
\end{displaymath}
If $G_{y^*}$ has a perfect matching $M$, then $y^*$ and $\chi^M$ 
are optimal solutions of the corresponding linear programs,
In particular, since, by definition of $\chi^M$,
\begin{displaymath}
   \chi^M_{uv} > 0 \Rightarrow y^*_u + y^*_v = w_{uv},
\end{displaymath} 
$M$ is a \emph{minimum weight perfect matching}.

\subsection*{The Primal-Dual Algorithm}
Let $G=(V,E)$ be a complete bipartite graph with bipartition 
$V=A \sqcup B$ and edge weights $w: E \rightarrow \setR$. 
The primal-dual algorithm is a method to find a minimum weight 
matching in $G$.
It starts from a feasible solution $y^*$ of the dual LP recalled 
before theorem~\ref{thr:optmat} (notice that a possible 
choice for an initial feasible solution can always be $y^*$ such 
that $y^*_v$ is equal to $0$ or to the smallest edge-weight in 
the graph in the case it is smaller than $0$) 
and it updates this solution until it is optimal. 
The last question we have to answer is: how do we update $y^*$
in such a way that our algorithm terminates correctly?

Let $G_{y^*}$ be as previously defined and let $M$ be a maximum 
cardinality matching in $G_{y^*}$. Turn $G_{y^*}$ into a directed
graph by orientating the matching edges from $A$ to $B$ and the
non-matching edges from $B$ to $A$ (call this new graph 
$\tilde{G}_{y^*}$). 
Let $M$ be a maximum cardinality matching in $G_{y^*}$.
Let $L \subseteq A \sqcup B$ be the set of vertices reachable 
from the set of exposed nodes of $B$ in the directed graph 
$\tilde{G}_{y^*}$.

\begin{claim}
   $G_{y^*}$ does not contain an edge $uv$ with 
   $u \in A \textbackslash L$ and $v \in B \cap L$.
\end{claim}

\begin{proof}
   Let $uv$ be such an edge. If $uv$ is a non-matching edge, then, 
   since $v \in L$, also $u$ has to be in $L$. Suppose now that 
   $uv$ is a matching edge. Since $v \in L$, there must be a 
   matching edge $\tilde{u}v$ with $\tilde{u} \in L$. This implies 
   that $u = \tilde{u}$ and this leads to a contradiction since 
   $u \not\in L$.
\end{proof}

We want to update our solution $y^*$ in such a way to obtain 
an optimal solution. Let 
$\delta = min\{w_{uv}-y^*_u-y^*_v$ $|$
$uv \in E, u \in A \textbackslash L, v \in B \cap L\}$.
If $y^*$ is not already optimal, then $\delta$ is strictly positive.
This allows us to construct a feasible solution $\tilde{y}$ 
in this way:
\begin{enumerate}[i)]
   \item $\tilde{y}_v = y^*_v + \delta$ if 
         $v \in A \textbackslash L$
   \item $\tilde{y}_v = y^*_v - \delta$ if 
         $v \in B \textbackslash L$
   \item $\tilde{y}_v = y^*_v$ otherwise
\end{enumerate}

This gives us a new graph $G_{\tilde{y}}$. 
Notice that the previous claim shows that all the edges of $M$ 
are in $G_{\tilde{y}}$. 
Furthermore, by the definition of $\delta$, at least one edge 
connecting a node in $A \textbackslash L$ to a vertex in 
$B \cap L$ becomes tight. Thus, $L$ is augmented.
This implies that every time we update our solution there 
are two things that can happen:
\begin{enumerate}[i)]
   \item $G_{\tilde{y}}$ contains an augmenting path with respect 
         to $M$. This allows us to increase our matching.
   \item The set $L$ increases.
\end{enumerate}
Since $V$ is a finite set and by the criterion of optimality 
showed before, we can assert that the Primal-Dual algorithm 
terminates correctly.

We conclude by resuming the algorithm, giving his running time 
and an example, which can be useful to understand better how 
the Primal-Dual algorithm works.

\begin{algorithm}
  \begin{tabbing}
     \\
     {\bf Input:}           \= $A$ $feasible$ $solution$ $y^*$   \\
     {\bf Initialisation:}  \= $Turn$ $G$ $into$ $a$ $directed$ 
                               $graph$                           \\  
     {\bf while}            \= $M$ $is$ $not$ $optimal$          \\                               
                            \> $Compute$ $L$                     \\
                            \> $Let$ $\delta = min\{w_{uv}
                                - y^*_u - y^*_v\}$               \\
                            \>$Update$ $y^*$                     \\
     {\bf return}           \= $M$  
  \end{tabbing}
\end{algorithm}

\begin{theorem}
   The Primal-Dual algorithm runs in $O(|V^2|)$ time.
\end{theorem}

\begin{proof}
   We can have at most $|V|$ many augmentation of $L$ 
   without augmenting $M$ and $M$ can be augmented at most $|V| / 2$ 
   times. 
\end{proof}

Figure~\ref{fig:egprimaldual} shows how to apply Primal-Dual 
algorithm on a complete bipartite graph. 

\begin{figure}[htbp]
   \begin{center}
      \includegraphics{figures/egprimaldual.pdf}
   \end{center}
 \caption{The edges of $G_{y^*}$ are coloured, red for 
 the matching edges, and green for the non-matching edges.}
 \label{fig:egprimaldual}
\end{figure}

\section*{Exercises} 
\begin{enumerate}
   \item   Let 
           \begin{displaymath}
              \begin{array}{c}
                 max \ c^T x \\
                 Ax \leq b
              \end{array}
           \end{displaymath}
           be a primal LP and let 
           \begin{displaymath}
              \begin{array}{c}
                 min \ b^T y \\
                 A^T y = c \\
                 y \geq 0
              \end{array}
           \end{displaymath}
           be his dual. Let $x^*$ and $y^*$ be primal and dual    
           feasible solutions respectively. Suppose they are
           both optimal. Is the following true: 
           if $a^T_i x^* = b_i$, then $y^*_i > 0$?
           
   \item   Find $G_{y^*}$ from the graph $G=(V,E)$ shown in
           figure~\ref{fig:exgraph} as described in the 
           criterion for the optimality of a minimum
           weight perfect graph.
          
           \begin{figure}[htbp]
             \begin{center}
               \includegraphics{figures/exgraph.pdf}
             \end{center}
               \caption{$G=(V,E)$}
               \label{fig:exgraph} 
           \end{figure} 
           
   \item   Are there rational numbers $y_1, y_2, y_3, y_4$ such 
           that 
           \begin{displaymath}
              \begin{array}{c}
                 y_1 + y_3 \leq 5 \\
                 y_1 + y_4 = 2 \\
                 y_2 + y_3 = 9 \\
                 y_2 + y_4 \leq 3
              \end{array}
           \end{displaymath}
           
   \item   Apply Primal-Dual Algorithm to the complete bipartite 
           graph of figure~\ref{fig:exprimaldualalgo} (we give 
           an initial feasible solution $y^*=(0,0,0,0)$).
           \begin{figure}[htbp]
             \begin{center}
               \includegraphics{figures/exprimaldualalgo.pdf}
             \end{center}
             \caption{}
               \label{fig:exprimaldualalgo}
           \end{figure}
           
\end{enumerate}


%%% Local Variables: 
%%% mode: latex
%%% TeX-master: "lecture"
%%% End: 
 





\bibliographystyle{abbrv}
\bibliography{mybib,papers,books,my_publications}



\end{document}

%%% Local Variables: 
%%% mode: latex
%%% TeX-master: t
%%% End: 
